\documentclass[11pt,fleqn]{book}

% ============================================================
% ORTHODROMIC INFRASTRUCTURE BLUEPRINT
% A Long-Term Architectural Specification for Civil Convergence
%
% Compile with LuaLaTeX or XeLaTeX.
% ============================================================

\usepackage{fontspec}
\setmainfont{TeX Gyre Pagella}

\usepackage{unicode-math}
\setmathfont{TeX Gyre Pagella Math}

\usepackage[
  letterpaper,
  inner=1.20in,
  outer=1.00in,
  top=1.00in,
  bottom=1.10in,
  headheight=15pt
]{geometry}

\usepackage{amsmath}
\usepackage{amssymb}
\usepackage{mathtools}
\usepackage{amsthm}
\usepackage{xcolor}
\usepackage{mdframed}
\usepackage{booktabs}
\usepackage{longtable}
\usepackage{array}
\usepackage{microtype}
\usepackage{graphicx}
\usepackage{tikz}
\usepackage{pgfplots}
\usepackage{fancyhdr}
\usepackage{titlesec}
\usepackage{hyperref}
\usepackage{cleveref}

\usetikzlibrary{
  arrows.meta,
  calc,
  intersections,
  positioning,
  decorations.pathreplacing,
  patterns,
  angles,
  quotes
}

\pgfplotsset{compat=1.18}

\hypersetup{
  colorlinks=true,
  linkcolor=black,
  citecolor=black,
  urlcolor=black,
  pdftitle={Orthodromic Infrastructure Blueprint},
  pdfauthor={Flyxion}
}

\setlength{\parindent}{1.4em}
\setlength{\parskip}{0.25em}
\setlength{\mathindent}{2em}

\numberwithin{equation}{chapter}

\pagestyle{fancy}
\fancyhf{}
\fancyhead[LE]{\small\thepage\qquad Orthodromic Infrastructure Blueprint}
\fancyhead[RO]{\small\leftmark\qquad\thepage}
\renewcommand{\headrulewidth}{0.4pt}

\titleformat{\chapter}[display]
  {\normalfont\huge\bfseries}
  {\chaptername\ \thechapter}
  {0.6em}
  {\Huge}

\titleformat{\section}
  {\normalfont\Large\bfseries}
  {\thesection}
  {0.8em}
  {}

\titleformat{\subsection}
  {\normalfont\large\bfseries}
  {\thesubsection}
  {0.8em}
  {}

% ============================================================
% SEMANTIC ENVIRONMENTS
% ============================================================

\newmdenv[
  topline=false,
  bottomline=false,
  rightline=false,
  leftline=true,
  linewidth=3pt,
  linecolor=black,
  innertopmargin=9pt,
  innerbottommargin=9pt,
  innerleftmargin=12pt,
  innerrightmargin=12pt,
  skipabove=10pt,
  skipbelow=10pt,
  backgroundcolor=gray!6
]{definitionEnv}

\newenvironment{definition}[1]
  {\begin{definitionEnv}
   \noindent\textbf{Definition: #1}\par\smallskip}
  {\end{definitionEnv}}

\newmdenv[
  topline=false,
  bottomline=false,
  rightline=false,
  leftline=true,
  linewidth=3pt,
  linecolor=gray!55,
  innertopmargin=9pt,
  innerbottommargin=9pt,
  innerleftmargin=12pt,
  innerrightmargin=12pt,
  skipabove=10pt,
  skipbelow=10pt,
  backgroundcolor=gray!2
]{principleEnv}

\newenvironment{principle}[1]
  {\begin{principleEnv}
   \noindent\textbf{Design Principle: #1}\par\smallskip\itshape}
  {\end{principleEnv}}

\newmdenv[
  topline=true,
  bottomline=true,
  rightline=true,
  leftline=true,
  linewidth=1pt,
  linecolor=gray!65,
  innertopmargin=9pt,
  innerbottommargin=9pt,
  innerleftmargin=12pt,
  innerrightmargin=12pt,
  skipabove=10pt,
  skipbelow=10pt,
  backgroundcolor=white
]{specificationEnv}

\newenvironment{specification}[1]
  {\begin{specificationEnv}
   \noindent\textbf{Specification: #1}\par\smallskip}
  {\end{specificationEnv}}

\newmdenv[
  topline=false,
  bottomline=false,
  rightline=false,
  leftline=true,
  linewidth=2pt,
  linecolor=black!80,
  innertopmargin=9pt,
  innerbottommargin=9pt,
  innerleftmargin=12pt,
  innerrightmargin=12pt,
  skipabove=10pt,
  skipbelow=10pt,
  backgroundcolor=gray!4
]{researchEnv}

\newenvironment{research}[1]
  {\begin{researchEnv}
   \noindent\textbf{Research Question: #1}\par\smallskip}
  {\end{researchEnv}}

\newmdenv[
  topline=true,
  bottomline=true,
  rightline=false,
  leftline=false,
  linewidth=0.6pt,
  linecolor=gray!55,
  innertopmargin=8pt,
  innerbottommargin=8pt,
  innerleftmargin=10pt,
  innerrightmargin=10pt,
  skipabove=10pt,
  skipbelow=10pt,
  backgroundcolor=gray!1
]{derivationEnv}

\newenvironment{derivation}[1]
  {\begin{derivationEnv}
   \noindent\textbf{Derivation: #1}\par\smallskip}
  {\end{derivationEnv}}

\newmdenv[
  topline=true,
  bottomline=true,
  rightline=true,
  leftline=true,
  linewidth=1.2pt,
  linecolor=black,
  innertopmargin=8pt,
  innerbottommargin=8pt,
  innerleftmargin=10pt,
  innerrightmargin=10pt,
  skipabove=10pt,
  skipbelow=10pt,
  backgroundcolor=gray!3
]{protocolEnv}

\newenvironment{protocol}[1]
  {\begin{protocolEnv}
   \noindent\textbf{Protocol: #1}\par\smallskip}
  {\end{protocolEnv}}

% ============================================================
% THEOREM ENVIRONMENTS
% ============================================================

\theoremstyle{plain}
\newtheorem{theorem}{Theorem}[chapter]
\newtheorem{proposition}[theorem]{Proposition}
\newtheorem{lemma}[theorem]{Lemma}
\newtheorem{corollary}[theorem]{Corollary}

\theoremstyle{definition}
\newtheorem{formaldefinition}[theorem]{Formal Definition}
\newtheorem{construction}[theorem]{Construction}

\theoremstyle{remark}
\newtheorem{remark}[theorem]{Remark}
\newtheorem{example}[theorem]{Example}

% ============================================================
% MATHEMATICAL COMMANDS
% ============================================================

\newcommand{\Sphere}{\mathbb{S}^{2}}
\newcommand{\Ball}{\mathbb{B}^{3}}
\newcommand{\R}{\mathbb{R}}
\newcommand{\N}{\mathbb{N}}
\newcommand{\Oset}{\mathcal{O}}
\newcommand{\Nset}{\mathcal{N}}
\newcommand{\Cset}{\mathcal{C}}
\newcommand{\Vset}{\mathcal{V}}
\newcommand{\Dset}{\mathcal{D}}
\newcommand{\Gset}{\mathcal{G}}
\newcommand{\Mset}{\mathcal{M}}
\newcommand{\Pset}{\mathcal{P}}
\newcommand{\Aset}{\mathcal{A}}
\newcommand{\dist}{\operatorname{dist}}
\newcommand{\argmin}{\operatorname*{arg\,min}}
\newcommand{\argmax}{\operatorname*{arg\,max}}
\newcommand{\Vor}{\operatorname{Vor}}
\newcommand{\Del}{\operatorname{Del}}
\newcommand{\conv}{\operatorname{conv}}
\newcommand{\area}{\operatorname{area}}
\newcommand{\length}{\operatorname{length}}
\newcommand{\diam}{\operatorname{diam}}
\newcommand{\sgn}{\operatorname{sgn}}
\newcommand{\id}{\operatorname{id}}
\newcommand{\Pareto}{\operatorname{Pareto}}
\newcommand{\esssup}{\operatorname*{ess\,sup}}
\newcommand{\suchthat}{\;\middle|\;}

% ============================================================
% DOCUMENT METADATA
% ============================================================

\title{
  \textbf{Orthodromic Infrastructure Blueprint}\\[0.5em]
  \Large A Long-Term Architectural Specification\\
  for Civil Convergence
}

\author{
  Flyxion\\
  \small Independent Researcher
}

\date{\today}

\begin{document}

\frontmatter

\maketitle

\begin{abstract}
\noindent
This monograph develops a mathematical and architectural specification for a planetary infrastructure system organized around a finite arrangement of orthodromes, or great circles, on the terrestrial sphere. The proposal is not presented as a prediction that a particular technology, political order, or economic institution will prevail. It is instead formulated as a protocol for preserving compatibility among civil works produced by different societies at different levels of technological capacity and at widely separated historical moments.

The central architectural thesis is that planetary geometry can serve as a stable intergenerational interface. The physical realization of a corridor may change from a surveyed reserve to a road, hydraulic conduit, electrical backbone, guided transportation line, geothermal installation, or linear mass-acceleration structure. The invariant element is not the machinery occupying the corridor but the geometric and topological relation among corridors, nodes, reserved volumes, and permissible interfaces. The system therefore separates spatial protocol from technological implementation.

The adopted framework consists of nine primary orthodromes whose intersections generate a global arrangement graph. Five orthodromes are associated with a reference equator and a family of polar great circles, while four additional inclined orthodromes remain subject to constrained optimization. Because the initial arrangement is not necessarily in general position, the specification develops a formal treatment of multiway intersections, antipodal vertices, edge multiplicities, spherical cells, and degenerate configurations.

The induced orthodromic arrangement is supplemented by two dual proximity structures. A spherical Voronoi tessellation assigns every surface point to its nearest infrastructure node under geodesic distance. A spherical Delaunay complex connects nodes whose service regions meet and thereby supplies a locally adaptive secondary network distinct from the fixed primary orthodromes. The monograph derives the relationship between these structures, establishes their duality under nondegeneracy assumptions, and explains how weighted and constrained variants can represent terrain cost, ecological exclusion, population demand, hydrological capacity, tectonic risk, and launch suitability.

The resulting blueprint is developmental rather than terminal. Regions may construct low-capability implementations while retaining compatibility with more advanced sections elsewhere. Models, survey monuments, roads, hydraulic works, electrical systems, and megastructures can all instantiate the same protocol at different scales. Civilizations separated by centuries may therefore converge upon a common architecture without requiring synchronized construction, uniform technology, or a permanently centralized authority.
\end{abstract}

\tableofcontents

\section*{Notation}

Symbols in this monograph are scoped to the chapter or section in which
they are introduced rather than reserved globally. Three reuses are
frequent enough to note explicitly, since none of the three co-occurs
within a single derivation:

\(R\) denotes the planetary radius in Parts~I--II (\(\Sphere_R\),
geodesic distance, corridor width), electrical resistance in the
Energy Transport chapter, a reliability function \(R(t)\) and a
weighted continuity measure in the resilience chapters of Part~V, and
a resilience functional \(\mathcal R\) later in the same Part.
\(C\) denotes a generic cost functional, the graph coupling matrix of
Chapter~6, and channel capacity in the Communication Networks chapter.
\(Q\) denotes hydraulic volumetric flow throughout Parts~II and~VI,
and is kept distinct from communication queue length, written
\(\mathcal Q_i(t)\), in the Communication Networks chapter.

Readers moving between chapters, or assembling appendices that cite
several chapters at once, should treat each symbol as local to its
section unless a chapter explicitly imports an earlier definition.

\mainmatter

% ============================================================
% PART I
% ============================================================

\part{Mathematical and Architectural Foundations}

\chapter{Purpose, Scope, and Semantic Discipline}

\section{The Blueprint as a Protocol}

The present work is a specification of a spatial protocol. It does not require that every corridor be built, that every node receive the same machinery, or that every region progress through the same technological sequence. It instead asks what common geometric and interface commitments would permit independently constructed infrastructure to become mutually compatible over very long intervals.

Ordinary infrastructure planning begins with a known project, an available technology, a budget, a jurisdiction, and a finite construction horizon. A planetary blueprint inverts this ordering. It begins with a persistent geometric substrate and asks which project families can be embedded into that substrate without foreclosing future development. The infrastructure is therefore not defined by the first machinery installed within it. A surveyed route, a pipeline, a railway, and an evacuated launch tube can be different temporal realizations of the same reserved corridor.

This distinction requires a strict semantic separation among mathematical definitions, architectural principles, normative specifications, and research questions. A definition introduces an object and determines how that object is to be interpreted. A principle states a reason for preferring one design direction over another but does not by itself establish compliance. A specification establishes a requirement whose satisfaction can in principle be tested. A research question identifies a parameter, construction, or theorem that remains unresolved.

\begin{principle}{Protocol Before Implementation}
The blueprint shall define invariant spatial relations, interface obligations, and compatibility conditions without prematurely fixing the technologies by which those obligations are realized. A compliant implementation is identified by the preservation of required relations rather than by resemblance to a particular historical machine.
\end{principle}

The statement that geometry is defined once while implementation is optimized indefinitely must not be interpreted as a claim that every initial geometric decision is beyond revision. It means instead that, once a version of the protocol is adopted for construction, its interfaces must remain sufficiently stable that later work can connect to earlier work. Versioning and revision are therefore themselves required parts of the protocol. The geometry is stable relative to an adopted protocol epoch, while the governance system retains a formally constrained procedure for amendments.

\section{Normative Vocabulary}

This monograph adopts a standards-oriented distinction among the verbs \emph{shall}, \emph{should}, and \emph{may}. The word \emph{shall} identifies a mandatory condition of compliance. The word \emph{should} identifies a recommendation that may be departed from when a documented regional justification exists. The word \emph{may} identifies an optional implementation choice.

A statement containing \emph{shall} must eventually admit an operational test. If no test can be specified, the statement belongs either among the design principles or among the research questions. This convention prevents the language of aspiration from being confused with the language of engineering obligation.

\begin{specification}{Testability of Normative Claims}
Every requirement formulated with the word \emph{shall} shall be accompanied, either locally or in a later compliance chapter, by a measurable predicate, a tolerance relation, or a documented decision procedure. A requirement that cannot be translated into such a test shall not be treated as a binding engineering specification.
\end{specification}

\section{Planning Horizon}

The relevant time horizon is deliberately longer than the expected life of any particular machine, institution, or political settlement. This changes the meaning of optimization. A conventional project minimizes cost relative to a near-term function. A multi-century project must also minimize irreversible obstruction of unknown future functions.

Let \(I_t\) denote the physical implementation occupying a reserved corridor at time \(t\). Let \(\Gamma\) denote the geometric corridor itself, and let \(\Pi\) denote the interface protocol governing permissible connections. The long-term architecture is represented schematically by

\begin{equation}
    (\Gamma,\Pi)
    \longrightarrow
    I_{t_1}
    \longrightarrow
    I_{t_2}
    \longrightarrow
    I_{t_3}
    \longrightarrow \cdots,
\end{equation}

where the implementations \(I_{t_k}\) may differ radically while remaining compatible with the same corridor and protocol.

The relevant preservation law is therefore not

\begin{equation}
    I_{t_{k+1}} = I_{t_k},
\end{equation}

but rather

\begin{equation}
    I_{t_{k+1}} \sim_{\Pi} I_{t_k},
\end{equation}

where \(\sim_{\Pi}\) denotes protocol compatibility. The exact definition of \(\sim_{\Pi}\) will depend on the subsystem being considered. For hydraulic systems it may involve flange geometry, operating pressure intervals, contamination barriers, and flow-direction conventions. For electrical systems it may involve voltage ranges, fault isolation, conversion interfaces, and synchronization rules. For spatial reservations it may involve clearance envelopes and access rights.

\begin{definition}{Intergenerational Compatibility}
Two temporally separated implementations \(I_a\) and \(I_b\) are intergenerationally compatible under protocol \(\Pi\), written
\[
I_a \sim_{\Pi} I_b,
\]
when there exists a finite sequence of permitted adapters, conversions, or replacement operations that joins \(I_a\) to \(I_b\) without violating the protected spatial and safety invariants of the corridor.
\end{definition}

This definition intentionally permits technological discontinuity. A future system need not preserve obsolete machinery. It must preserve either direct interoperability or a feasible transition path. The protected object is therefore the capacity for continuation rather than the persistence of hardware.

\section{Scale of Realization}

The blueprint admits multiple scales of physical embodiment. A classroom model, a regional survey, a reduced-scale hydraulic demonstrator, a transportation corridor, and a planetary network may all instantiate related incidence relations. This does not mean that every property is scale invariant. Gravity, material strength, heat transfer, fluid turbulence, and atmospheric drag introduce scale-dependent effects. The scale invariance lies instead in the combinatorial organization of modules and interfaces.

Let \(X\) be a prototype and \(Y\) a full-scale realization. A geometric similarity map

\begin{equation}
    f_\lambda : X \to Y,
    \qquad
    f_\lambda(x)=\lambda x,
\end{equation}

preserves angles and incidence while multiplying lengths, areas, and volumes by different powers of \(\lambda\):

\begin{equation}
    L_Y = \lambda L_X,
    \qquad
    A_Y = \lambda^2 A_X,
    \qquad
    V_Y = \lambda^3 V_X.
\end{equation}

The differing exponents show why a physical prototype cannot be interpreted by simple visual enlargement. Structural mass grows as \(\lambda^3\), while cross-sectional support area commonly grows as \(\lambda^2\). The characteristic stress generated by self-weight therefore often grows proportionally to \(\lambda\). Consequently, the modular geometry may scale while the material realization must change.

\begin{derivation}{Square--Cube Scaling Constraint}
Suppose a structural member has characteristic length \(L\), density \(\rho\), and a geometrically similar cross-section. Its mass satisfies
\[
m \propto \rho L^3.
\]
Its weight is therefore
\[
W \propto \rho gL^3.
\]
If the supporting cross-sectional area satisfies
\[
A \propto L^2,
\]
then the compressive stress generated by self-weight obeys
\[
\sigma = \frac{W}{A}
\propto
\frac{\rho gL^3}{L^2}
=
\rho gL.
\]
Thus a geometrically similar enlargement by a factor \(\lambda\) multiplies the self-weight stress by approximately \(\lambda\). Structural similarity alone does not imply mechanical similarity.
\end{derivation}

The specification therefore distinguishes \emph{topological modularity}, which can be preserved exactly, from \emph{mechanical scaling}, which requires new materials, geometries, active supports, or distributed load paths.

\chapter{The Planetary Sphere and Orthodromic Geometry}

\section{Reference Sphere}

The first mathematical approximation treats the planet as a sphere of radius \(R>0\). The reference surface is

\begin{equation}
    \Sphere_R
    =
    \left\{
    x\in\R^3
    \suchthat
    \|x\|=R
    \right\}.
\end{equation}

When no ambiguity arises, the unit sphere

\begin{equation}
    \Sphere
    =
    \left\{
    u\in\R^3
    \suchthat
    \|u\|=1
    \right\}
\end{equation}

will be used for derivations. Physical distances are recovered by multiplication by \(R\).

A real planet is neither a perfect sphere nor a static surface. Earth is more accurately represented by a reference ellipsoid together with a time-dependent geoid and local terrain model. Nevertheless, the sphere is the correct initial object for defining great circles, antipodal symmetry, spherical Voronoi cells, and the abstract incidence graph. Later implementations may replace spherical geodesics by ellipsoidal geodesics or by terrain-constrained routes while retaining the spherical protocol as the global reference layer.

\begin{definition}{Reference and Realization Layers}
The \emph{reference layer} is the idealized spherical or ellipsoidal geometry in which orthodromes, nodes, cells, and interfaces are defined. The \emph{realization layer} is the terrain-dependent physical route by which a regional implementation approximates the corresponding reference object.
\end{definition}

This distinction permits a corridor to remain globally identifiable even where mountains, oceans, faults, settlements, or protected ecosystems require local displacement.

\section{Great Circles}

A great circle is the intersection of \(\Sphere_R\) with a plane passing through the origin. Let \(n\in\Sphere\) be a unit normal vector. The associated great circle is

\begin{equation}
    O(n)
    =
    \left\{
    x\in\Sphere_R
    \suchthat
    n\cdot x=0
    \right\}.
\end{equation}

Because \(n\) and \(-n\) determine the same plane,

\begin{equation}
    O(n)=O(-n).
\end{equation}

Thus the parameter space of unoriented great circles is not \(\Sphere\) itself but the real projective plane

\begin{equation}
    \mathbb{RP}^{2}
    =
    \Sphere/\{n\sim -n\}.
\end{equation}

This observation is important for optimization. Searching over all great circles means searching over plane orientations, and each orientation is represented twice on the unit sphere unless antipodal normals are identified.

\begin{definition}{Orthodrome}
An orthodrome is an unoriented great circle \(O(n)\subset\Sphere_R\), represented by an equivalence class \([n]\in\mathbb{RP}^{2}\).
\end{definition}

\section{Intersection of Two Orthodromes}

Let \(O(n_a)\) and \(O(n_b)\) be distinct great circles. If \(n_a\) and \(n_b\) are not parallel, their defining planes intersect in a line through the origin. The direction of this line is given by the cross product

\begin{equation}
    q_{ab}
    =
    n_a\times n_b.
\end{equation}

Normalizing gives

\begin{equation}
    \widehat q_{ab}
    =
    \frac{n_a\times n_b}
         {\|n_a\times n_b\|}.
\end{equation}

The two intersection points on \(\Sphere_R\) are therefore

\begin{equation}
    p_{ab}^{\pm}
    =
    \pm R
    \frac{n_a\times n_b}
         {\|n_a\times n_b\|}.
\end{equation}

\begin{proposition}
Two distinct great circles intersect in exactly two antipodal points.
\end{proposition}

\begin{proof}
Each great circle is the intersection of the sphere with a two-dimensional linear subspace of \(\R^3\). Two distinct two-dimensional linear subspaces through the origin intersect in a one-dimensional linear subspace. A one-dimensional linear subspace is a line
\[
L=\{tq\mid t\in\R\}
\]
for some nonzero vector \(q\). The condition \(\|tq\|=R\) has exactly two solutions,
\[
t=\pm\frac{R}{\|q\|}.
\]
Hence \(L\cap\Sphere_R\) consists of two antipodal points.
\end{proof}

If \(n_a\) and \(n_b\) are parallel, the two orthodromes coincide and do not constitute distinct elements of the arrangement. A valid orthodromic family therefore excludes duplicate projective normals.

\section{Geodesic Distance}

For \(x,y\in\Sphere_R\), define the central angle

\begin{equation}
    \theta(x,y)
    =
    \arccos\left(
      \frac{x\cdot y}{R^2}
    \right),
    \qquad
    0\leq\theta\leq\pi.
\end{equation}

The spherical geodesic distance is

\begin{equation}
    d_{\Sphere_R}(x,y)
    =
    R\theta(x,y)
    =
    R\arccos\left(
      \frac{x\cdot y}{R^2}
    \right).
\end{equation}

\begin{theorem}
For non-antipodal points \(x,y\in\Sphere_R\), the shorter arc of the unique great circle through \(x\) and \(y\) minimizes surface length among all piecewise smooth curves joining them.
\end{theorem}

\begin{proof}
Let \(\gamma:[0,1]\to\Sphere_R\) be a piecewise smooth curve from \(x\) to \(y\). Normalize by writing \(u=\gamma/R\), so \(u(t)\in\Sphere\). The length of \(\gamma\) is
\[
L(\gamma)
=
R\int_0^1\|u'(t)\|\,dt.
\]
The angular displacement between \(u(0)\) and \(u(1)\) is
\[
\theta
=
\arccos(u(0)\cdot u(1)).
\]
Any curve on the unit sphere connecting these endpoints must accumulate total tangent rotation at least \(\theta\). More formally, partition the curve into sufficiently short segments and apply the spherical triangle inequality to the induced angular metric. The sum of the segment angles is at least the direct angular distance:
\[
\sum_{k}
\arccos\bigl(u(t_k)\cdot u(t_{k+1})\bigr)
\geq
\arccos\bigl(u(0)\cdot u(1)\bigr).
\]
Passing to the limit gives
\[
\int_0^1\|u'(t)\|\,dt\geq\theta.
\]
Thus
\[
L(\gamma)\geq R\theta.
\]
Equality is attained by the constant-speed shorter great-circle arc.
\end{proof}

The theorem justifies orthodromes as minimum-distance routes on an ideal sphere. It does not prove that an actual pipeline, road, or launch corridor should follow the exact great circle at every point. Physical routing introduces additional costs.

Let \(c(x,\dot x)\) denote a local cost density incorporating slope, excavation, ecology, geology, legal access, and other constraints. A realized path \(\gamma\) may instead minimize

\begin{equation}
    \mathcal{L}_c[\gamma]
    =
    \int_0^1
    c\bigl(\gamma(t),\dot\gamma(t)\bigr)
    \,dt.
\end{equation}

The reference orthodrome supplies the global alignment target, while the realized route solves a constrained variational problem around it.

\section{Distance from a Point to an Orthodrome}

Let \(O(n)\) be a great circle with unit normal \(n\), and let \(x\in\Sphere_R\). The signed angular latitude of \(x\) relative to the plane of \(O(n)\) is

\begin{equation}
    \alpha_n(x)
    =
    \arcsin\left(
      \frac{n\cdot x}{R}
    \right).
\end{equation}

The unsigned geodesic distance from \(x\) to the great circle is

\begin{equation}
    d_{\Sphere_R}\bigl(x,O(n)\bigr)
    =
    R
    \left|
    \arcsin\left(
      \frac{n\cdot x}{R}
    \right)
    \right|.
\end{equation}

\begin{derivation}{Point-to-Orthodrome Distance}
Write
\[
x=x_{\parallel}+x_{\perp},
\]
where
\[
x_{\perp}=(n\cdot x)n
\]
is normal to the defining plane and
\[
x_{\parallel}=x-(n\cdot x)n
\]
lies in the plane. The nearest point on the great circle is obtained by radially normalizing \(x_{\parallel}\):
\[
y
=
R\frac{x_{\parallel}}{\|x_{\parallel}\|}.
\]
The angle between \(x\) and the plane is \(\alpha\), where
\[
\sin\alpha
=
\frac{\|x_{\perp}\|}{R}
=
\frac{|n\cdot x|}{R}.
\]
Therefore
\[
\alpha
=
\left|
\arcsin\left(\frac{n\cdot x}{R}\right)
\right|,
\]
and the corresponding surface distance is \(R\alpha\).
\end{derivation}

This formula makes it possible to define a reserved corridor of angular half-width \(\varepsilon\).

\begin{definition}{Orthodromic Corridor Tube}
For an orthodrome \(O(n)\) and width parameter \(0<\varepsilon<\pi R/2\), the associated spherical corridor tube is
\[
\mathcal{T}_{\varepsilon}(O(n))
=
\left\{
x\in\Sphere_R
\suchthat
d_{\Sphere_R}\bigl(x,O(n)\bigr)
\leq\varepsilon
\right\}.
\]
Equivalently,
\[
\mathcal{T}_{\varepsilon}(O(n))
=
\left\{
x\in\Sphere_R
\suchthat
|n\cdot x|
\leq
R\sin\left(\frac{\varepsilon}{R}\right)
\right\}.
\]
\end{definition}

The corridor is therefore a spherical band centered on the great circle. The reference line may be infinitesimally thin, but the protected right-of-way has finite width.

\begin{derivation}{Area of a Spherical Corridor Tube}
Let the angular half-width be
\[
\delta=\frac{\varepsilon}{R}.
\]
Choose coordinates in which the great circle is the equator. The corridor then occupies latitudes
\[
-\delta\leq\varphi\leq\delta.
\]
The area element on \(\Sphere_R\) is
\[
dA=R^2\cos\varphi\,d\varphi\,d\lambda.
\]
Integrating over longitude \(0\leq\lambda<2\pi\) gives
\begin{align}
A_{\mathrm{tube}}
&=
\int_0^{2\pi}
\int_{-\delta}^{\delta}
R^2\cos\varphi
\,d\varphi\,d\lambda
\\
&=
2\pi R^2
\left[
\sin\varphi
\right]_{-\delta}^{\delta}
\\
&=
4\pi R^2\sin\delta.
\end{align}
Thus
\[
A_{\mathrm{tube}}
=
4\pi R^2
\sin\left(\frac{\varepsilon}{R}\right).
\]
For \(\varepsilon\ll R\),
\[
\sin\left(\frac{\varepsilon}{R}\right)
\approx
\frac{\varepsilon}{R},
\]
so
\[
A_{\mathrm{tube}}
\approx
4\pi R\varepsilon.
\]
This agrees with the planar intuition that area is approximately circumference \(2\pi R\) multiplied by total width \(2\varepsilon\).
\end{derivation}

\begin{specification}{Reference Alignment and Realization Deviation}
Every primary corridor shall be associated with a unique reference orthodrome and a declared admissible tube \(\mathcal{T}_{\varepsilon}(O)\). A realized segment may depart from the centerline when required by terrain, geology, ecology, settlement avoidance, or structural safety, but it shall remain within its declared tube except at formally registered bypass sections. A bypass shall include entry and exit interfaces that preserve the identity and continuity of the reference corridor.
\end{specification}

\chapter{Arrangements of Great Circles}

\section{The Orthodromic Arrangement}

Let

\begin{equation}
    \Oset
    =
    \{O_1,O_2,\dots,O_n\}
\end{equation}

be a finite set of distinct great circles. Their union

\begin{equation}
    \mathscr{A}(\Oset)
    =
    \bigcup_{i=1}^{n}O_i
\end{equation}

partitions the sphere into vertices, edges, and two-dimensional cells.

A point \(p\in\Sphere_R\) is an arrangement vertex when it belongs to at least two orthodromes. Its multiplicity is

\begin{equation}
    m(p)
    =
    \#\{O_i\in\Oset\mid p\in O_i\}.
\end{equation}

An arrangement is in \emph{general position} when no three great circles pass through the same point. Since every intersection occurs antipodally, this means every vertex has multiplicity two.

The proposed network is not necessarily in general position. A family of polar meridian great circles all meets at the north and south poles. These are multiway intersections and must be represented as single physical nodes with multiplicity greater than two.

\begin{definition}{Arrangement Vertex Set}
The vertex set of the orthodromic arrangement is
\[
V(\Oset)
=
\left\{
p\in\Sphere_R
\suchthat
m(p)\geq 2
\right\}.
\]
Each geometric point appears once, regardless of how many orthodromes pass through it.
\end{definition}

\section{Edges and Cells}

Fix an orthodrome \(O_i\). Remove from it all arrangement vertices lying on \(O_i\). The remaining set is a disjoint union of open arcs. The closures of these arcs are the arrangement edges supported by \(O_i\).

Let

\begin{equation}
    r_i
    =
    \#\bigl(V(\Oset)\cap O_i\bigr)
\end{equation}

be the number of distinct vertices on \(O_i\). If \(r_i\geq 2\), then the circle is divided into exactly \(r_i\) closed arcs.

\begin{lemma}
A circle containing \(r\geq 2\) marked points is divided into exactly \(r\) arcs.
\end{lemma}

\begin{proof}
Order the marked points cyclically as
\[
p_1,p_2,\dots,p_r.
\]
There is one arc from \(p_k\) to \(p_{k+1}\) for each \(k<r\), together with one final arc from \(p_r\) to \(p_1\). Hence the number of arcs equals \(r\).
\end{proof}

Therefore the total number of arrangement edges, counted as geometric arcs supported by individual orthodromes, is

\begin{equation}
    E
    =
    \sum_{i=1}^{n}r_i.
\end{equation}

The same count can be expressed using vertex multiplicities.

\begin{proposition}
For an arrangement of distinct great circles with at least two intersection points on each circle,
\[
E
=
\sum_{p\in V(\Oset)}m(p).
\]
\end{proposition}

\begin{proof}
The quantity
\[
\sum_{i=1}^{n}r_i
\]
counts incidences \((O_i,p)\) such that \(p\in O_i\). Counting the same incidence relation by vertices gives
\[
\sum_{p\in V(\Oset)}m(p).
\]
The two sums are therefore equal.
\end{proof}

Let \(F\) denote the number of connected components of

\begin{equation}
    \Sphere_R\setminus \mathscr{A}(\Oset).
\end{equation}

These components are the open arrangement cells.

Assuming the induced graph is connected, Euler's formula on the sphere gives

\begin{equation}
    V-E+F=2.
\end{equation}

Hence

\begin{equation}
    F=2-V+E.
\end{equation}

Substituting the multiplicity formula gives

\begin{equation}
    F
    =
    2
    -
    \#V(\Oset)
    +
    \sum_{p\in V(\Oset)}m(p).
\end{equation}

Equivalently,

\begin{equation}
    F
    =
    2
    +
    \sum_{p\in V(\Oset)}
    \bigl(m(p)-1\bigr).
\end{equation}

\begin{theorem}{Cell Count with Multiway Intersections}
For a connected arrangement of distinct great circles,
\[
F
=
2
+
\sum_{p\in V(\Oset)}
\bigl(m(p)-1\bigr).
\]
\end{theorem}

\begin{proof}
Euler's relation gives
\[
F=2-V+E.
\]
Since
\[
V=\sum_{p\in V(\Oset)}1
\]
and
\[
E=\sum_{p\in V(\Oset)}m(p),
\]
we obtain
\begin{align}
F
&=
2
-
\sum_{p\in V(\Oset)}1
+
\sum_{p\in V(\Oset)}m(p)
\\
&=
2
+
\sum_{p\in V(\Oset)}
\bigl(m(p)-1\bigr).
\end{align}
\end{proof}

This formula is useful because it remains valid when multiple orthodromes concur.

\section{General-Position Counts}

Suppose \(n\) great circles are in general position. Every pair intersects in two antipodal points, and no intersection point belongs to a third circle.

The number of unordered pairs is

\begin{equation}
    \binom{n}{2}.
\end{equation}

Each pair contributes two vertices, so

\begin{equation}
    V
    =
    2\binom{n}{2}
    =
    n(n-1).
\end{equation}

Each great circle intersects each of the other \(n-1\) circles twice, producing

\begin{equation}
    r_i=2(n-1)
\end{equation}

vertices on each circle. Therefore

\begin{equation}
    E
    =
    n\cdot 2(n-1)
    =
    2n(n-1).
\end{equation}

Euler's formula gives

\begin{align}
    F
    &=
    2-V+E
    \\
    &=
    2-n(n-1)+2n(n-1)
    \\
    &=
    n^2-n+2.
\end{align}

\begin{corollary}
An arrangement of \(n\) great circles in general position has
\[
V=n(n-1),
\qquad
E=2n(n-1),
\qquad
F=n^2-n+2.
\]
\end{corollary}

For \(n=9\), the maximum general-position counts are

\begin{equation}
    V=72,
    \qquad
    E=144,
    \qquad
    F=74.
\end{equation}

These numbers are upper bounds for arrangements in which multiway concurrence is allowed. When several circles pass through the same antipodal pair, vertices and cells merge.

\section{A Family with Polar Concurrence}

Suppose \(q\) distinct meridian great circles all pass through a common polar axis. They meet at two antipodal points \(P_N\) and \(P_S\), each of multiplicity \(q\). Add one equatorial great circle \(O_E\), which intersects each meridian at a distinct antipodal pair of equatorial points.

The polar vertices contribute two geometric vertices. The equator contributes \(2q\) distinct meridian intersections. Thus

\begin{equation}
    V=2+2q.
\end{equation}

Each meridian contains the two poles and two equatorial intersections, so it contains four distinct vertices and contributes four edges. The equator contains \(2q\) vertices and contributes \(2q\) edges. Therefore

\begin{equation}
    E=4q+2q=6q.
\end{equation}

Euler's formula gives

\begin{equation}
    F
    =
    2-(2+2q)+6q
    =
    4q.
\end{equation}

For \(q=4\),

\begin{equation}
    V=10,
    \qquad
    E=24,
    \qquad
    F=16.
\end{equation}

This is the arrangement produced by one equator and four polar great circles before the inclined orthodromes are introduced.

\begin{remark}
A ``meridian'' in geographic language is often described as a semicircle from pole to pole. A great circle contains a meridian and its antipodal continuation. Consequently, longitudes separated by \(180^\circ\) belong to the same meridian great circle. A specification involving four meridian orthodromes must therefore identify four distinct planes through the polar axis rather than four isolated longitude rays.
\end{remark}

\section{Adding a Generic Inclined Orthodrome}

Suppose a new great circle \(O_{n+1}\) is added to an existing arrangement and does not pass through any existing vertex. It intersects each of the \(n\) existing circles in two points. If all \(2n\) new points are distinct, the new circle is divided into \(2n\) arcs.

Each new arc passes through the interior of one existing cell and divides that cell into two. The number of cells therefore increases by \(2n\).

\begin{proposition}
If a new great circle is added in general position relative to an existing arrangement of \(n\) great circles, then
\[
F_{n+1}=F_n+2n.
\]
\end{proposition}

\begin{proof}
The new circle contains \(2n\) distinct intersection points and is therefore divided into \(2n\) arcs. Each arc lies in the interior of an old cell except at its endpoints. Introducing the arc splits that cell into two. Hence every arc increases the number of cells by one, yielding an increase of \(2n\).
\end{proof}

Starting with \(F_1=2\), repeated application gives

\begin{align}
F_n
&=
2+\sum_{k=1}^{n-1}2k
\\
&=
2+n(n-1)
\\
&=
n^2-n+2.
\end{align}

This provides an inductive derivation of the general-position cell formula.

\section{The Nine-Orthodrome Realization}

The adopted architecture begins with five reference orthodromes and four flexible inclined orthodromes. Let

\begin{equation}
    \Oset_{\mathrm{fixed}}
    =
    \{O_1,O_2,O_3,O_4,O_5\},
\end{equation}

and let

\begin{equation}
    \Oset_{\mathrm{flex}}
    =
    \{O_6,O_7,O_8,O_9\}.
\end{equation}

The complete family is

\begin{equation}
    \Oset
    =
    \Oset_{\mathrm{fixed}}
    \cup
    \Oset_{\mathrm{flex}}.
\end{equation}

The reference equator may be represented by the normal

\begin{equation}
    n_1
    =
    \begin{pmatrix}
      0\\0\\1
    \end{pmatrix}.
\end{equation}

A polar great circle with azimuth parameter \(\lambda\) has a horizontal normal

\begin{equation}
    n(\lambda)
    =
    \begin{pmatrix}
      \cos\lambda\\
      \sin\lambda\\
      0
    \end{pmatrix}.
\end{equation}

Four symmetrically spaced polar great circles may be represented by

\begin{equation}
    \lambda_k
    =
    \lambda_0+\frac{k\pi}{4},
    \qquad
    k=0,1,2,3.
\end{equation}

The spacing is \(\pi/4\), not \(\pi/2\), because normals differing by \(\pi\) represent the same great circle. Thus four distinct projective orientations are distributed over a half-turn.

The flexible inclined orthodromes can be parameterized by unit normals

\begin{equation}
    n_i
    =
    \begin{pmatrix}
      \cos\beta_i\cos\lambda_i\\
      \cos\beta_i\sin\lambda_i\\
      \sin\beta_i
    \end{pmatrix},
    \qquad
    i=6,7,8,9,
\end{equation}

where \(\beta_i\) is the latitude of the pole of the great circle and \(\lambda_i\) is its azimuth. Because \(n_i\sim -n_i\), the parameter pairs satisfy a projective identification.

\begin{definition}{Nine-Orthodrome Configuration Space}
The unconstrained configuration space of the four flexible orthodromes is
\[
\mathfrak{C}
=
\left(\mathbb{RP}^{2}\right)^4
\setminus \Delta,
\]
where \(\Delta\) denotes configurations containing duplicate orthodromes or other forbidden degeneracies.
\end{definition}

The dimension of \(\mathbb{RP}^{2}\) is two. Four freely chosen inclined orthodromes therefore contribute eight continuous degrees of freedom before symmetry, exclusion, and optimization constraints are imposed.

If the four inclined orthodromes are required to arise from a single seed orthodrome under a cyclic symmetry group, their degrees of freedom can be reduced. Let \(R\in SO(3)\) be a rotation satisfying

\begin{equation}
    R^4=I.
\end{equation}

Choose one normal \(n_6\), and define

\begin{equation}
    n_{6+k}=R^k n_6,
    \qquad
    k=0,1,2,3.
\end{equation}

The flexible family is then generated by one seed and one prescribed symmetry operation. This reduces arbitrariness but may also prevent adaptation to population, geography, or geology.

\begin{research}{Symmetry Versus Planetary Fit}
Should the four inclined orthodromes be constrained to an exact finite symmetry group, or should symmetry appear only as a regularization term within a broader optimization? Exact symmetry simplifies modular replication and yields uniform cells, but a geographically optimized configuration may require controlled asymmetry.
\end{research}

\chapter{Spherical Service Geometry}

\section{Why the Orthodromic Arrangement Is Not Sufficient}

The orthodromic arrangement specifies primary corridors and their intersections. It does not by itself assign every point on the planetary surface to a service node, determine which node is closest to a settlement, or identify locally efficient secondary connections among nodes.

These functions require a proximity geometry. Given a finite node set

\begin{equation}
    \Nset
    =
    \{p_1,p_2,\dots,p_m\}
    \subset\Sphere_R,
\end{equation}

the spherical Voronoi tessellation partitions the planet into nearest-node service regions. Its dual, the spherical Delaunay complex, connects nodes whose service regions meet.

The orthodromic graph and the Delaunay graph must not be conflated. The orthodromic graph is fixed by the chosen great circles. The Delaunay graph is induced by the node locations and may include edges that do not lie on primary orthodromes. It is therefore suitable for secondary distribution, regional redundancy, local roads, feeder pipes, and alternative routing.

\section{Spherical Voronoi Cells}

Let the geodesic distance on \(\Sphere_R\) be denoted by \(d_S\). For each node \(p_i\), define its Voronoi cell by

\begin{equation}
    \Vor(p_i)
    =
    \left\{
    x\in\Sphere_R
    \suchthat
    d_S(x,p_i)
    \leq
    d_S(x,p_j)
    \text{ for all }j
    \right\}.
\end{equation}

A point belongs to the cell of \(p_i\) when \(p_i\) is one of its nearest nodes.

Because

\begin{equation}
    d_S(x,p_i)
    =
    R\arccos\left(
      \frac{x\cdot p_i}{R^2}
    \right),
\end{equation}

and because \(\arccos\) is strictly decreasing on \([-1,1]\), the inequality

\begin{equation}
    d_S(x,p_i)
    \leq
    d_S(x,p_j)
\end{equation}

is equivalent to

\begin{equation}
    x\cdot p_i
    \geq
    x\cdot p_j.
\end{equation}

Thus

\begin{equation}
    \Vor(p_i)
    =
    \left\{
    x\in\Sphere_R
    \suchthat
    x\cdot(p_i-p_j)\geq 0
    \text{ for all }j
    \right\}.
\end{equation}

Each pairwise inequality defines a hemisphere bounded by a great circle.

\begin{theorem}
Every unweighted spherical Voronoi cell is an intersection of closed hemispheres and is therefore spherically convex.
\end{theorem}

\begin{proof}
For fixed \(i\) and \(j\), the condition
\[
d_S(x,p_i)\leq d_S(x,p_j)
\]
is equivalent to
\[
x\cdot(p_i-p_j)\geq 0.
\]
The set
\[
H_{ij}
=
\{x\in\Sphere_R\mid x\cdot(p_i-p_j)\geq 0\}
\]
is a closed hemisphere. Therefore
\[
\Vor(p_i)
=
\bigcap_{j\neq i}H_{ij}.
\]
An intersection of spherically convex sets is spherically convex.
\end{proof}

\section{Spherical Bisectors}

The bisector between nodes \(p_i\) and \(p_j\) is the set

\begin{equation}
    B_{ij}
    =
    \left\{
    x\in\Sphere_R
    \suchthat
    d_S(x,p_i)=d_S(x,p_j)
    \right\}.
\end{equation}

Using the dot-product form,

\begin{equation}
    B_{ij}
    =
    \left\{
    x\in\Sphere_R
    \suchthat
    x\cdot(p_i-p_j)=0
    \right\}.
\end{equation}

This is a great circle whose defining plane is perpendicular to \(p_i-p_j\).

\begin{proposition}
The geodesic bisector of two distinct non-antipodal sites on a sphere is a great circle.
\end{proposition}

\begin{proof}
Equality of distances is equivalent to
\[
x\cdot p_i=x\cdot p_j,
\]
and therefore to
\[
x\cdot(p_i-p_j)=0.
\]
This is the equation of a plane through the origin. Its intersection with the sphere is a great circle.
\end{proof}

If \(p_j=-p_i\), every point on the great circle orthogonal to \(p_i\) is equidistant, but the two associated Voronoi hemispheres remain well defined. Antipodal site pairs require care because many shortest paths between them exist.

\section{Voronoi Vertices}

A generic Voronoi vertex is equidistant from three sites \(p_i,p_j,p_k\). It satisfies

\begin{equation}
    x\cdot(p_i-p_j)=0
\end{equation}

and

\begin{equation}
    x\cdot(p_i-p_k)=0.
\end{equation}

Hence \(x\) is parallel to

\begin{equation}
    (p_i-p_j)\times(p_i-p_k).
\end{equation}

The two antipodal candidates are

\begin{equation}
    x_{\pm}
    =
    \pm R
    \frac{
      (p_i-p_j)\times(p_i-p_k)
    }{
      \|
      (p_i-p_j)\times(p_i-p_k)
      \|
    }.
\end{equation}

Only a candidate whose distance to \(p_i,p_j,p_k\) is no greater than its distance to every other site is an actual Voronoi vertex.

\begin{definition}{Empty Spherical Circumdisk}
Let \(p_i,p_j,p_k\in\Sphere_R\). A spherical disk \(D(c,\rho)\) with center \(c\in\Sphere_R\) and angular radius \(\rho\) is a circumdisk of the three sites when
\[
d_S(c,p_i)
=
d_S(c,p_j)
=
d_S(c,p_k)
=
\rho.
\]
It is empty relative to \(\Nset\) when
\[
d_S(c,p_\ell)\geq \rho
\]
for every other site \(p_\ell\in\Nset\).
\end{definition}

A Voronoi vertex is the center of an empty spherical circumdisk through at least three sites.

\section{The Spherical Delaunay Complex}

The Delaunay graph is the adjacency dual of the Voronoi tessellation.

\begin{definition}{Spherical Delaunay Adjacency}
Two sites \(p_i\) and \(p_j\) are Delaunay adjacent when their Voronoi cells share a one-dimensional boundary arc:
\[
\dim\bigl(
\Vor(p_i)\cap\Vor(p_j)
\bigr)=1.
\]
The resulting graph is denoted
\[
\Del(\Nset).
\]
\end{definition}

Under generic conditions, three sites form a Delaunay triangle when their Voronoi cells meet at a common Voronoi vertex.

\begin{definition}{Spherical Delaunay Triangle}
A triple \((p_i,p_j,p_k)\) defines a spherical Delaunay triangle when there exists an empty spherical circumdisk whose boundary contains \(p_i,p_j,p_k\).
\end{definition}

\begin{theorem}{Voronoi--Delaunay Duality}
Assume that the site set contains no degeneracy in which four or more sites lie on the boundary of the same empty spherical disk. Then every Voronoi edge corresponds to one Delaunay edge, every Voronoi vertex corresponds to one Delaunay triangular face, and every Voronoi cell corresponds to one Delaunay vertex.
\end{theorem}

\begin{proof}
A Voronoi cell is indexed by one site, so the correspondence between Voronoi cells and Delaunay vertices is immediate.

Suppose \(\Vor(p_i)\) and \(\Vor(p_j)\) share an edge. Every point \(x\) in the interior of that edge satisfies
\[
d_S(x,p_i)=d_S(x,p_j)
\]
and
\[
d_S(x,p_i)<d_S(x,p_k)
\]
for every other site \(p_k\). A sufficiently small spherical disk centered at \(x\) and expanded until it reaches \(p_i\) and \(p_j\) has those sites on its boundary and contains no other site in its interior. Thus \(p_i\) and \(p_j\) satisfy the empty-disk condition and are joined by a Delaunay edge.

Conversely, suppose \(p_i\) and \(p_j\) are joined by a Delaunay edge. Then there exists an empty spherical disk having both sites on its boundary. Its center is equidistant from \(p_i\) and \(p_j\) and no closer to any other site. Hence the center lies in
\[
\Vor(p_i)\cap\Vor(p_j).
\]
Under nondegeneracy this intersection contains a one-dimensional arc rather than only an isolated higher-order coincidence, so the cells share a Voronoi edge.

Now let \(x\) be a Voronoi vertex. Under the stated nondegeneracy assumption, exactly three cells meet at \(x\), corresponding to sites \(p_i,p_j,p_k\). The spherical disk centered at \(x\) with radius
\[
\rho=d_S(x,p_i)
\]
has all three sites on its boundary and no site in its interior. Therefore the triple forms a Delaunay triangle.

Conversely, if \(p_i,p_j,p_k\) admit an empty circumdisk with center \(x\), then \(x\) is equidistant from the three sites and no farther from them than from any other site. Hence
\[
x\in
\Vor(p_i)\cap
\Vor(p_j)\cap
\Vor(p_k).
\]
Under nondegeneracy this intersection is a single Voronoi vertex.
\end{proof}

\section{Euler Relations for the Spherical Delaunay Triangulation}

Suppose the Delaunay complex is a triangulation of the entire sphere with \(m\) vertices, \(e\) edges, and \(f\) triangular faces. Euler's relation gives

\begin{equation}
    m-e+f=2.
\end{equation}

Each triangular face has three edges, while each edge belongs to two faces. Therefore

\begin{equation}
    3f=2e.
\end{equation}

Substituting

\begin{equation}
    f=\frac{2e}{3}
\end{equation}

into Euler's relation gives

\begin{align}
m-e+\frac{2e}{3}
&=2,
\\
m-\frac{e}{3}
&=2,
\\
e
&=3m-6.
\end{align}

Then

\begin{equation}
    f
    =
    \frac{2}{3}(3m-6)
    =
    2m-4.
\end{equation}

\begin{corollary}
A nondegenerate spherical Delaunay triangulation with \(m\geq 3\) sites has
\[
e=3m-6
\]
edges and
\[
f=2m-4
\]
triangular faces.
\end{corollary}

The average Delaunay degree is

\begin{equation}
    \overline d
    =
    \frac{2e}{m}
    =
    6-\frac{12}{m}.
\end{equation}

Thus the average degree approaches six for large node sets. This is the spherical analogue of the local hexagonal tendency of planar Voronoi systems, modified by the global curvature of the sphere.

\begin{derivation}{Curvature Defect in a Spherical Triangulation}
For each Delaunay vertex \(p_i\), let \(d_i\) denote its degree. Since
\[
\sum_{i=1}^{m}d_i=2e,
\]
and
\[
e=3m-6,
\]
we obtain
\[
\sum_{i=1}^{m}d_i
=
6m-12.
\]
Therefore
\[
\sum_{i=1}^{m}(6-d_i)
=
6m-(6m-12)
=
12.
\]
The total combinatorial degree defect is always twelve. A perfectly hexagonal triangulation with \(d_i=6\) everywhere is impossible on the sphere. Positive curvature requires a total deficit equivalent, for example, to twelve degree-five vertices.
\end{derivation}

This result is architecturally significant. A globally uniform network cannot reproduce an infinite planar hexagonal lattice without defects. The sphere forces exceptional nodes. These exceptional nodes may be interpreted as high-centrality hubs, geometric singularities, or places where local modular rules must be modified.

\section{Convex-Hull Construction of the Delaunay Complex}

The spherical Delaunay triangulation can be constructed using the convex hull of the site vectors in \(\R^3\). Let

\begin{equation}
    P
    =
    \conv\{p_1,p_2,\dots,p_m\}.
\end{equation}

Assume the origin lies strictly inside \(P\), and assume no four sites lie on a common supporting plane of \(P\). Then each facet of \(P\) is triangular. Radially projecting each facet to \(\Sphere_R\) produces a spherical triangle.

\begin{theorem}
Under the stated assumptions, the radial projections of the facets of \(P\) form the spherical Delaunay triangulation of the sites.
\end{theorem}

\begin{proof}
Let a triangular facet have vertices \(p_i,p_j,p_k\). Because it is a supporting facet, there exists a nonzero vector \(c\) and scalar \(h>0\) such that
\[
c\cdot p_i
=
c\cdot p_j
=
c\cdot p_k
=
h
\]
and
\[
c\cdot p_\ell
\leq h
\]
for every other site \(p_\ell\).

Normalize \(c\) to a point on the sphere:
\[
\widehat c
=
R\frac{c}{\|c\|}.
\]
Since
\[
\widehat c\cdot p_i
=
\widehat c\cdot p_j
=
\widehat c\cdot p_k,
\]
the geodesic distances from \(\widehat c\) to the three facet vertices are equal. For every other site,
\[
\widehat c\cdot p_\ell
\leq
\widehat c\cdot p_i.
\]
Because \(\arccos\) is decreasing,
\[
d_S(\widehat c,p_\ell)
\geq
d_S(\widehat c,p_i).
\]
Thus the spherical circumdisk centered at \(\widehat c\) through \(p_i,p_j,p_k\) is empty. The triple is therefore Delaunay.

Conversely, suppose \(p_i,p_j,p_k\) admit an empty spherical circumdisk centered at \(\widehat c\). Then
\[
\widehat c\cdot p_i
=
\widehat c\cdot p_j
=
\widehat c\cdot p_k
\]
and
\[
\widehat c\cdot p_\ell
\leq
\widehat c\cdot p_i
\]
for all other sites. Hence the plane
\[
\widehat c\cdot x
=
\widehat c\cdot p_i
\]
supports the convex hull at the three sites. Under nondegeneracy, they form a triangular facet.
\end{proof}

This construction gives a practical computational method. Once the orthodromic intersections are calculated as unit vectors, a three-dimensional convex-hull algorithm can recover the spherical Delaunay complex.

\section{The Orthodromic Graph and the Delaunay Graph}

Let

\begin{equation}
    G_O=(V,E_O)
\end{equation}

denote the graph induced by the orthodromic arrangement, and let

\begin{equation}
    G_D=(V,E_D)
\end{equation}

denote the Delaunay graph of the same node set.

The graphs generally differ. An edge in \(E_O\) connects consecutive intersection nodes along a primary orthodrome. An edge in \(E_D\) connects nodes with adjacent nearest-service regions. Some edges may belong to both sets:

\begin{equation}
    E_O\cap E_D\neq\varnothing.
\end{equation}

The complete architectural graph may be defined as

\begin{equation}
    G_A
    =
    \bigl(
      V,
      E_O\cup E_D\cup E_R
    \bigr),
\end{equation}

where \(E_R\) denotes additional regionally authorized links.

Primary orthodromic edges possess protocol priority because they realize the fixed global geometry. Delaunay edges provide adaptive local connectivity. Regional edges accommodate geography, settlements, existing infrastructure, or political boundaries.

\begin{definition}{Three-Layer Connectivity}
The architectural network consists of a primary geodesic layer \(E_O\), a secondary proximity layer \(E_D\), and a tertiary realization layer \(E_R\). The primary layer preserves global identity, the secondary layer minimizes local service separation, and the tertiary layer accommodates contingent regional requirements.
\end{definition}

\begin{principle}{Fixed Skeleton, Adaptive Tissue}
The orthodromic arrangement should act as a stable planetary skeleton, while the Delaunay and regional layers act as adaptive connective tissue. Global convergence does not require every local route to be a great-circle arc. It requires local routes to preserve reliable access to the invariant primary skeleton.
\end{principle}

\section{A Spanner Interpretation}

A network is useful when graph travel distances do not greatly exceed direct geodesic distances. For nodes \(p_i,p_j\), let \(d_G(p_i,p_j)\) be the length of the shortest network path, with each edge weighted by its spherical length. Define the stretch ratio

\begin{equation}
    \sigma_G(p_i,p_j)
    =
    \frac{
      d_G(p_i,p_j)
    }{
      d_S(p_i,p_j)
    }.
\end{equation}

The global stretch is

\begin{equation}
    \sigma(G)
    =
    \sup_{i\neq j}
    \sigma_G(p_i,p_j).
\end{equation}

A small stretch indicates that the network approximates direct travel efficiently.

The Delaunay graph often performs well as a geometric spanner in Euclidean settings. On the sphere, analogous behavior depends on site distribution and maximum cell diameter. Rather than assuming a universal constant, this specification treats spherical stretch as a measurable design criterion.

\begin{specification}{Network Stretch Evaluation}
Every proposed orthodromic configuration shall be evaluated for weighted shortest-path stretch on the combined graph \(G_O\cup G_D\). The evaluation shall report both worst-case stretch and demand-weighted mean stretch. A configuration shall not be selected solely on visual symmetry if it produces excessive detours between major service regions.
\end{specification}

\section{Voronoi Coverage Radius}

The covering radius of a node set is

\begin{equation}
    \rho_{\mathrm{cov}}(\Nset)
    =
    \max_{x\in\Sphere_R}
    \min_{p_i\in\Nset}
    d_S(x,p_i).
\end{equation}

It is the greatest distance from any surface point to its nearest node. Equivalently, it is the maximum circumradius among Voronoi vertices.

A small covering radius gives uniform access. A configuration with visually regular corridors may nevertheless contain large underserved cells. The Voronoi tessellation makes this visible.

The separation radius is

\begin{equation}
    \rho_{\mathrm{sep}}(\Nset)
    =
    \frac{1}{2}
    \min_{i\neq j}
    d_S(p_i,p_j).
\end{equation}

The mesh ratio is

\begin{equation}
    \eta(\Nset)
    =
    \frac{
      \rho_{\mathrm{cov}}(\Nset)
    }{
      \rho_{\mathrm{sep}}(\Nset)
    }.
\end{equation}

A low mesh ratio indicates that the sites are neither excessively clustered nor separated by large holes.

\begin{principle}{Uniformity Without Forced Uniformity}
The network should avoid both severe node clustering and extensive service voids. Exact equality of Voronoi areas is not required, because population, ecology, geology, and hydrology are nonuniform. Nevertheless, deviations from uniformity should arise from declared planning objectives rather than accidental geometric imbalance.
\end{principle}

\section{Weighted Voronoi Tessellations}

Geodesic nearness is not the only relevant notion of accessibility. A node may possess greater hydraulic capacity, more reliable energy, lower elevation cost, or greater political accessibility. Such differences can be represented by weights.

An additively weighted spherical Voronoi cell may be defined by

\begin{equation}
    \Vor_w(p_i)
    =
    \left\{
    x\in\Sphere_R
    \suchthat
    d_S(x,p_i)-w_i
    \leq
    d_S(x,p_j)-w_j
    \text{ for all }j
    \right\}.
\end{equation}

A positive \(w_i\) expands the effective service region of \(p_i\), while a negative weight contracts it.

A multiplicatively weighted cell may be defined by

\begin{equation}
    \Vor_\mu(p_i)
    =
    \left\{
    x\in\Sphere_R
    \suchthat
    \mu_i d_S(x,p_i)
    \leq
    \mu_j d_S(x,p_j)
    \text{ for all }j
    \right\}.
\end{equation}

Here \(\mu_i>1\) can represent increased effective travel cost or reduced reliability.

Weighted bisectors are generally not great circles. Consequently, weighted cells need not be bounded by orthodromic arcs. This is acceptable because they represent service domains rather than primary corridor geometry.

\begin{research}{Capacity-Weighted Service Domains}
Which weighting model best represents a node whose service capability is limited simultaneously by water throughput, electrical capacity, transportation volume, emergency redundancy, and maintenance latency? An additive scalar weight is simple but may collapse physically distinct constraints into a single number. A vector-valued or multi-commodity service geometry may be required.
\end{research}

\section{Anisotropic and Terrain-Constrained Distance}

A mountain range, ocean trench, conflict boundary, ecological reserve, or unstable substrate can make travel direction-dependent. Replace the spherical metric by a Finsler-type cost

\begin{equation}
    \mathcal{D}(x,y)
    =
    \inf_{\gamma(0)=x,\;\gamma(1)=y}
    \int_0^1
    F\bigl(\gamma(t),\dot\gamma(t)\bigr)
    \,dt,
\end{equation}

where \(F(x,v)\) gives the cost of motion through \(x\) in direction \(v\).

The corresponding generalized Voronoi cell is

\begin{equation}
    \Vor_F(p_i)
    =
    \left\{
    x
    \suchthat
    \mathcal{D}(x,p_i)
    \leq
    \mathcal{D}(x,p_j)
    \text{ for all }j
    \right\}.
\end{equation}

If \(F(x,v)\neq F(x,-v)\), then travel cost is directional. This can model prevailing winds, ocean currents, elevation-dependent pumping, gravity-assisted freight, or one-way launch trajectories.

The dual graph of such a generalized Voronoi diagram need not be an ordinary Delaunay triangulation. It is more properly understood as a cost-adjacency graph. The classical spherical Delaunay complex remains the geometric reference, while anisotropic variants guide regional realization.

\chapter{Optimization of the Inclined Orthodromes}

\section{From Arbitrary Placement to Explicit Objectives}

The four inclined orthodromes constitute the principal continuous freedom in the adopted nine-circle realization. Their selection cannot be justified by symmetry alone if the network is expected to interact with population, geology, ecology, hydrology, and launch dynamics.

Let

\begin{equation}
    \theta
    =
    (\beta_6,\lambda_6,\dots,\beta_9,\lambda_9)
\end{equation}

be an eight-parameter representation of the flexible normals. The resulting orthodrome family is denoted by \(\Oset(\theta)\), the arrangement node set by \(\Nset(\theta)\), and the combined graph by \(G_A(\theta)\).

The optimization problem is not intrinsically scalar. Let the objective vector be

\begin{equation}
    \mathbf{J}(\theta)
    =
    \begin{pmatrix}
      J_{\mathrm{coverage}}(\theta)\\
      J_{\mathrm{tectonic}}(\theta)\\
      J_{\mathrm{population}}(\theta)\\
      J_{\mathrm{launch}}(\theta)\\
      J_{\mathrm{ecology}}(\theta)\\
      J_{\mathrm{hydrology}}(\theta)\\
      J_{\mathrm{symmetry}}(\theta)\\
      J_{\mathrm{stretch}}(\theta)
    \end{pmatrix}.
\end{equation}

Some components are to be maximized and others minimized. A scalarized objective can be written

\begin{equation}
    J_w(\theta)
    =
    \sum_{k=1}^{q}s_k w_kJ_k(\theta),
\end{equation}

where \(w_k\geq 0\) and \(s_k\in\{+1,-1\}\) records whether the corresponding criterion is beneficial or costly.

The optimal configuration is then

\begin{equation}
    \theta^\star
    \in
    \argmax_{\theta\in\Theta_{\mathrm{adm}}}
    J_w(\theta),
\end{equation}

where \(\Theta_{\mathrm{adm}}\) is the admissible parameter space after excluding duplicates, prohibited intersections, excessive corridor overlap, and other constraints.

\begin{principle}{Pareto Transparency}
No weighted optimum should be presented as uniquely dictated by geometry when its location depends on political or ethical weights. The protocol should preserve the Pareto frontier and disclose the weights or priority rules by which a particular configuration is selected.
\end{principle}

\section{Pareto Optimality}

Let \(\mathbf{C}(\theta)\in\R^q\) be a vector of costs, all of which are to be minimized. A configuration \(\theta_1\) dominates \(\theta_2\) when

\begin{equation}
    C_k(\theta_1)
    \leq
    C_k(\theta_2)
\end{equation}

for every \(k\), with strict inequality for at least one component.

A configuration is Pareto optimal when no admissible configuration dominates it. Here \(\mathbf{C}(\theta')\prec\mathbf{C}(\theta)\) denotes exactly the dominance relation just defined: componentwise inequality in every cost with strict inequality in at least one.

\begin{definition}{Pareto Set}
The Pareto set is
\[
\mathcal{P}
=
\left\{
\theta\in\Theta_{\mathrm{adm}}
\suchthat
\nexists\,
\theta'
\in\Theta_{\mathrm{adm}}
\text{ with }
\mathbf{C}(\theta')
\prec
\mathbf{C}(\theta)
\right\}.
\]
\end{definition}

The image

\begin{equation}
    \mathbf{C}(\mathcal{P})
\end{equation}

is the Pareto frontier. It represents the irreducible trade-offs among criteria. For example, reducing tectonic crossings may increase distance from population centers. Improving equatorial launch orientation may worsen freshwater coverage at high latitudes. Increasing symmetry may force routes through ecologically sensitive regions.

This is a first instance of Pareto dominance, specialized to the eight-parameter orthodrome configuration space \(\Theta_{\mathrm{adm}}\). Chapter~\ref{chap:multi-objective-optimization} restates the same dominance order over the general admissible function space \(\mathcal A\), so that every subsequent infrastructure functional (routing, capacity, resilience) inherits this construction rather than a fresh one.

\section{Coverage Functional}

Let the nearest-node distance field be

\begin{equation}
    r_\theta(x)
    =
    \min_{p\in\Nset(\theta)}
    d_S(x,p).
\end{equation}

The worst-case coverage cost is

\begin{equation}
    C_{\infty}(\theta)
    =
    \max_{x\in\Sphere_R}
    r_\theta(x).
\end{equation}

An average coverage cost relative to surface measure \(dA\) is

\begin{equation}
    C_1(\theta)
    =
    \frac{1}{4\pi R^2}
    \int_{\Sphere_R}
    r_\theta(x)\,dA.
\end{equation}

A quadratic form is

\begin{equation}
    C_2(\theta)
    =
    \frac{1}{4\pi R^2}
    \int_{\Sphere_R}
    r_\theta(x)^2\,dA.
\end{equation}

The worst-case functional protects remote locations, while the average functionals reward broad overall access. They need not choose the same configuration.

If \(w_P(x)\) is a population density or projected future settlement density, define the demand-weighted coverage functional

\begin{equation}
    C_P(\theta)
    =
    \frac{
      \int_{\Sphere_R}
      w_P(x)r_\theta(x)^2\,dA
    }{
      \int_{\Sphere_R}
      w_P(x)\,dA
    }.
\end{equation}

A future-oriented blueprint should not simply substitute present population density for \(w_P\). Present population reflects current coastlines, political boundaries, and historical infrastructure. If the network is intended to influence settlement over centuries, demand and infrastructure co-evolve.

\section{Corridor Accessibility Functional}

Distance to nodes and distance to corridors answer different questions. Let

\begin{equation}
    a_\theta(x)
    =
    \min_{O\in\Oset(\theta)}
    d_S(x,O).
\end{equation}

The mean corridor-access cost is

\begin{equation}
    A_2(\theta)
    =
    \frac{
      \int_{\Sphere_R}
      w_A(x)a_\theta(x)^2\,dA
    }{
      \int_{\Sphere_R}
      w_A(x)\,dA
    }.
\end{equation}

This functional measures how far settlements, agricultural zones, or potential inland cities lie from the nearest primary orthodrome.

Because \(a_\theta\) is a minimum over smooth distance functions, it is continuous but generally not differentiable where two orthodromes are equally near. Numerical optimization may therefore require nonsmooth methods, stochastic search, or smooth approximations.

One smooth approximation to the minimum is the soft minimum

\begin{equation}
    \operatorname{smin}_{\tau}
    (z_1,\dots,z_n)
    =
    -\tau
    \log\left(
      \sum_{i=1}^{n}
      e^{-z_i/\tau}
    \right),
\end{equation}

where \(\tau>0\). As \(\tau\to 0^+\),

\begin{equation}
    \operatorname{smin}_{\tau}
    (z_1,\dots,z_n)
    \to
    \min_i z_i.
\end{equation}

This permits gradient-based exploratory optimization while preserving the exact minimum for final evaluation.

\section{Symmetry Regularization}

Let \(H\subset SO(3)\) be a target symmetry group. A configuration is exactly invariant under \(H\) when

\begin{equation}
    h\Oset=\Oset
\end{equation}

for every \(h\in H\).

When exact invariance is too restrictive, define a symmetry defect. Let \(N\) be the matrix whose columns are chosen representative normals. For each \(h\in H\), compare the transformed normal set \(hN\) with the original set after the best permutation and sign choices.

A possible defect is

\begin{equation}
    C_{\mathrm{sym}}(\theta)
    =
    \frac{1}{|H|}
    \sum_{h\in H}
    \min_{\pi,\epsilon}
    \sum_{i}
    \left\|
      hn_i
      -
      \epsilon_i n_{\pi(i)}
    \right\|^2,
\end{equation}

where \(\pi\) ranges over permutations and each \(\epsilon_i\in\{-1,+1\}\) accounts for projective equivalence.

Exact symmetry gives

\begin{equation}
    C_{\mathrm{sym}}=0.
\end{equation}

A small nonzero value represents controlled asymmetry.

\section{Degeneracy Penalties}

Certain configurations produce nearly coincident orthodromes or unstable clusters of intersections. Let the projective angular separation between normals be

\begin{equation}
    \delta_{ij}
    =
    \arccos\bigl(|n_i\cdot n_j|\bigr).
\end{equation}

The absolute value identifies \(n\) and \(-n\). Duplicate orthodromes satisfy \(\delta_{ij}=0\).

A repulsion penalty may be defined as

\begin{equation}
    C_{\mathrm{dup}}(\theta)
    =
    \sum_{i<j}
    \psi(\delta_{ij}),
\end{equation}

where \(\psi(\delta)\to\infty\) as \(\delta\to 0\). For example,

\begin{equation}
    \psi(\delta)
    =
    \frac{1}{(\delta+\varepsilon_0)^p}
\end{equation}

with \(p>0\) and small regularizer \(\varepsilon_0\).

Multiway intersections may be desirable at selected hubs but undesirable elsewhere. Let

\begin{equation}
    \chi_{\mathrm{hub}}(p)
\end{equation}

identify permitted high-multiplicity nodes. A concurrence penalty can be defined by measuring the proximity of pairwise intersection points generated by different circle pairs and penalizing clusters outside designated hub regions.

\section{Topological Stability}

As the orthodrome parameters vary, the arrangement graph remains combinatorially unchanged until a degeneracy occurs, such as three circles becoming concurrent. The configuration space is therefore partitioned into regions within which the graph topology is constant.

Let \(\mathfrak{D}\subset\mathfrak{C}\) denote the discriminant set of degenerate configurations. Then

\begin{equation}
    \mathfrak{C}\setminus\mathfrak{D}
\end{equation}

is divided into topological chambers. Moving continuously within one chamber changes edge lengths, cell areas, and node positions without changing the incidence graph.

\begin{principle}{Topology Before Fine Position}
Optimization should distinguish changes that merely deform an existing arrangement from changes that alter the incidence type. Fine geometric optimization may proceed within a stable topological chamber, while transitions across the discriminant set should be treated as explicit architectural revisions.
\end{principle}

\begin{research}{Classification of Nine-Circle Chambers}
What are the combinatorially distinct arrangement types obtainable from the adopted five fixed and four flexible orthodromes? Which chambers preserve the intended correspondence with the supercube, triteck, Hoberman-sphere, or exploded-cube incidence structures?
\end{research}

\section{A Constrained Variational Statement}

The flexible configuration can be written as the solution of a constrained optimization problem:

\begin{equation}
    \theta^\star
    \in
    \argmin_{\theta\in\Theta}
    \mathcal{F}(\theta),
\end{equation}

where

\begin{align}
\mathcal{F}(\theta)
&=
\alpha_{\mathrm{cov}}C_{\mathrm{cov}}(\theta)
+
\alpha_{\mathrm{access}}C_{\mathrm{access}}(\theta)
+
\alpha_{\mathrm{tect}}C_{\mathrm{tect}}(\theta)
\\
&\quad
+
\alpha_{\mathrm{eco}}C_{\mathrm{eco}}(\theta)
+
\alpha_{\mathrm{launch}}C_{\mathrm{launch}}(\theta)
+
\alpha_{\mathrm{stretch}}C_{\mathrm{stretch}}(\theta)
\\
&\quad
+
\alpha_{\mathrm{sym}}C_{\mathrm{sym}}(\theta)
+
\alpha_{\mathrm{deg}}C_{\mathrm{deg}}(\theta).
\end{align}

The admissible set is

\begin{equation}
    \Theta_{\mathrm{adm}}
    =
    \left\{
    \theta\in\Theta
    \suchthat
    g_a(\theta)\leq 0,
    \;
    h_b(\theta)=0
    \right\},
\end{equation}

where \(g_a\) are inequality constraints and \(h_b\) are equality constraints.

Examples of equality constraints include exact rotational relations among the four inclined orthodromes. Examples of inequality constraints include minimum projective angular separation, maximum overlap with protected ecological zones, minimum Delaunay angle quality, and upper bounds on network stretch.

The Lagrangian is

\begin{equation}
    \mathscr{L}
    (\theta,\lambda,\mu)
    =
    \mathcal{F}(\theta)
    +
    \sum_a\lambda_a g_a(\theta)
    +
    \sum_b\mu_b h_b(\theta),
\end{equation}

with multipliers satisfying

\begin{equation}
    \lambda_a\geq 0.
\end{equation}

At a regular local optimum, the Karush--Kuhn--Tucker conditions require

\begin{equation}
    \nabla_\theta
    \mathscr{L}
    (\theta^\star,\lambda^\star,\mu^\star)
    =
    0,
\end{equation}

together with primal feasibility,

\begin{equation}
    g_a(\theta^\star)\leq 0,
    \qquad
    h_b(\theta^\star)=0,
\end{equation}

dual feasibility,

\begin{equation}
    \lambda_a^\star\geq 0,
\end{equation}

and complementary slackness,

\begin{equation}
    \lambda_a^\star
    g_a(\theta^\star)
    =
    0.
\end{equation}

These conditions do not guarantee a global optimum in a nonconvex configuration space. They nevertheless clarify which constraints actively determine a candidate solution.

\begin{specification}{Reproducibility of Orthodrome Selection}
Any adopted configuration of the four inclined orthodromes shall be accompanied by its parameterization, objective definitions, constraint set, data epoch, numerical method, initialization strategy, and sensitivity analysis. A map alone shall not constitute a sufficient justification of placement.
\end{specification}

\section{Sensitivity and Long-Horizon Robustness}

An optimum computed from present-day data may be fragile. Let \(z\) denote uncertain exogenous data, including future population distributions, sea level, ecological boundaries, political exclusions, energy technology, and launch requirements. The objective becomes

\begin{equation}
    \mathcal{F}(\theta;z).
\end{equation}

A nominal optimum solves

\begin{equation}
    \theta^\star(z_0)
    \in
    \argmin_\theta
    \mathcal{F}(\theta;z_0)
\end{equation}

for one assumed scenario \(z_0\). A robust optimum instead minimizes worst-case cost over an uncertainty set \(\mathcal{Z}\):

\begin{equation}
    \theta^\star_{\mathrm{rob}}
    \in
    \argmin_{\theta}
    \sup_{z\in\mathcal{Z}}
    \mathcal{F}(\theta;z).
\end{equation}

A stochastic formulation minimizes expected cost:

\begin{equation}
    \theta^\star_{\mathrm{stoch}}
    \in
    \argmin_{\theta}
    \mathbb{E}_{z}
    \left[
      \mathcal{F}(\theta;z)
    \right].
\end{equation}

A risk-sensitive formulation may combine expectation with a tail-risk measure:

\begin{equation}
    \theta^\star
    \in
    \argmin_{\theta}
    \left\{
      \mathbb{E}
      [\mathcal{F}(\theta;z)]
      +
      \kappa
      \operatorname{CVaR}_{\alpha}
      \bigl(
        \mathcal{F}(\theta;z)
      \bigr)
    \right\}.
\end{equation}

For a blueprint intended to persist across centuries, robustness may be more important than nominal optimality. A slightly inferior present-day configuration may be preferable if its performance remains acceptable across a wide range of futures.

\begin{principle}{Long-Horizon Minimax Regret}
The placement of primary geometry should avoid decisions whose apparent present advantage depends on narrow assumptions about future population, technology, climate, or political organization. Where uncertainty cannot be reduced, the preferred configuration should minimize irreversible regret rather than maximize short-term efficiency.
\end{principle}

\chapter{Primary Cells, Service Cells, and Regional Modules}

\section{Two Distinct Tessellations}

The great-circle arrangement produces one tessellation, while the node-based Voronoi construction produces another.

The orthodromic cells are the connected components of

\begin{equation}
    \Sphere_R
    \setminus
    \bigcup_{i=1}^{9}O_i.
\end{equation}

Their boundaries lie on primary corridors. They express the global geometric skeleton.

The Voronoi cells are the nearest-node regions

\begin{equation}
    \Vor(p_i).
\end{equation}

Their boundaries lie on node bisectors. They express service responsibility and local accessibility.

These tessellations should not be expected to coincide. An orthodromic cell may intersect several Voronoi cells, and a Voronoi cell may cross several orthodromic cells.

Their common refinement is formed by all nonempty intersections

\begin{equation}
    Q_{\alpha i}
    =
    C_\alpha
    \cap
    \Vor(p_i),
\end{equation}

where \(C_\alpha\) is an orthodromic cell.

\begin{definition}{Operational Microcell}
An operational microcell is a connected component of
\[
C_\alpha\cap\Vor(p_i)
\]
for some orthodromic cell \(C_\alpha\) and some service node \(p_i\). It records both the primary geometric region and the nearest-node assignment.
\end{definition}

The common refinement permits regional planning to answer two questions simultaneously: which primary corridor boundaries structure the region, and which node is responsible for its nearest service interface?

\section{Cell Area}

For a spherical polygon with interior angles

\begin{equation}
    \alpha_1,\dots,\alpha_k,
\end{equation}

the spherical excess is

\begin{equation}
    \mathcal{E}
    =
    \sum_{j=1}^{k}\alpha_j
    -
    (k-2)\pi.
\end{equation}

On a sphere of radius \(R\), its area is

\begin{equation}
    A
    =
    R^2\mathcal{E}.
\end{equation}

\begin{theorem}{Girard's Formula for a Spherical Triangle}
For a spherical triangle with angles \(A,B,C\),
\[
\area(\triangle)
=
R^2(A+B+C-\pi).
\]
\end{theorem}

\begin{proof}
Consider a spherical lune of angle \(A\). A lune of angle \(2\pi\) covers the sphere twice, so its area is proportional to \(A\):
\[
\area(\text{lune}_A)
=
2R^2A.
\]
Construct the three lunes associated with the triangle's angles \(A,B,C\). Together they cover the triangle and its antipodal triangle three times in total, while covering the remainder of the sphere once. Therefore
\[
2R^2(A+B+C)
=
4\pi R^2
+
2\area(\triangle).
\]
Solving gives
\[
\area(\triangle)
=
R^2(A+B+C-\pi).
\]
\end{proof}

A spherical polygon can be triangulated from an interior point, yielding the general spherical-excess formula.

Cell area provides one measure of uniformity. However, equal area alone does not guarantee good shape. A long narrow cell may have the same area as a compact cell but produce much worse access.

\section{Cell Diameter and Aspect}

For a cell \(C\subset\Sphere_R\), define its geodesic diameter by

\begin{equation}
    \diam_S(C)
    =
    \sup_{x,y\in C}
    d_S(x,y).
\end{equation}

Define its inradius relative to a chosen center \(c\in C\) as the largest \(\rho\) for which

\begin{equation}
    B_S(c,\rho)\subseteq C.
\end{equation}

A compactness ratio may be written

\begin{equation}
    \kappa(C)
    =
    \frac{
      \diam_S(C)
    }{
      2\rho_{\mathrm{in}}(C)
    }.
\end{equation}

Values near one indicate compact cells, while large values indicate elongated or irregular cells.

For a Voronoi cell, the generating site is a natural center, though it need not maximize the inradius. For an orthodromic cell, a spherical Chebyshev center may be defined by

\begin{equation}
    c^\star
    \in
    \argmax_{c\in C}
    \min_{x\in\partial C}
    d_S(c,x).
\end{equation}

\begin{specification}{Cell Quality Reporting}
The geometric evaluation of a candidate orthodrome arrangement shall report cell area, geodesic diameter, inradius, compactness, and adjacency degree. No claim of uniform global coverage shall be based on area alone.
\end{specification}

\section{Centroidal Voronoi Relaxation}

Given a density function \(\rho(x)\geq 0\), the mass of a Voronoi cell is

\begin{equation}
    M_i
    =
    \int_{\Vor(p_i)}
    \rho(x)\,dA.
\end{equation}

An extrinsic centroid vector is

\begin{equation}
    \bar c_i
    =
    \frac{
      \int_{\Vor(p_i)}
      \rho(x)x\,dA
    }{
      \int_{\Vor(p_i)}
      \rho(x)\,dA
    }.
\end{equation}

Because \(\bar c_i\) generally lies inside the sphere rather than on it, project it radially:

\begin{equation}
    c_i
    =
    R\frac{\bar c_i}{\|\bar c_i\|}.
\end{equation}

A centroidal spherical Voronoi tessellation satisfies

\begin{equation}
    p_i=c_i
\end{equation}

for every site.

The Lloyd iteration repeatedly replaces each site with the projected centroid of its cell. This tends to equalize weighted service regions.

The orthodromic nodes cannot generally move independently because their positions are constrained as intersections of great circles. Nevertheless, centroidal error can be used as an optimization criterion:

\begin{equation}
    C_{\mathrm{CVT}}(\theta)
    =
    \sum_i
    d_S\bigl(p_i(\theta),c_i(\theta)\bigr)^2.
\end{equation}

A small value means the constrained orthodromic nodes approximate the centroids of their service regions.

\begin{research}{Constrained Centroidal Orthodromic Design}
Can the four inclined orthodromes be chosen so that their induced intersection nodes approximate a density-weighted centroidal Voronoi tessellation while preserving the intended incidence structure? This is a constrained inverse problem because node positions cannot be varied independently.
\end{research}

\section{Delaunay Angle Quality}

A Delaunay triangulation is often preferred because it avoids unnecessarily small angles in the planar case. On the sphere, triangle quality can likewise be assessed by its minimum angle.

For a spherical triangle \(T\) with angles \(A_T,B_T,C_T\), define

\begin{equation}
    q_{\angle}(T)
    =
    \min\{A_T,B_T,C_T\}.
\end{equation}

The global minimum-angle quality is

\begin{equation}
    Q_{\angle}
    =
    \min_{T\in\Del(\Nset)}
    q_{\angle}(T).
\end{equation}

Small values indicate needle-like triangles and poor local conditioning.

A shape functional may also use edge lengths \(a,b,c\) and area \(A_T\). One normalized quantity is

\begin{equation}
    q_T
    =
    \frac{
      4\sqrt{3}\,A_T
    }{
      a^2+b^2+c^2
    },
\end{equation}

interpreted locally when triangles are small relative to \(R\). For larger spherical triangles, a curvature-corrected quality measure should be used.

\section{Delaunay Paths as Secondary Corridors}

Suppose a region lies far from a primary orthodrome but near two or more nodes. Delaunay edges can provide secondary feeder routes. Because two Delaunay-adjacent nodes possess neighboring service territories, a connecting route crosses no intervening node's domain in the ordinary nearest-site geometry.

This gives a hierarchy:

\begin{equation}
    \text{local settlement}
    \longrightarrow
    \text{Voronoi service node}
    \longrightarrow
    \text{Delaunay feeder edge}
    \longrightarrow
    \text{orthodromic backbone}.
\end{equation}

The hierarchy prevents the primary network from being expected to solve every local routing problem.

\begin{protocol}{Primary and Secondary Route Identity}
A route lying on an orthodromic edge shall be registered as a primary corridor segment. A route connecting Delaunay-adjacent nodes but not lying within a primary orthodromic tube shall be registered as a secondary feeder. Secondary feeder status shall not grant authority to redefine the reference geometry of the primary network.
\end{protocol}

\section{Redundancy and Failure}

Let the operational graph be

\begin{equation}
    G=(V,E).
\end{equation}

A graph is \(k\)-edge-connected when at least \(k\) edges must be removed to disconnect it. It is \(k\)-vertex-connected when at least \(k\) vertices must be removed to disconnect it.

For infrastructure, edge connectivity measures resilience to corridor failures, while vertex connectivity measures resilience to node failures.

The local edge connectivity between nodes \(u\) and \(v\) is the maximum number of edge-disjoint \(u\)-to-\(v\) paths. By Menger's theorem, this equals the minimum number of edges whose removal separates \(u\) from \(v\).

\begin{theorem}[Edge Form of Menger's Theorem]
For distinct vertices \(u\) and \(v\) in a finite graph, the maximum number of pairwise edge-disjoint paths from \(u\) to \(v\) equals the minimum number of edges whose removal disconnects \(u\) from \(v\).
\end{theorem}

The theorem motivates the use of Delaunay and regional links as redundancy layers. A pure orthodromic arrangement may have high-degree nodes yet still contain vulnerable arcs. Secondary edges can increase local connectivity without altering the primary protocol.

\begin{specification}{Failure-Mode Connectivity}
Every critical node pair shall be evaluated for edge-disjoint and vertex-disjoint path counts under the combined primary and secondary graph. Hydraulic, electrical, transportation, communications, and launch functions shall be evaluated separately because an edge that is redundant for one subsystem may not carry another.
\end{specification}

\section{Multi-Commodity Graphs}

The physical corridor graph is not identical to any one subsystem graph. Let

\begin{equation}
    G^{(w)}
    =
    (V,E^{(w)})
\end{equation}

be the water graph,

\begin{equation}
    G^{(e)}
    =
    (V,E^{(e)})
\end{equation}

the electrical graph,

\begin{equation}
    G^{(t)}
    =
    (V,E^{(t)})
\end{equation}

the transportation graph, and

\begin{equation}
    G^{(\ell)}
    =
    (V,E^{(\ell)})
\end{equation}

the launch graph.

All are subgraphs of the reserved spatial graph

\begin{equation}
    G^{(\mathrm{space})}.
\end{equation}

Thus

\begin{equation}
    E^{(w)},
    E^{(e)},
    E^{(t)},
    E^{(\ell)}
    \subseteq
    E^{(\mathrm{space})}.
\end{equation}

Co-location means that multiple subsystem edges occupy a common spatial corridor. It does not mean they share a physical conduit or failure boundary.

\begin{definition}{Spatial Co-location Without Functional Identification}
Two subsystem edges are co-located when their protected volumes lie within the same registered corridor tube. They remain functionally distinct unless the protocol explicitly declares a coupling interface.
\end{definition}

This distinction is essential for preventing cascading failure. A water rupture must not automatically destroy electrical transmission. An electrical fault must not disable communications redundancy. A launch-track incident must not contaminate freshwater conduits.

\begin{principle}{Shared Alignment, Separated Failure Domains}
The economic and ecological benefits of co-location should be obtained without collapsing independent systems into a single failure domain. Spatial proximity shall be balanced against physical separation, shielding, access independence, and emergency isolation.
\end{principle}

%%%%%%%%%%%%%%%%%%%%%%%%%%%%%%%%%%%%%%%%%%%%%%%%%%%%%%%%%%%%%%%%%%%%%%
\part{The Geometry of Infrastructure Corridors}

\chapter{The Corridor as a Geometric Object}

\section{From Curves to Corridors}

The orthodromic network introduced in the previous chapters defines the
one-dimensional geometric skeleton of the planetary infrastructure.
However, no physical engineering system possesses zero thickness.
Roads occupy land, pipelines possess diameter, electrical systems require
clearance, and maintenance operations demand accessible working volume.

The primary mathematical object is therefore not the orthodrome itself,
but a finite neighborhood surrounding it.  The orthodrome serves only as
a geometric reference from which the physical corridor is generated.

This distinction mirrors the relationship between an abstract graph and
its embedding in Euclidean space.  The graph determines connectivity,
while the embedding determines geometry.

\begin{definition}{Reference Curve}

Let

\[
\gamma:[0,L]\rightarrow\Sphere_R
\]

be a continuously differentiable curve parameterized by arc length.

The curve satisfies

\[
\left\|
\frac{d\gamma}{ds}
\right\|
=1.
\]

When $\gamma$ is an orthodrome, it is a geodesic of the sphere.

\end{definition}

Throughout this monograph the parameter $s$ denotes arc length measured
along the reference curve.

%%%%%%%%%%%%%%%%%%%%%%%%%%%%%%%%%%%%%%%%%%%%%%%%%%%%%%%%%%%%%%%

\section{Moving Frames Along the Corridor}

Every point of the reference curve possesses a natural orthogonal frame.

Define the tangent vector

\[
T(s)
=
\frac{d\gamma}{ds}.
\]

Since

\[
\gamma(s)\cdot\gamma(s)=R^2,
\]

differentiating gives

\[
2\gamma\cdot T=0.
\]

Hence

\[
\gamma\cdot T=0.
\]

Therefore the tangent vector is everywhere tangent to the planetary
surface.

Next define the outward planetary normal

\[
N(s)
=
\frac{\gamma(s)}{R}.
\]

Immediately,

\[
N\cdot T=0.
\]

Finally define the binormal direction

\[
B=T\times N.
\]

Consequently,

\[
(T,N,B)
\]

forms an orthonormal basis at every point of the corridor.

%%%%%%%%%%%%%%%%%%%%%%%%%%%%%%%%%%%%%%%%%%%%%%%%%%%%%%%%%%%%%%%

\begin{theorem}

The moving frame

\[
(T,N,B)
\]

is orthonormal.

\end{theorem}

\begin{proof}

Since

\[
\|T\|=1
\]

by construction and

\[
\|N\|=1
\]

because

\[
\|N\|
=
\frac{\|\gamma\|}{R}
=
1,
\]

it remains only to establish orthogonality.

We already have

\[
N\cdot T=0.
\]

Now

\[
B=T\times N,
\]

and every cross product is orthogonal to each of its factors.

Hence

\[
B\cdot T=0,
\]

and

\[
B\cdot N=0.
\]

Finally,

\[
\|B\|
=
\|T\|
\|N\|
=
1.
\]

Therefore

\[
(T,N,B)
\]

is an orthonormal frame.

\end{proof}

%%%%%%%%%%%%%%%%%%%%%%%%%%%%%%%%%%%%%%%%%%%%%%%%%%%%%%%%%%%%%%%

\section{Curvature of Great Circles}

Unlike arbitrary curves, orthodromes possess zero intrinsic curvature on
the sphere.

Their geodesic curvature satisfies

\[
\kappa_g=0.
\]

The total curvature arises entirely from the curvature of the planetary
surface.

%%%%%%%%%%%%%%%%%%%%%%%%%%%%%%%%%%%%%%%%%%%%%%%%%%%%%%%%%%%%%%%

\begin{derivation}{Curvature of a Great Circle}

Parameterize the equator by

\[
\gamma(\theta)
=
R
\begin{pmatrix}
\cos\theta\\
\sin\theta\\
0
\end{pmatrix}.
\]

Since

\[
s=R\theta,
\]

we obtain

\[
\theta=\frac{s}{R}.
\]

Hence

\[
\gamma(s)
=
R
\begin{pmatrix}
\cos(s/R)\\
\sin(s/R)\\
0
\end{pmatrix}.
\]

Differentiating,

\[
T
=
\frac{d\gamma}{ds}
=
\begin{pmatrix}
-\sin(s/R)\\
\cos(s/R)\\
0
\end{pmatrix}.
\]

Differentiating once more,

\[
\frac{dT}{ds}
=
-\frac1R
\begin{pmatrix}
\cos(s/R)\\
\sin(s/R)\\
0
\end{pmatrix}.
\]

Therefore

\[
\left\|
\frac{dT}{ds}
\right\|
=
\frac1R.
\]

Thus

\[
\boxed{\kappa=\frac1R.}
\]

\end{derivation}

This result is fundamental because every sufficiently long pipeline,
railway, transmission line, or launch structure inherits the same
planetary curvature.

%%%%%%%%%%%%%%%%%%%%%%%%%%%%%%%%%%%%%%%%%%%%%%%%%%%%%%%%%%%%%%%

\section{The Corridor Tube}

Engineering systems occupy finite spatial extent.

Rather than identifying the corridor with the reference curve, define it
as a tubular neighborhood surrounding that curve.

\begin{definition}{Infrastructure Corridor}

Let

\[
\gamma:[0,L]\rightarrow\Sphere_R
\]

be a reference orthodrome.

For corridor half-width $w$ define

\[
\mathcal C_w
=
\left\{
x\in\Sphere_R
:
d(x,\gamma)\le w
\right\},
\]

where $d$ denotes geodesic distance.

\end{definition}

The corridor therefore possesses finite area rather than zero measure.

%%%%%%%%%%%%%%%%%%%%%%%%%%%%%%%%%%%%%%%%%%%%%%%%%%%%%%%%%%%%%%%

\section{Subsystem Occupancy}

Every engineering discipline occupies only part of the available
corridor.

Let

\[
X_i\subseteq\mathcal C_w
\]

denote the spatial region occupied by subsystem $i$.

Examples include hydraulic transport, electrical transmission,
communications, transportation, maintenance, ecological buffering,
inspection access, and future reserved capacity.

%%%%%%%%%%%%%%%%%%%%%%%%%%%%%%%%%%%%%%%%%%%%%%%%%%%%%%%%%%%%%%%

\begin{definition}{Occupancy}

The total occupied region is

\[
X
=
\bigcup_iX_i.
\]

The occupancy fraction is

\[
u
=
\frac{\operatorname{Area}(X)}
{\operatorname{Area}(\mathcal C_w)}.
\]

\end{definition}

Clearly

\[
0
\le
u
\le
1.
\]

%%%%%%%%%%%%%%%%%%%%%%%%%%%%%%%%%%%%%%%%%%%%%%%%%%%%%%%%%%%%%%%

\begin{theorem}

If

\[
u<1,
\]

then additional infrastructure may be incorporated without increasing
the protected corridor width.

\end{theorem}

\begin{proof}

Since

\[
u<1,
\]

we have

\[
\operatorname{Area}(X)
<
\operatorname{Area}(\mathcal C_w).
\]

Hence

\[
\mathcal C_w\setminus X
\neq\emptyset.
\]

Therefore unused corridor capacity remains available for future
construction.

\end{proof}

%%%%%%%%%%%%%%%%%%%%%%%%%%%%%%%%%%%%%%%%%%%%%%%%%%%%%%%%%%%%%%%

\section{Reserved Capacity}

Infrastructure planning often optimizes for immediate utilization.

The present blueprint instead recognizes unused corridor capacity as a
long-term strategic resource.

A reserved region

\[
\mathcal U
=
\mathcal C_w\setminus X
\]

is not wasted space.

Rather, it functions as future optionality, permitting later generations
to introduce technologies that cannot presently be anticipated without
requiring acquisition of entirely new rights-of-way.

Consequently, corridor utilization should not be maximized
indiscriminately.

%%%%%%%%%%%%%%%%%%%%%%%%%%%%%%%%%%%%%%%%%%%%%%%%%%%%%%%%%%%%%%%

\begin{principle}{Protocol Before Occupancy}

The geometric identity of the corridor precedes every engineering
realization placed within it.

Roads, pipelines, transmission lines, communication systems, and launch
facilities are transient occupants of a permanent geometric protocol.

The protocol persists while individual technologies evolve.

\end{principle}

%%%%%%%%%%%%%%%%%%%%%%%%%%%%%%%%%%%%%%%%%%%%%%%%%%%%%%%%%%%%%%%

\section{Cross-Sections}

For every arc-length coordinate $s$, define the local cross-section

\[
\Sigma(s)
\]

as the intersection of the corridor with the plane normal to the tangent
vector

\[
T(s).
\]

Each subsystem occupies a measurable subset

\[
\Sigma_i(s)
\subseteq
\Sigma(s).
\]

Hence

\[
\Sigma(s)
=
\bigcup_i\Sigma_i(s).
\]

The cross-sectional decomposition allows engineering disciplines to be
treated independently while preserving global geometric compatibility.

Later chapters will derive optimal cross-sectional layouts for hydraulic,
electrical, transportation, communication, maintenance, and launch
systems using this common geometric framework.

%%%%%%%%%%%%%%%%%%%%%%%%%%%%%%%%%%%%%%%%%%%%%%%%%%%%%%%%%%%%%%%%%%%%%%
\chapter{Cross-Sectional Geometry}

\section{The Corridor as a Finite Resource}

The previous chapter established that every infrastructure corridor is
represented by a finite tubular neighborhood surrounding a reference
orthodrome. We now examine the internal geometry of that neighborhood.

Unlike the reference curve itself, the corridor possesses finite
cross-sectional area. Every physical system occupying the corridor must
therefore compete for limited geometric space. Water conduits,
electrical transmission systems, transportation infrastructure,
communications, maintenance access, drainage systems, ecological
buffers, and reserved future capacity all require mutually compatible
allocation within a common spatial envelope.

The central mathematical problem is therefore one of geometric
partitioning.

Rather than designing independent rights-of-way for each engineering
discipline, the blueprint treats the corridor as a conserved geometric
resource whose allocation evolves over time while preserving global
compatibility.

%%%%%%%%%%%%%%%%%%%%%%%%%%%%%%%%%%%%%%%%%%%%%%%%%%%%%%%%%%%%%%%

\section{Cross-Sectional Coordinates}

Fix a point

\[
\gamma(s)
\]

along the reference orthodrome.

The orthonormal frame

\[
(T,N,B)
\]

constructed previously determines a natural coordinate system within the
cross section.

The tangent direction

\[
T
\]

is longitudinal.

The vectors

\[
N
\quad\text{and}\quad
B
\]

span the transverse plane.

Consequently every point in a sufficiently small neighborhood of the
corridor may be written

\[
x(s,u,v)
=
\gamma(s)
+
uN(s)
+
vB(s),
\]

where

\[
(u,v)
\]

are local transverse coordinates.

When the corridor width is small compared with planetary radius this
representation is accurate to first order.

%%%%%%%%%%%%%%%%%%%%%%%%%%%%%%%%%%%%%%%%%%%%%%%%%%%%%%%%%%%%%%%

\begin{definition}{Cross Section}

The cross section at position $s$ is

\[
\Sigma(s)
=
\left\{
(u,v):
(u,v)
\in
D(s)
\right\},
\]

where

\[
D(s)
\subset\mathbb R^2
\]

is a bounded planar region.

\end{definition}

Throughout the remainder of this monograph the exact shape of
\(D(s)\) remains implementation-dependent.

Only its topology and interface structure are fixed.

%%%%%%%%%%%%%%%%%%%%%%%%%%%%%%%%%%%%%%%%%%%%%%%%%%%%%%%%%%%%%%%

\section{Cross-Sectional Measure}

The usable geometric capacity of a corridor is quantified by its
cross-sectional area.

Define

\[
A(s)
=
\operatorname{Area}
(
\Sigma(s)
).
\]

If the corridor possesses constant width,

\[
A(s)=A_0.
\]

Otherwise,

\[
A=A(s)
\]

may vary continuously along the route in response to local geological
constraints.

%%%%%%%%%%%%%%%%%%%%%%%%%%%%%%%%%%%%%%%%%%%%%%%%%%%%%%%%%%%%%%%

The total geometric volume occupied by the corridor equals

\[
V
=
\int_0^L
A(s)\,ds.
\]

Consequently,

\[
V
\]

is the total three-dimensional resource available for every present and
future engineering system.

%%%%%%%%%%%%%%%%%%%%%%%%%%%%%%%%%%%%%%%%%%%%%%%%%%%%%%%%%%%%%%%

\begin{derivation}{Volume of a Uniform Corridor}

Suppose

\[
A(s)=A_0.
\]

Then

\[
V
=
\int_0^L
A_0\,ds.
\]

Since

\[
A_0
\]

is constant,

\[
V
=
A_0
\int_0^Lds.
\]

Therefore

\[
\boxed{
V=A_0L.
}
\]

Thus corridor volume scales linearly with corridor length.

\end{derivation}

%%%%%%%%%%%%%%%%%%%%%%%%%%%%%%%%%%%%%%%%%%%%%%%%%%%%%%%%%%%%%%%

\section{Geometric Allocation}

Suppose

\[
m
\]

independent engineering systems occupy the corridor.

Each subsystem receives a measurable region

\[
\Sigma_i(s)
\subseteq
\Sigma(s).
\]

These satisfy

\[
\Sigma_i
\cap
\Sigma_j
=
\emptyset,
\qquad
i\neq j,
\]

unless explicit interface coupling is declared.

Furthermore,

\[
\bigcup_{i=1}^{m}
\Sigma_i
\subseteq
\Sigma.
\]

%%%%%%%%%%%%%%%%%%%%%%%%%%%%%%%%%%%%%%%%%%%%%%%%%%%%%%%%%%%%%%%

\begin{definition}{Allocation Function}

The allocation function

\[
a_i(s)
=
\operatorname{Area}
(
\Sigma_i(s)
)
\]

assigns cross-sectional area to subsystem

\[
i.
\]

Hence

\[
a_i(s)\ge0.
\]

The occupied area becomes

\[
A_{\rm occ}(s)
=
\sum_i
a_i(s).
\]

\end{definition}

%%%%%%%%%%%%%%%%%%%%%%%%%%%%%%%%%%%%%%%%%%%%%%%%%%%%%%%%%%%%%%%

Immediately,

\[
A_{\rm occ}(s)
\le
A(s).
\]

Equality corresponds to complete occupation.

Strict inequality corresponds to reserved future capacity.

%%%%%%%%%%%%%%%%%%%%%%%%%%%%%%%%%%%%%%%%%%%%%%%%%%%%%%%%%%%%%%%

\begin{theorem}

The allocation functions satisfy

\[
\sum_i
a_i(s)
\le
A(s).
\]

\end{theorem}

\begin{proof}

Since

\[
\Sigma_i
\]

are pairwise disjoint,

measure additivity gives

\[
\operatorname{Area}
\left(
\bigcup_i
\Sigma_i
\right)
=
\sum_i
a_i.
\]

Because

\[
\bigcup_i
\Sigma_i
\subseteq
\Sigma,
\]

monotonicity of measure implies

\[
\sum_i
a_i
=
\operatorname{Area}
\left(
\bigcup_i
\Sigma_i
\right)
\le
\operatorname{Area}
(\Sigma).
\]

Therefore

\[
\boxed{
\sum_i
a_i(s)
\le
A(s).
}
\]

\end{proof}

%%%%%%%%%%%%%%%%%%%%%%%%%%%%%%%%%%%%%%%%%%%%%%%%%%%%%%%%%%%%%%%

\section{Reserved Capacity}

Define

\[
a_R(s)
=
A(s)
-
A_{\rm occ}(s).
\]

This quantity measures unused corridor capacity.

Unlike conventional engineering optimization,
the blueprint does not regard

\[
a_R=0
\]

as necessarily desirable.

Instead,

\[
a_R
\]

represents optionality available to future generations.

%%%%%%%%%%%%%%%%%%%%%%%%%%%%%%%%%%%%%%%%%%%%%%%%%%%%%%%%%%%%%%%

\begin{definition}{Future Capacity Ratio}

The future capacity ratio is

\[
r_F(s)
=
\frac{
a_R(s)
}{
A(s)
}.
\]

Clearly,

\[
0
\le
r_F
\le
1.
\]

\end{definition}

High values indicate significant geometric flexibility.

Low values indicate heavily saturated corridors.

%%%%%%%%%%%%%%%%%%%%%%%%%%%%%%%%%%%%%%%%%%%%%%%%%%%%%%%%%%%%%%%

\section{Interface Boundaries}

Adjacent infrastructure systems require physical interfaces.

For any two neighboring allocations,

\[
\Sigma_i
\quad\text{and}\quad
\Sigma_j,
\]

their common boundary is

\[
\Gamma_{ij}
=
\partial\Sigma_i
\cap
\partial\Sigma_j.
\]

%%%%%%%%%%%%%%%%%%%%%%%%%%%%%%%%%%%%%%%%%%%%%%%%%%%%%%%%%%%%%%%

The interface length

\[
L_{ij}
=
\operatorname{Length}
(
\Gamma_{ij}
)
\]

plays an important engineering role.

Large interface lengths facilitate maintenance access,
shared cooling,
and coordinated construction.

Conversely,
hazardous subsystems may deliberately minimize

\[
L_{ij}
\]

through geometric separation.

%%%%%%%%%%%%%%%%%%%%%%%%%%%%%%%%%%%%%%%%%%%%%%%%%%%%%%%%%%%%%%%

\section{Adjacency Graph}

The cross section naturally defines a graph.

Each subsystem becomes a vertex.

Two vertices are adjacent whenever their allocated regions share a common
boundary.

%%%%%%%%%%%%%%%%%%%%%%%%%%%%%%%%%%%%%%%%%%%%%%%%%%%%%%%%%%%%%%%

\begin{definition}{Cross-Sectional Adjacency Graph}

Let

\[
G_\Sigma
=
(V_\Sigma,E_\Sigma)
\]

where

\[
V_\Sigma
=
\{
\Sigma_i
\},
\]

and

\[
(\Sigma_i,\Sigma_j)
\in
E_\Sigma
\]

whenever

\[
\Gamma_{ij}
\neq
\emptyset.
\]

\end{definition}

This graph captures geometric relationships independently of the
particular engineering discipline.

%%%%%%%%%%%%%%%%%%%%%%%%%%%%%%%%%%%%%%%%%%%%%%%%%%%%%%%%%%%%%%%

\section{Voronoi Allocation Within the Corridor}

Cross-sectional allocation need not be performed manually.

Suppose subsystem centers are represented by points

\[
p_1,\ldots,p_n
\]

inside

\[
\Sigma.
\]

The Voronoi partition is

\[
V_i
=
\{
x\in\Sigma:
\|x-p_i\|
\le
\|x-p_j\|,
\;
\forall j
\}.
\]

Each subsystem automatically receives the region nearest its assigned
center.

%%%%%%%%%%%%%%%%%%%%%%%%%%%%%%%%%%%%%%%%%%%%%%%%%%%%%%%%%%%%%%%

This construction possesses several desirable properties.

First,
every point belongs to exactly one subsystem.

Second,
boundaries lie along perpendicular bisectors.

Third,
small perturbations of subsystem centers produce continuous changes in
allocation geometry.

%%%%%%%%%%%%%%%%%%%%%%%%%%%%%%%%%%%%%%%%%%%%%%%%%%%%%%%%%%%%%%%

\begin{theorem}

The Voronoi regions form a complete partition of the cross section.

\end{theorem}

\begin{proof}

For every point

\[
x
\in
\Sigma,
\]

the finite set of distances

\[
\{
\|x-p_i\|
\}
\]

contains at least one minimum.

Assign

\[
x
\]

to any minimizing subsystem.

Hence

\[
\bigcup_iV_i
=
\Sigma.
\]

Distinct Voronoi interiors cannot overlap because two different strict
minimum distances cannot occur simultaneously.

Therefore

\[
V_i^\circ
\cap
V_j^\circ
=
\emptyset,
\qquad
i\neq j.
\]

Thus

\[
\{V_i\}
\]

forms a partition.

\end{proof}

%%%%%%%%%%%%%%%%%%%%%%%%%%%%%%%%%%%%%%%%%%%%%%%%%%%%%%%%%%%%%%%

\section{Delaunay Connectivity}

The dual graph of the Voronoi partition defines natural subsystem
adjacency.

Subsystem centers become vertices.

Edges connect centers whose Voronoi regions share a boundary.

This Delaunay graph minimizes unnecessarily elongated neighbor
relationships and supplies a natural communication topology for
inspection, monitoring, and maintenance routing within the corridor.

Consequently the same Voronoi--Delaunay duality introduced for planetary
geometry reappears at the scale of a single corridor cross section.

This recurrence of the same mathematical structure across multiple
orders of magnitude illustrates one of the central themes of the present
specification: infrastructure organization is governed not by scale but
by invariant geometric relationships. The same constructions that define
global service regions and node connectivity can be applied recursively
to the internal organization of individual corridors, yielding a
hierarchy of compatible geometric decompositions from planetary
arrangements down to subsystem layouts.

%%%%%%%%%%%%%%%%%%%%%%%%%%%%%%%%%%%%%%%%%%%%%%%%%%%%%%%%%%%%%%%%%%%%%%
\chapter{Hydraulic Transport}

\section{Introduction}

Water is the oldest and most fundamental transported commodity.
Long before electrical grids, railways, fiber optics, or atmospheric
launch systems, civilizations constructed hydraulic infrastructure to
redistribute spatially heterogeneous resources. The orthodromic
framework therefore treats hydraulic transport as the canonical example
of corridor utilization.

The objective of the present chapter is not to prescribe one particular
pipeline technology or pumping architecture. Instead, the aim is to
derive the governing mathematical equations that every hydraulic
implementation must satisfy.

Throughout this chapter water is modeled as a continuum occupying a
connected graph of conduits embedded within the corridor network.

%%%%%%%%%%%%%%%%%%%%%%%%%%%%%%%%%%%%%%%%%%%%%%%%%%%%%%%%%%%%%%%

\section{Conservation of Mass}

Consider a control volume

\[
\Omega
\]

containing a portion of a pipeline.

Let

\[
\rho(x,t)
\]

denote the fluid density.

The total mass contained within the control volume is

\[
M(t)
=
\int_{\Omega}
\rho
\,dV.
\]

Mass can change only through transport across the boundary.

If

\[
\mathbf v
\]

denotes the velocity field,

the outward mass flux becomes

\[
\rho\mathbf v.
\]

Hence

\[
\frac{d}{dt}
\int_{\Omega}
\rho\,dV
=
-
\int_{\partial\Omega}
\rho
\mathbf v
\cdot
\mathbf n
\,dS.
\]

%%%%%%%%%%%%%%%%%%%%%%%%%%%%%%%%%%%%%%%%%%%%%%%%%%%%%%%%%%%%%%%

Applying the divergence theorem,

\[
\int_{\partial\Omega}
\rho
\mathbf v
\cdot
\mathbf n
\,dS
=
\int_{\Omega}
\nabla\cdot
(\rho\mathbf v)
\,dV.
\]

Therefore

\[
\int_\Omega
\left(
\frac{\partial\rho}{\partial t}
+
\nabla\cdot
(\rho\mathbf v)
\right)
dV
=
0.
\]

Since the control volume is arbitrary,

\[
\boxed{
\frac{\partial\rho}{\partial t}
+
\nabla\cdot(\rho\mathbf v)
=
0.
}
\]

%%%%%%%%%%%%%%%%%%%%%%%%%%%%%%%%%%%%%%%%%%%%%%%%%%%%%%%%%%%%%%%

\begin{definition}{Continuity Equation}

The equation

\[
\frac{\partial\rho}{\partial t}
+
\nabla\cdot(\rho\mathbf v)
=
0
\]

is called the continuity equation.

Every hydraulic system satisfying conservation of mass obeys this
equation independently of pipe material, pump design, or corridor
geometry.

\end{definition}

%%%%%%%%%%%%%%%%%%%%%%%%%%%%%%%%%%%%%%%%%%%%%%%%%%%%%%%%%%%%%%%

\section{Incompressible Approximation}

For freshwater transported under ordinary engineering conditions,

density variations are extremely small.

Consequently,

\[
\rho
=
\rho_0,
\]

where

\[
\rho_0
\]

is constant.

The continuity equation simplifies immediately.

%%%%%%%%%%%%%%%%%%%%%%%%%%%%%%%%%%%%%%%%%%%%%%%%%%%%%%%%%%%%%%%

\begin{derivation}

Substituting

\[
\rho=\rho_0,
\]

gives

\[
\frac{\partial\rho_0}{\partial t}
+
\nabla\cdot(\rho_0\mathbf v)
=
0.
\]

Since

\[
\rho_0
\]

is constant,

\[
\rho_0
\nabla\cdot\mathbf v
=
0.
\]

Hence

\[
\boxed{
\nabla\cdot\mathbf v=0.
}
\]

\end{derivation}

The velocity field is therefore divergence free.

%%%%%%%%%%%%%%%%%%%%%%%%%%%%%%%%%%%%%%%%%%%%%%%%%%%%%%%%%%%%%%%

\section{Volume Flow Rate}

Consider a pipe with cross-sectional area

\[
A.
\]

Suppose velocity is approximately uniform.

The volume flow rate becomes

\[
Q
=
Av.
\]

Its SI units are

\[
{\rm m^3\,s^{-1}}.
\]

The corresponding mass flow rate equals

\[
\dot m
=
\rho Q.
\]

%%%%%%%%%%%%%%%%%%%%%%%%%%%%%%%%%%%%%%%%%%%%%%%%%%%%%%%%%%%%%%%

\begin{theorem}

For incompressible steady flow,

the volume flow rate remains constant along every connected pipe segment.

\end{theorem}

\begin{proof}

Let

\[
A_1,
v_1
\]

represent one cross section,

and

\[
A_2,
v_2
\]

another.

Mass conservation requires

\[
\rho A_1v_1
=
\rho A_2v_2.
\]

Since

\[
\rho
\]

is constant,

\[
A_1v_1
=
A_2v_2.
\]

Therefore

\[
Q_1
=
Q_2.
\]

\end{proof}

%%%%%%%%%%%%%%%%%%%%%%%%%%%%%%%%%%%%%%%%%%%%%%%%%%%%%%%%%%%%%%%

\section{Bernoulli's Equation}

The total mechanical energy per unit volume consists of

pressure,

kinetic energy,

and gravitational potential.

These contributions are

\[
P,
\]

\[
\frac12\rho v^2,
\]

and

\[
\rho gz.
\]

%%%%%%%%%%%%%%%%%%%%%%%%%%%%%%%%%%%%%%%%%%%%%%%%%%%%%%%%%%%%%%%

Ignoring viscous losses,

their sum remains constant along a streamline.

%%%%%%%%%%%%%%%%%%%%%%%%%%%%%%%%%%%%%%%%%%%%%%%%%%%%%%%%%%%%%%%

\begin{theorem}[Bernoulli]

For steady incompressible inviscid flow,

\[
\boxed{
P
+
\frac12\rho v^2
+
\rho gz
=
{\rm constant}.
}
\]

\end{theorem}

%%%%%%%%%%%%%%%%%%%%%%%%%%%%%%%%%%%%%%%%%%%%%%%%%%%%%%%%%%%%%%%

\begin{proof}

Begin with Euler's equation,

\[
\rho
(\mathbf v\cdot\nabla)\mathbf v
=
-
\nabla P
+
\rho\mathbf g.
\]

Using the vector identity

\[
(\mathbf v\cdot\nabla)\mathbf v
=
\nabla
\left(
\frac12v^2
\right)
-
\mathbf v
\times
(\nabla\times\mathbf v),
\]

and restricting attention to motion along a streamline,

the rotational term vanishes.

Hence

\[
\nabla
\left(
\frac12\rho v^2
+
P
+
\rho gz
\right)
=
0.
\]

Therefore

\[
P
+
\frac12\rho v^2
+
\rho gz
=
{\rm constant}.
\]

\end{proof}

%%%%%%%%%%%%%%%%%%%%%%%%%%%%%%%%%%%%%%%%%%%%%%%%%%%%%%%%%%%%%%%

\section{Hydraulic Head}

Rather than energy density,

civil engineering often measures

head.

Divide Bernoulli's equation by

\[
\rho g.
\]

One obtains

\[
\boxed{
\frac{P}{\rho g}
+
\frac{v^2}{2g}
+
z
=
H.
}
\]

The three terms correspond respectively to

pressure head,

velocity head,

and

elevation head.

%%%%%%%%%%%%%%%%%%%%%%%%%%%%%%%%%%%%%%%%%%%%%%%%%%%%%%%%%%%%%%%

\section{Head Loss}

Real pipelines possess friction.

Consequently,

Bernoulli's equation becomes

\[
\frac{P_1}{\rho g}
+
\frac{v_1^2}{2g}
+
z_1
=
\frac{P_2}{\rho g}
+
\frac{v_2^2}{2g}
+
z_2
+
h_f.
\]

The quantity

\[
h_f
\]

is the head loss.

%%%%%%%%%%%%%%%%%%%%%%%%%%%%%%%%%%%%%%%%%%%%%%%%%%%%%%%%%%%%%%%

For turbulent flow,

the Darcy--Weisbach relation gives

\[
\boxed{
h_f
=
f
\frac{L}{D}
\frac{v^2}{2g}.
}
\]

Here

\[
f
\]

is the friction factor,

\[
L
\]

the pipe length,

and

\[
D
\]

its diameter.

%%%%%%%%%%%%%%%%%%%%%%%%%%%%%%%%%%%%%%%%%%%%%%%%%%%%%%%%%%%%%%%

\section{Pump Power}

Suppose a pump raises water through head

\[
H.
\]

The work required each second equals

\[
\rho gQH.
\]

Hence

\[
\boxed{
P_{\rm pump}
=
\rho gQH.
}
\]

Accounting for efficiency,

\[
\eta,
\]

gives

\[
P_{\rm electrical}
=
\frac{\rho gQH}{\eta}.
\]

%%%%%%%%%%%%%%%%%%%%%%%%%%%%%%%%%%%%%%%%%%%%%%%%%%%%%%%%%%%%%%%

\section{Gravity Recovery}

Descending water possesses recoverable potential energy.

Suppose water falls through elevation

\[
\Delta h.
\]

The available power becomes

\[
\boxed{
P
=
\rho gQ\Delta h.
}
\]

If turbines recover a fraction

\[
\eta_t,
\]

then

\[
P_{\rm recovered}
=
\eta_t
\rho gQ\Delta h.
\]

This immediately suggests a design principle.

%%%%%%%%%%%%%%%%%%%%%%%%%%%%%%%%%%%%%%%%%%%%%%%%%%%%%%%%%%%%%%%

\begin{principle}{Energy Recovery}

Hydraulic transport should not be regarded solely as an energy consumer.

Where topography permits,

descending flow should regenerate electrical energy through integrated
turbine systems,

partially offsetting the pumping requirements of ascending sections.

\end{principle}

%%%%%%%%%%%%%%%%%%%%%%%%%%%%%%%%%%%%%%%%%%%%%%%%%%%%%%%%%%%%%%%

\section{Hydraulic Networks}

The corridor does not contain isolated pipes.

Instead,

the system forms a graph

\[
G=(V,E).
\]

Each edge represents a pipeline.

Each vertex represents a switching node.

%%%%%%%%%%%%%%%%%%%%%%%%%%%%%%%%%%%%%%%%%%%%%%%%%%%%%%%%%%%%%%%

Assign flow

\[
Q_e
\]

to every edge.

At every interior node,

mass conservation requires

\[
\boxed{
\sum_{e\in{\rm in}}
Q_e
=
\sum_{e\in{\rm out}}
Q_e.
}
\]

%%%%%%%%%%%%%%%%%%%%%%%%%%%%%%%%%%%%%%%%%%%%%%%%%%%%%%%%%%%%%%%

This equation is the graph-theoretic analogue of the continuity
equation.

%%%%%%%%%%%%%%%%%%%%%%%%%%%%%%%%%%%%%%%%%%%%%%%%%%%%%%%%%%%%%%%

\section{Incidence Matrix Formulation}

Let

\[
B
\]

be the incidence matrix of the hydraulic graph.

Collect all edge flows into the vector

\[
q.
\]

External withdrawals and injections form

\[
b.
\]

Then conservation at every node becomes

\[
\boxed{
Bq=b.
}
\]

This compact matrix equation governs every hydraulic network regardless
of its size.

%%%%%%%%%%%%%%%%%%%%%%%%%%%%%%%%%%%%%%%%%%%%%%%%%%%%%%%%%%%%%%%

\section{Optimization of Water Routing}

Suppose each edge possesses energy cost

\[
c_e(Q_e).
\]

The total operating cost becomes

\[
J(Q)
=
\sum_e
c_e(Q_e).
\]

The routing problem is

\[
\min_Q
J(Q)
\]

subject to

\[
Bq=b.
\]

%%%%%%%%%%%%%%%%%%%%%%%%%%%%%%%%%%%%%%%%%%%%%%%%%%%%%%%%%%%%%%%

This constrained optimization determines the energetically optimal
distribution of water throughout the orthodromic network while
respecting conservation of mass.

Later chapters will couple this hydraulic optimization to electrical
generation, geothermal extraction, and transportation so that all
corridor subsystems become components of a unified multi-commodity
network.

%%%%%%%%%%%%%%%%%%%%%%%%%%%%%%%%%%%%%%%%%%%%%%%%%%%%%%%%%%%%%%%%%%%%%%
\chapter{Energy Transport and Electrical Infrastructure}

\section{Introduction}

The hydraulic system developed in the previous chapter transports
material. Electrical infrastructure transports energy. Although the
physical mechanisms differ, both systems satisfy the same underlying
conservation principles. One conserves mass while the other conserves
charge. Both may therefore be described within a common mathematical
framework.

The purpose of this chapter is to derive the governing equations of
electrical transport from first principles before embedding them into
the orthodromic corridor network.

Throughout this chapter the electrical system is treated as a continuum
at large scales and as a weighted graph at network scales.

%%%%%%%%%%%%%%%%%%%%%%%%%%%%%%%%%%%%%%%%%%%%%%%%%%%%%%%%%%%%%%%

\section{Conservation of Charge}

Let

\[
\rho_e(x,t)
\]

denote the electric charge density.

The total charge contained within a control volume

\[
\Omega
\]

is

\[
Q(t)
=
\int_\Omega
\rho_e\,dV.
\]

Charge can change only through current crossing the boundary.

Let

\[
\mathbf J
\]

denote current density.

Then

\[
\frac{d}{dt}
\int_\Omega
\rho_e\,dV
=
-
\int_{\partial\Omega}
\mathbf J\cdot\mathbf n\,dS.
\]

Applying the divergence theorem gives

\[
\int_\Omega
\left(
\frac{\partial\rho_e}{\partial t}
+
\nabla\cdot\mathbf J
\right)
dV
=
0.
\]

Since the control volume is arbitrary,

\[
\boxed{
\frac{\partial\rho_e}{\partial t}
+
\nabla\cdot\mathbf J
=
0.
}
\]

%%%%%%%%%%%%%%%%%%%%%%%%%%%%%%%%%%%%%%%%%%%%%%%%%%%%%%%%%%%%%%%

\begin{definition}{Electrical Continuity Equation}

The equation

\[
\frac{\partial\rho_e}{\partial t}
+
\nabla\cdot\mathbf J
=
0
\]

expresses conservation of electric charge.

No electrical network satisfying classical electromagnetism can violate
this relation.

\end{definition}

%%%%%%%%%%%%%%%%%%%%%%%%%%%%%%%%%%%%%%%%%%%%%%%%%%%%%%%%%%%%%%%

\section{Steady-State Networks}

Electrical transmission systems normally operate near steady state.

Consequently,

\[
\frac{\partial\rho_e}{\partial t}
\approx0,
\]

and therefore

\[
\nabla\cdot\mathbf J
=
0.
\]

Current neither accumulates nor disappears inside transmission lines.

%%%%%%%%%%%%%%%%%%%%%%%%%%%%%%%%%%%%%%%%%%%%%%%%%%%%%%%%%%%%%%%

\section{Ohm's Law}

For ordinary conductors,

current density is proportional to electric field.

\[
\boxed{
\mathbf J
=
\sigma
\mathbf E,
}
\]

where

\[
\sigma
\]

is electrical conductivity.

Since

\[
\mathbf E
=
-\nabla V,
\]

we obtain

\[
\boxed{
\mathbf J
=
-
\sigma
\nabla V.
}
\]

Current therefore flows down the voltage gradient.

%%%%%%%%%%%%%%%%%%%%%%%%%%%%%%%%%%%%%%%%%%%%%%%%%%%%%%%%%%%%%%%

\section{Potential Equation}

Substituting Ohm's law into charge conservation yields

\[
\nabla\cdot
(
\sigma\nabla V
)
=
0.
\]

If conductivity is constant,

\[
\sigma=\sigma_0,
\]

then

\[
\sigma_0
\nabla^2V
=
0.
\]

Hence

\[
\boxed{
\nabla^2V=0.
}
\]

The electrical potential therefore satisfies Laplace's equation.

%%%%%%%%%%%%%%%%%%%%%%%%%%%%%%%%%%%%%%%%%%%%%%%%%%%%%%%%%%%%%%%

\begin{remark}

Hydraulic head satisfies equations closely analogous to electrical
potential.

The correspondence

\[
V
\longleftrightarrow
H
\]

allows many hydraulic optimization methods to be translated directly
into electrical network analysis.

\end{remark}

%%%%%%%%%%%%%%%%%%%%%%%%%%%%%%%%%%%%%%%%%%%%%%%%%%%%%%%%%%%%%%%

\section{Electrical Power}

Suppose current

\[
I
\]

flows through voltage difference

\[
V.
\]

Power is

\[
\boxed{
P=VI.
}
\]

Using Ohm's law,

\[
V=IR,
\]

one immediately obtains

\[
P
=
I^2R.
\]

Equivalently,

\[
P
=
\frac{V^2}{R}.
\]

%%%%%%%%%%%%%%%%%%%%%%%%%%%%%%%%%%%%%%%%%%%%%%%%%%%%%%%%%%%%%%%

\section{Transmission Losses}

Energy is dissipated by resistance.

Suppose a transmission line has resistance

\[
R.
\]

The thermal loss is

\[
\boxed{
P_{\rm loss}
=
I^2R.
}
\]

This quadratic dependence immediately explains why long-distance
transmission employs high voltages.

%%%%%%%%%%%%%%%%%%%%%%%%%%%%%%%%%%%%%%%%%%%%%%%%%%%%%%%%%%%%%%%

\begin{derivation}

Power delivered is

\[
P
=
VI.
\]

For fixed transmitted power,

\[
I
=
\frac PV.
\]

Substituting into the loss equation,

\[
P_{\rm loss}
=
\left(
\frac PV
\right)^2
R.
\]

Hence

\[
\boxed{
P_{\rm loss}
=
\frac{P^2R}{V^2}.
}
\]

Therefore increasing voltage decreases resistive loss quadratically.

\end{derivation}

%%%%%%%%%%%%%%%%%%%%%%%%%%%%%%%%%%%%%%%%%%%%%%%%%%%%%%%%%%%%%%%

\section{Electrical Graphs}

Let

\[
G=(V,E)
\]

represent the electrical backbone.

Vertices correspond to substations.

Edges correspond to transmission corridors.

Assign current

\[
I_e
\]

to every edge.

Kirchhoff's Current Law states

\[
\boxed{
\sum_{\rm in}I
=
\sum_{\rm out}I.
}
\]

This is precisely the graph analogue of charge conservation.

%%%%%%%%%%%%%%%%%%%%%%%%%%%%%%%%%%%%%%%%%%%%%%%%%%%%%%%%%%%%%%%

\section{Incidence Matrix}

Let

\[
B
\]

be the graph incidence matrix.

Collect all edge currents into

\[
i.
\]

Generator injections and loads form

\[
s.
\]

Then

\[
\boxed{
Bi=s.
}
\]

This equation is structurally identical to the hydraulic conservation
equation

\[
Bq=b.
\]

%%%%%%%%%%%%%%%%%%%%%%%%%%%%%%%%%%%%%%%%%%%%%%%%%%%%%%%%%%%%%%%

\section{Electrical Resistance Optimization}

Suppose edge

\[
e
\]

possesses resistance

\[
R_e.
\]

The total dissipation is

\[
J(i)
=
\sum_e
R_eI_e^2.
\]

The optimal current distribution minimizes

\[
J(i)
\]

subject to

\[
Bi=s.
\]

%%%%%%%%%%%%%%%%%%%%%%%%%%%%%%%%%%%%%%%%%%%%%%%%%%%%%%%%%%%%%%%

\begin{theorem}

The electrical operating state minimizes total resistive dissipation
subject to current conservation.

\end{theorem}

\begin{proof}

Introduce Lagrange multipliers

\[
\lambda.
\]

The Lagrangian is

\[
L
=
\sum_e
R_eI_e^2
+
\lambda^T(Bi-s).
\]

Differentiating,

\[
2R_eI_e
+
(B^T\lambda)_e
=
0.
\]

Therefore

\[
I_e
=
-
\frac{
(B^T\lambda)_e
}
{2R_e}.
\]

Substituting into the conservation equations determines the unique
minimum-energy solution.

\end{proof}

%%%%%%%%%%%%%%%%%%%%%%%%%%%%%%%%%%%%%%%%%%%%%%%%%%%%%%%%%%%%%%%

\section{Hydraulic--Electrical Duality}

The previous chapter developed hydraulic transport.

This chapter developed electrical transport.

Their governing equations possess striking structural similarity.

Hydraulics satisfies

\[
Bq=b,
\]

while electricity satisfies

\[
Bi=s.
\]

Hydraulic energy loss scales approximately with

\[
Q^2,
\]

while electrical dissipation scales with

\[
I^2.
\]

Pressure head corresponds to voltage.

Flow rate corresponds to current.

Hydraulic resistance corresponds to electrical resistance.

These analogies are summarized by the correspondence

\[
\begin{array}{ccc}
\textbf{Hydraulics}
&
&
\textbf{Electricity}
\\[1ex]
Q
&
\leftrightarrow
&
I
\\
H
&
\leftrightarrow
&
V
\\
R_h
&
\leftrightarrow
&
R_e
\\
BQ=b
&
\leftrightarrow
&
BI=s.
\end{array}
\]

This duality is not merely pedagogical.

It permits common optimization algorithms to solve both systems.

%%%%%%%%%%%%%%%%%%%%%%%%%%%%%%%%%%%%%%%%%%%%%%%%%%%%%%%%%%%%%%%

\section{Energy Storage Along Corridors}

Electrical demand and renewable generation fluctuate continuously.

Consequently, storage becomes an essential corridor function.

Let

\[
E(t)
\]

denote stored energy.

If charging power is

\[
P_c
\]

and discharge power is

\[
P_d,
\]

then

\[
\boxed{
\frac{dE}{dt}
=
\eta_cP_c
-
\frac{P_d}{\eta_d},
}
\]

where

\[
\eta_c,
\eta_d
\]

are charging and discharge efficiencies.

The storage state satisfies

\[
0
\le
E
\le
E_{\max}.
\]

Storage devices therefore become dynamical state variables rather than
mere network components.

%%%%%%%%%%%%%%%%%%%%%%%%%%%%%%%%%%%%%%%%%%%%%%%%%%%%%%%%%%%%%%%

\section{Coupling with Hydraulic Infrastructure}

Electrical and hydraulic systems are not independent.

Pumps consume electricity.

Hydroelectric turbines produce electricity.

Consequently,

the hydraulic flow vector

\[
Q
\]

and electrical current vector

\[
I
\]

must be coupled.

A pumping station contributes

\[
P_{\rm pump}
=
\rho gQH,
\]

to the electrical load vector.

A hydroelectric station contributes

\[
P_{\rm hydro}
=
\eta
\rho gQ\Delta h,
\]

to the generator vector.

The two infrastructure systems therefore exchange energy continuously.

%%%%%%%%%%%%%%%%%%%%%%%%%%%%%%%%%%%%%%%%%%%%%%%%%%%%%%%%%%%%%%%

\section{Unified Conservation Structure}

Hydraulic transport conserves mass.

Electrical transport conserves charge.

Both obey graph conservation laws.

Consequently they may be combined into a single state vector

\[
x
=
\begin{pmatrix}
Q\\
I\\
E
\end{pmatrix},
\]

whose evolution satisfies

\[
\boxed{
\frac{dx}{dt}
=
F(x,u),
}
\]

where

\[
u
\]

represents controllable infrastructure decisions such as pumping,
generation, switching, storage dispatch, and demand management.

Later chapters enlarge this state vector to include
transportation, communications, thermal systems, and launch
infrastructure, yielding a unified mathematical description of the
entire orthodromic corridor network.

%%%%%%%%%%%%%%%%%%%%%%%%%%%%%%%%%%%%%%%%%%%%%%%%%%%%%%%%%%%%%%%%%%%%%%
\chapter{Transportation Dynamics}

\section{Introduction}

The movement of matter and the movement of people are governed by many of
the same mathematical principles. A freight train, a convoy of trucks,
an autonomous vehicle platoon, and a stream of passengers all represent
transported quantities moving through a constrained network.

Unlike electrical current, transportation systems cannot generally be
treated as continuous media under all operating conditions. Individual
vehicles remain discrete objects. Nevertheless, when the number of
vehicles becomes sufficiently large, continuum models provide remarkably
accurate approximations. Such models permit transportation to be
integrated naturally with the hydraulic and electrical systems developed
in the preceding chapters.

Throughout this chapter the transported quantity will be referred to as
traffic density, although the same mathematical framework applies to
cargo, autonomous freight, pedestrians, rail vehicles, maintenance
equipment, and other mobile assets.

%%%%%%%%%%%%%%%%%%%%%%%%%%%%%%%%%%%%%%%%%%%%%%%%%%%%%%%%%%%%%%%

\section{Traffic Density}

Consider a corridor parameterized by arc length

\[
s\in[0,L].
\]

Let

\[
\rho(s,t)
\]

denote the average number of vehicles per unit length.

Its units are

\[
{\rm vehicles\,m^{-1}}.
\]

Similarly define

\[
v(s,t)
\]

as the average vehicle speed.

The corresponding traffic flux is

\[
q(s,t)
=
\rho(s,t)v(s,t).
\]

Flux measures the number of vehicles crossing a fixed point per unit
time.

%%%%%%%%%%%%%%%%%%%%%%%%%%%%%%%%%%%%%%%%%%%%%%%%%%%%%%%%%%%%%%%

\section{Conservation of Vehicles}

Consider an interval

\[
[a,b]
\]

along the corridor.

The total number of vehicles contained within this interval equals

\[
N(t)
=
\int_a^b
\rho(s,t)\,ds.
\]

Vehicles are neither created nor destroyed inside the interval.

They enter through one boundary and leave through the other.

Consequently,

\[
\frac{dN}{dt}
=
q(a,t)
-
q(b,t).
\]

Using the Fundamental Theorem of Calculus,

\[
\frac{d}{dt}
\int_a^b
\rho\,ds
=
-
\int_a^b
\frac{\partial q}{\partial s}\,ds.
\]

Since the interval is arbitrary,

\[
\boxed{
\frac{\partial\rho}{\partial t}
+
\frac{\partial q}{\partial s}
=
0.
}
\]

%%%%%%%%%%%%%%%%%%%%%%%%%%%%%%%%%%%%%%%%%%%%%%%%%%%%%%%%%%%%%%%

\begin{definition}{Traffic Continuity Equation}

The equation

\[
\frac{\partial\rho}{\partial t}
+
\frac{\partial q}{\partial s}
=
0
\]

expresses conservation of vehicles along a transportation corridor.

It is mathematically identical in form to the continuity equations for
mass and electric charge derived previously.

\end{definition}

%%%%%%%%%%%%%%%%%%%%%%%%%%%%%%%%%%%%%%%%%%%%%%%%%%%%%%%%%%%%%%%

\section{Fundamental Diagram}

Vehicle speed depends upon traffic density.

When density is low,

vehicles travel near the design speed

\[
v_{\max}.
\]

As density increases,

interactions between vehicles reduce average velocity.

A simple constitutive relation is

\[
\boxed{
v(\rho)
=
v_{\max}
\left(
1-
\frac{\rho}{\rho_{\max}}
\right),
}
\]

where

\[
\rho_{\max}
\]

is the jam density.

Substituting,

\[
q(\rho)
=
\rho
v_{\max}
\left(
1-
\frac{\rho}{\rho_{\max}}
\right).
\]

This quadratic relation is known as the Greenshields fundamental
diagram.

%%%%%%%%%%%%%%%%%%%%%%%%%%%%%%%%%%%%%%%%%%%%%%%%%%%%%%%%%%%%%%%

\begin{derivation}{Maximum Throughput}

Differentiate

\[
q(\rho)
=
v_{\max}
\left(
\rho
-
\frac{\rho^2}{\rho_{\max}}
\right).
\]

Hence

\[
\frac{dq}{d\rho}
=
v_{\max}
\left(
1-
2\frac{\rho}{\rho_{\max}}
\right).
\]

Setting

\[
\frac{dq}{d\rho}=0
\]

gives

\[
\rho
=
\frac{\rho_{\max}}2.
\]

Substituting,

\[
q_{\max}
=
v_{\max}
\frac{\rho_{\max}}4.
\]

Thus maximum corridor throughput occurs at one-half of jam density.

\end{derivation}

%%%%%%%%%%%%%%%%%%%%%%%%%%%%%%%%%%%%%%%%%%%%%%%%%%%%%%%%%%%%%%%

\section{The Lighthill--Whitham--Richards Equation}

Substituting the constitutive law

\[
q=q(\rho)
\]

into the conservation equation yields

\[
\boxed{
\frac{\partial\rho}{\partial t}
+
\frac{\partial q(\rho)}{\partial s}
=
0.
}
\]

Applying the chain rule,

\[
\frac{\partial q}{\partial s}
=
\frac{dq}{d\rho}
\frac{\partial\rho}{\partial s},
\]

gives

\[
\boxed{
\frac{\partial\rho}{\partial t}
+
c(\rho)
\frac{\partial\rho}{\partial s}
=
0,
}
\]

where

\[
c(\rho)
=
\frac{dq}{d\rho}
\]

is the propagation speed of traffic disturbances.

%%%%%%%%%%%%%%%%%%%%%%%%%%%%%%%%%%%%%%%%%%%%%%%%%%%%%%%%%%%%%%%

\section{Traffic Waves}

A local increase in density propagates through the corridor.

Unlike vehicles,

the disturbance itself may travel upstream.

Suppose two traffic states satisfy

\[
(\rho_1,q_1)
\]

and

\[
(\rho_2,q_2).
\]

The interface separating them moves with speed

\[
\boxed{
w
=
\frac{q_2-q_1}
{\rho_2-\rho_1}.
}
\]

This is the Rankine--Hugoniot condition.

Traffic congestion therefore behaves mathematically as a moving shock.

%%%%%%%%%%%%%%%%%%%%%%%%%%%%%%%%%%%%%%%%%%%%%%%%%%%%%%%%%%%%%%%

\section{Transportation Graphs}

Planetary transportation consists of corridors joined at nodes.

Represent the system as

\[
G=(V,E).
\]

Each edge carries traffic flow

\[
q_e.
\]

Each node satisfies conservation.

Thus

\[
\boxed{
\sum_{\rm in}
q_e
=
\sum_{\rm out}
q_e
+
d_v,
}
\]

where

\[
d_v
\]

is local demand.

%%%%%%%%%%%%%%%%%%%%%%%%%%%%%%%%%%%%%%%%%%%%%%%%%%%%%%%%%%%%%%%

\section{Incidence Matrix Formulation}

Collect edge flows into

\[
q.
\]

Let

\[
B
\]

denote the incidence matrix.

Demand becomes

\[
d.
\]

The network equations become

\[
\boxed{
Bq=d.
}
\]

This is structurally identical to the hydraulic equation

\[
BQ=b
\]

and the electrical equation

\[
BI=s.
\]

%%%%%%%%%%%%%%%%%%%%%%%%%%%%%%%%%%%%%%%%%%%%%%%%%%%%%%%%%%%%%%%

\section{Travel Time}

Each edge possesses length

\[
L_e
\]

and average speed

\[
v_e.
\]

The travel time is

\[
\boxed{
T_e
=
\frac{L_e}{v_e}.
}
\]

The travel time of an entire route equals

\[
T
=
\sum_e
T_e.
\]

When speed depends upon congestion,

\[
v_e=v_e(\rho),
\]

travel time becomes nonlinear.

%%%%%%%%%%%%%%%%%%%%%%%%%%%%%%%%%%%%%%%%%%%%%%%%%%%%%%%%%%%%%%%

\section{Shortest Paths}

Assign edge weights

\[
w_e=T_e.
\]

The optimal route minimizes

\[
\boxed{
T
=
\sum_{e\in P}
w_e,
}
\]

where

\[
P
\]

is a path joining origin and destination.

When congestion changes,

weights evolve dynamically.

Consequently,

optimal routing becomes a time-dependent optimization problem.

%%%%%%%%%%%%%%%%%%%%%%%%%%%%%%%%%%%%%%%%%%%%%%%%%%%%%%%%%%%%%%%

\section{Transportation Capacity}

Each corridor possesses finite carrying capacity.

Let

\[
C_e
\]

denote the maximum sustainable flow.

Then

\[
0
\le
q_e
\le
C_e.
\]

Capacity itself depends upon lane count,

vehicle spacing,

speed limits,

control systems,

weather,

maintenance,

and infrastructure condition.

%%%%%%%%%%%%%%%%%%%%%%%%%%%%%%%%%%%%%%%%%%%%%%%%%%%%%%%%%%%%%%%

\section{Rail Systems}

Rail transport differs from highway traffic primarily through stronger
constraints on vehicle trajectories.

Every train follows a predetermined guideway.

Passing opportunities occur only at designated sidings or multiple-track
segments.

Mathematically,

the conservation equation remains unchanged.

Only the admissible state space differs.

Let

\[
x_i(t)
\]

denote the position of train

\[
i.
\]

Safety requires

\[
x_{i+1}-x_i
\ge
d_{\rm safe},
\]

where

\[
d_{\rm safe}
\]

is the minimum permitted headway.

Consequently,

rail capacity becomes a constrained optimization problem over allowable
vehicle trajectories.

%%%%%%%%%%%%%%%%%%%%%%%%%%%%%%%%%%%%%%%%%%%%%%%%%%%%%%%%%%%%%%%

\section{Freight Flow}

Instead of counting vehicles,

consider transported mass.

Let

\[
m(s,t)
\]

denote freight density.

If

\[
v_f
\]

is freight velocity,

then

\[
f
=
mv_f
\]

is freight flux.

The governing equation becomes

\[
\boxed{
\frac{\partial m}{\partial t}
+
\frac{\partial f}{\partial s}
=
0.
}
\]

Thus freight transport satisfies exactly the same conservation law as
ordinary traffic.

%%%%%%%%%%%%%%%%%%%%%%%%%%%%%%%%%%%%%%%%%%%%%%%%%%%%%%%%%%%%%%%

\section{Passenger Flow}

Passenger transport differs only by replacing freight density with human
occupancy.

Define

\[
p(s,t)
\]

as passenger density.

Then

\[
J_p
=
pv.
\]

Hence

\[
\boxed{
\frac{\partial p}{\partial t}
+
\frac{\partial J_p}{\partial s}
=
0.
}
\]

This equation governs passenger movement through stations, tunnels,
corridors, and evacuation routes.

%%%%%%%%%%%%%%%%%%%%%%%%%%%%%%%%%%%%%%%%%%%%%%%%%%%%%%%%%%%%%%%

\section{Unified Transportation State}

Road vehicles,

rail traffic,

maintenance equipment,

construction machinery,

freight,

and passengers all satisfy conservation equations.

Collect these quantities into

\[
x_T
=
\begin{pmatrix}
\rho\\
m\\
p
\end{pmatrix}.
\]

Their dynamics satisfy

\[
\boxed{
\frac{\partial x_T}{\partial t}
+
\nabla\cdot
F_T(x_T)
=
S_T,
}
\]

where

\[
S_T
\]

represents loading,

unloading,

boarding,

maintenance operations,

or emergency interventions.

%%%%%%%%%%%%%%%%%%%%%%%%%%%%%%%%%%%%%%%%%%%%%%%%%%%%%%%%%%%%%%%

\section{Relationship to the Other Corridor Systems}

At this stage the mathematical similarity among all infrastructure
systems becomes apparent.

Hydraulics conserves mass.

Electricity conserves charge.

Transportation conserves vehicles,

cargo,

or passengers.

Each produces a conservation law of the form

\[
\frac{\partial u}{\partial t}
+
\nabla\cdot F(u)
=
S.
\]

The specific physical interpretation changes,

but the underlying mathematical structure remains invariant.

This observation motivates the next chapter, in which all corridor
subsystems will be combined into a unified multi-commodity dynamical
system governed by a single vector conservation equation rather than by
independent discipline-specific models.

%%%%%%%%%%%%%%%%%%%%%%%%%%%%%%%%%%%%%%%%%%%%%%%%%%%%%%%%%%%%%%%%%%%%%%
\chapter{Unified Multi-Commodity Dynamics}

\section{Introduction}

The previous chapters derived governing equations for hydraulic,
electrical, and transportation systems independently.

Although these systems differ physically, each obeys a conservation law.
Water conserves mass. Electrical systems conserve charge. Transportation
conserves vehicles, passengers, or freight.

This repeated mathematical structure suggests that the distinction
between infrastructure disciplines is secondary. The primary object is
the conservation law itself.

Infrastructure may therefore be regarded as a collection of interacting
conserved quantities evolving on a common geometric substrate.

%%%%%%%%%%%%%%%%%%%%%%%%%%%%%%%%%%%%%%%%%%%%%%%%%%%%%%%%%%%%%%%

\section{Commodity Space}

Suppose the corridor supports

\[
n
\]

independent commodities.

Examples include

\[
\begin{aligned}
u_1 &= \text{water},\\
u_2 &= \text{electrical energy},\\
u_3 &= \text{rail vehicles},\\
u_4 &= \text{road traffic},\\
u_5 &= \text{communications},\\
u_6 &= \text{thermal energy},\\
u_7 &= \text{hydrogen},\\
u_8 &= \text{maintenance resources}.
\end{aligned}
\]

Collect these into the state vector

\[
u(x,t)
=
\begin{pmatrix}
u_1\\
u_2\\
\vdots\\
u_n
\end{pmatrix}.
\]

The state of the entire infrastructure network is therefore represented
by a single vector-valued field.

%%%%%%%%%%%%%%%%%%%%%%%%%%%%%%%%%%%%%%%%%%%%%%%%%%%%%%%%%%%%%%%

\section{General Conservation Law}

Each commodity possesses a flux

\[
F_i(u).
\]

Collect these into

\[
F(u)
=
\begin{pmatrix}
F_1\\
F_2\\
\vdots\\
F_n
\end{pmatrix}.
\]

External production and consumption become

\[
S(u).
\]

The governing equation becomes

\[
\boxed{
\frac{\partial u}{\partial t}
+
\nabla\cdot F(u)
=
S(u).
}
\]

This equation will serve as the fundamental dynamical equation of the
present blueprint.

%%%%%%%%%%%%%%%%%%%%%%%%%%%%%%%%%%%%%%%%%%%%%%%%%%%%%%%%%%%%%%%

\begin{definition}{Infrastructure State}

The pair

\[
(u,F)
\]

consisting of the commodity field and its associated flux field is called
an infrastructure state.

\end{definition}

%%%%%%%%%%%%%%%%%%%%%%%%%%%%%%%%%%%%%%%%%%%%%%%%%%%%%%%%%%%%%%%

\section{Linear Coupling}

Infrastructure systems rarely evolve independently.

Electricity powers pumps.

Transportation delivers repair crews.

Communications control substations.

Water cools generators.

Represent these interactions through a coupling matrix

\[
C
=
(c_{ij}).
\]

Entry

\[
c_{ij}
\]

measures the influence of commodity

\[
j
\]

upon commodity

\[
i.
\]

The source term becomes

\[
S
=
Cu
+
b.
\]

Consequently

\[
\boxed{
\frac{\partial u}{\partial t}
+
\nabla\cdot F(u)
=
Cu+b.
}
\]

%%%%%%%%%%%%%%%%%%%%%%%%%%%%%%%%%%%%%%%%%%%%%%%%%%%%%%%%%%%%%%%

\section{Interpretation of the Coupling Matrix}

The diagonal entries

\[
c_{ii}
\]

represent internal production or decay.

The off-diagonal entries

\[
c_{ij},
\qquad
i\neq j,
\]

represent interactions between infrastructure systems.

For example,

\[
c_{\text{water},\text{electric}}
>0
\]

models electrically driven pumping.

Similarly,

\[
c_{\text{electric},\text{water}}
>0
\]

models hydroelectric generation.

The matrix therefore represents the infrastructure ecology.

%%%%%%%%%%%%%%%%%%%%%%%%%%%%%%%%%%%%%%%%%%%%%%%%%%%%%%%%%%%%%%%

\begin{definition}{Infrastructure Ecology}

The weighted directed graph induced by the coupling matrix

\[
C
\]

is called the infrastructure ecology.

Vertices correspond to conserved commodities.

Edges correspond to energetic, informational, or operational
dependencies.

\end{definition}

%%%%%%%%%%%%%%%%%%%%%%%%%%%%%%%%%%%%%%%%%%%%%%%%%%%%%%%%%%%%%%%

\section{Equilibrium States}

A stationary infrastructure satisfies

\[
\frac{\partial u}{\partial t}
=
0.
\]

Hence

\[
\nabla\cdot F(u)
=
Cu+b.
\]

If spatial gradients also vanish,

then

\[
Cu+b=0.
\]

When

\[
C
\]

is invertible,

\[
\boxed{
u^*
=
-C^{-1}b.
}
\]

Thus equilibrium infrastructure is determined by solving a linear system.

%%%%%%%%%%%%%%%%%%%%%%%%%%%%%%%%%%%%%%%%%%%%%%%%%%%%%%%%%%%%%%%

\section{Linear Stability}

Suppose

\[
u^*
\]

is an equilibrium.

Introduce a perturbation

\[
u=u^*+\delta u.
\]

Linearizing,

\[
\frac{\partial(\delta u)}{\partial t}
=
C\,\delta u.
\]

Solutions are

\[
\delta u(t)
=
e^{Ct}\delta u(0).
\]

Therefore the eigenvalues of

\[
C
\]

determine stability.

%%%%%%%%%%%%%%%%%%%%%%%%%%%%%%%%%%%%%%%%%%%%%%%%%%%%%%%%%%%%%%%

\begin{theorem}

If every eigenvalue of

\[
C
\]

has strictly negative real part,

the equilibrium is asymptotically stable.

\end{theorem}

\begin{proof}

Diagonalize

\[
C=P\Lambda P^{-1}.
\]

Then

\[
e^{Ct}
=
Pe^{\Lambda t}P^{-1}.
\]

Since every exponential

\[
e^{\lambda_i t}
\]

decays whenever

\[
\operatorname{Re}(\lambda_i)<0,
\]

all perturbations converge exponentially toward equilibrium.

\end{proof}

%%%%%%%%%%%%%%%%%%%%%%%%%%%%%%%%%%%%%%%%%%%%%%%%%%%%%%%%%%%%%%%

\section{Infrastructure Capacity}

Every commodity possesses finite carrying capacity.

Let

\[
K_i
\]

denote the capacity associated with commodity

\[
i.
\]

Normalize the state by

\[
\phi_i
=
\frac{u_i}{K_i}.
\]

Consequently,

\[
0
\le
\phi_i
\le
1
\]

measures infrastructure utilization independently of engineering units.

This normalization permits water flow, electrical transmission,
transportation throughput, and communication bandwidth to be compared
within a common mathematical framework.

%%%%%%%%%%%%%%%%%%%%%%%%%%%%%%%%%%%%%%%%%%%%%%%%%%%%%%%%%%%%%%%

\section{The Universal Infrastructure Equation}

The previous chapters now collapse into a single mathematical object.

\[
\boxed{
\frac{\partial u}{\partial t}
+
\nabla\cdot F(u)
=
Cu+b.
}
\]

The geometry established in Parts I and II determines where the state
exists.

The conservation laws determine how commodities move.

The coupling matrix determines how the infrastructure systems interact.

The remainder of this monograph will study particular choices of the
flux function

\[
F,
\]

the coupling matrix

\[
C,
\]

and the optimization functionals governing their design.

%%%%%%%%%%%%%%%%%%%%%%%%%%%%%%%%%%%%%%%%%%%%%%%%%%%%%%%%%%%%%%%%%%%%%%
\part{Variational Infrastructure Theory}

%%%%%%%%%%%%%%%%%%%%%%%%%%%%%%%%%%%%%%%%%%%%%%%%%%%%%%%%%%%%%%%%%%%%%%
\chapter{Infrastructure Functionals}

\section{Introduction}

The previous chapters established that hydraulic transport,
electrical transmission, transportation systems, and other corridor
commodities satisfy conservation laws. These equations describe the
evolution of infrastructure states.

A complementary question remains.

Among all admissible infrastructure states, why does one particular state
occur rather than another?

The answer lies in optimization.

Many physical systems evolve toward configurations minimizing energy,
travel time, dissipation, or construction cost. Others maximize
throughput, reliability, accessibility, or resilience.

Rather than treating each optimization independently, the present
chapter introduces a common mathematical language through the theory of
functionals.

%%%%%%%%%%%%%%%%%%%%%%%%%%%%%%%%%%%%%%%%%%%%%%%%%%%%%%%%%%%%%%%

\section{Function Spaces}

Let

\[
\Omega
\]

denote the planetary infrastructure domain.

An infrastructure configuration is represented by a vector-valued field

\[
u:\Omega\rightarrow\mathbb R^n.
\]

The collection of all admissible configurations forms a function space

\[
\mathcal A.
\]

Every optimization considered throughout this monograph is therefore
performed over

\[
\mathcal A,
\]

rather than over finite-dimensional vectors.

%%%%%%%%%%%%%%%%%%%%%%%%%%%%%%%%%%%%%%%%%%%%%%%%%%%%%%%%%%%%%%%

\begin{definition}{Admissible Configuration}

An admissible configuration is a field

\[
u\in\mathcal A
\]

satisfying

\begin{enumerate}
\item conservation laws,
\item boundary conditions,
\item capacity constraints,
\item geometric compatibility,
\item physical feasibility.
\end{enumerate}

\end{definition}

%%%%%%%%%%%%%%%%%%%%%%%%%%%%%%%%%%%%%%%%%%%%%%%%%%%%%%%%%%%%%%%

\section{Infrastructure Functionals}

Unlike ordinary functions,

\[
f:\mathbb R^n\rightarrow\mathbb R,
\]

a functional assigns a real number to an entire field.

\[
J:
\mathcal A
\rightarrow
\mathbb R.
\]

Examples include

construction cost,

energy expenditure,

travel time,

environmental impact,

maintenance burden,

and resilience.

%%%%%%%%%%%%%%%%%%%%%%%%%%%%%%%%%%%%%%%%%%%%%%%%%%%%%%%%%%%%%%%

\begin{definition}{Infrastructure Functional}

A functional is any mapping

\[
J[u]
\]

assigning a real number to every admissible infrastructure state.

\end{definition}

%%%%%%%%%%%%%%%%%%%%%%%%%%%%%%%%%%%%%%%%%%%%%%%%%%%%%%%%%%%%%%%

\section{Integral Functionals}

Most engineering objectives possess the form

\[
\boxed{
J[u]
=
\int_\Omega
L(x,u,\nabla u)\,dV,
}
\]

where

\[
L
\]

is called the infrastructure density.

The quantity

\[
L
\]

may represent energetic cost,

construction expense,

carbon footprint,

or any other local objective.

%%%%%%%%%%%%%%%%%%%%%%%%%%%%%%%%%%%%%%%%%%%%%%%%%%%%%%%%%%%%%%%

\section{Examples}

Construction cost:

\[
J_C
=
\int_\Omega
c(x)\,dV.
\]

Energy dissipation:

\[
J_E
=
\int_\Omega
R(x)
I(x)^2\,dV.
\]

Hydraulic pumping:

\[
J_H
=
\int_\Omega
\rho gQH\,dV.
\]

Travel time:

\[
J_T
=
\int_\Omega
\frac1{v(x)}\,ds.
\]

%%%%%%%%%%%%%%%%%%%%%%%%%%%%%%%%%%%%%%%%%%%%%%%%%%%%%%%%%%%%%%%

\section{The Principle of Optimal Infrastructure}

Physical infrastructure rarely occupies arbitrary configurations.

Instead it tends toward states minimizing one or more objective
functionals while respecting conservation laws.

%%%%%%%%%%%%%%%%%%%%%%%%%%%%%%%%%%%%%%%%%%%%%%%%%%%%%%%%%%%%%%%

\begin{principle}{Optimal Infrastructure Principle}

Among all admissible configurations,

the realized infrastructure minimizes an objective functional,

\[
\boxed{
u^\star
=
\operatorname*{arg\,min}_{u\in\mathcal A}
J[u].
}
\]

\end{principle}

%%%%%%%%%%%%%%%%%%%%%%%%%%%%%%%%%%%%%%%%%%%%%%%%%%%%%%%%%%%%%%%

\section{Variations}

Suppose

\[
u
\]

is perturbed slightly.

Introduce

\[
u_\varepsilon
=
u
+
\varepsilon\eta,
\]

where

\[
\eta
\]

is an arbitrary smooth perturbation satisfying the boundary conditions.

The variation of

\[
J
\]

is

\[
\delta J
=
\left.
\frac d{d\varepsilon}
J[u+\varepsilon\eta]
\right|_{\varepsilon=0}.
\]

%%%%%%%%%%%%%%%%%%%%%%%%%%%%%%%%%%%%%%%%%%%%%%%%%%%%%%%%%%%%%%%

\begin{definition}{Stationary Infrastructure}

An admissible configuration is stationary if

\[
\boxed{
\delta J=0
}
\]

for every admissible perturbation.

\end{definition}

%%%%%%%%%%%%%%%%%%%%%%%%%%%%%%%%%%%%%%%%%%%%%%%%%%%%%%%%%%%%%%%

\section{Derivation of the Euler--Lagrange Equation}

Suppose

\[
J[u]
=
\int_\Omega
L(u,\nabla u)\,dV.
\]

Substitute

\[
u+\varepsilon\eta.
\]

Differentiating,

\[
\delta J
=
\int_\Omega
\left(
\frac{\partial L}{\partial u}\eta
+
\frac{\partial L}{\partial(\nabla u)}
\cdot
\nabla\eta
\right)
dV.
\]

Integrating the second term by parts,

\[
\delta J
=
\int_\Omega
\left(
\frac{\partial L}{\partial u}
-
\nabla\cdot
\frac{\partial L}
{\partial(\nabla u)}
\right)
\eta
\,dV.
\]

Since

\[
\eta
\]

is arbitrary,

the Fundamental Lemma of the Calculus of Variations gives

\[
\boxed{
\frac{\partial L}{\partial u}
-
\nabla\cdot
\left(
\frac{\partial L}
{\partial(\nabla u)}
\right)
=
0.
}
\]

%%%%%%%%%%%%%%%%%%%%%%%%%%%%%%%%%%%%%%%%%%%%%%%%%%%%%%%%%%%%%%%

\begin{theorem}{Euler--Lagrange Equation}

Every stationary infrastructure configuration satisfies the
Euler--Lagrange equation.

\end{theorem}

%%%%%%%%%%%%%%%%%%%%%%%%%%%%%%%%%%%%%%%%%%%%%%%%%%%%%%%%%%%%%%%

\section{Interpretation}

The Euler--Lagrange equation is remarkable because it is independent of
the engineering discipline.

Different choices of

\[
L
\]

produce different governing equations.

Hydraulic equilibrium,

electrical equilibrium,

minimal travel paths,

optimal corridor shapes,

centroidal Voronoi tessellations,

elastic bridge forms,

and diffusion processes all emerge from different infrastructure
densities.

Consequently, the diversity of engineering systems reflects not
different mathematics, but different choices of objective functional.

%%%%%%%%%%%%%%%%%%%%%%%%%%%%%%%%%%%%%%%%%%%%%%%%%%%%%%%%%%%%%%%

\section{Examples}

For

\[
L
=
\frac12
|\nabla u|^2,
\]

the Euler--Lagrange equation becomes

\[
\nabla^2u=0,
\]

recovering Laplace's equation.

For

\[
L
=
\sqrt{1+|\nabla u|^2},
\]

one obtains the minimal surface equation.

For

\[
L
=
\rho gQH,
\]

the stationary equations describe optimal hydraulic transport.

For

\[
L
=
RI^2,
\]

one recovers the minimum-dissipation principle governing electrical
networks.

%%%%%%%%%%%%%%%%%%%%%%%%%%%%%%%%%%%%%%%%%%%%%%%%%%%%%%%%%%%%%%%

\section{Toward Multiple Objectives}

Real infrastructure is never optimized for only one quantity.

Construction cost,

energy,

ecology,

maintainability,

safety,

robustness,

and accessibility compete simultaneously.

The next chapter therefore extends the single functional

\[
J
\]

to vector-valued objective functionals and develops the theory of
multi-objective optimization and Pareto-efficient infrastructure.

%%%%%%%%%%%%%%%%%%%%%%%%%%%%%%%%%%%%%%%%%%%%%%%%%%%%%%%%%%%%%%%%%%%%%%
\chapter{Multi-Objective Optimization}
\label{chap:multi-objective-optimization}

\section{Introduction}

The previous chapter formulated infrastructure design as the
minimization of a single functional

\[
J[u].
\]

Such formulations are mathematically elegant but rarely describe real
engineering projects.

Every infrastructure decision simultaneously affects construction cost,
operating cost, environmental impact, accessibility, reliability,
maintainability, resilience, and future expandability.

These objectives generally conflict.

Reducing construction cost may reduce redundancy.

Increasing redundancy may increase environmental impact.

Improving accessibility may require greater capital investment.

Consequently there rarely exists a configuration minimizing every
objective simultaneously.

The mathematical problem is therefore not single-objective
optimization but multi-objective optimization.

%%%%%%%%%%%%%%%%%%%%%%%%%%%%%%%%%%%%%%%%%%%%%%%%%%%%%%%%%%%%%%%

\section{Objective Vectors}

Suppose

\[
m
\]

performance measures are identified.

Each defines a functional

\[
J_i:\mathcal A\rightarrow\mathbb R.
\]

Collectively,

\[
J(u)
=
\begin{pmatrix}
J_1(u)\\
J_2(u)\\
\vdots\\
J_m(u)
\end{pmatrix}.
\]

Rather than minimizing a scalar,

the optimization now seeks desirable points within objective space.

%%%%%%%%%%%%%%%%%%%%%%%%%%%%%%%%%%%%%%%%%%%%%%%%%%%%%%%%%%%%%%%

\section{Examples}

Typical objectives include

\[
J_C
\]

construction cost,

\[
J_E
\]

operating energy,

\[
J_M
\]

maintenance burden,

\[
J_R
\]

network resilience,

\[
J_A
\]

average accessibility,

\[
J_P
\]

population coverage,

\[
J_L
\]

ecological disturbance,

and

\[
J_F
\]

future expansion capacity.

Each objective may itself be represented as an integral functional over
the corridor network.

%%%%%%%%%%%%%%%%%%%%%%%%%%%%%%%%%%%%%%%%%%%%%%%%%%%%%%%%%%%%%%%

\section{Partial Ordering}

The dominance order was introduced earlier in a narrower setting, over
the eight-parameter orthodrome configuration space
\(\Theta_{\mathrm{adm}}\). The present section restates it over the
general admissible function space \(\mathcal A\), so that every
functional introduced in the remainder of this Part is ordered by the
same relation rather than by a locally redefined one.

Unlike ordinary optimization,

objective vectors possess no natural total ordering.

Instead define

\[
u_1
\preceq
u_2
\]

whenever

\[
J_i(u_1)
\le
J_i(u_2)
\]

for every objective

\[
i,
\]

with strict inequality for at least one objective.

%%%%%%%%%%%%%%%%%%%%%%%%%%%%%%%%%%%%%%%%%%%%%%%%%%%%%%%%%%%%%%%

\begin{definition}{Pareto Dominance}

Configuration

\[
u_1
\]

dominates

\[
u_2
\]

if

\[
u_1
\preceq
u_2.
\]

\end{definition}

%%%%%%%%%%%%%%%%%%%%%%%%%%%%%%%%%%%%%%%%%%%%%%%%%%%%%%%%%%%%%%%

\begin{definition}{Pareto Optimality}

An admissible configuration

\[
u^\star
\]

is Pareto optimal if no admissible configuration dominates it.

The collection of all Pareto-optimal configurations is called the
Pareto frontier.

\end{definition}

%%%%%%%%%%%%%%%%%%%%%%%%%%%%%%%%%%%%%%%%%%%%%%%%%%%%%%%%%%%%%%%

\section{Weighted Scalarization}

One approach converts multiple objectives into a single functional.

Assign positive weights

\[
w_i>0.
\]

Define

\[
J_w
=
\sum_i
w_iJ_i.
\]

The optimization becomes

\[
u^\star
=
\operatorname*{arg\,min}
J_w.
\]

Different choices of

\[
w_i
\]

select different regions of the Pareto frontier.

%%%%%%%%%%%%%%%%%%%%%%%%%%%%%%%%%%%%%%%%%%%%%%%%%%%%%%%%%%%%%%%

\begin{theorem}

Every minimizer of a positively weighted objective is Pareto optimal.

\end{theorem}

\begin{proof}

Suppose

\[
u^\star
\]

were dominated by another admissible configuration

\[
v.
\]

Then

\[
J_i(v)
\le
J_i(u^\star)
\]

for every objective,

with strict inequality for at least one.

Since every weight is positive,

\[
\sum_i
w_iJ_i(v)
<
\sum_i
w_iJ_i(u^\star),
\]

contradicting the optimality of

\[
u^\star.
\]

Therefore

\[
u^\star
\]

must be Pareto optimal.

\end{proof}

%%%%%%%%%%%%%%%%%%%%%%%%%%%%%%%%%%%%%%%%%%%%%%%%%%%%%%%%%%%%%%%

\section{Infrastructure Trade-offs}

The geometry of the Pareto frontier reveals engineering compromises.

Regions of large curvature correspond to situations in which small
improvements in one objective require substantial sacrifice in another.

Flat regions indicate relatively inexpensive improvements.

Sharp corners identify fundamentally incompatible objectives.

The frontier therefore serves not merely as an optimization result but
as a geometric description of design difficulty.

%%%%%%%%%%%%%%%%%%%%%%%%%%%%%%%%%%%%%%%%%%%%%%%%%%%%%%%%%%%%%%%

\section{Capacity as an Objective}

Previous chapters introduced corridor utilization

\[
\phi_i
=
\frac{u_i}{K_i}.
\]

Future expansion is measured by

\[
1-\phi_i.
\]

Consequently,

future adaptability itself becomes an optimization objective.

Unlike conventional engineering practice,

unused capacity is not interpreted as inefficiency but as preserved
optionality.

Define

\[
J_F
=
-
\int_\Omega
(1-\phi)
\,dV.
\]

Minimizing

\[
J_F
\]

is therefore equivalent to maximizing long-term expansion potential.

%%%%%%%%%%%%%%%%%%%%%%%%%%%%%%%%%%%%%%%%%%%%%%%%%%%%%%%%%%%%%%%

\section{Robust Optimization}

Infrastructure parameters are never known exactly.

Demand fluctuates.

Climate changes.

Population migrates.

Technology evolves.

Suppose

\[
\theta
\]

represents uncertain parameters belonging to an uncertainty set

\[
\Theta.
\]

The robust optimization problem becomes

\[
\boxed{
\min_u
\;
\max_{\theta\in\Theta}
J(u,\theta).
}
\]

The resulting infrastructure minimizes worst-case performance rather
than nominal performance.

%%%%%%%%%%%%%%%%%%%%%%%%%%%%%%%%%%%%%%%%%%%%%%%%%%%%%%%%%%%%%%%

\section{Adaptive Infrastructure}

Rather than optimizing only for present conditions,

consider infrastructure capable of evolving.

Let

\[
u(t)
\]

denote the infrastructure state through time.

The optimization objective becomes

\[
J
=
\int_0^T
L(u,\dot u,t)\,dt.
\]

The solution is no longer a single configuration but an optimal
development trajectory.

This formulation naturally accommodates phased construction, population
growth, technological change, and future corridor expansion.

%%%%%%%%%%%%%%%%%%%%%%%%%%%%%%%%%%%%%%%%%%%%%%%%%%%%%%%%%%%%%%%

\section{Discussion}

Single-objective optimization produces one mathematically optimal
solution.

Multi-objective optimization produces an entire family of efficient
solutions.

Selecting among these solutions requires additional criteria that are
not purely mathematical.

In the following chapter we therefore introduce infrastructure
governance as a constrained decision process operating over the Pareto
frontier rather than over individual engineering variables.

%%%%%%%%%%%%%%%%%%%%%%%%%%%%%%%%%%%%%%%%%%%%%%%%%%%%%%%%%%%%%%%%%%%%%%
\chapter{Convexity, Duality, and Existence}

\section{Introduction}

The preceding chapters formulated infrastructure design as a
variational optimization problem. Before such formulations can be used
for engineering or planning, one must establish three fundamental
questions.

First, does an optimal infrastructure configuration exist?

Second, if a solution exists, is it unique?

Third, how do geometric and physical constraints modify the optimum?

These questions belong to the mathematical theory of optimization rather
than to any individual engineering discipline.

%%%%%%%%%%%%%%%%%%%%%%%%%%%%%%%%%%%%%%%%%%%%%%%%%%%%%%%%%%%%%%%

\section{Convex Sets}

The admissible configurations introduced previously form a subset

\[
\mathcal A
\subseteq
X,
\]

where

\[
X
\]

is an appropriate function space.

%%%%%%%%%%%%%%%%%%%%%%%%%%%%%%%%%%%%%%%%%%%%%%%%%%%%%%%%%%%%%%%

\begin{definition}{Convex Set}

A subset

\[
C\subseteq X
\]

is convex if

\[
\lambda u
+
(1-\lambda)v
\in C
\]

whenever

\[
u,v\in C
\]

and

\[
0\le\lambda\le1.
\]

\end{definition}

Convexity expresses the principle that mixtures of admissible
configurations remain admissible.

%%%%%%%%%%%%%%%%%%%%%%%%%%%%%%%%%%%%%%%%%%%%%%%%%%%%%%%%%%%%%%%

\section{Convex Functionals}

Infrastructure objectives frequently possess convex structure.

%%%%%%%%%%%%%%%%%%%%%%%%%%%%%%%%%%%%%%%%%%%%%%%%%%%%%%%%%%%%%%%

\begin{definition}{Convex Functional}

A functional

\[
J:\mathcal A\rightarrow\mathbb R
\]

is convex if

\[
J(\lambda u+(1-\lambda)v)
\le
\lambda J(u)
+
(1-\lambda)J(v)
\]

for every

\[
u,v\in\mathcal A.
\]

If strict inequality holds whenever

\[
u\neq v,
\]

the functional is strictly convex.

\end{definition}

%%%%%%%%%%%%%%%%%%%%%%%%%%%%%%%%%%%%%%%%%%%%%%%%%%%%%%%%%%%%%%%

Strict convexity has an important practical consequence.

%%%%%%%%%%%%%%%%%%%%%%%%%%%%%%%%%%%%%%%%%%%%%%%%%%%%%%%%%%%%%%%

\begin{theorem}

A strictly convex functional possesses at most one minimizer.

\end{theorem}

\begin{proof}

Assume

\[
u_1
\]

and

\[
u_2
\]

are distinct minimizers.

Then

\[
J(u_1)
=
J(u_2)
=
J_{\min}.
\]

Strict convexity implies

\[
J
\left(
\frac{u_1+u_2}{2}
\right)
<
\frac12
J(u_1)
+
\frac12
J(u_2)
=
J_{\min},
\]

contradicting minimality.

Therefore

\[
u_1=u_2.
\]

\end{proof}

%%%%%%%%%%%%%%%%%%%%%%%%%%%%%%%%%%%%%%%%%%%%%%%%%%%%%%%%%%%%%%%

\section{Coercivity}

Not every optimization problem possesses a solution.

A functional may decrease indefinitely without approaching a minimum.

To exclude such pathological behavior we introduce coercivity.

%%%%%%%%%%%%%%%%%%%%%%%%%%%%%%%%%%%%%%%%%%%%%%%%%%%%%%%%%%%%%%%

\begin{definition}{Coercive Functional}

A functional

\[
J
\]

is coercive if

\[
J(u)
\rightarrow
+\infty
\]

whenever

\[
\|u\|
\rightarrow
\infty.
\]

\end{definition}

Coercivity prevents minimizing sequences from escaping to infinity.

%%%%%%%%%%%%%%%%%%%%%%%%%%%%%%%%%%%%%%%%%%%%%%%%%%%%%%%%%%%%%%%

\section{Existence of Optimal Infrastructure}

The following theorem summarizes the basic existence principle.

%%%%%%%%%%%%%%%%%%%%%%%%%%%%%%%%%%%%%%%%%%%%%%%%%%%%%%%%%%%%%%%

\begin{theorem}

Suppose

\[
\mathcal A
\]

is nonempty and closed.

Suppose

\[
J
\]

is coercive and lower semicontinuous.

Then

\[
J
\]

possesses at least one minimizer.

\end{theorem}

\begin{proof}

Choose a minimizing sequence

\[
u_n.
\]

Coercivity implies

\[
\{u_n\}
\]

is bounded.

Hence a convergent subsequence exists.

Closedness of

\[
\mathcal A
\]

ensures the limit remains admissible.

Lower semicontinuity yields

\[
J(u^\star)
\le
\liminf
J(u_n),
\]

showing

\[
u^\star
\]

attains the minimum.

\end{proof}

%%%%%%%%%%%%%%%%%%%%%%%%%%%%%%%%%%%%%%%%%%%%%%%%%%%%%%%%%%%%%%%

\section{Equality Constraints}

Infrastructure systems are rarely unconstrained.

Mass conservation,

charge conservation,

budget limitations,

land availability,

and corridor width all impose equality constraints.

Let

\[
g_i(u)=0,
\qquad
i=1,\ldots,m.
\]

The optimization problem becomes

\[
\min_u
J(u)
\]

subject to

\[
g_i(u)=0.
\]

%%%%%%%%%%%%%%%%%%%%%%%%%%%%%%%%%%%%%%%%%%%%%%%%%%%%%%%%%%%%%%%

\section{Lagrange Multipliers}

Introduce multipliers

\[
\lambda_i.
\]

The Lagrangian is

\[
\boxed{
\mathcal L(u,\lambda)
=
J(u)
+
\sum_{i=1}^m
\lambda_i
g_i(u).
}
\]

Stationary points satisfy

\[
\frac{\partial\mathcal L}{\partial u}=0,
\]

together with

\[
g_i(u)=0.
\]

%%%%%%%%%%%%%%%%%%%%%%%%%%%%%%%%%%%%%%%%%%%%%%%%%%%%%%%%%%%%%%%

\begin{remark}

In engineering applications the Lagrange multipliers often possess
direct physical interpretations.

A multiplier associated with a budget constraint measures the marginal
benefit of additional funding.

A multiplier associated with corridor width measures the marginal value
of additional right-of-way.

A multiplier associated with electrical capacity measures the value of
additional transmission capability.

Thus dual variables are not merely mathematical artifacts but quantities
with operational significance.

\end{remark}

%%%%%%%%%%%%%%%%%%%%%%%%%%%%%%%%%%%%%%%%%%%%%%%%%%%%%%%%%%%%%%%

\section{Inequality Constraints}

Many engineering constraints possess the form

\[
h_j(u)\le0.
\]

Examples include

\[
\phi_i\le1,
\]

maximum allowable stress,

minimum safety distance,

maximum environmental impact,

and available construction budgets.

%%%%%%%%%%%%%%%%%%%%%%%%%%%%%%%%%%%%%%%%%%%%%%%%%%%%%%%%%%%%%%%

The optimization problem becomes

\[
\min
J(u)
\]

subject to

\[
g_i(u)=0,
\]

and

\[
h_j(u)\le0.
\]

%%%%%%%%%%%%%%%%%%%%%%%%%%%%%%%%%%%%%%%%%%%%%%%%%%%%%%%%%%%%%%%

\section{Karush--Kuhn--Tucker Conditions}

Necessary conditions for constrained optimality are summarized by the
Karush--Kuhn--Tucker equations.

Stationarity requires

\[
\nabla_u
\mathcal L
=
0.
\]

Primal feasibility requires

\[
g_i(u)=0,
\]

and

\[
h_j(u)\le0.
\]

Dual feasibility requires

\[
\mu_j\ge0.
\]

Finally,

\[
\boxed{
\mu_jh_j(u)=0.
}
\]

This complementary slackness condition implies that every inactive
constraint possesses zero multiplier.

%%%%%%%%%%%%%%%%%%%%%%%%%%%%%%%%%%%%%%%%%%%%%%%%%%%%%%%%%%%%%%%

\section{Sensitivity Analysis}

Suppose the optimization depends upon a parameter

\[
\alpha.
\]

Examples include

population,

construction cost,

climate,

or energy demand.

The optimal objective therefore becomes

\[
J^\star(\alpha).
\]

Its derivative

\[
\frac{dJ^\star}{d\alpha}
\]

measures the sensitivity of optimal infrastructure performance to
changes in environmental conditions.

Sensitivity analysis therefore identifies which assumptions most
strongly influence the resulting design.

%%%%%%%%%%%%%%%%%%%%%%%%%%%%%%%%%%%%%%%%%%%%%%%%%%%%%%%%%%%%%%%

\section{Discussion}

The preceding results establish that infrastructure optimization is not
merely an algorithmic exercise but a well-posed mathematical problem.

Existence theorems guarantee that optimal configurations are attainable.

Convexity identifies situations in which those configurations are
unique.

Duality assigns physical meaning to engineering constraints.

Sensitivity analysis quantifies robustness under changing assumptions.

These mathematical results provide the foundation upon which the
remaining chapters build. Subsequent discussions of governance,
expansion, resilience, and adaptive planning will assume that the
underlying optimization problems satisfy the conditions established in
this chapter.

%%%%%%%%%%%%%%%%%%%%%%%%%%%%%%%%%%%%%%%%%%%%%%%%%%%%%%%%%%%%%%%%%%%%%%
\part{Optimal Control and Adaptive Infrastructure}

%%%%%%%%%%%%%%%%%%%%%%%%%%%%%%%%%%%%%%%%%%%%%%%%%%%%%%%%%%%%%%%%%%%%%%
\chapter{Dynamic Infrastructure Systems}

\section{Introduction}

The preceding chapters treated infrastructure primarily as a static
optimization problem. Given a collection of admissible configurations,
the objective was to determine an optimal equilibrium satisfying the
appropriate physical and geometric constraints.

Real infrastructure, however, is never static.

Cities expand.

Populations migrate.

Technology improves.

Climate changes.

Demand fluctuates on timescales ranging from seconds to centuries.

Infrastructure must therefore be regarded not as a fixed object but as a
dynamical system whose state evolves continuously through time.

The mathematical problem changes accordingly. Instead of selecting an
optimal configuration, we seek an optimal trajectory.

%%%%%%%%%%%%%%%%%%%%%%%%%%%%%%%%%%%%%%%%%%%%%%%%%%%%%%%%%%%%%%%

\section{Infrastructure State Space}

Let

\[
X
\]

denote the infrastructure state space introduced in previous chapters.

Every point

\[
x\in X
\]

represents a complete description of the infrastructure at one instant.

A typical state vector may be written

\[
x=
\begin{pmatrix}
u\\
E\\
Q\\
I\\
\rho\\
m
\end{pmatrix},
\]

where

\[
u
\]

represents commodity distributions,

\[
E
\]

stored energy,

\[
Q
\]

hydraulic flows,

\[
I
\]

electrical currents,

\[
\rho
\]

traffic density,

and

\[
m
\]

freight density.

The exact composition of

\[
x
\]

depends upon the engineering model under consideration.

%%%%%%%%%%%%%%%%%%%%%%%%%%%%%%%%%%%%%%%%%%%%%%%%%%%%%%%%%%%%%%%

\section{Control Variables}

Not every aspect of the system is directly controllable.

Instead, planners manipulate a smaller collection of control variables.

Let

\[
u_c(t)
\]

denote the control vector.

Typical controls include

\[
\begin{aligned}
&
\text{pump outputs},\\
&
\text{generator dispatch},\\
&
\text{switching operations},\\
&
\text{traffic signalling},\\
&
\text{maintenance scheduling},\\
&
\text{construction investment},\\
&
\text{storage charging},
\end{aligned}
\]

and other operational decisions.

%%%%%%%%%%%%%%%%%%%%%%%%%%%%%%%%%%%%%%%%%%%%%%%%%%%%%%%%%%%%%%%

\section{State Equations}

The infrastructure evolves according to

\[
\boxed{
\frac{dx}{dt}
=
f(x,u_c,t).
}
\]

The vector field

\[
f
\]

combines

conservation laws,

network coupling,

physical constraints,

and external forcing.

%%%%%%%%%%%%%%%%%%%%%%%%%%%%%%%%%%%%%%%%%%%%%%%%%%%%%%%%%%%%%%%

\begin{definition}{Infrastructure Trajectory}

A trajectory is a continuously differentiable mapping

\[
x:[0,T]\rightarrow X
\]

satisfying

\[
\frac{dx}{dt}
=
f(x,u_c,t).
\]

\end{definition}

%%%%%%%%%%%%%%%%%%%%%%%%%%%%%%%%%%%%%%%%%%%%%%%%%%%%%%%%%%%%%%%

\section{Reachability}

Not every state can necessarily be reached from every other state.

Suppose

\[
x_0
\]

is the initial infrastructure state.

The reachable set is

\[
\mathcal R(x_0,T)
=
\{
x(T):
x(t)
\text{ satisfies the state equations}
\}.
\]

Reachability quantifies what infrastructure transformations are
physically achievable within a finite planning horizon.

%%%%%%%%%%%%%%%%%%%%%%%%%%%%%%%%%%%%%%%%%%%%%%%%%%%%%%%%%%%%%%%

\section{Trajectory Optimization}

The objective is no longer a static functional.

Instead,

optimization occurs over trajectories.

Let

\[
L(x,u_c,t)
\]

denote the instantaneous operating cost.

The total objective becomes

\[
\boxed{
J
=
\int_0^T
L(x,u_c,t)\,dt
+
\Phi(x(T)),
}
\]

where

\[
\Phi
\]

is a terminal cost.

Examples include

remaining construction obligations,

future maintenance liabilities,

or resilience targets at the end of the planning horizon.

%%%%%%%%%%%%%%%%%%%%%%%%%%%%%%%%%%%%%%%%%%%%%%%%%%%%%%%%%%%%%%%

\section{The Principle of Optimal Operation}

%%%%%%%%%%%%%%%%%%%%%%%%%%%%%%%%%%%%%%%%%%%%%%%%%%%%%%%%%%%%%%%

\begin{principle}{Optimal Infrastructure Operation}

Among all admissible control histories,

the realized infrastructure trajectory minimizes

\[
J.
\]

\end{principle}

%%%%%%%%%%%%%%%%%%%%%%%%%%%%%%%%%%%%%%%%%%%%%%%%%%%%%%%%%%%%%%%

\section{The Hamiltonian}

Introduce the adjoint variable

\[
\lambda(t).
\]

Define the Hamiltonian

\[
\boxed{
H(x,u_c,\lambda,t)
=
L(x,u_c,t)
+
\lambda^Tf(x,u_c,t).
}
\]

The Hamiltonian combines

instantaneous operating cost

with

future consequences encoded by the adjoint variables.

%%%%%%%%%%%%%%%%%%%%%%%%%%%%%%%%%%%%%%%%%%%%%%%%%%%%%%%%%%%%%%%

\section{Pontryagin's Maximum Principle}

Necessary conditions for optimality are obtained by considering small
variations in the control history.

%%%%%%%%%%%%%%%%%%%%%%%%%%%%%%%%%%%%%%%%%%%%%%%%%%%%%%%%%%%%%%%

\begin{theorem}{Pontryagin Maximum Principle}

If

\[
u_c^\star(t)
\]

is optimal,

then there exists an adjoint trajectory

\[
\lambda(t)
\]

satisfying

\[
\boxed{
\frac{dx}{dt}
=
\frac{\partial H}{\partial\lambda},
}
\]

\[
\boxed{
\frac{d\lambda}{dt}
=
-
\frac{\partial H}{\partial x},
}
\]

together with the stationarity condition

\[
\boxed{
\frac{\partial H}{\partial u_c}
=
0.
}
\]

\end{theorem}

%%%%%%%%%%%%%%%%%%%%%%%%%%%%%%%%%%%%%%%%%%%%%%%%%%%%%%%%%%%%%%%

\section{Interpretation}

The state equations propagate forward through time.

The adjoint equations propagate backward.

The optimal control lies at the intersection of these two dynamical
systems.

This forward--backward structure appears throughout engineering,
economics, robotics, and aerospace guidance.

Infrastructure planning is no exception.

%%%%%%%%%%%%%%%%%%%%%%%%%%%%%%%%%%%%%%%%%%%%%%%%%%%%%%%%%%%%%%%

\section{Feedback Control}

Open-loop optimization assumes future disturbances are known.

Real infrastructure instead operates under continual uncertainty.

Consequently,

controls should depend upon the observed system state.

Define

\[
u_c
=
K(x).
\]

Substituting into the state equation gives

\[
\frac{dx}{dt}
=
f(x,K(x),t).
\]

The resulting closed-loop system automatically responds to changing
conditions.

%%%%%%%%%%%%%%%%%%%%%%%%%%%%%%%%%%%%%%%%%%%%%%%%%%%%%%%%%%%%%%%

\section{Lyapunov Stability}

Infrastructure should recover from small disturbances.

Let

\[
x^\star
\]

be an equilibrium.

Suppose

\[
V(x)
\]

is a continuously differentiable function satisfying

\[
V(x^\star)=0,
\]

and

\[
V(x)>0
\]

for

\[
x\neq x^\star.
\]

%%%%%%%%%%%%%%%%%%%%%%%%%%%%%%%%%%%%%%%%%%%%%%%%%%%%%%%%%%%%%%%

\begin{definition}{Lyapunov Function}

Such a function

\[
V
\]

is called a Lyapunov function if

\[
\frac{dV}{dt}
\le0.
\]

\end{definition}

%%%%%%%%%%%%%%%%%%%%%%%%%%%%%%%%%%%%%%%%%%%%%%%%%%%%%%%%%%%%%%%

\begin{theorem}

If

\[
\frac{dV}{dt}<0
\]

throughout a neighborhood of

\[
x^\star,
\]

then

\[
x^\star
\]

is asymptotically stable.

\end{theorem}

%%%%%%%%%%%%%%%%%%%%%%%%%%%%%%%%%%%%%%%%%%%%%%%%%%%%%%%%%%%%%%%

\section{Infrastructure Resilience}

The previous theorem provides a natural mathematical definition of
resilience.

%%%%%%%%%%%%%%%%%%%%%%%%%%%%%%%%%%%%%%%%%%%%%%%%%%%%%%%%%%%%%%%

\begin{definition}{Infrastructure Resilience}

The resilience of an infrastructure system is its ability to return
toward an admissible operating region following perturbation.

Mathematically,

resilience corresponds to asymptotic stability of the controlled
dynamical system.

\end{definition}

This definition emphasizes recovery rather than resistance.

A resilient infrastructure need not avoid disturbance entirely.

Instead,

it restores acceptable operation after disturbance.

%%%%%%%%%%%%%%%%%%%%%%%%%%%%%%%%%%%%%%%%%%%%%%%%%%%%%%%%%%%%%%%

\section{Time Scales}

Infrastructure evolves simultaneously on multiple temporal scales.

Electrical switching occurs within milliseconds.

Traffic control evolves over seconds.

Reservoir operation changes over hours.

Maintenance planning spans months.

Urban expansion occurs over decades.

Planetary infrastructure develops over centuries.

Introduce characteristic times

\[
\tau_1
<
\tau_2
<
\cdots
<
\tau_n.
\]

Each subsystem contributes dynamics on one or more characteristic
timescales.

This multiscale structure motivates hierarchical control strategies in
which fast controllers stabilize local processes while slower planning
layers modify long-term infrastructure geometry.

%%%%%%%%%%%%%%%%%%%%%%%%%%%%%%%%%%%%%%%%%%%%%%%%%%%%%%%%%%%%%%%

\section{Toward Predictive Control}

The mathematical framework developed in this chapter transforms
infrastructure planning into a dynamical optimization problem.

Rather than repeatedly reacting to disturbances,

future controllers should anticipate them.

Chapter~\ref{chap:model-predictive-control} takes up Model Predictive Control, in which finite
horizon optimization is solved repeatedly using updated observations,
allowing the infrastructure to adapt continuously while remaining
consistent with the long-term geometric protocol established throughout
this monograph.

%%%%%%%%%%%%%%%%%%%%%%%%%%%%%%%%%%%%%%%%%%%%%%%%%%%%%%%%%%%%%%%%%%%%%%
\chapter{Model Predictive Control}
\label{chap:model-predictive-control}

\section{Introduction}

The previous chapter formulated infrastructure operation as an optimal
control problem over a finite planning horizon. Although this framework
provides an ideal mathematical description, practical infrastructure
systems rarely possess perfect knowledge of future conditions.

Electrical demand fluctuates.

Weather changes.

Traffic patterns evolve.

Construction schedules shift.

Equipment fails.

Population distributions change.

Consequently, a controller based upon a single optimization performed at
the beginning of the planning horizon rapidly becomes obsolete.

Model Predictive Control addresses this problem by repeatedly solving a
finite-horizon optimization using continually updated observations.

Rather than computing one permanent control law, the controller
constructs a sequence of locally optimal decisions whose repeated
application approximates long-term optimal operation.

%%%%%%%%%%%%%%%%%%%%%%%%%%%%%%%%%%%%%%%%%%%%%%%%%%%%%%%%%%%%%%%

\section{Prediction Horizon}

Let

\[
t
\]

denote the current time.

Choose a prediction horizon

\[
T_p>0.
\]

The controller predicts infrastructure evolution over the interval

\[
[t,t+T_p].
\]

The optimization problem is

\[
\min_{u_c}
\;
\int_t^{t+T_p}
L(x,u_c,\tau)
\,d\tau
+
\Phi(x(t+T_p)).
\]

Only the first portion of the optimal control trajectory is actually
implemented.

%%%%%%%%%%%%%%%%%%%%%%%%%%%%%%%%%%%%%%%%%%%%%%%%%%%%%%%%%%%%%%%

\section{Rolling Optimization}

After a short interval

\[
\Delta t,
\]

new observations become available.

The optimization is solved again using the updated infrastructure state.

Thus,

\[
x(t)
\rightarrow
x(t+\Delta t)
\rightarrow
x(t+2\Delta t)
\]

produces a sequence of optimization problems rather than one permanent
solution.

%%%%%%%%%%%%%%%%%%%%%%%%%%%%%%%%%%%%%%%%%%%%%%%%%%%%%%%%%%%%%%%

\begin{definition}{Rolling Horizon}

A rolling horizon is the sequence of prediction intervals

\[
[t_k,t_k+T_p],
\]

where

\[
t_k=k\Delta t.
\]

Each interval begins with the observed infrastructure state rather than
the previously predicted state.

\end{definition}

%%%%%%%%%%%%%%%%%%%%%%%%%%%%%%%%%%%%%%%%%%%%%%%%%%%%%%%%%%%%%%%

\section{State Estimation}

Optimization requires an estimate of the current system state.

Suppose

\[
y(t)
\]

represents measurements collected from sensors.

These may include

reservoir levels,

electrical frequencies,

traffic densities,

pipeline pressures,

weather observations,

and equipment status.

The estimated infrastructure state is

\[
\hat{x}(t).
\]

%%%%%%%%%%%%%%%%%%%%%%%%%%%%%%%%%%%%%%%%%%%%%%%%%%%%%%%%%%%%%%%

Prediction then begins from

\[
\hat{x}(t),
\]

rather than from an idealized mathematical state.

%%%%%%%%%%%%%%%%%%%%%%%%%%%%%%%%%%%%%%%%%%%%%%%%%%%%%%%%%%%%%%%

\section{Disturbances}

Infrastructure is continually influenced by external forcing.

Represent disturbances by

\[
d(t).
\]

The state equation becomes

\[
\boxed{
\frac{dx}{dt}
=
f(x,u_c,d,t).
}
\]

Examples include

unexpected demand,

storms,

equipment failures,

traffic incidents,

wildfire,

or seismic activity.

%%%%%%%%%%%%%%%%%%%%%%%%%%%%%%%%%%%%%%%%%%%%%%%%%%%%%%%%%%%%%%%

\section{Constraint Satisfaction}

Model Predictive Control naturally incorporates engineering constraints.

Typical examples include

\[
Q_{\max},
\]

maximum hydraulic flow,

\[
I_{\max},
\]

maximum transmission current,

\[
\rho_{\max},
\]

maximum traffic density,

and

\[
E_{\max},
\]

maximum storage capacity.

The optimization therefore becomes

\[
\min
J
\]

subject to

\[
\frac{dx}{dt}
=
f(x,u_c,d,t),
\]

together with all operational constraints.

%%%%%%%%%%%%%%%%%%%%%%%%%%%%%%%%%%%%%%%%%%%%%%%%%%%%%%%%%%%%%%%

\section{Terminal Constraints}

Long-term stability cannot be guaranteed by local optimization alone.

Introduce a terminal admissible region

\[
\mathcal X_f.
\]

Require

\[
x(t+T_p)
\in
\mathcal X_f.
\]

The optimization therefore preserves the ability to continue operating
after the prediction horizon has ended.

%%%%%%%%%%%%%%%%%%%%%%%%%%%%%%%%%%%%%%%%%%%%%%%%%%%%%%%%%%%%%%%

\begin{remark}

Terminal constraints prevent controllers from achieving artificially low
short-term costs by exhausting resources immediately before the end of
the prediction interval.

\end{remark}

%%%%%%%%%%%%%%%%%%%%%%%%%%%%%%%%%%%%%%%%%%%%%%%%%%%%%%%%%%%%%%%

\section{Recursive Feasibility}

A predictive controller should never produce a solution that prevents
future optimization.

%%%%%%%%%%%%%%%%%%%%%%%%%%%%%%%%%%%%%%%%%%%%%%%%%%%%%%%%%%%%%%%

\begin{definition}{Recursive Feasibility}

A Model Predictive Controller is recursively feasible if every feasible
optimization produces a successor state from which the next optimization
problem is also feasible.

\end{definition}

Recursive feasibility guarantees that the controller does not guide the
infrastructure into operational dead ends.

%%%%%%%%%%%%%%%%%%%%%%%%%%%%%%%%%%%%%%%%%%%%%%%%%%%%%%%%%%%%%%%

\section{Distributed Predictive Control}

Planetary infrastructure cannot realistically be controlled from a
single location.

Instead, each orthodromic node performs local optimization while
communicating with neighboring nodes.

Suppose

\[
x_i
\]

is the state of node

\[
i.
\]

Its local optimization becomes

\[
\min
J_i
\]

subject to

\[
\frac{dx_i}{dt}
=
f_i(x_i,u_i,x_{\mathcal N(i)}),
\]

where

\[
x_{\mathcal N(i)}
\]

denotes neighboring node states.

Distributed optimization greatly reduces computational complexity while
maintaining global coordination.

%%%%%%%%%%%%%%%%%%%%%%%%%%%%%%%%%%%%%%%%%%%%%%%%%%%%%%%%%%%%%%%

\section{Hierarchical Predictive Control}

The orthodromic architecture naturally supports hierarchical decision
making.

Local controllers regulate equipment operating on timescales of seconds.

Regional controllers coordinate corridor flows over hours or days.

Planetary controllers optimize long-term development over decades.

Formally, let

\[
u^{(1)},
u^{(2)},
\ldots,
u^{(k)}
\]

denote control variables operating on progressively longer temporal
scales.

Each layer solves

\[
\min
J_i
\]

subject to decisions inherited from higher planning levels.

%%%%%%%%%%%%%%%%%%%%%%%%%%%%%%%%%%%%%%%%%%%%%%%%%%%%%%%%%%%%%%%

\section{Protocol Preservation}

Unlike conventional infrastructure optimization, every controller in the
present framework operates under an additional geometric constraint.

The orthodromic protocol established in Part~I is treated as invariant.

Controllers may alter

flows,

routing,

capacity allocation,

maintenance schedules,

generation,

and construction sequencing.

They may not alter the underlying geometric protocol except through
explicit protocol revision procedures.

%%%%%%%%%%%%%%%%%%%%%%%%%%%%%%%%%%%%%%%%%%%%%%%%%%%%%%%%%%%%%%%

\begin{principle}{Protocol Preservation}

Operational optimization shall preserve the geometric compatibility of
the orthodromic network.

Control modifies infrastructure utilization rather than infrastructure
identity.

\end{principle}

%%%%%%%%%%%%%%%%%%%%%%%%%%%%%%%%%%%%%%%%%%%%%%%%%%%%%%%%%%%%%%%

\section{Infrastructure Inertia}

Infrastructure possesses enormous physical persistence.

Pipelines remain for decades.

Rail corridors may survive centuries.

Urban rights-of-way can remain recognizable for millennia.

Consequently, changing infrastructure geometry incurs substantial cost.

Introduce an inertia functional

\[
I[x]
=
\int_\Omega
\alpha(x)
\,
\|x-x_{\rm ref}\|^2
\,dV,
\]

where

\[
x_{\rm ref}
\]

denotes the existing infrastructure configuration.

Large values of

\[
\alpha
\]

represent highly persistent assets.

Small values represent rapidly reconfigurable systems.

%%%%%%%%%%%%%%%%%%%%%%%%%%%%%%%%%%%%%%%%%%%%%%%%%%%%%%%%%%%%%%%

The optimization objective becomes

\[
J_{\rm total}
=
J
+
I.
\]

Thus, controllers naturally prefer modifications that preserve existing
high-value infrastructure unless substantial long-term benefit justifies
their replacement.

%%%%%%%%%%%%%%%%%%%%%%%%%%%%%%%%%%%%%%%%%%%%%%%%%%%%%%%%%%%%%%%

\section{Adaptive Capacity}

Demand grows through time.

Instead of treating corridor capacity as fixed, let

\[
K(t)
\]

represent available capacity.

Expansion projects satisfy

\[
\frac{dK}{dt}
=
g(u_c,t),
\]

where

\[
g
\]

depends upon construction investment and resource availability.

Infrastructure therefore evolves through coupled operational and
construction dynamics.

%%%%%%%%%%%%%%%%%%%%%%%%%%%%%%%%%%%%%%%%%%%%%%%%%%%%%%%%%%%%%%%

\section{Discussion}

Model Predictive Control transforms infrastructure from a passive network
into an adaptive dynamical system.

Repeated optimization permits continual adjustment to changing operating
conditions while preserving long-term geometric compatibility.

The introduction of protocol preservation and infrastructure inertia
extends classical predictive control by distinguishing between transient
operational decisions and persistent geometric commitments.

This distinction is developed further in the next chapter, which introduces
multiscale planning, in which infrastructure evolves simultaneously on
operational, developmental, and civilizational timescales.

%%%%%%%%%%%%%%%%%%%%%%%%%%%%%%%%%%%%%%%%%%%%%%%%%%%%%%%%%%%%%%%%%%%%%%
\chapter{Multiscale Planning and Civilizational Timescales}

\section{Introduction}

Infrastructure exists simultaneously on many temporal scales.

Electrical protection systems respond within milliseconds.

Traffic control systems operate over seconds and minutes.

Reservoir management evolves over hours and days.

Maintenance scheduling spans months.

Construction projects require years.

Urban growth unfolds over decades.

Civilizations reshape continental infrastructure over centuries.

Conventional optimization frequently treats these processes
independently.

The orthodromic framework instead regards them as coupled components of
a single multiscale dynamical system.

%%%%%%%%%%%%%%%%%%%%%%%%%%%%%%%%%%%%%%%%%%%%%%%%%%%%%%%%%%%%%%%

\section{Temporal Hierarchy}

Let

\[
\tau_1
<
\tau_2
<
\cdots
<
\tau_N
\]

denote characteristic timescales.

Associate with each scale a state

\[
x^{(k)}(t),
\]

where

\[
k=1,\ldots,N.
\]

The complete infrastructure state becomes

\[
\boxed{
X
=
\left(
x^{(1)},
x^{(2)},
\ldots,
x^{(N)}
\right).
}
\]

%%%%%%%%%%%%%%%%%%%%%%%%%%%%%%%%%%%%%%%%%%%%%%%%%%%%%%%%%%%%%%%

\begin{definition}{Temporal Layer}

A temporal layer consists of all infrastructure variables whose
characteristic evolution occurs on approximately the same timescale.

\end{definition}

%%%%%%%%%%%%%%%%%%%%%%%%%%%%%%%%%%%%%%%%%%%%%%%%%%%%%%%%%%%%%%%

For example,

\[
x^{(1)}
\]

may describe electrical switching,

while

\[
x^{(N)}
\]

describes the evolution of planetary corridors.

%%%%%%%%%%%%%%%%%%%%%%%%%%%%%%%%%%%%%%%%%%%%%%%%%%%%%%%%%%%%%%%

\section{Nested Dynamics}

Fast variables respond almost instantaneously to slow variables.

Conversely,

slow variables evolve according to long-term averages of fast dynamics.

Accordingly,

the governing equations become

\[
\frac{dx^{(k)}}{dt}
=
f_k
\left(
x^{(1)},
\ldots,
x^{(N)}
\right).
\]

Each temporal layer therefore depends upon every other layer.

%%%%%%%%%%%%%%%%%%%%%%%%%%%%%%%%%%%%%%%%%%%%%%%%%%%%%%%%%%%%%%%

When

\[
\tau_i
\ll
\tau_j,
\]

the fast subsystem rapidly approaches local equilibrium while the slow
subsystem remains nearly constant.

This separation of timescales permits approximate decoupling.

%%%%%%%%%%%%%%%%%%%%%%%%%%%%%%%%%%%%%%%%%%%%%%%%%%%%%%%%%%%%%%%

\section{Slow Manifolds}

Suppose

\[
\tau_1
\ll
\tau_N.
\]

Fast dynamics satisfy

\[
\frac{dx^{(1)}}{dt}
=
f_1(x^{(1)},x^{(N)}).
\]

Holding

\[
x^{(N)}
\]

constant,

the fast subsystem converges toward

\[
x^{(1)}
=
h(x^{(N)}).
\]

The graph of

\[
h
\]

defines the slow manifold.

Consequently,

the long-term infrastructure dynamics reduce to

\[
\boxed{
\frac{dx^{(N)}}{dt}
=
f_N
\left(
h(x^{(N)}),
x^{(N)}
\right).
}
\]

%%%%%%%%%%%%%%%%%%%%%%%%%%%%%%%%%%%%%%%%%%%%%%%%%%%%%%%%%%%%%%%

\section{Infrastructure Memory}

Every construction project modifies future optimization.

Completed infrastructure constrains subsequent development.

Let

\[
M(t)
\]

denote accumulated infrastructure memory.

Examples include

existing corridors,

rights-of-way,

bridges,

tunnels,

reservoirs,

urban development,

and legal agreements.

The memory state evolves according to

\[
\boxed{
\frac{dM}{dt}
=
G(x,u_c).
}
\]

Unlike operational variables,

memory is only weakly reversible.

%%%%%%%%%%%%%%%%%%%%%%%%%%%%%%%%%%%%%%%%%%%%%%%%%%%%%%%%%%%%%%%

\begin{remark}

Infrastructure memory differs fundamentally from operational state.

Operational variables describe current utilization.

Memory describes persistent structural commitments.

\end{remark}

%%%%%%%%%%%%%%%%%%%%%%%%%%%%%%%%%%%%%%%%%%%%%%%%%%%%%%%%%%%%%%%

\section{Path Dependence}

Optimization therefore depends not only upon the present state but also
upon historical development.

Let

\[
\gamma
=
\{
x(t):
0\le t\le T
\}
\]

denote the development trajectory.

Rather than

\[
J(x(T)),
\]

consider

\[
\boxed{
J[\gamma].
}
\]

Infrastructure optimization therefore becomes a problem over histories
rather than instantaneous configurations.

%%%%%%%%%%%%%%%%%%%%%%%%%%%%%%%%%%%%%%%%%%%%%%%%%%%%%%%%%%%%%%%

\section{Irreversible Investment}

Many infrastructure decisions are effectively irreversible.

Excavated tunnels,

deep foundations,

major dams,

and protected corridors cannot be abandoned without substantial loss.

Introduce an irreversible investment variable

\[
I(t).
\]

Its evolution satisfies

\[
\frac{dI}{dt}
\ge
0.
\]

Hence

\[
I
\]

is monotone increasing.

%%%%%%%%%%%%%%%%%%%%%%%%%%%%%%%%%%%%%%%%%%%%%%%%%%%%%%%%%%%%%%%

\begin{definition}{Committed Capital}

Committed capital consists of infrastructure investment that cannot be
economically recovered through ordinary operation.

\end{definition}

%%%%%%%%%%%%%%%%%%%%%%%%%%%%%%%%%%%%%%%%%%%%%%%%%%%%%%%%%%%%%%%

This monotonicity distinguishes infrastructure development from ordinary
inventory management.

%%%%%%%%%%%%%%%%%%%%%%%%%%%%%%%%%%%%%%%%%%%%%%%%%%%%%%%%%%%%%%%

\section{Infrastructure Momentum}

Changes in infrastructure geometry require sustained effort.

Define geometric momentum

\[
P_G
=
\beta
\frac{dx_G}{dt},
\]

where

\[
x_G
\]

represents geometric configuration and

\[
\beta
\]

is a persistence coefficient.

Large

\[
\beta
\]

corresponds to massive civil works.

Small

\[
\beta
\]

corresponds to software-defined infrastructure such as traffic routing
or communication protocols.

%%%%%%%%%%%%%%%%%%%%%%%%%%%%%%%%%%%%%%%%%%%%%%%%%%%%%%%%%%%%%%%

The associated kinetic functional is

\[
K_G
=
\frac12
\beta
\left\|
\frac{dx_G}{dt}
\right\|^2.
\]

Rapid geometric reconfiguration therefore incurs increasing cost.

%%%%%%%%%%%%%%%%%%%%%%%%%%%%%%%%%%%%%%%%%%%%%%%%%%%%%%%%%%%%%%%

\section{Development Trajectories}

Infrastructure planning becomes a trajectory optimization problem.

Choose

\[
x_G(t)
\]

to minimize

\[
\boxed{
J
=
\int_0^T
\left(
L
+
K_G
\right)
dt.
}
\]

The first term measures operational performance.

The second penalizes unnecessary geometric restructuring.

%%%%%%%%%%%%%%%%%%%%%%%%%%%%%%%%%%%%%%%%%%%%%%%%%%%%%%%%%%%%%%%

\section{Generational Optimization}

Suppose planning occurs over generations.

Let

\[
T_g
\]

denote one generation.

Infrastructure decisions made during generation

\[
n
\]

become inherited constraints during generation

\[
n+1.
\]

Accordingly,

\[
x_{n+1}(0)
=
x_n(T_g).
\]

Optimization therefore proceeds recursively across generations.

%%%%%%%%%%%%%%%%%%%%%%%%%%%%%%%%%%%%%%%%%%%%%%%%%%%%%%%%%%%%%%%

\section{Protocol Evolution}

Although operational control preserves the orthodromic protocol,
exceptional circumstances may justify protocol revision.

Let

\[
\Pi(t)
\]

denote the governing protocol.

Ordinary operation satisfies

\[
\frac{d\Pi}{dt}
=
0.
\]

Protocol revision occurs only through discrete events,

\[
\Pi
\rightarrow
\Pi'.
\]

%%%%%%%%%%%%%%%%%%%%%%%%%%%%%%%%%%%%%%%%%%%%%%%%%%%%%%%%%%%%%%%

\begin{principle}{Protocol Stability}

The threshold for modifying the governing protocol shall be
substantially higher than the threshold for modifying operational
policies.

Protocols define long-term interoperability.

Policies define temporary implementation.

\end{principle}

%%%%%%%%%%%%%%%%%%%%%%%%%%%%%%%%%%%%%%%%%%%%%%%%%%%%%%%%%%%%%%%

\section{Civilizational Continuity}

The objective of infrastructure planning extends beyond immediate
performance.

Instead,

planning should maximize continuity of service across multiple
generations.

Define the continuity functional

\[
C(T)
=
\frac1T
\int_0^T
R(t)\,dt,
\]

where

\[
R(t)
\]

measures infrastructure reliability.

As

\[
T
\rightarrow
\infty,
\]

the asymptotic value

\[
C_\infty
=
\lim_{T\rightarrow\infty}
C(T)
\]

quantifies long-term civilizational continuity.

%%%%%%%%%%%%%%%%%%%%%%%%%%%%%%%%%%%%%%%%%%%%%%%%%%%%%%%%%%%%%%%

\section{Discussion}

The introduction of temporal hierarchy fundamentally changes the nature
of infrastructure optimization.

Instead of constructing an optimal static network,

the planner constructs an optimal sequence of adaptations extending over
multiple temporal scales.

Operational decisions,

maintenance,

construction,

and protocol evolution become components of one continuous dynamical
process.

Resilience and network recovery are taken up next, showing how infrastructure can preserve continuity in
the presence of failures, natural disasters, and technological change.

%%%%%%%%%%%%%%%%%%%%%%%%%%%%%%%%%%%%%%%%%%%%%%%%%%%%%%%%%%%%%%%%%%%%%%
\chapter{Resilience, Recovery, and Continuation}

\section{Introduction}

Infrastructure systems are designed not merely to perform efficiently
under ordinary conditions but to continue operating under extraordinary
conditions.

Floods,

earthquakes,

wildfire,

component failure,

cyber attacks,

human error,

and long-term environmental change all perturb the infrastructure away
from its nominal operating state.

Traditional engineering often measures resilience through redundancy,
reserve capacity, or restoration time.

While these quantities remain useful, they describe only particular
manifestations of a more general mathematical property.

Within the present framework, resilience is understood as the existence
of admissible recovery trajectories.

The central question therefore becomes not whether failure occurs, but
whether the system possesses a continuous path back toward acceptable
operation.

%%%%%%%%%%%%%%%%%%%%%%%%%%%%%%%%%%%%%%%%%%%%%%%%%%%%%%%%%%%%%%%

\section{Operational State Space}

Let

\[
X
\]

denote the infrastructure state space.

Every point

\[
x\in X
\]

represents one possible operational state.

Not every state is acceptable.

Instead define an admissible operating region

\[
\mathcal A
\subseteq
X.
\]

States outside

\[
\mathcal A
\]

represent degraded or failed operation.

%%%%%%%%%%%%%%%%%%%%%%%%%%%%%%%%%%%%%%%%%%%%%%%%%%%%%%%%%%%%%%%

\begin{definition}{Operational Admissibility}

A state

\[
x
\]

is operationally admissible whenever every required infrastructure
service satisfies its prescribed performance constraints.

Symbolically,

\[
x\in\mathcal A.
\]

\end{definition}

%%%%%%%%%%%%%%%%%%%%%%%%%%%%%%%%%%%%%%%%%%%%%%%%%%%%%%%%%%%%%%%

\section{Perturbations}

Failures are modeled as perturbations of the operational state.

Let

\[
\delta x
\]

denote a disturbance.

The perturbed state becomes

\[
x'
=
x+\delta x.
\]

Typical perturbations include

loss of transmission lines,

pipeline rupture,

bridge closure,

communication failure,

unexpected demand,

or environmental disruption.

%%%%%%%%%%%%%%%%%%%%%%%%%%%%%%%%%%%%%%%%%%%%%%%%%%%%%%%%%%%%%%%

\section{Recovery Trajectories}

Suppose

\[
x'
\notin
\mathcal A.
\]

Recovery consists of constructing a trajectory

\[
\gamma(t),
\qquad
0\le t\le T,
\]

satisfying

\[
\gamma(0)=x',
\]

and

\[
\gamma(T)\in\mathcal A.
\]

%%%%%%%%%%%%%%%%%%%%%%%%%%%%%%%%%%%%%%%%%%%%%%%%%%%%%%%%%%%%%%%

\begin{definition}{Recovery Path}

A recovery path is any admissible trajectory connecting a disturbed state
to an operationally admissible state.

\end{definition}

%%%%%%%%%%%%%%%%%%%%%%%%%%%%%%%%%%%%%%%%%%%%%%%%%%%%%%%%%%%%%%%

\section{Continuation Principle}

Infrastructure recovery should preserve as much existing structure as
possible.

Accordingly, among all recovery trajectories we prefer those minimizing
structural modification.

Let

\[
J_R[\gamma]
=
\int_0^T
C(\gamma,\dot\gamma)\,dt,
\]

where

\[
C
\]

measures recovery effort.

%%%%%%%%%%%%%%%%%%%%%%%%%%%%%%%%%%%%%%%%%%%%%%%%%%%%%%%%%%%%%%%

\begin{principle}{Continuation Principle}

Among all admissible recovery trajectories, preferred recovery minimizes
cumulative intervention while restoring acceptable operation.

\end{principle}

%%%%%%%%%%%%%%%%%%%%%%%%%%%%%%%%%%%%%%%%%%%%%%%%%%%%%%%%%%%%%%%

\section{Recovery Distance}

Introduce a functional measuring the minimum effort required to recover
from disturbance.

%%%%%%%%%%%%%%%%%%%%%%%%%%%%%%%%%%%%%%%%%%%%%%%%%%%%%%%%%%%%%%%

\begin{definition}{Recovery Distance}

The recovery distance from

\[
x
\]

to the admissible region is

\[
\boxed{
D_R(x)
=
\inf_{\gamma}
J_R[\gamma],
}
\]

where the infimum is taken over every recovery trajectory satisfying

\[
\gamma(0)=x,
\qquad
\gamma(T)\in\mathcal A.
\]

\end{definition}

Clearly,

\[
D_R(x)=0
\]

whenever

\[
x\in\mathcal A.
\]

%%%%%%%%%%%%%%%%%%%%%%%%%%%%%%%%%%%%%%%%%%%%%%%%%%%%%%%%%%%%%%%

\section{Resilience Basin}

Recovery is not always possible.

Define

\[
\mathcal B_R
=
\left\{
x:
D_R(x)<\infty
\right\}.
\]

%%%%%%%%%%%%%%%%%%%%%%%%%%%%%%%%%%%%%%%%%%%%%%%%%%%%%%%%%%%%%%%

\begin{definition}{Resilience Basin}

The resilience basin consists of every disturbed state from which an
admissible recovery trajectory exists.

\end{definition}

States outside

\[
\mathcal B_R
\]

require external reconstruction rather than ordinary recovery.

%%%%%%%%%%%%%%%%%%%%%%%%%%%%%%%%%%%%%%%%%%%%%%%%%%%%%%%%%%%%%%%

\section{Graceful Degradation}

Complete service preservation is not always achievable.

Instead,

systems may temporarily operate at reduced capacity while preserving
critical functionality.

Let

\[
S_i(x)
\]

denote the performance level of service

\[
i.
\]

Define normalized performance

\[
0
\le
\sigma_i
=
\frac{S_i}{S_i^{\max}}
\le
1.
\]

%%%%%%%%%%%%%%%%%%%%%%%%%%%%%%%%%%%%%%%%%%%%%%%%%%%%%%%%%%%%%%%

Rather than requiring

\[
\sigma_i=1,
\]

during disturbance,

we require

\[
\sigma_i
\ge
\sigma_i^{\min}.
\]

The infrastructure therefore degrades continuously rather than
catastrophically.

%%%%%%%%%%%%%%%%%%%%%%%%%%%%%%%%%%%%%%%%%%%%%%%%%%%%%%%%%%%%%%%

\section{Critical Services}

Not every service possesses equal importance.

Assign weights

\[
w_i>0.
\]

Define the weighted continuity measure

\[
\boxed{
R(x)
=
\sum_i
w_i
\sigma_i(x).
}
\]

Life-support systems,

emergency communications,

water supply,

and medical logistics may therefore receive substantially greater weight
than less critical services.

%%%%%%%%%%%%%%%%%%%%%%%%%%%%%%%%%%%%%%%%%%%%%%%%%%%%%%%%%%%%%%%

\section{Redundancy}

Suppose two independent subsystems provide the same service.

Let

\[
p_1,
\qquad
p_2
\]

denote their probabilities of failure.

If failures are statistically independent,

the probability that both fail simultaneously equals

\[
p_1p_2.
\]

Thus redundancy reduces failure probability multiplicatively rather than
additively.

%%%%%%%%%%%%%%%%%%%%%%%%%%%%%%%%%%%%%%%%%%%%%%%%%%%%%%%%%%%%%%%

\section{Network Connectivity}

Let

\[
G=(V,E)
\]

represent an infrastructure graph.

Failures remove vertices or edges.

Connectivity therefore changes over time.

%%%%%%%%%%%%%%%%%%%%%%%%%%%%%%%%%%%%%%%%%%%%%%%%%%%%%%%%%%%%%%%

\begin{definition}{Operational Connectivity}

A graph is operationally connected if every demand node remains connected
to at least one admissible supply node.

\end{definition}

%%%%%%%%%%%%%%%%%%%%%%%%%%%%%%%%%%%%%%%%%%%%%%%%%%%%%%%%%%%%%%%

Connectivity therefore depends upon functionality rather than purely
graph-theoretic adjacency.

%%%%%%%%%%%%%%%%%%%%%%%%%%%%%%%%%%%%%%%%%%%%%%%%%%%%%%%%%%%%%%%

\section{Minimum Recovery Sets}

Suppose

\[
F
\subseteq
E
\]

is the failed component set.

A repair set

\[
R
\subseteq
F
\]

is sufficient whenever restoring every component of

\[
R
\]

returns the infrastructure to admissible operation.

%%%%%%%%%%%%%%%%%%%%%%%%%%%%%%%%%%%%%%%%%%%%%%%%%%%%%%%%%%%%%%%

\begin{definition}{Minimal Recovery Set}

A recovery set is minimal if no proper subset is sufficient to restore
admissible operation.

\end{definition}

%%%%%%%%%%%%%%%%%%%%%%%%%%%%%%%%%%%%%%%%%%%%%%%%%%%%%%%%%%%%%%%

Identifying minimal recovery sets is generally a combinatorial
optimization problem.

%%%%%%%%%%%%%%%%%%%%%%%%%%%%%%%%%%%%%%%%%%%%%%%%%%%%%%%%%%%%%%%

\section{Recovery Scheduling}

Suppose repairs occur sequentially.

Let

\[
r_i
\]

denote repair action

\[
i.
\]

Its duration is

\[
\tau_i.
\]

Recovery planning minimizes

\[
\boxed{
J
=
\sum_i
w_i
\tau_i,
}
\]

subject to precedence,

resource,

and safety constraints.

%%%%%%%%%%%%%%%%%%%%%%%%%%%%%%%%%%%%%%%%%%%%%%%%%%%%%%%%%%%%%%%

\section{Resilience Functional}

Combining operational continuity and recovery effort yields the
resilience functional

\[
\boxed{
\mathcal R
=
\int_0^T
R(t)\,dt
-
\lambda
J_R.
}
\]

Large values correspond to systems that both preserve service during
disturbance and recover efficiently afterward.

%%%%%%%%%%%%%%%%%%%%%%%%%%%%%%%%%%%%%%%%%%%%%%%%%%%%%%%%%%%%%%%

\begin{theorem}

Suppose recovery effort is finite and the operational continuity
functional is bounded below.

Then

\[
\mathcal R
\]

is bounded below.

\end{theorem}

\begin{proof}

Since

\[
R(t)\ge0,
\]

its time integral is non-negative.

If

\[
J_R<\infty,
\]

then

\[
\lambda J_R
\]

is finite.

Therefore

\[
\mathcal R
>
-\infty.
\]

Thus the resilience functional is well-defined for every recoverable
trajectory.

\end{proof}

%%%%%%%%%%%%%%%%%%%%%%%%%%%%%%%%%%%%%%%%%%%%%%%%%%%%%%%%%%%%%%%

\section{Discussion}

Resilience is frequently described as the ability to withstand
disturbance.

The present formulation adopts a broader mathematical perspective.

A resilient infrastructure is one for which admissible recovery
trajectories exist and can be realized with finite effort while
maintaining acceptable continuity of essential services.

The emphasis therefore shifts from resistance to continuation.

Failures become temporary excursions through state space rather than
terminal conditions.

The next chapter develops this geometric viewpoint further by studying
the topology of infrastructure networks, introducing connectivity
invariants, cut sets, cycles, and multiscale graph structures that
govern both ordinary operation and recovery.

%%%%%%%%%%%%%%%%%%%%%%%%%%%%%%%%%%%%%%%%%%%%%%%%%%%%%%%%%%%%%%%%%%%%%%
\chapter{Topology and Geometry of Infrastructure Networks}

\section{Introduction}

The preceding chapters described infrastructure as a dynamical system
whose state evolves through time.

Those dynamics, however, are constrained by an underlying geometric
structure.

Pipelines follow corridors.

Transmission lines terminate at substations.

Roads intersect only at designated junctions.

Railways form branching and cyclic networks.

The topology of these networks determines not merely where commodities
can travel, but which failures are recoverable, which regions remain
connected, and which expansions preserve long-term continuity.

Accordingly, this chapter develops the mathematical topology of
infrastructure networks independently of any particular engineering
system.

%%%%%%%%%%%%%%%%%%%%%%%%%%%%%%%%%%%%%%%%%%%%%%%%%%%%%%%%%%%%%%%

\section{Infrastructure Graphs}

Let

\[
G=(V,E)
\]

denote an infrastructure graph.

The vertex set

\[
V
\]

represents operational nodes,

including

generation,

storage,

distribution,

intersections,

and terminal facilities.

The edge set

\[
E
\]

represents physical corridors,

including

roads,

railways,

pipelines,

transmission lines,

fiber-optic cables,

shipping lanes,

and tunnels.

%%%%%%%%%%%%%%%%%%%%%%%%%%%%%%%%%%%%%%%%%%%%%%%%%%%%%%%%%%%%%%%

Each edge possesses associated attributes,

including

length,

capacity,

construction cost,

maintenance cost,

travel time,

and reliability.

Formally,

\[
a:E\rightarrow\mathbb R^m.
\]

%%%%%%%%%%%%%%%%%%%%%%%%%%%%%%%%%%%%%%%%%%%%%%%%%%%%%%%%%%%%%%%

\section{Paths}

A path consists of an ordered sequence of adjacent edges.

%%%%%%%%%%%%%%%%%%%%%%%%%%%%%%%%%%%%%%%%%%%%%%%%%%%%%%%%%%%%%%%

\begin{definition}{Infrastructure Path}

A path from

\[
u
\]

to

\[
v
\]

is a sequence

\[
(v_0,v_1,\ldots,v_n)
\]

such that

\[
v_0=u,
\qquad
v_n=v,
\]

and

\[
(v_i,v_{i+1})\in E
\]

for every

\[
i.
\]

\end{definition}

%%%%%%%%%%%%%%%%%%%%%%%%%%%%%%%%%%%%%%%%%%%%%%%%%%%%%%%%%%%%%%%

A path corresponds to one feasible transportation route.

%%%%%%%%%%%%%%%%%%%%%%%%%%%%%%%%%%%%%%%%%%%%%%%%%%%%%%%%%%%%%%%

\section{Cycles}

%%%%%%%%%%%%%%%%%%%%%%%%%%%%%%%%%%%%%%%%%%%%%%%%%%%%%%%%%%%%%%%

\begin{definition}{Cycle}

A cycle is a path whose initial and terminal vertices coincide while no
other vertex is repeated.

\end{definition}

Cycles provide alternative routing during failures.

Networks without cycles possess no intrinsic redundancy.

%%%%%%%%%%%%%%%%%%%%%%%%%%%%%%%%%%%%%%%%%%%%%%%%%%%%%%%%%%%%%%%

\section{Trees}

%%%%%%%%%%%%%%%%%%%%%%%%%%%%%%%%%%%%%%%%%%%%%%%%%%%%%%%%%%%%%%%

\begin{definition}{Infrastructure Tree}

An infrastructure tree is a connected graph containing no cycles.

\end{definition}

Trees minimize construction cost but possess minimal redundancy.

Failure of one edge disconnects the network.

%%%%%%%%%%%%%%%%%%%%%%%%%%%%%%%%%%%%%%%%%%%%%%%%%%%%%%%%%%%%%%%

\section{Connectivity}

%%%%%%%%%%%%%%%%%%%%%%%%%%%%%%%%%%%%%%%%%%%%%%%%%%%%%%%%%%%%%%%

\begin{definition}{Connected Network}

An infrastructure graph is connected if every pair of vertices is joined
by at least one path.

\end{definition}

Connectivity is the minimum topological requirement for integrated
operation.

%%%%%%%%%%%%%%%%%%%%%%%%%%%%%%%%%%%%%%%%%%%%%%%%%%%%%%%%%%%%%%%

\section{Connected Components}

Suppose failures partition the graph.

Each maximal connected subgraph forms a connected component.

Let

\[
c(G)
\]

denote the number of connected components.

A resilient infrastructure seeks to minimize growth of

\[
c(G)
\]

under disturbance.

%%%%%%%%%%%%%%%%%%%%%%%%%%%%%%%%%%%%%%%%%%%%%%%%%%%%%%%%%%%%%%%

\section{Edge Cuts}

%%%%%%%%%%%%%%%%%%%%%%%%%%%%%%%%%%%%%%%%%%%%%%%%%%%%%%%%%%%%%%%

\begin{definition}{Edge Cut}

An edge cut is a subset

\[
F\subseteq E
\]

whose removal disconnects the graph.

\end{definition}

Edge cuts identify vulnerable infrastructure corridors.

%%%%%%%%%%%%%%%%%%%%%%%%%%%%%%%%%%%%%%%%%%%%%%%%%%%%%%%%%%%%%%%

\section{Vertex Cuts}

%%%%%%%%%%%%%%%%%%%%%%%%%%%%%%%%%%%%%%%%%%%%%%%%%%%%%%%%%%%%%%%

\begin{definition}{Vertex Cut}

A vertex cut is a subset

\[
S\subseteq V
\]

whose removal disconnects the graph.

\end{definition}

Major substations,

bridges,

harbors,

or intermodal terminals frequently correspond to high-value vertex cuts.

%%%%%%%%%%%%%%%%%%%%%%%%%%%%%%%%%%%%%%%%%%%%%%%%%%%%%%%%%%%%%%%

\section{Bridge Edges}

%%%%%%%%%%%%%%%%%%%%%%%%%%%%%%%%%%%%%%%%%%%%%%%%%%%%%%%%%%%%%%%

\begin{definition}{Bridge}

An edge is a bridge if its removal increases the number of connected
components.

\end{definition}

Bridges represent single points of structural failure.

%%%%%%%%%%%%%%%%%%%%%%%%%%%%%%%%%%%%%%%%%%%%%%%%%%%%%%%%%%%%%%%

\begin{theorem}

Every bridge belongs to no cycle.

\end{theorem}

\begin{proof}

Suppose a bridge belonged to a cycle.

Removing the bridge would leave the remaining portion of the cycle as an
alternative path.

The graph would therefore remain connected,

contradicting the definition of a bridge.

Hence every bridge lies outside every cycle.

\end{proof}

%%%%%%%%%%%%%%%%%%%%%%%%%%%%%%%%%%%%%%%%%%%%%%%%%%%%%%%%%%%%%%%

\section{Cycle Space}

The collection of all cycles forms a vector space over

\[
\mathbb Z_2.
\]

Its dimension equals

\[
\boxed{
|E|-|V|+c(G).
}
\]

This quantity measures independent routing alternatives within the
network.

Larger cycle spaces generally imply greater routing flexibility.

%%%%%%%%%%%%%%%%%%%%%%%%%%%%%%%%%%%%%%%%%%%%%%%%%%%%%%%%%%%%%%%

\section{Incidence Matrix}

Number the vertices

\[
v_1,\ldots,v_n,
\]

and edges

\[
e_1,\ldots,e_m.
\]

The incidence matrix

\[
B
\]

is defined by

\[
B_{ij}
=
\begin{cases}
1,&e_j\text{ enters }v_i,\\
-1,&e_j\text{ leaves }v_i,\\
0,&\text{otherwise.}
\end{cases}
\]

Conservation of flow becomes

\[
Bf=s,
\]

where

\[
f
\]

is the edge-flow vector and

\[
s
\]

contains external sources and sinks.

%%%%%%%%%%%%%%%%%%%%%%%%%%%%%%%%%%%%%%%%%%%%%%%%%%%%%%%%%%%%%%%

\section{Graph Laplacian}

Define

\[
L
=
BB^T.
\]

The graph Laplacian governs diffusion,

electrical potential,

heat transport,

consensus dynamics,

and many optimization problems.

%%%%%%%%%%%%%%%%%%%%%%%%%%%%%%%%%%%%%%%%%%%%%%%%%%%%%%%%%%%%%%%

\begin{theorem}

The multiplicity of the eigenvalue

\[
0
\]

of

\[
L
\]

equals the number of connected components of

\[
G.
\]

\end{theorem}

This theorem connects topology directly with spectral properties.

%%%%%%%%%%%%%%%%%%%%%%%%%%%%%%%%%%%%%%%%%%%%%%%%%%%%%%%%%%%%%%%

\section{Algebraic Connectivity}

Let

\[
\lambda_2
\]

denote the second-smallest eigenvalue of the Laplacian.

%%%%%%%%%%%%%%%%%%%%%%%%%%%%%%%%%%%%%%%%%%%%%%%%%%%%%%%%%%%%%%%

\begin{definition}{Algebraic Connectivity}

The quantity

\[
\lambda_2
\]

is called the algebraic connectivity of the network.

\end{definition}

Large values of

\[
\lambda_2
\]

indicate highly connected infrastructure with many alternative routing
paths.

Small values indicate structural vulnerability.

%%%%%%%%%%%%%%%%%%%%%%%%%%%%%%%%%%%%%%%%%%%%%%%%%%%%%%%%%%%%%%%

\section{Topological Resilience}

Topology provides a purely structural measure of resilience.

Define

\[
T(G)
=
\alpha\lambda_2
+
\beta\dim(\mathcal C)
-
\gamma b,
\]

where

\[
\dim(\mathcal C)
\]

is the cycle-space dimension,

\[
b
\]

is the number of bridges,

and

\[
\alpha,\beta,\gamma>0.
\]

Large values of

\[
T(G)
\]

correspond to structurally robust infrastructure.

%%%%%%%%%%%%%%%%%%%%%%%%%%%%%%%%%%%%%%%%%%%%%%%%%%%%%%%%%%%%%%%

\section{Hierarchical Networks}

Real infrastructure is neither purely planar nor purely homogeneous.

Instead,

regional subnetworks connect to continental corridors,

which themselves connect through global hubs.

Represent this hierarchy by a sequence

\[
G_1
\subseteq
G_2
\subseteq
\cdots
\subseteq
G_n.
\]

Each level inherits connectivity from lower layers while introducing
longer-range coordination.

%%%%%%%%%%%%%%%%%%%%%%%%%%%%%%%%%%%%%%%%%%%%%%%%%%%%%%%%%%%%%%%

\section{Discussion}

Topology determines what infrastructure is capable of doing before any
flow equations are solved.

Conservation laws govern movement through a network.

Topology governs whether such movement is possible at all.

The graph-theoretic quantities introduced in this chapter—connectivity,
cuts, cycles, bridges, Laplacian spectra, and hierarchical
decompositions—provide structural invariants that remain meaningful
across transportation, electrical, hydraulic, communication, and
logistics systems alike.

In the following chapter these topological invariants are combined with
continuous geometry to study embedded infrastructure networks on curved
surfaces, culminating in a differential-geometric treatment of
large-scale planetary corridor systems.

%%%%%%%%%%%%%%%%%%%%%%%%%%%%%%%%%%%%%%%%%%%%%%%%%%%%%%%%%%%%%%%%%%%%%%
\chapter{Graphs Embedded on Manifolds}

\section{Introduction}

Infrastructure networks are simultaneously discrete and continuous.

The discrete structure consists of nodes and corridors represented by a
graph.

The continuous structure consists of the physical surface upon which the
network is embedded.

Neither description alone is sufficient.

Graph theory captures connectivity but ignores geography.

Differential geometry describes the underlying surface but ignores the
combinatorial organization of infrastructure.

The objective of this chapter is to unify these viewpoints by studying
graphs embedded on smooth manifolds.

%%%%%%%%%%%%%%%%%%%%%%%%%%%%%%%%%%%%%%%%%%%%%%%%%%%%%%%%%%%%%%%

\section{The Underlying Manifold}

Let

\[
(M,g)
\]

be a connected Riemannian manifold.

Throughout this work,

\[
M
\]

represents the physical domain occupied by the infrastructure.

Examples include

continental surfaces,

river basins,

coastal regions,

or the Earth's surface itself.

The metric tensor

\[
g
\]

determines distances,

angles,

areas,

and geodesics.

%%%%%%%%%%%%%%%%%%%%%%%%%%%%%%%%%%%%%%%%%%%%%%%%%%%%%%%%%%%%%%%

\section{Embedded Graphs}

Let

\[
G=(V,E)
\]

be an infrastructure graph.

An embedding is a mapping

\[
\iota:G\rightarrow M
\]

satisfying

\[
\iota(V)\subset M
\]

and

\[
\iota(E)
\]

is a continuous curve joining its endpoints.

Thus every infrastructure corridor becomes a curve on the manifold.

%%%%%%%%%%%%%%%%%%%%%%%%%%%%%%%%%%%%%%%%%%%%%%%%%%%%%%%%%%%%%%%

\begin{definition}{Embedded Infrastructure Network}

An embedded infrastructure network consists of a graph

\[
G
\]

together with an embedding

\[
\iota:G\rightarrow M.
\]

\end{definition}

%%%%%%%%%%%%%%%%%%%%%%%%%%%%%%%%%%%%%%%%%%%%%%%%%%%%%%%%%%%%%%%

\section{Edge Length}

Every embedded corridor inherits a geometric length.

If

\[
\gamma:[0,1]\rightarrow M
\]

parameterizes an edge,

its length is

\[
\boxed{
L(\gamma)
=
\int_0^1
\sqrt{
g(\dot\gamma,\dot\gamma)
}
\,dt.
}
\]

Construction cost,

travel time,

and maintenance often depend upon this quantity.

%%%%%%%%%%%%%%%%%%%%%%%%%%%%%%%%%%%%%%%%%%%%%%%%%%%%%%%%%%%%%%%

\section{Geodesic Corridors}

Among all curves joining two nodes,

geodesics minimize local distance.

%%%%%%%%%%%%%%%%%%%%%%%%%%%%%%%%%%%%%%%%%%%%%%%%%%%%%%%%%%%%%%%

\begin{definition}{Geodesic Corridor}

A corridor is geodesic if its embedding satisfies the geodesic equation

\[
\nabla_{\dot\gamma}\dot\gamma=0.
\]

\end{definition}

On a sphere,

geodesic corridors coincide with great-circle segments.

%%%%%%%%%%%%%%%%%%%%%%%%%%%%%%%%%%%%%%%%%%%%%%%%%%%%%%%%%%%%%%%

\section{Terrain Potential}

Shortest distance is rarely optimal.

Elevation,

geology,

population,

environmental protection,

and political boundaries all influence routing.

Introduce a scalar terrain potential

\[
\Phi:M\rightarrow\mathbb R.
\]

Define the weighted corridor functional

\[
J(\gamma)
=
\int_\gamma
\Phi\,ds.
\]

Optimal routing therefore minimizes weighted rather than purely
geometric distance.

%%%%%%%%%%%%%%%%%%%%%%%%%%%%%%%%%%%%%%%%%%%%%%%%%%%%%%%%%%%%%%%

\section{Weighted Geodesics}

The Euler--Lagrange equations applied to

\[
J(\gamma)
\]

produce weighted geodesics.

These curves balance geometric efficiency against environmental and
economic cost.

Classical shortest paths arise when

\[
\Phi\equiv1.
\]

%%%%%%%%%%%%%%%%%%%%%%%%%%%%%%%%%%%%%%%%%%%%%%%%%%%%%%%%%%%%%%%

\section{Curvature and Routing}

The geometry of the underlying manifold influences optimal network
structure.

Positive curvature causes nearby geodesics to converge.

Negative curvature causes them to diverge.

Consequently,

terrain curvature influences corridor spacing,

junction density,

and redundancy.

%%%%%%%%%%%%%%%%%%%%%%%%%%%%%%%%%%%%%%%%%%%%%%%%%%%%%%%%%%%%%%%

\section{Discrete Curvature}

The graph itself possesses intrinsic geometry.

One useful measure is based upon neighboring connectivity.

Let

\[
\kappa(e)
\]

denote the curvature assigned to edge

\[
e.
\]

Positive values indicate locally redundant regions.

Negative values indicate bottlenecks.

Near-zero values correspond to approximately flat network geometry.

%%%%%%%%%%%%%%%%%%%%%%%%%%%%%%%%%%%%%%%%%%%%%%%%%%%%%%%%%%%%%%%

Although several mathematical definitions exist,

the essential idea is identical:

curvature measures how strongly neighboring routes reinforce one another.

%%%%%%%%%%%%%%%%%%%%%%%%%%%%%%%%%%%%%%%%%%%%%%%%%%%%%%%%%%%%%%%

\section{Continuous--Discrete Coupling}

The embedded network therefore possesses two distinct notions of
curvature.

The manifold contributes geometric curvature,

while the graph contributes combinatorial curvature.

Together they determine infrastructure behavior.

A mountain ridge may force high manifold curvature while the graph
remains richly interconnected.

Conversely,

a perfectly flat region may contain severe graph bottlenecks due to
limited bridges or legal constraints.

%%%%%%%%%%%%%%%%%%%%%%%%%%%%%%%%%%%%%%%%%%%%%%%%%%%%%%%%%%%%%%%

\section{Energy of an Embedded Network}

The total geometric energy is

\[
\boxed{
E(G)
=
\sum_{e\in E}
\int_e
\Phi
\,ds.
}
\]

This quantity combines geometry,

terrain,

and network topology into a single variational functional.

%%%%%%%%%%%%%%%%%%%%%%%%%%%%%%%%%%%%%%%%%%%%%%%%%%%%%%%%%%%%%%%

\section{Embedded Flow Equations}

Let

\[
f_e
\]

denote the flow along edge

\[
e.
\]

Mass conservation remains

\[
Bf=s,
\]

while geometric routing minimizes

\[
E(G).
\]

Infrastructure optimization therefore consists of simultaneously solving

the discrete conservation equations,

continuous geometric optimization,

and operational constraints.

%%%%%%%%%%%%%%%%%%%%%%%%%%%%%%%%%%%%%%%%%%%%%%%%%%%%%%%%%%%%%%%

\section{Scale Separation}

At local scales,

individual corridors follow terrain in great detail.

At continental scales,

only coarse geometric structure remains important.

Accordingly,

introduce a family of embedded graphs

\[
G_\epsilon,
\]

indexed by spatial resolution

\[
\epsilon.
\]

As

\[
\epsilon
\rightarrow0,
\]

the graph approaches detailed engineering geometry.

As

\[
\epsilon
\rightarrow\infty,
\]

it converges toward a coarse strategic network.

Thus infrastructure admits simultaneous microscopic and macroscopic
descriptions.

%%%%%%%%%%%%%%%%%%%%%%%%%%%%%%%%%%%%%%%%%%%%%%%%%%%%%%%%%%%%%%%

\section{The Embedded Network Principle}

\begin{principle}{Embedded Network Principle}

Infrastructure optimization is the simultaneous optimization of

\[
\begin{aligned}
&
\text{topology},\\
&
\text{geometry},\\
&
\text{physical transport},\\
&
\text{operational control}.
\end{aligned}
\]

No single aspect determines optimal infrastructure independently of the
others.

\end{principle}

%%%%%%%%%%%%%%%%%%%%%%%%%%%%%%%%%%%%%%%%%%%%%%%%%%%%%%%%%%%%%%%

\section{Discussion}

The mathematical framework developed in this chapter unifies graph
theory and differential geometry.

Infrastructure is represented as a graph embedded within a continuous
manifold, allowing discrete connectivity and continuous spatial
constraints to be studied within one formalism.

Subsequent chapters will extend this framework by introducing stochastic
geometry and uncertainty, recognizing that terrain, demand, failures,
and future development cannot be known exactly. The resulting theory
treats infrastructure not as a fixed geometric object, but as a family
of probabilistic embedded networks evolving through space and time.

%%%%%%%%%%%%%%%%%%%%%%%%%%%%%%%%%%%%%%%%%%%%%%%%%%%%%%%%%%%%%%%%%%%%%%
\part{Stochastic Infrastructure Theory}

%%%%%%%%%%%%%%%%%%%%%%%%%%%%%%%%%%%%%%%%%%%%%%%%%%%%%%%%%%%%%%%%%%%%%%
\chapter{Uncertainty, Random Fields, and Probabilistic Infrastructure}

\section{Introduction}

The mathematical framework developed thus far has treated infrastructure
as a deterministic system.

Given the network geometry,

physical laws,

boundary conditions,

and operating policy,

the future evolution of the system is assumed to be uniquely determined.

Real infrastructure operates under substantially greater uncertainty.

Electrical demand fluctuates.

Transportation patterns evolve.

Weather changes.

Construction costs vary.

Component failures occur unpredictably.

Political priorities shift.

Future technologies alter system capabilities.

Consequently, the infrastructure state should be regarded not as a
single trajectory but as a probability distribution over possible
trajectories.

This chapter develops the mathematical foundations of stochastic
infrastructure systems.

%%%%%%%%%%%%%%%%%%%%%%%%%%%%%%%%%%%%%%%%%%%%%%%%%%%%%%%%%%%%%%%

\section{Probability Space}

Let

\[
(\Omega,\mathcal F,\mathbb P)
\]

be a probability space.

Here,

\[
\Omega
\]

denotes the collection of all possible realizations of future
conditions,

\[
\mathcal F
\]

is the associated sigma-algebra,

and

\[
\mathbb P
\]

assigns probabilities to measurable events.

Every infrastructure quantity now becomes a random variable defined on
this space.

%%%%%%%%%%%%%%%%%%%%%%%%%%%%%%%%%%%%%%%%%%%%%%%%%%%%%%%%%%%%%%%

\section{Random Infrastructure States}

The infrastructure state becomes

\[
X(t,\omega),
\]

where

\[
\omega\in\Omega.
\]

Instead of predicting one future state,

we seek the distribution

\[
\mathbb P(X(t)\in A)
\]

for every measurable region

\[
A\subseteq X.
\]

%%%%%%%%%%%%%%%%%%%%%%%%%%%%%%%%%%%%%%%%%%%%%%%%%%%%%%%%%%%%%%%

\begin{definition}{Random Infrastructure Process}

A random infrastructure process is a measurable mapping

\[
X:[0,T]\times\Omega\rightarrow\mathcal X.
\]

\end{definition}

%%%%%%%%%%%%%%%%%%%%%%%%%%%%%%%%%%%%%%%%%%%%%%%%%%%%%%%%%%%%%%%

\section{Random Fields}

Infrastructure parameters frequently vary continuously over space.

Examples include

population density,

wind speed,

solar irradiance,

soil stability,

construction cost,

and flood probability.

Represent these by a random field

\[
\Phi(x,\omega),
\]

defined on the underlying manifold.

%%%%%%%%%%%%%%%%%%%%%%%%%%%%%%%%%%%%%%%%%%%%%%%%%%%%%%%%%%%%%%%

Its expectation is

\[
\mathbb E[\Phi(x)].
\]

Its covariance is

\[
C(x,y)
=
\mathbb E
\left[
(\Phi(x)-\mu(x))
(\Phi(y)-\mu(y))
\right].
\]

%%%%%%%%%%%%%%%%%%%%%%%%%%%%%%%%%%%%%%%%%%%%%%%%%%%%%%%%%%%%%%%

\section{Random Demand}

Suppose demand evolves as

\[
D(t,\omega).
\]

Its expected value is

\[
\bar D(t)
=
\mathbb E[D(t)].
\]

Its variance is

\[
\operatorname{Var}(D)
=
\mathbb E[(D-\bar D)^2].
\]

Infrastructure planning must therefore satisfy both expected demand and
demand variability.

%%%%%%%%%%%%%%%%%%%%%%%%%%%%%%%%%%%%%%%%%%%%%%%%%%%%%%%%%%%%%%%

\section{Failure Processes}

Component failures occur randomly.

Let

\[
N(t)
\]

denote the cumulative number of failures.

A common approximation models

\[
N(t)
\]

as a Poisson process satisfying

\[
\mathbb P(N(t+h)-N(t)=1)
=
\lambda h
+
o(h).
\]

Here,

\[
\lambda
\]

is the failure intensity.

%%%%%%%%%%%%%%%%%%%%%%%%%%%%%%%%%%%%%%%%%%%%%%%%%%%%%%%%%%%%%%%

\section{Reliability}

Suppose

\[
T
\]

denotes component lifetime.

The reliability function is

\[
R(t)
=
\mathbb P(T>t).
\]

Its derivative satisfies

\[
f(t)
=
-
\frac{dR}{dt},
\]

where

\[
f
\]

is the lifetime density.

%%%%%%%%%%%%%%%%%%%%%%%%%%%%%%%%%%%%%%%%%%%%%%%%%%%%%%%%%%%%%%%

The hazard rate is

\[
\boxed{
h(t)
=
\frac{f(t)}{R(t)}.
}
\]

Large hazard rates indicate rapidly increasing failure probability.

%%%%%%%%%%%%%%%%%%%%%%%%%%%%%%%%%%%%%%%%%%%%%%%%%%%%%%%%%%%%%%%

\section{Expected Cost}

Optimization now minimizes expected rather than deterministic cost.

Define

\[
\boxed{
J
=
\mathbb E
\left[
\int_0^T
L(X,u,t)
dt
+
\Phi(X(T))
\right].
}
\]

The expectation is taken over every admissible realization of future
uncertainty.

%%%%%%%%%%%%%%%%%%%%%%%%%%%%%%%%%%%%%%%%%%%%%%%%%%%%%%%%%%%%%%%

\section{Chance Constraints}

Deterministic constraints are often unnecessarily conservative.

Instead,

permit rare violations.

Suppose

\[
g(X)\le0
\]

is required.

Replace this by

\[
\boxed{
\mathbb P(g(X)\le0)
\ge
1-\varepsilon.
}
\]

The parameter

\[
\varepsilon
\]

represents acceptable risk.

%%%%%%%%%%%%%%%%%%%%%%%%%%%%%%%%%%%%%%%%%%%%%%%%%%%%%%%%%%%%%%%

\section{Scenario Optimization}

Analytical probability distributions are rarely available.

Instead,

construct representative scenarios

\[
\omega_1,\ldots,\omega_N.
\]

Approximate the expectation by

\[
\mathbb E[J]
\approx
\frac1N
\sum_{i=1}^N
J(\omega_i).
\]

Scenario optimization converts stochastic optimization into a finite
family of deterministic problems.

%%%%%%%%%%%%%%%%%%%%%%%%%%%%%%%%%%%%%%%%%%%%%%%%%%%%%%%%%%%%%%%

\section{Risk Measures}

Expected value alone cannot distinguish stable systems from highly
volatile ones.

Introduce the variance

\[
\sigma^2
=
\operatorname{Var}(J).
\]

More generally,

let

\[
\rho(J)
\]

be a coherent risk measure.

Optimization becomes

\[
\boxed{
\min
\left(
\mathbb E[J]
+
\lambda
\rho(J)
\right).
}
\]

Large

\[
\lambda
\]

produces increasingly risk-averse infrastructure.

%%%%%%%%%%%%%%%%%%%%%%%%%%%%%%%%%%%%%%%%%%%%%%%%%%%%%%%%%%%%%%%

\section{Stochastic Resilience}

The resilience functional introduced previously also becomes random.

Define

\[
\mathcal R(\omega).
\]

Its expected value is

\[
\mathbb E[\mathcal R].
\]

Its variance measures uncertainty in recovery capability.

Planning therefore seeks systems possessing both large expected
resilience and small resilience variability.

%%%%%%%%%%%%%%%%%%%%%%%%%%%%%%%%%%%%%%%%%%%%%%%%%%%%%%%%%%%%%%%

\section{Adaptive Learning}

Probability distributions should evolve as new observations become
available.

Let

\[
\mathcal D_t
\]

denote all observations collected before time

\[
t.
\]

The conditional distribution becomes

\[
\mathbb P(X|\mathcal D_t).
\]

Infrastructure planning therefore continually updates its probabilistic
model rather than relying upon fixed assumptions.

%%%%%%%%%%%%%%%%%%%%%%%%%%%%%%%%%%%%%%%%%%%%%%%%%%%%%%%%%%%%%%%

\section{The Principle of Probabilistic Sufficiency}

\begin{principle}{Probabilistic Sufficiency}

Infrastructure should not be optimized for one predicted future.

Instead, it should perform satisfactorily across the entire probability
distribution of plausible futures.

\end{principle}

%%%%%%%%%%%%%%%%%%%%%%%%%%%%%%%%%%%%%%%%%%%%%%%%%%%%%%%%%%%%%%%

\section{Discussion}

Deterministic optimization identifies one ideal infrastructure under one
ideal future.

Stochastic optimization instead evaluates families of possible futures,
balancing efficiency against uncertainty.

Probability therefore becomes an intrinsic component of infrastructure
mathematics rather than merely a correction applied after deterministic
planning.

Robust optimization is developed next: designs are
constructed to remain admissible under worst-case perturbations,
providing a bridge between probabilistic planning and guaranteed system
performance.

%%%%%%%%%%%%%%%%%%%%%%%%%%%%%%%%%%%%%%%%%%%%%%%%%%%%%%%%%%%%%%%%%%%%%%
\chapter{Robust Optimization and Adversarial Infrastructure}

\section{Introduction}

The preceding chapter incorporated uncertainty through probability.

Probabilistic optimization is appropriate when future conditions can be
described by meaningful probability distributions.

Many important infrastructure problems, however, lack reliable
statistical information.

Extreme weather,

geopolitical disruption,

cyber attacks,

technological discontinuities,

and unprecedented demographic change may possess no trustworthy
probability model.

Rather than optimizing expected performance,

one instead seeks infrastructure that performs satisfactorily under every
admissible realization of uncertainty.

This philosophy leads naturally to robust optimization.

%%%%%%%%%%%%%%%%%%%%%%%%%%%%%%%%%%%%%%%%%%%%%%%%%%%%%%%%%%%%%%%

\section{Uncertainty Sets}

Instead of assigning probabilities,

assume uncertain parameters belong to a prescribed uncertainty set

\[
\Theta.
\]

Each

\[
\theta\in\Theta
\]

represents one admissible realization of future conditions.

Examples include

future demand,

construction costs,

climate,

equipment degradation,

fuel prices,

or environmental regulations.

%%%%%%%%%%%%%%%%%%%%%%%%%%%%%%%%%%%%%%%%%%%%%%%%%%%%%%%%%%%%%%%

\begin{definition}{Uncertainty Set}

An uncertainty set is a compact subset

\[
\Theta\subseteq\mathbb R^m
\]

containing every admissible realization of uncertain parameters.

\end{definition}

%%%%%%%%%%%%%%%%%%%%%%%%%%%%%%%%%%%%%%%%%%%%%%%%%%%%%%%%%%%%%%%

\section{Worst-Case Optimization}

Suppose

\[
J(x,\theta)
\]

depends upon uncertain parameters.

The robust optimization problem becomes

\[
\boxed{
\min_x
\;
\max_{\theta\in\Theta}
J(x,\theta).
}
\]

The outer optimization chooses infrastructure.

The inner optimization represents the most adverse admissible
environment.

%%%%%%%%%%%%%%%%%%%%%%%%%%%%%%%%%%%%%%%%%%%%%%%%%%%%%%%%%%%%%%%

\section{Minimax Interpretation}

The optimization can be interpreted as a strategic interaction.

The planner selects

\[
x.
\]

Nature subsequently selects

\[
\theta.
\]

Infrastructure is therefore designed assuming that every uncertainty acts
in the least favorable admissible manner.

%%%%%%%%%%%%%%%%%%%%%%%%%%%%%%%%%%%%%%%%%%%%%%%%%%%%%%%%%%%%%%%

\begin{principle}{Robust Design}

Preferred infrastructure minimizes the greatest admissible loss rather
than maximizing average performance.

\end{principle}

%%%%%%%%%%%%%%%%%%%%%%%%%%%%%%%%%%%%%%%%%%%%%%%%%%%%%%%%%%%%%%%

\section{Robust Feasibility}

Suppose operational constraints satisfy

\[
g(x,\theta)\le0.
\]

A deterministic design verifies the constraint for one value of

\[
\theta.
\]

A robust design requires

\[
\boxed{
g(x,\theta)\le0,
\qquad
\forall\theta\in\Theta.
}
\]

Thus every admissible uncertainty preserves operational feasibility.

%%%%%%%%%%%%%%%%%%%%%%%%%%%%%%%%%%%%%%%%%%%%%%%%%%%%%%%%%%%%%%%

\section{Performance Regret}

Absolute optimality may be impossible under uncertainty.

Instead compare a robust solution with the hypothetical optimum that
would have been chosen had the future been known.

Define the regret

\[
\boxed{
R(x,\theta)
=
J(x,\theta)
-
J^\star(\theta),
}
\]

where

\[
J^\star(\theta)
\]

is the optimal cost for realization

\[
\theta.
\]

%%%%%%%%%%%%%%%%%%%%%%%%%%%%%%%%%%%%%%%%%%%%%%%%%%%%%%%%%%%%%%%

Infrastructure design may therefore minimize

maximum regret rather than maximum cost.

%%%%%%%%%%%%%%%%%%%%%%%%%%%%%%%%%%%%%%%%%%%%%%%%%%%%%%%%%%%%%%%

\section{Robust Corridor Capacity}

Suppose corridor demand satisfies

\[
D\in[D_{\min},D_{\max}].
\]

Capacity

\[
K
\]

must satisfy

\[
K
\ge
D_{\max}.
\]

This conservative design sacrifices efficiency to guarantee feasibility.

%%%%%%%%%%%%%%%%%%%%%%%%%%%%%%%%%%%%%%%%%%%%%%%%%%%%%%%%%%%%%%%

More generally,

capacity optimization becomes

\[
\min_K
C(K)
\]

subject to

\[
K
\ge
D(\theta),
\qquad
\forall\theta\in\Theta.
\]

%%%%%%%%%%%%%%%%%%%%%%%%%%%%%%%%%%%%%%%%%%%%%%%%%%%%%%%%%%%%%%%

\section{Budget Uncertainty}

Construction costs frequently vary.

Suppose

\[
c_i(\theta)
\]

is the uncertain cost of project

\[
i.
\]

Total investment becomes

\[
\sum_i
c_i(\theta)x_i.
\]

Robust budgeting requires

\[
\sum_i
c_i(\theta)x_i
\le
B,
\qquad
\forall\theta\in\Theta.
\]

%%%%%%%%%%%%%%%%%%%%%%%%%%%%%%%%%%%%%%%%%%%%%%%%%%%%%%%%%%%%%%%

\section{Failure Tolerance}

Suppose at most

\[
k
\]

components fail simultaneously.

Let

\[
F_k
\]

denote every failure set containing at most

\[
k
\]

components.

%%%%%%%%%%%%%%%%%%%%%%%%%%%%%%%%%%%%%%%%%%%%%%%%%%%%%%%%%%%%%%%

A robust infrastructure satisfies operational admissibility after every
failure in

\[
F_k.
\]

%%%%%%%%%%%%%%%%%%%%%%%%%%%%%%%%%%%%%%%%%%%%%%%%%%%%%%%%%%%%%%%

\begin{definition}{\(k\)-Failure Robustness}

An infrastructure is

\[
k
\]

-failure robust if every admissible failure involving at most

\[
k
\]

components admits continued operation within prescribed performance
limits.

\end{definition}

%%%%%%%%%%%%%%%%%%%%%%%%%%%%%%%%%%%%%%%%%%%%%%%%%%%%%%%%%%%%%%%

\section{Adversarial Routing}

Congestion,

failures,

or attacks may deliberately attempt to reduce network performance.

Suppose an adversary removes

\[
r
\]

edges.

The resulting optimization becomes

\[
\min_x
\;
\max_{F\subseteq E,\;|F|=r}
J(x,F).
\]

Robust corridor placement therefore anticipates targeted disruption.

%%%%%%%%%%%%%%%%%%%%%%%%%%%%%%%%%%%%%%%%%%%%%%%%%%%%%%%%%%%%%%%

\section{Robustness Margin}

Define

\[
\delta^\star
\]

as the largest uncertainty magnitude preserving admissible operation.

%%%%%%%%%%%%%%%%%%%%%%%%%%%%%%%%%%%%%%%%%%%%%%%%%%%%%%%%%%%%%%%

\begin{definition}{Robustness Margin}

The robustness margin is

\[
\boxed{
\delta^\star
=
\sup
\{
\delta:
x
\text{ remains admissible for all }
\|\theta\|\le\delta
\}.
}
\]

\end{definition}

Large margins indicate designs capable of tolerating substantial
uncertainty before failure.

%%%%%%%%%%%%%%%%%%%%%%%%%%%%%%%%%%%%%%%%%%%%%%%%%%%%%%%%%%%%%%%

\section{Trade-offs}

Robustness is rarely free.

Increasing reserve capacity,

additional redundancy,

and stronger components increase construction and maintenance costs.

Consequently,

robust optimization naturally extends the Pareto framework developed
earlier.

Typical competing objectives include

construction cost,

operating efficiency,

environmental impact,

and robustness margin.

%%%%%%%%%%%%%%%%%%%%%%%%%%%%%%%%%%%%%%%%%%%%%%%%%%%%%%%%%%%%%%%

\section{Adaptive Robustness}

Static robustness prepares for uncertainty before operation begins.

Adaptive robustness additionally permits future modification.

Let

\[
u(t)
\]

represent adaptive control.

Optimization becomes

\[
\min_u
\;
\max_{\theta}
J(u,\theta).
\]

Thus operational flexibility becomes an explicit component of robust
design rather than an afterthought.

%%%%%%%%%%%%%%%%%%%%%%%%%%%%%%%%%%%%%%%%%%%%%%%%%%%%%%%%%%%%%%%

\section{Discussion}

Robust optimization complements probabilistic optimization rather than
replacing it.

Probabilistic methods exploit information when reliable statistical
models exist.

Robust methods remain applicable when uncertainty is bounded but poorly
characterized.

Together they provide complementary mathematical frameworks for
long-term infrastructure planning.

The final chapter of this Part synthesizes deterministic, probabilistic, and
robust optimization into a unified theory of adaptive infrastructure,
where planning, operation, uncertainty, and resilience are treated as
different aspects of one mathematical system.

%%%%%%%%%%%%%%%%%%%%%%%%%%%%%%%%%%%%%%%%%%%%%%%%%%%%%%%%%%%%%%%%%%%%%%
\chapter{A Unified Theory of Adaptive Infrastructure}

\section{Introduction}

The preceding chapters have developed several apparently distinct
mathematical theories.

Infrastructure was first formulated as a variational optimization
problem.

Dynamics were introduced through systems of differential equations.

Optimal control extended these dynamics to decision making.

Graph theory characterized network topology.

Differential geometry described embedded corridor systems.

Probability incorporated uncertainty.

Robust optimization addressed worst-case behavior.

Despite their differing mathematical language, these theories describe
different aspects of a single object.

The purpose of this chapter is to construct a unified mathematical
framework in which each earlier formulation appears as a special case.

%%%%%%%%%%%%%%%%%%%%%%%%%%%%%%%%%%%%%%%%%%%%%%%%%%%%%%%%%%%%%%%

\section{The Infrastructure System}

Let

\[
\mathcal I
=
(M,G,X,U,\Theta),
\]

where

\[
M
\]

is the underlying manifold,

\[
G
\]

is the embedded infrastructure graph,

\[
X
\]

is the state space,

\[
U
\]

is the admissible control space,

and

\[
\Theta
\]

represents uncertainty.

%%%%%%%%%%%%%%%%%%%%%%%%%%%%%%%%%%%%%%%%%%%%%%%%%%%%%%%%%%%%%%%

Every infrastructure problem considered in this monograph is obtained by
restricting one or more components of

\[
\mathcal I.
\]

%%%%%%%%%%%%%%%%%%%%%%%%%%%%%%%%%%%%%%%%%%%%%%%%%%%%%%%%%%%%%%%

\section{State Evolution}

The complete infrastructure evolves according to

\[
\boxed{
\frac{dx}{dt}
=
f(x,u,\theta,t).
}
\]

The function

\[
f
\]

includes

physical transport,

operational decisions,

external disturbances,

and network coupling.

%%%%%%%%%%%%%%%%%%%%%%%%%%%%%%%%%%%%%%%%%%%%%%%%%%%%%%%%%%%%%%%

\section{Infrastructure Histories}

A complete infrastructure history is a mapping

\[
\Gamma
=
(x,u):
[0,T]
\rightarrow
X\times U.
\]

Rather than studying isolated states,

the mathematical object of interest becomes

the entire trajectory.

%%%%%%%%%%%%%%%%%%%%%%%%%%%%%%%%%%%%%%%%%%%%%%%%%%%%%%%%%%%%%%%

\begin{definition}{Infrastructure History}

An infrastructure history consists of every admissible state and control
pair throughout the planning interval.

\end{definition}

%%%%%%%%%%%%%%%%%%%%%%%%%%%%%%%%%%%%%%%%%%%%%%%%%%%%%%%%%%%%%%%

\section{Universal Objective Functional}

Every optimization problem considered previously can be written as

\[
\boxed{
J[\Gamma]
=
\mathbb E
\left[
\int_0^T
L(x,u,\theta,t)
dt
+
\Phi(x(T))
\right].
}
\]

Different mathematical theories correspond to different choices of

\[
L,
\]

\[
\Phi,
\]

and

\[
\Theta.
\]

%%%%%%%%%%%%%%%%%%%%%%%%%%%%%%%%%%%%%%%%%%%%%%%%%%%%%%%%%%%%%%%

\section{Special Cases}

When

\[
\Theta
=
\{\theta_0\},
\]

the theory reduces to deterministic optimization.

When

\[
\Theta
\]

possesses a probability measure,

the theory becomes stochastic optimization.

When only an uncertainty set is specified,

the expectation is replaced by

\[
\max_{\theta\in\Theta},
\]

yielding robust optimization.

%%%%%%%%%%%%%%%%%%%%%%%%%%%%%%%%%%%%%%%%%%%%%%%%%%%%%%%%%%%%%%%

\section{Infrastructure Constraints}

Every admissible history satisfies

\[
\frac{dx}{dt}
=
f(x,u,\theta,t),
\]

together with

capacity constraints,

conservation laws,

operational policies,

and geometric compatibility.

Collect these into the admissible history set

\[
\mathcal H.
\]

Optimization therefore becomes

\[
\boxed{
\min_{\Gamma\in\mathcal H}
J[\Gamma].
}
\]

%%%%%%%%%%%%%%%%%%%%%%%%%%%%%%%%%%%%%%%%%%%%%%%%%%%%%%%%%%%%%%%

\section{Geometry and Topology}

The manifold

\[
M
\]

determines

distance,

curvature,

and terrain.

The graph

\[
G
\]

determines

connectivity,

cycles,

cuts,

and routing alternatives.

Operational optimization therefore occurs simultaneously over

continuous geometry

and

discrete topology.

%%%%%%%%%%%%%%%%%%%%%%%%%%%%%%%%%%%%%%%%%%%%%%%%%%%%%%%%%%%%%%%

\section{Temporal Structure}

Infrastructure variables evolve on multiple timescales.

Represent this hierarchy by

\[
x
=
(x_f,x_s),
\]

where

\[
x_f
\]

contains fast operational variables and

\[
x_s
\]

contains slowly varying structural variables.

The coupled dynamics satisfy

\[
\frac{dx_f}{dt}
=
f_f(x_f,x_s,u),
\]

and

\[
\frac{dx_s}{dt}
=
f_s(x_f,x_s,u).
\]

%%%%%%%%%%%%%%%%%%%%%%%%%%%%%%%%%%%%%%%%%%%%%%%%%%%%%%%%%%%%%%%

Fast control therefore operates within an evolving geometric framework.

%%%%%%%%%%%%%%%%%%%%%%%%%%%%%%%%%%%%%%%%%%%%%%%%%%%%%%%%%%%%%%%

\section{Resilience}

Failures correspond to perturbations

\[
x
\rightarrow
x+\delta x.
\]

Recovery consists of constructing another admissible history

\[
\Gamma'
\]

connecting the disturbed state back into the admissible operating
region.

Thus resilience is naturally expressed within the same mathematical
framework.

%%%%%%%%%%%%%%%%%%%%%%%%%%%%%%%%%%%%%%%%%%%%%%%%%%%%%%%%%%%%%%%

\section{Learning}

Infrastructure models improve through observation.

Let

\[
\mathcal D_t
\]

denote accumulated observations.

The uncertainty model evolves according to

\[
\Theta_t
=
F(\Theta_{t-1},\mathcal D_t).
\]

Planning therefore becomes increasingly informed through continued
operation.

%%%%%%%%%%%%%%%%%%%%%%%%%%%%%%%%%%%%%%%%%%%%%%%%%%%%%%%%%%%%%%%

\section{Adaptation}

Operational controls respond continuously.

Structural variables evolve gradually.

Protocols evolve only occasionally.

This hierarchy may be written

\[
u
\prec
x_s
\prec
\Pi,
\]

where

\[
\Pi
\]

denotes the governing protocol.

Operational decisions therefore occur much more frequently than protocol
revision.

%%%%%%%%%%%%%%%%%%%%%%%%%%%%%%%%%%%%%%%%%%%%%%%%%%%%%%%%%%%%%%%

\section{Unified Optimization Principle}

\begin{principle}{Unified Infrastructure Principle}

An adaptive infrastructure system seeks an admissible history that

simultaneously

preserves physical feasibility,

maintains operational continuity,

minimizes cumulative cost,

remains resilient under disturbance,

and adapts to uncertain future conditions.

\end{principle}

%%%%%%%%%%%%%%%%%%%%%%%%%%%%%%%%%%%%%%%%%%%%%%%%%%%%%%%%%%%%%%%

\section{Mathematical Interpretation}

The mathematical object studied throughout this monograph is therefore
not

a graph,

a manifold,

a differential equation,

or a variational functional individually.

Instead,

it is the admissible history of an embedded network evolving under
physical law, operational control, uncertainty, and long-term
adaptation.

Every previous chapter contributes one projection of this more general
structure.

%%%%%%%%%%%%%%%%%%%%%%%%%%%%%%%%%%%%%%%%%%%%%%%%%%%%%%%%%%%%%%%

\section{Toward Infrastructure Science}

The framework developed throughout this volume is intentionally
independent of any specific engineering discipline.

Electrical grids,

water systems,

transportation,

communications,

energy distribution,

and logistics all appear as particular realizations of the same
mathematical formalism.

The common structure is

geometry,

topology,

conservation,

optimization,

control,

uncertainty,

adaptation,

and continuation.

%%%%%%%%%%%%%%%%%%%%%%%%%%%%%%%%%%%%%%%%%%%%%%%%%%%%%%%%%%%%%%%

\section{Concluding Remarks}

The mathematical theory developed in this monograph provides a common
language for describing large-scale infrastructure systems across both
space and time.

Beginning with orthodromic geometry and proceeding through variational
principles, network topology, optimal control, stochastic analysis, and
robust optimization, each chapter has expanded the scope of the theory
without abandoning its foundational concepts.

The resulting framework does not prescribe one optimal infrastructure.
Instead, it provides the mathematical tools required to analyze,
compare, optimize, and adapt infrastructure systems under realistic
physical, operational, and environmental constraints.

Future developments may extend this theory toward autonomous planning,
planetary-scale resource allocation, interplanetary logistics, and
self-adapting infrastructure capable of continual optimization over
civilizational timescales.

%%%%%%%%%%%%%%%%%%%%%%%%%%%%%%%%%%%%%%%%%%%%%%%%%%%%%%%%%%%%%%%%%%%%%%
\part{Applications and Case Studies}

%%%%%%%%%%%%%%%%%%%%%%%%%%%%%%%%%%%%%%%%%%%%%%%%%%%%%%%%%%%%%%%%%%%%%%
\chapter{Electrical Power Networks}
\label{chap:electrical-power-networks}

\section{Introduction}

Electrical transmission systems provide one of the clearest
manifestations of the mathematical framework developed throughout this
monograph.

They possess

continuous physical dynamics,

network topology,

capacity constraints,

optimal control,

stochastic disturbances,

and long-term infrastructure evolution.

Consequently they provide an ideal first application.

%%%%%%%%%%%%%%%%%%%%%%%%%%%%%%%%%%%%%%%%%%%%%%%%%%%%%%%%%%%%%%%

\section{Network Representation}

Represent the electrical grid by the graph

\[
G=(V,E).
\]

Vertices represent

generators,

substations,

storage facilities,

and loads.

Edges represent transmission corridors.

Each transmission line possesses

resistance,

reactance,

thermal capacity,

construction cost,

maintenance cost,

and reliability.

%%%%%%%%%%%%%%%%%%%%%%%%%%%%%%%%%%%%%%%%%%%%%%%%%%%%%%%%%%%%%%%

\section{Power Balance}

Let

\[
P_i
\]

denote net power injection at node

\[
i.
\]

Kirchhoff's Current Law becomes

\[
\boxed{
\sum_j
P_{ij}
=
P_i.
}
\]

Power is therefore conserved at every node.

%%%%%%%%%%%%%%%%%%%%%%%%%%%%%%%%%%%%%%%%%%%%%%%%%%%%%%%%%%%%%%%

\section{DC Power Flow}

Under the standard DC approximation,

line flow satisfies

\[
\boxed{
P_{ij}
=
B_{ij}
(\theta_i-\theta_j),
}
\]

where

\[
B_{ij}
\]

is line susceptance and

\[
\theta_i
\]

is voltage phase angle.

Collecting every node yields

\[
L\theta=P,
\]

where

\[
L
\]

is the weighted graph Laplacian.

%%%%%%%%%%%%%%%%%%%%%%%%%%%%%%%%%%%%%%%%%%%%%%%%%%%%%%%%%%%%%%%

\section{Capacity Constraints}

Every transmission corridor possesses thermal limits.

Accordingly,

\[
|P_{ij}|
\le
P_{ij}^{\max}.
\]

Violation produces overheating,

sag,

or equipment damage.

%%%%%%%%%%%%%%%%%%%%%%%%%%%%%%%%%%%%%%%%%%%%%%%%%%%%%%%%%%%%%%%

\section{Generation Constraints}

Each generator satisfies

\[
P_i^{\min}
\le
P_i
\le
P_i^{\max}.
\]

Ramp-rate limitations additionally require

\[
\left|
\frac{dP_i}{dt}
\right|
\le
R_i.
\]

%%%%%%%%%%%%%%%%%%%%%%%%%%%%%%%%%%%%%%%%%%%%%%%%%%%%%%%%%%%%%%%

\section{Economic Dispatch}

Suppose generator

\[
i
\]

possesses operating cost

\[
C_i(P_i).
\]

The dispatch problem becomes

\[
\boxed{
\min
\sum_i
C_i(P_i),
}
\]

subject to

power balance,

network constraints,

and generation limits.

%%%%%%%%%%%%%%%%%%%%%%%%%%%%%%%%%%%%%%%%%%%%%%%%%%%%%%%%%%%%%%%

\section{Optimal Power Flow}

Including transmission constraints yields

\[
\boxed{
\min
J(P,\theta)
}
\]

subject to

\[
L\theta=P,
\]

generation limits,

transmission limits,

and voltage constraints.

Optimal Power Flow is therefore a direct specialization of the unified
optimization framework developed earlier.

%%%%%%%%%%%%%%%%%%%%%%%%%%%%%%%%%%%%%%%%%%%%%%%%%%%%%%%%%%%%%%%

\section{Energy Storage}

Introduce storage state

\[
E(t).
\]

Storage dynamics satisfy

\[
\boxed{
\frac{dE}{dt}
=
P_{\rm charge}
-
P_{\rm discharge}.
}
\]

Storage introduces temporal coupling into network optimization.

%%%%%%%%%%%%%%%%%%%%%%%%%%%%%%%%%%%%%%%%%%%%%%%%%%%%%%%%%%%%%%%

\section{Renewable Generation}

Wind and solar generation satisfy

\[
P_R(t,\omega),
\]

where

\[
\omega
\]

represents environmental uncertainty.

Optimization therefore naturally becomes stochastic,

as developed in Part~VI.

%%%%%%%%%%%%%%%%%%%%%%%%%%%%%%%%%%%%%%%%%%%%%%%%%%%%%%%%%%%%%%%

\section{Grid Stability}

Let

\[
\omega_i
\]

denote frequency deviation at node

\[
i.
\]

The swing equation is

\[
M_i
\frac{d^2\delta_i}{dt^2}
+
D_i
\frac{d\delta_i}{dt}
=
P_m-P_e.
\]

This dynamical equation couples directly to the optimal control
framework developed previously.

%%%%%%%%%%%%%%%%%%%%%%%%%%%%%%%%%%%%%%%%%%%%%%%%%%%%%%%%%%%%%%%

\section{Contingency Analysis}

Power systems are commonly designed according to the

\[
N-1
\]

criterion.

%%%%%%%%%%%%%%%%%%%%%%%%%%%%%%%%%%%%%%%%%%%%%%%%%%%%%%%%%%%%%%%

\begin{definition}{\(N-1\) Security}

A power system satisfies

\[
N-1
\]

security if continued admissible operation remains possible after the
failure of any single transmission element or generator.

\end{definition}

This is precisely the notion of

\[
k
\]

-failure robustness introduced in the previous chapter with

\[
k=1.
\]

%%%%%%%%%%%%%%%%%%%%%%%%%%%%%%%%%%%%%%%%%%%%%%%%%%%%%%%%%%%%%%%

\section{Adaptive Grid Operation}

Demand,

renewable generation,

and market prices change continually.

Model Predictive Control therefore repeatedly solves

Optimal Power Flow,

implements the first control interval,

observes the updated system,

and repeats.

Electrical grid management thus becomes a direct application of adaptive
optimal control.

%%%%%%%%%%%%%%%%%%%%%%%%%%%%%%%%%%%%%%%%%%%%%%%%%%%%%%%%%%%%%%%

\section{Infrastructure Geometry}

Transmission corridors are embedded on the Earth's surface.

Routing therefore minimizes

\[
\int_\gamma
\Phi
\,ds,
\]

where

\[
\Phi
\]

incorporates terrain,

construction cost,

environmental impact,

and land acquisition.

Thus geometric optimization determines where transmission lines should
be built,

while network optimization determines how they should operate.

%%%%%%%%%%%%%%%%%%%%%%%%%%%%%%%%%%%%%%%%%%%%%%%%%%%%%%%%%%%%%%%

\section{Discussion}

Electrical transmission systems illustrate every major mathematical idea
developed in this volume.

Graph theory determines connectivity.

Differential equations govern physical dynamics.

Variational principles determine routing.

Optimization allocates generation.

Control theory regulates operation.

Probability models renewable uncertainty.

Robust optimization ensures contingency performance.

Consequently,

electrical infrastructure serves as an archetypal example of adaptive
network optimization.

An analogous formulation for hydraulic
networks follows, demonstrating that the same mathematical framework extends
naturally from electrical energy transport to water distribution and
fluid infrastructure.

%%%%%%%%%%%%%%%%%%%%%%%%%%%%%%%%%%%%%%%%%%%%%%%%%%%%%%%%%%%%%%%%%%%%%%
\chapter{Hydraulic Networks}

\section{Introduction}

Water distribution systems,

irrigation networks,

stormwater systems,

and long-distance aqueducts provide another canonical example of
networked infrastructure.

Although the transported commodity differs from electrical energy,
the underlying mathematical framework remains remarkably similar.

Nodes represent reservoirs,

junctions,

pumping stations,

treatment facilities,

and consumers.

Edges represent pipes,

canals,

tunnels,

or rivers.

The governing equations differ,

but the topology,

optimization,

control,

and resilience principles developed throughout this monograph remain
unchanged.

%%%%%%%%%%%%%%%%%%%%%%%%%%%%%%%%%%%%%%%%%%%%%%%%%%%%%%%%%%%%%%%

\section{Network Representation}

Represent the hydraulic system by

\[
G=(V,E).
\]

Each edge possesses

diameter,

length,

roughness,

capacity,

construction cost,

maintenance cost,

and elevation profile.

Each node possesses

pressure,

storage,

and demand.

%%%%%%%%%%%%%%%%%%%%%%%%%%%%%%%%%%%%%%%%%%%%%%%%%%%%%%%%%%%%%%%

\section{Mass Conservation}

Let

\[
Q_{ij}
\]

denote volumetric flow.

Mass conservation requires

\[
\boxed{
\sum_j
Q_{ij}
=
D_i,
}
\]

where

\[
D_i
\]

is net demand at node

\[
i.
\]

This equation is mathematically identical to Kirchhoff's Current Law in
electrical networks.

%%%%%%%%%%%%%%%%%%%%%%%%%%%%%%%%%%%%%%%%%%%%%%%%%%%%%%%%%%%%%%%

\section{Energy Conservation}

Pressure losses satisfy

\[
\boxed{
H_i-H_j
=
h_{ij}(Q_{ij}),
}
\]

where

\[
H
\]

is hydraulic head.

The function

\[
h_{ij}
\]

depends upon pipe geometry and friction.

%%%%%%%%%%%%%%%%%%%%%%%%%%%%%%%%%%%%%%%%%%%%%%%%%%%%%%%%%%%%%%%

\section{Pipe Resistance}

Under the Darcy--Weisbach equation,

\[
h_f
=
f
\frac{L}{D}
\frac{v^2}{2g},
\]

where

\[
f
\]

is the friction factor,

\[
L
\]

the pipe length,

\[
D
\]

its diameter,

and

\[
v
\]

the average velocity.

%%%%%%%%%%%%%%%%%%%%%%%%%%%%%%%%%%%%%%%%%%%%%%%%%%%%%%%%%%%%%%%

\section{Pump Dynamics}

Pumps increase hydraulic head.

Let

\[
H_p(Q)
\]

denote the pump curve.

The governing equation becomes

\[
H_{\rm out}
=
H_{\rm in}
+
H_p(Q).
\]

Pump scheduling therefore becomes a control variable within the unified
optimization framework.

%%%%%%%%%%%%%%%%%%%%%%%%%%%%%%%%%%%%%%%%%%%%%%%%%%%%%%%%%%%%%%%

\section{Reservoir Dynamics}

Suppose a storage reservoir contains volume

\[
V(t).
\]

Its dynamics satisfy

\[
\boxed{
\frac{dV}{dt}
=
Q_{\rm in}
-
Q_{\rm out}.
}
\]

Reservoirs therefore introduce temporal coupling analogous to battery
storage in electrical systems.

%%%%%%%%%%%%%%%%%%%%%%%%%%%%%%%%%%%%%%%%%%%%%%%%%%%%%%%%%%%%%%%

\section{Pressure Constraints}

Operational safety requires

\[
H_{\min}
\le
H_i
\le
H_{\max}.
\]

Low pressures threaten service continuity.

High pressures increase leakage and pipe failure.

%%%%%%%%%%%%%%%%%%%%%%%%%%%%%%%%%%%%%%%%%%%%%%%%%%%%%%%%%%%%%%%

\section{Leakage}

Real hydraulic systems experience distributed loss.

Represent leakage by

\[
L_i(H_i).
\]

Mass conservation therefore becomes

\[
\sum_j
Q_{ij}
=
D_i
+
L_i(H_i).
\]

Leakage transforms conservation into a dissipative transport process.

%%%%%%%%%%%%%%%%%%%%%%%%%%%%%%%%%%%%%%%%%%%%%%%%%%%%%%%%%%%%%%%

\section{Optimal Pump Scheduling}

Let

\[
P_k(t)
\]

denote electrical power consumed by pump

\[
k.
\]

Optimization minimizes

\[
\boxed{
J
=
\int_0^T
\sum_k
c(t)
P_k(t)
\,dt,
}
\]

subject to

hydraulic dynamics,

pressure constraints,

and reservoir requirements.

%%%%%%%%%%%%%%%%%%%%%%%%%%%%%%%%%%%%%%%%%%%%%%%%%%%%%%%%%%%%%%%

\section{Water Quality}

Transported commodities are not always conserved perfectly.

Let

\[
C(x,t)
\]

denote contaminant concentration.

Its evolution satisfies the advection--diffusion equation

\[
\boxed{
\frac{\partial C}{\partial t}
+
u\cdot\nabla C
=
D\nabla^2C
-
R(C),
}
\]

where

\[
R(C)
\]

represents chemical decay.

Hydraulic optimization therefore couples flow and water quality.

%%%%%%%%%%%%%%%%%%%%%%%%%%%%%%%%%%%%%%%%%%%%%%%%%%%%%%%%%%%%%%%

\section{Drought and Uncertainty}

Available inflow satisfies

\[
Q_{\rm in}(t,\omega).
\]

Reservoir operation therefore becomes a stochastic optimal control
problem.

Chance constraints may require

\[
\mathbb P
(
V(t)\ge V_{\min}
)
\ge
0.99.
\]

%%%%%%%%%%%%%%%%%%%%%%%%%%%%%%%%%%%%%%%%%%%%%%%%%%%%%%%%%%%%%%%

\section{Failure Recovery}

Hydraulic failures include

pipe rupture,

pump failure,

reservoir contamination,

and valve malfunction.

Recovery consists of constructing admissible rerouting trajectories while
maintaining minimum pressure and water quality throughout the affected
region.

The resilience framework developed previously applies without
modification.

%%%%%%%%%%%%%%%%%%%%%%%%%%%%%%%%%%%%%%%%%%%%%%%%%%%%%%%%%%%%%%%

\section{Geometry}

Pipelines are embedded curves on the Earth's surface.

Construction minimizes

\[
\int_\gamma
\Phi
\,ds,
\]

where

\[
\Phi
\]

includes

terrain,

excavation cost,

geology,

land ownership,

and environmental impact.

The geometric optimization problem is therefore identical to that for
electrical transmission corridors.

%%%%%%%%%%%%%%%%%%%%%%%%%%%%%%%%%%%%%%%%%%%%%%%%%%%%%%%%%%%%%%%

\section{Discussion}

Hydraulic systems illustrate the universality of the mathematical
framework developed throughout this monograph.

Replacing voltage by pressure,

current by flow,

and batteries by reservoirs changes the governing physics but leaves the
underlying mathematical structure essentially unchanged.

The conservation laws,

graph topology,

variational optimization,

control,

stochastic planning,

and resilience theory all transfer directly between infrastructure
domains.

The same framework is applied next to transportation
networks, where the transported quantity is no longer energy or water
but vehicles, passengers, and freight.

%%%%%%%%%%%%%%%%%%%%%%%%%%%%%%%%%%%%%%%%%%%%%%%%%%%%%%%%%%%%%%%%%%%%%%
\chapter{Transportation Networks}

\section{Introduction}

Transportation networks differ fundamentally from electrical and
hydraulic systems because the transported commodity possesses individual
decision-making capability.

Electric current follows physical laws.

Water follows pressure gradients.

Drivers, passengers, and freight operators select routes in response to
travel time, congestion, cost, safety, and policy.

Consequently, transportation systems couple physical infrastructure with
distributed decision processes.

The resulting mathematical framework combines conservation laws,
continuum mechanics, optimization, and game theory.

%%%%%%%%%%%%%%%%%%%%%%%%%%%%%%%%%%%%%%%%%%%%%%%%%%%%%%%%%%%%%%%

\section{Network Representation}

Represent the transportation system by

\[
G=(V,E).
\]

Vertices correspond to

intersections,

stations,

terminals,

ports,

distribution centers,

and airports.

Edges represent

roads,

railways,

shipping corridors,

air routes,

or pedestrian pathways.

Each edge possesses

length,

speed limit,

capacity,

construction cost,

maintenance cost,

and reliability.

%%%%%%%%%%%%%%%%%%%%%%%%%%%%%%%%%%%%%%%%%%%%%%%%%%%%%%%%%%%%%%%

\section{Traffic Density}

Let

\[
\rho(x,t)
\]

denote vehicle density.

Traffic flow satisfies

\[
q=\rho v,
\]

where

\[
v
\]

is average speed.

%%%%%%%%%%%%%%%%%%%%%%%%%%%%%%%%%%%%%%%%%%%%%%%%%%%%%%%%%%%%%%%

\section{Conservation of Vehicles}

Vehicle conservation is expressed by

\[
\boxed{
\frac{\partial\rho}{\partial t}
+
\nabla\cdot(\rho v)
=
s,
}
\]

where

\[
s
\]

represents vehicle sources and sinks.

This equation is directly analogous to the conservation laws developed
earlier for energy and hydraulic transport.

%%%%%%%%%%%%%%%%%%%%%%%%%%%%%%%%%%%%%%%%%%%%%%%%%%%%%%%%%%%%%%%

\section{Fundamental Diagram}

Traffic speed decreases as density increases.

Let

\[
v=v(\rho).
\]

Then

\[
q(\rho)
=
\rho v(\rho).
\]

A common approximation is

\[
v(\rho)
=
v_f
\left(
1-
\frac{\rho}{\rho_{\max}}
\right),
\]

where

\[
v_f
\]

is the free-flow speed.

%%%%%%%%%%%%%%%%%%%%%%%%%%%%%%%%%%%%%%%%%%%%%%%%%%%%%%%%%%%%%%%

\section{The Lighthill--Whitham--Richards Model}

Substituting

\[
q=q(\rho)
\]

into the conservation equation gives

\[
\boxed{
\frac{\partial\rho}{\partial t}
+
\frac{\partial q(\rho)}{\partial x}
=
0.
}
\]

This nonlinear conservation law predicts shock waves,
queue formation,
and congestion propagation.

%%%%%%%%%%%%%%%%%%%%%%%%%%%%%%%%%%%%%%%%%%%%%%%%%%%%%%%%%%%%%%%

\section{Node Conservation}

At every intersection,

incoming and outgoing flow satisfy

\[
\sum_j
q_{ji}
-
\sum_k
q_{ik}
=
d_i,
\]

where

\[
d_i
\]

represents local demand.

%%%%%%%%%%%%%%%%%%%%%%%%%%%%%%%%%%%%%%%%%%%%%%%%%%%%%%%%%%%%%%%

\section{Capacity Constraints}

Each corridor satisfies

\[
q
\le
q_{\max}.
\]

Beyond this limit,

additional demand produces congestion rather than increased throughput.

%%%%%%%%%%%%%%%%%%%%%%%%%%%%%%%%%%%%%%%%%%%%%%%%%%%%%%%%%%%%%%%

\section{Travel Time}

Suppose edge

\[
e
\]

has travel time

\[
T_e(q).
\]

Typically,

\[
\frac{dT_e}{dq}
>
0.
\]

Congestion therefore increases travel time nonlinearly.

%%%%%%%%%%%%%%%%%%%%%%%%%%%%%%%%%%%%%%%%%%%%%%%%%%%%%%%%%%%%%%%

\section{Shortest Paths}

In uncongested networks,

routing minimizes

\[
\sum_e
T_e.
\]

On curved surfaces,

the route additionally minimizes

\[
\int_\gamma
\Phi
\,ds,
\]

where

\[
\Phi
\]

represents terrain,

construction,

or policy costs.

%%%%%%%%%%%%%%%%%%%%%%%%%%%%%%%%%%%%%%%%%%%%%%%%%%%%%%%%%%%%%%%

\section{User Equilibrium}

Drivers select routes independently.

%%%%%%%%%%%%%%%%%%%%%%%%%%%%%%%%%%%%%%%%%%%%%%%%%%%%%%%%%%%%%%%

\begin{definition}{Wardrop Equilibrium}

A transportation network is in user equilibrium when no traveler can
reduce travel time by unilaterally changing routes.

\end{definition}

This equilibrium resembles a Nash equilibrium in game theory.

%%%%%%%%%%%%%%%%%%%%%%%%%%%%%%%%%%%%%%%%%%%%%%%%%%%%%%%%%%%%%%%

\section{System Optimum}

The socially optimal routing minimizes

\[
\boxed{
\min
\sum_e
q_e
T_e(q_e).
}
\]

Unlike user equilibrium,

system optimum minimizes total travel time across all travelers.

%%%%%%%%%%%%%%%%%%%%%%%%%%%%%%%%%%%%%%%%%%%%%%%%%%%%%%%%%%%%%%%

\section{Traffic Control}

Traffic signals,

variable speed limits,

ramp metering,

and reversible lanes become control variables

\[
u(t).
\]

Model Predictive Control repeatedly optimizes

signal timing,

lane allocation,

and routing

using continually updated traffic observations.

%%%%%%%%%%%%%%%%%%%%%%%%%%%%%%%%%%%%%%%%%%%%%%%%%%%%%%%%%%%%%%%

\section{Public Transportation}

Transit systems introduce scheduled services.

Suppose

\[
h_i
\]

is the headway of route

\[
i.
\]

Average waiting time satisfies

\[
W_i
=
\frac{h_i}{2},
\]

assuming random passenger arrivals.

Optimization balances

vehicle frequency,

operating cost,

and passenger delay.

%%%%%%%%%%%%%%%%%%%%%%%%%%%%%%%%%%%%%%%%%%%%%%%%%%%%%%%%%%%%%%%

\section{Freight Logistics}

Freight transport introduces inventory dynamics.

Let

\[
I(t)
\]

denote inventory at a logistics node.

Its evolution satisfies

\[
\boxed{
\frac{dI}{dt}
=
Q_{\rm in}
-
Q_{\rm out}
-
D.
}
\]

Transportation and storage therefore become coupled optimization
problems.

%%%%%%%%%%%%%%%%%%%%%%%%%%%%%%%%%%%%%%%%%%%%%%%%%%%%%%%%%%%%%%%

\section{Network Resilience}

Transportation resilience includes

bridge failures,

road closures,

rail disruptions,

weather events,

and accidents.

Recovery seeks alternate routes minimizing disruption while preserving
critical mobility.

Graph connectivity,

cycle space,

and redundancy therefore become essential structural properties.

%%%%%%%%%%%%%%%%%%%%%%%%%%%%%%%%%%%%%%%%%%%%%%%%%%%%%%%%%%%%%%%

\section{Autonomous Transportation}

Autonomous vehicles provide continuously adjustable routing.

Instead of static route assignment,

routing becomes a distributed optimal control problem.

Each vehicle solves

\[
\min
J_i,
\]

subject to

global coordination constraints,

thereby coupling transportation optimization with communication
networks.

%%%%%%%%%%%%%%%%%%%%%%%%%%%%%%%%%%%%%%%%%%%%%%%%%%%%%%%%%%%%%%%

\section{Discussion}

Transportation systems complete the triad of canonical infrastructure
examples.

Electrical networks transport energy.

Hydraulic networks transport fluids.

Transportation networks transport people and goods.

Although the governing constitutive equations differ,

each domain shares the same mathematical foundations:

graph topology,

conservation laws,

embedded geometry,

optimization,

control,

uncertainty,

and resilience.

These ideas extend next to communication
networks, where the transported quantity is information rather than
matter or energy, demonstrating once again that the mathematical
framework remains unchanged while the physical interpretation evolves.

%%%%%%%%%%%%%%%%%%%%%%%%%%%%%%%%%%%%%%%%%%%%%%%%%%%%%%%%%%%%%%%%%%%%%%
\chapter{Communication Networks}

\section{Introduction}

Communication networks transport information rather than matter or
energy.

Unlike electrical, hydraulic, and transportation systems, the conserved
quantity is information flow.

Messages,

voice,

video,

sensor measurements,

control commands,

and distributed computations all traverse the same underlying network.

Although the governing physical mechanisms differ, the mathematical
structure developed throughout this monograph remains unchanged.

Nodes correspond to devices capable of generating, processing, storing,
or forwarding information.

Edges correspond to communication channels.

Routing,

capacity,

optimization,

control,

uncertainty,

and resilience retain precisely the same mathematical interpretation as
in previous applications.

%%%%%%%%%%%%%%%%%%%%%%%%%%%%%%%%%%%%%%%%%%%%%%%%%%%%%%%%%%%%%%%

\section{Network Representation}

Represent the communication infrastructure by

\[
G=(V,E).
\]

Vertices represent

routers,

switches,

base stations,

satellites,

servers,

data centers,

and user devices.

Edges represent

fiber-optic links,

wireless channels,

microwave relays,

undersea cables,

or satellite connections.

%%%%%%%%%%%%%%%%%%%%%%%%%%%%%%%%%%%%%%%%%%%%%%%%%%%%%%%%%%%%%%%

\section{Information Flow}

Let

\[
F_{ij}
\]

denote information flow along edge

\[
(i,j).
\]

Information conservation requires

\[
\boxed{
\sum_jF_{ji}
-
\sum_kF_{ik}
=
S_i,
}
\]

where

\[
S_i
\]

represents data generated or consumed at node

\[
i.
\]

This equation has the same algebraic form as conservation of charge,
water, and traffic.

%%%%%%%%%%%%%%%%%%%%%%%%%%%%%%%%%%%%%%%%%%%%%%%%%%%%%%%%%%%%%%%

\section{Bandwidth Constraints}

Every communication channel possesses finite capacity.

Let

\[
C_{ij}
\]

denote channel bandwidth.

Then

\[
\boxed{
F_{ij}
\le
C_{ij}.
}
\]

Exceeding channel capacity produces congestion,
delay,
or packet loss.

%%%%%%%%%%%%%%%%%%%%%%%%%%%%%%%%%%%%%%%%%%%%%%%%%%%%%%%%%%%%%%%

\section{Latency}

Each communication edge possesses propagation delay

\[
\tau_e.
\]

The end-to-end latency of a path

\[
\gamma
\]

is

\[
\boxed{
T(\gamma)
=
\sum_{e\in\gamma}
\tau_e.
}
\]

Optimization therefore seeks low-latency routes while respecting
capacity constraints.

%%%%%%%%%%%%%%%%%%%%%%%%%%%%%%%%%%%%%%%%%%%%%%%%%%%%%%%%%%%%%%%

\section{Queue Dynamics}

Routers temporarily store packets awaiting transmission.

Let

\[
\mathcal Q_i(t)
\]

denote the queue length at node

\[
i.
\]

Its evolution satisfies

\[
\boxed{
\frac{d\mathcal Q_i}{dt}
=
A_i
-
D_i,
}
\]

where

\[
A_i
\]

is the arrival rate and

\[
D_i
\]

the departure rate.

Stable operation requires bounded queue lengths.

%%%%%%%%%%%%%%%%%%%%%%%%%%%%%%%%%%%%%%%%%%%%%%%%%%%%%%%%%%%%%%%

\section{Routing}

Suppose

\[
\Gamma
\]

denotes the collection of feasible communication paths.

Routing determines

\[
\gamma^\star
=
\operatorname*{arg\,min}
J(\gamma),
\]

where the objective may include

latency,

bandwidth utilization,

energy consumption,

reliability,

or congestion.

%%%%%%%%%%%%%%%%%%%%%%%%%%%%%%%%%%%%%%%%%%%%%%%%%%%%%%%%%%%%%%%

\section{Network Utility Maximization}

Many communication systems optimize an aggregate utility function.

Let

\[
U_i(F_i)
\]

be the utility of flow

\[
F_i.
\]

The optimization problem becomes

\[
\boxed{
\max
\sum_i
U_i(F_i),
}
\]

subject to

capacity and conservation constraints.

This is mathematically dual to the cost-minimization problems developed
in earlier chapters.

%%%%%%%%%%%%%%%%%%%%%%%%%%%%%%%%%%%%%%%%%%%%%%%%%%%%%%%%%%%%%%%

\section{Reliability}

Communication links may fail due to

hardware faults,

fiber damage,

weather,

or interference.

Let

\[
p_e
\]

denote the probability that edge

\[
e
\]

remains operational.

The reliability of a communication path is

\[
R(\gamma)
=
\prod_{e\in\gamma}
p_e,
\]

assuming statistically independent failures.

%%%%%%%%%%%%%%%%%%%%%%%%%%%%%%%%%%%%%%%%%%%%%%%%%%%%%%%%%%%%%%%

\section{Redundant Routing}

To improve resilience,

multiple disjoint paths may simultaneously connect the same endpoints.

Let

\[
\gamma_1,\ldots,\gamma_k
\]

be edge-disjoint paths.

Failure of any single path therefore does not interrupt communication.

This is a direct application of the redundancy principles developed in
Part~V.

%%%%%%%%%%%%%%%%%%%%%%%%%%%%%%%%%%%%%%%%%%%%%%%%%%%%%%%%%%%%%%%

\section{Wireless Networks}

Wireless communication introduces spatial dependence.

Signal power satisfies approximately

\[
P(r)
\propto
r^{-\alpha},
\]

where

\[
\alpha
\]

is the path-loss exponent.

Communication quality therefore depends explicitly upon the geometry of
the embedded network.

%%%%%%%%%%%%%%%%%%%%%%%%%%%%%%%%%%%%%%%%%%%%%%%%%%%%%%%%%%%%%%%

\section{Satellite Networks}

Satellite communication extends the underlying manifold beyond the
Earth's surface.

The embedded graph now occupies

three-dimensional space,

with nodes following time-dependent orbital trajectories.

Routing therefore becomes both spatially and temporally dynamic.

%%%%%%%%%%%%%%%%%%%%%%%%%%%%%%%%%%%%%%%%%%%%%%%%%%%%%%%%%%%%%%%

\section{Adaptive Network Control}

Modern communication systems continuously adjust

routing,

transmission power,

coding,

channel allocation,

and congestion control.

Model Predictive Control repeatedly updates these decisions using
real-time measurements of network state.

Operational adaptation therefore follows exactly the framework developed
in Part~IV.

%%%%%%%%%%%%%%%%%%%%%%%%%%%%%%%%%%%%%%%%%%%%%%%%%%%%%%%%%%%%%%%

\section{Cyber Resilience}

Communication infrastructure must remain operational despite

equipment failures,

malicious attacks,

software defects,

and routing instability.

Recovery consists of restoring admissible communication paths while
preserving critical services.

The resilience basin,

recovery trajectories,

and robust optimization framework developed earlier apply without
modification.

%%%%%%%%%%%%%%%%%%%%%%%%%%%%%%%%%%%%%%%%%%%%%%%%%%%%%%%%%%%%%%%

\section{Information Geometry}

Unlike previous infrastructure domains,

communication networks transport representations rather than physical
substances.

Consequently,

network design frequently seeks not only efficient transport but also
efficient distribution of information.

Latency,

throughput,

reliability,

and redundancy together determine the informational geometry of the
network.

%%%%%%%%%%%%%%%%%%%%%%%%%%%%%%%%%%%%%%%%%%%%%%%%%%%%%%%%%%%%%%%

\section{Discussion}

Communication systems demonstrate that the unified mathematical
framework extends naturally beyond physical transport.

Electrical systems transport energy.

Hydraulic systems transport fluids.

Transportation systems transport passengers and freight.

Communication systems transport information.

Each domain differs in its governing constitutive equations, yet each is
described by the same mathematical language of

graph topology,

embedded geometry,

conservation,

capacity,

optimization,

control,

uncertainty,

and resilience.

Supply chains and
logistics are next, where multiple commodities interact simultaneously across
large-scale economic networks, providing the first truly
multi-commodity application of the unified theory.

%%%%%%%%%%%%%%%%%%%%%%%%%%%%%%%%%%%%%%%%%%%%%%%%%%%%%%%%%%%%%%%%%%%%%%
\chapter{Supply Chains and Logistics}

\section{Introduction}

Supply chains differ from the infrastructure systems examined in the
preceding chapters in one important respect.

Electrical networks transport a single commodity.

Hydraulic networks transport a single fluid.

Transportation systems primarily transport vehicles, passengers, and
freight.

Communication networks transport information.

Modern logistics systems simultaneously transport many interacting
commodities while coordinating production, storage, transformation, and
distribution across multiple temporal and spatial scales.

Consequently, supply chains constitute the first genuinely
multi-commodity application of the unified infrastructure theory.

%%%%%%%%%%%%%%%%%%%%%%%%%%%%%%%%%%%%%%%%%%%%%%%%%%%%%%%%%%%%%%%

\section{Network Representation}

Represent the logistics system by

\[
G=(V,E).
\]

Vertices correspond to

raw-material suppliers,

factories,

warehouses,

distribution centers,

retail facilities,

ports,

and customers.

Edges represent transportation corridors connecting these facilities.

%%%%%%%%%%%%%%%%%%%%%%%%%%%%%%%%%%%%%%%%%%%%%%%%%%%%%%%%%%%%%%%

\section{Commodity Space}

Suppose

\[
m
\]

distinct commodities are transported.

Define the commodity vector

\[
\mathbf q
=
(q^{(1)},q^{(2)},\ldots,q^{(m)}).
\]

Each component represents the flow of one material through the network.

%%%%%%%%%%%%%%%%%%%%%%%%%%%%%%%%%%%%%%%%%%%%%%%%%%%%%%%%%%%%%%%

\section{Multi-Commodity Conservation}

At every node,

each commodity satisfies

\[
\boxed{
\sum_j
q_{ji}^{(k)}
-
\sum_\ell
q_{i\ell}^{(k)}
=
s_i^{(k)},
}
\]

for every commodity

\[
k.
\]

The conservation laws therefore become a coupled system rather than a
single equation.

%%%%%%%%%%%%%%%%%%%%%%%%%%%%%%%%%%%%%%%%%%%%%%%%%%%%%%%%%%%%%%%

\section{Production}

Factories transform one collection of commodities into another.

Represent production by

\[
\mathbf y
=
T(\mathbf x),
\]

where

\[
T
\]

is the production operator.

Linear approximations take the form

\[
\boxed{
\mathbf y
=
A\mathbf x,
}
\]

where

\[
A
\]

is the production matrix.

%%%%%%%%%%%%%%%%%%%%%%%%%%%%%%%%%%%%%%%%%%%%%%%%%%%%%%%%%%%%%%%

\section{Bills of Materials}

Suppose

\[
a_{ij}
\]

denotes the quantity of component

\[
i
\]

required to produce one unit of product

\[
j.
\]

The matrix

\[
A=(a_{ij})
\]

defines the bill of materials.

Production planning therefore becomes constrained by

\[
A\mathbf p
\le
\mathbf x.
\]

%%%%%%%%%%%%%%%%%%%%%%%%%%%%%%%%%%%%%%%%%%%%%%%%%%%%%%%%%%%%%%%

\section{Inventory Dynamics}

Let

\[
I_i^{(k)}(t)
\]

denote inventory of commodity

\[
k
\]

at facility

\[
i.
\]

Then

\[
\boxed{
\frac{dI_i^{(k)}}{dt}
=
Q_{\rm in}^{(k)}
-
Q_{\rm out}^{(k)}
+
P_i^{(k)}
-
D_i^{(k)},
}
\]

where

\[
P_i^{(k)}
\]

represents production and

\[
D_i^{(k)}
\]

customer demand.

%%%%%%%%%%%%%%%%%%%%%%%%%%%%%%%%%%%%%%%%%%%%%%%%%%%%%%%%%%%%%%%

\section{Capacity Constraints}

Facilities possess finite storage.

Accordingly,

\[
0
\le
I_i^{(k)}
\le
I_{\max}^{(k)}.
\]

Similarly,

transportation corridors satisfy

\[
q^{(k)}
\le
q_{\max}^{(k)}.
\]

%%%%%%%%%%%%%%%%%%%%%%%%%%%%%%%%%%%%%%%%%%%%%%%%%%%%%%%%%%%%%%%

\section{Lead Times}

Every shipment experiences transportation delay.

Let

\[
\tau_e
\]

denote the lead time associated with edge

\[
e.
\]

Delivery at time

\[
t
\]

therefore depends upon decisions made at

\[
t-\tau_e.
\]

Logistics naturally introduces delayed dynamical systems.

%%%%%%%%%%%%%%%%%%%%%%%%%%%%%%%%%%%%%%%%%%%%%%%%%%%%%%%%%%%%%%%

\section{Demand Uncertainty}

Customer demand satisfies

\[
D(t,\omega),
\]

where

\[
\omega
\]

represents uncertainty.

Inventory policies therefore solve stochastic optimization problems of
the form

\[
\boxed{
\min
\mathbb E[J].
}
\]

%%%%%%%%%%%%%%%%%%%%%%%%%%%%%%%%%%%%%%%%%%%%%%%%%%%%%%%%%%%%%%%

\section{Bullwhip Effect}

Small demand fluctuations frequently amplify upstream.

Let

\[
\sigma_D
\]

denote customer-demand variance and

\[
\sigma_O
\]

order variance.

The amplification ratio is

\[
B
=
\frac{\sigma_O^2}
{\sigma_D^2}.
\]

Values

\[
B>1
\]

indicate the bullwhip effect.

Control policies seek to minimize

\[
B.
\]

%%%%%%%%%%%%%%%%%%%%%%%%%%%%%%%%%%%%%%%%%%%%%%%%%%%%%%%%%%%%%%%

\section{Facility Location}

Suppose candidate warehouse locations are

\[
x_1,\ldots,x_n.
\]

Facility placement minimizes

\[
\boxed{
J
=
\sum_i
c_i
+
\sum_j
d(x_j,\mathcal W),
}
\]

where

\[
\mathcal W
\]

denotes the chosen warehouse set.

This is a geometric optimization problem embedded within the underlying
transportation manifold.

%%%%%%%%%%%%%%%%%%%%%%%%%%%%%%%%%%%%%%%%%%%%%%%%%%%%%%%%%%%%%%%

\section{Routing}

Vehicle routing generalizes shortest-path problems.

Let

\[
\Gamma
\]

denote the set of feasible delivery routes.

Optimization seeks

\[
\Gamma^\star
=
\operatorname*{arg\,min}
J(\Gamma),
\]

subject to

vehicle capacity,

delivery windows,

driver constraints,

and inventory requirements.

%%%%%%%%%%%%%%%%%%%%%%%%%%%%%%%%%%%%%%%%%%%%%%%%%%%%%%%%%%%%%%%

\section{Integrated Logistics}

Transportation,

inventory,

production,

and facility planning cannot generally be optimized independently.

Instead,

they satisfy coupled equations,

and the complete optimization becomes

\[
\boxed{
\min_{\Gamma}
J_{\rm transport}
+
J_{\rm inventory}
+
J_{\rm production}
+
J_{\rm facilities}.
}
\]

This is one of the clearest examples of the unified optimization
principle developed earlier.

%%%%%%%%%%%%%%%%%%%%%%%%%%%%%%%%%%%%%%%%%%%%%%%%%%%%%%%%%%%%%%%

\section{Resilience}

Supply-chain disruptions include

factory shutdowns,

port closures,

transport interruptions,

component shortages,

and geopolitical instability.

Recovery consists of identifying alternative suppliers,

rerouting transportation,

reallocating inventories,

and modifying production schedules while maintaining essential service.

The continuation framework developed previously provides a natural
mathematical description of these adaptive recovery processes.

%%%%%%%%%%%%%%%%%%%%%%%%%%%%%%%%%%%%%%%%%%%%%%%%%%%%%%%%%%%%%%%

\section{Discussion}

Supply chains unify nearly every mathematical idea introduced throughout
this monograph.

Graph theory describes connectivity.

Geometry determines facility placement and transportation corridors.

Conservation laws govern inventories and material flow.

Transformation operators model manufacturing.

Optimal control schedules production and transportation.

Stochastic optimization addresses uncertain demand.

Robust optimization prepares for disruption.

Resilience theory guides recovery.

Unlike the previous application chapters, logistics demonstrates that
modern infrastructure is not merely a transport network but an
interconnected system in which transportation, storage,
transformation, and information continuously interact across multiple
timescales.

This perspective extends further in the next chapter's examination of
\emph{integrated infrastructure systems}, in which electrical,
hydraulic, transportation, communication, and logistics networks are
coupled into a single interdependent mathematical system.

%%%%%%%%%%%%%%%%%%%%%%%%%%%%%%%%%%%%%%%%%%%%%%%%%%%%%%%%%%%%%%%%%%%%%%
\chapter{Integrated Infrastructure Systems}

\section{Introduction}

No major infrastructure system operates in isolation.

Electrical power supplies communication networks.

Communication systems coordinate transportation.

Transportation distributes fuel for electrical generation.

Water systems require electrical power for pumping and treatment.

Supply chains depend simultaneously upon transportation,
communications,
energy,
and water.

The mathematical object of interest is therefore not an isolated
network, but an interconnected collection of mutually dependent
networks.

This chapter develops a unified mathematical framework for such
\emph{network-of-networks} systems.

%%%%%%%%%%%%%%%%%%%%%%%%%%%%%%%%%%%%%%%%%%%%%%%%%%%%%%%%%%%%%%%

\section{Infrastructure Layers}

Suppose

\[
\mathcal N
=
\{G_1,G_2,\ldots,G_n\},
\]

where each

\[
G_i
\]

represents an infrastructure layer.

Typical examples include

\[
G_E
\]

for electrical power,

\[
G_W
\]

for water,

\[
G_T
\]

for transportation,

\[
G_C
\]

for communications,

and

\[
G_L
\]

for logistics.

%%%%%%%%%%%%%%%%%%%%%%%%%%%%%%%%%%%%%%%%%%%%%%%%%%%%%%%%%%%%%%%

\section{Interdependency}

Infrastructure layers interact through dependency mappings

\[
\boxed{
D_{ij}:G_i
\rightarrow
G_j.
}
\]

For example,

electrical substations provide power to communication towers,

while communication networks coordinate electrical dispatch.

The dependency graph itself therefore forms another network.

%%%%%%%%%%%%%%%%%%%%%%%%%%%%%%%%%%%%%%%%%%%%%%%%%%%%%%%%%%%%%%%

\section{The Supra-Graph}

Define the integrated infrastructure graph

\[
\boxed{
\mathcal G
=
(V,E,E_D),
}
\]

where

\[
E
\]

contains ordinary infrastructure connections,

and

\[
E_D
\]

contains inter-layer dependency edges.

This structure is commonly called a supra-graph.

%%%%%%%%%%%%%%%%%%%%%%%%%%%%%%%%%%%%%%%%%%%%%%%%%%%%%%%%%%%%%%%

\section{Layer States}

Let

\[
x_i(t)
\]

denote the state vector of infrastructure layer

\[
i.
\]

The complete infrastructure state becomes

\[
\boxed{
X(t)
=
(x_1,x_2,\ldots,x_n).
}
\]

%%%%%%%%%%%%%%%%%%%%%%%%%%%%%%%%%%%%%%%%%%%%%%%%%%%%%%%%%%%%%%%

\section{Coupled Dynamics}

Individual infrastructure layers satisfy

\[
\dot x_i
=
f_i(x_i,u_i).
\]

Interdependence modifies these equations to

\[
\boxed{
\dot x_i
=
f_i
(x_i,u_i)
+
\sum_{j\neq i}
C_{ij}(x_j),
}
\]

where

\[
C_{ij}
\]

represents coupling between infrastructure systems.

%%%%%%%%%%%%%%%%%%%%%%%%%%%%%%%%%%%%%%%%%%%%%%%%%%%%%%%%%%%%%%%

\section{Dependency Matrix}

Coupling strengths may be collected into the matrix

\[
\boxed{
A
=
(a_{ij}),
}
\]

where

\[
a_{ij}>0
\]

indicates that infrastructure

\[
i
\]

depends upon infrastructure

\[
j.
\]

The dependency matrix characterizes the architecture of the entire
civil infrastructure system.

%%%%%%%%%%%%%%%%%%%%%%%%%%%%%%%%%%%%%%%%%%%%%%%%%%%%%%%%%%%%%%%

\section{Failure Propagation}

Suppose

\[
F_i(t)
\]

denotes failure within infrastructure layer

\[
i.
\]

A simple propagation model is

\[
\boxed{
\frac{dF_i}{dt}
=
-\alpha_iF_i
+
\sum_j
\beta_{ij}F_j,
}
\]

where

\[
\alpha_i
\]

represents local recovery,

and

\[
\beta_{ij}
\]

represents failure transmission.

Failure therefore behaves as a dynamical process on the dependency
network.

%%%%%%%%%%%%%%%%%%%%%%%%%%%%%%%%%%%%%%%%%%%%%%%%%%%%%%%%%%%%%%%

\section{Cascade Thresholds}

The integrated system may exhibit cascading failures.

Let

\[
\rho
\]

denote the spectral radius of the dependency matrix.

Large values of

\[
\rho
\]

indicate stronger coupling and increased susceptibility to cascading
disruptions.

%%%%%%%%%%%%%%%%%%%%%%%%%%%%%%%%%%%%%%%%%%%%%%%%%%%%%%%%%%%%%%%

\begin{theorem}[Qualitative Cascade Criterion]

If recovery dominates dependency propagation,

small localized failures remain localized.

If dependency propagation dominates recovery,

localized failures may propagate through multiple infrastructure layers.

\end{theorem}

%%%%%%%%%%%%%%%%%%%%%%%%%%%%%%%%%%%%%%%%%%%%%%%%%%%%%%%%%%%%%%%

\section{Resource Allocation}

Suppose resilience resources satisfy

\[
R_{\rm total}.
\]

Allocation variables

\[
r_i
\]

must satisfy

\[
\sum_i
r_i
=
R_{\rm total}.
\]

The optimization problem becomes

\[
\boxed{
\min
J(X,r),
}
\]

subject to coupled infrastructure dynamics.

%%%%%%%%%%%%%%%%%%%%%%%%%%%%%%%%%%%%%%%%%%%%%%%%%%%%%%%%%%%%%%%

\section{Joint Control}

Each infrastructure possesses independent controls

\[
u_i.
\]

Integrated management introduces the joint control vector

\[
\boxed{
u
=
(u_1,u_2,\ldots,u_n).
}
\]

Optimal operation minimizes

\[
J(X,u),
\]

rather than optimizing each infrastructure independently.

%%%%%%%%%%%%%%%%%%%%%%%%%%%%%%%%%%%%%%%%%%%%%%%%%%%%%%%%%%%%%%%

\section{Shared Information}

Integrated management requires continual exchange of information between
infrastructure layers.

Measurements

\[
y_i(t)
\]

from each subsystem are collected into

\[
Y(t)
=
(y_1,\ldots,y_n).
\]

State estimation therefore becomes a distributed inference problem over
the entire infrastructure network.

%%%%%%%%%%%%%%%%%%%%%%%%%%%%%%%%%%%%%%%%%%%%%%%%%%%%%%%%%%%%%%%

\section{Resilience of Interdependent Systems}

Recovery is no longer purely local.

Restoring electrical service may restore communications.

Restoring communications may accelerate transportation recovery.

Transportation recovery may restore fuel delivery.

Infrastructure repair therefore becomes a coupled optimization problem
rather than a collection of isolated repairs.

%%%%%%%%%%%%%%%%%%%%%%%%%%%%%%%%%%%%%%%%%%%%%%%%%%%%%%%%%%%%%%%

\section{Hierarchical Coordination}

Integrated infrastructure naturally operates across several levels.

Operational controllers regulate individual facilities.

Regional controllers coordinate infrastructure sectors.

National controllers allocate strategic resources.

Long-term planning modifies infrastructure architecture itself.

These levels correspond directly to the multiscale planning hierarchy
developed earlier in this volume.

%%%%%%%%%%%%%%%%%%%%%%%%%%%%%%%%%%%%%%%%%%%%%%%%%%%%%%%%%%%%%%%

\section{Digital Twins}

Modern infrastructure increasingly maintains a continuously updated
computational representation of the physical system.

Let

\[
\hat X(t)
\]

denote the estimated system state.

Measurements continuously update

\[
\hat X(t),
\]

while predictive models forecast future evolution.

Planning therefore becomes optimization over both the physical system
and its digital representation.

%%%%%%%%%%%%%%%%%%%%%%%%%%%%%%%%%%%%%%%%%%%%%%%%%%%%%%%%%%%%%%%

\section{Discussion}

Integrated infrastructure systems represent the culmination of the
mathematical framework developed throughout this monograph.

Electrical,
hydraulic,
transportation,
communication,
and logistics networks become components of a single coupled dynamical
system evolving over a shared geometric substrate.

Optimization,
control,
resilience,
and adaptation are therefore properties of the infrastructure ensemble
rather than of its individual parts.

This framework extends, finally, beyond terrestrial systems
to planetary and interplanetary infrastructure, where the underlying
geometry, communication delays, and resource constraints introduce
entirely new mathematical challenges while preserving the same unified
theory.

%%%%%%%%%%%%%%%%%%%%%%%%%%%%%%%%%%%%%%%%%%%%%%%%%%%%%%%%%%%%%%%%%%%%%%
\chapter{Planetary and Interplanetary Infrastructure}

\section{Introduction}

The mathematical framework developed throughout this volume is
independent of the scale upon which infrastructure operates.

Whether the network spans a city,
a continent,
or an entire planet,
its essential components remain

geometry,

topology,

transport,

storage,

control,

optimization,

uncertainty,

and resilience.

Extending infrastructure beyond Earth introduces new physical
constraints but does not alter the underlying mathematical principles.

Long communication delays,
orbital dynamics,
limited launch capacity,
radiation,
and autonomous operation become additional constraints within the same
unified framework.

%%%%%%%%%%%%%%%%%%%%%%%%%%%%%%%%%%%%%%%%%%%%%%%%%%%%%%%%%%%%%%%

\section{Infrastructure Manifold}

Unlike terrestrial systems, whose geometry is largely constrained to the
Earth's surface, planetary infrastructure occupies multiple embedded
manifolds.

Let

\[
M
=
M_E
\cup
M_O
\cup
M_L
\cup
M_M,
\]

where

\[
M_E
\]

denotes the Earth's surface,

\[
M_O
\]

Earth orbit,

\[
M_L
\]

the lunar surface,

and

\[
M_M
\]

the Martian surface.

Additional celestial bodies may be incorporated without altering the
formalism.

%%%%%%%%%%%%%%%%%%%%%%%%%%%%%%%%%%%%%%%%%%%%%%%%%%%%%%%%%%%%%%%

\section{Interplanetary Graph}

Define

\[
G=(V,E),
\]

where vertices represent

cities,

launch facilities,

orbital stations,

surface habitats,

mining installations,

communication relays,

and scientific outposts.

Edges represent

roads,

railways,

launch trajectories,

orbital transfers,

communication links,

and cargo routes.

Unlike terrestrial infrastructure, some edges exist only during
specific launch windows or orbital alignments.

%%%%%%%%%%%%%%%%%%%%%%%%%%%%%%%%%%%%%%%%%%%%%%%%%%%%%%%%%%%%%%%

\section{Time-Dependent Connectivity}

Interplanetary connectivity varies continuously with orbital motion.

Accordingly,

\[
E=E(t).
\]

The infrastructure graph therefore becomes a time-dependent graph,

requiring routing algorithms over evolving topologies rather than fixed
networks.

%%%%%%%%%%%%%%%%%%%%%%%%%%%%%%%%%%%%%%%%%%%%%%%%%%%%%%%%%%%%%%%

\section{Orbital Transfer Costs}

Each interplanetary trajectory possesses an associated transfer cost.

Let

\[
\Delta v(\gamma)
\]

denote the required velocity change along trajectory

\[
\gamma.
\]

A simplified transport objective becomes

\[
\boxed{
J(\gamma)
=
\alpha
\Delta v(\gamma)
+
\beta
T(\gamma)
+
\gamma
R(\gamma),
}
\]

where

\[
T
\]

is travel time and

\[
R
\]

represents mission risk.

%%%%%%%%%%%%%%%%%%%%%%%%%%%%%%%%%%%%%%%%%%%%%%%%%%%%%%%%%%%%%%%

\section{Launch Windows}

Certain transportation corridors exist only during finite intervals.

Let

\[
W_i
=
[t_i^{(1)},t_i^{(2)}].
\]

A trajectory is admissible only if departure occurs within an
appropriate launch window.

Planning therefore becomes a hybrid optimization problem over both space
and time.

%%%%%%%%%%%%%%%%%%%%%%%%%%%%%%%%%%%%%%%%%%%%%%%%%%%%%%%%%%%%%%%

\section{Communication Delay}

Unlike terrestrial communication,

signals require significant propagation time.

Let

\[
d(x,y)
\]

be the spatial separation between two nodes.

Communication latency satisfies

\[
\boxed{
\tau
=
\frac{d(x,y)}{c},
}
\]

where

\[
c
\]

is the speed of light.

Control architectures must therefore tolerate unavoidable delay.

%%%%%%%%%%%%%%%%%%%%%%%%%%%%%%%%%%%%%%%%%%%%%%%%%%%%%%%%%%%%%%%

\section{Autonomous Operation}

Because direct intervention is impossible during long communication
delays,

remote infrastructure must operate autonomously.

Control inputs satisfy

\[
u(t)
=
\pi(\hat x(t)),
\]

where

\[
\pi
\]

is a locally implemented control policy.

Autonomy therefore becomes a mathematical necessity rather than a matter
of convenience.

%%%%%%%%%%%%%%%%%%%%%%%%%%%%%%%%%%%%%%%%%%%%%%%%%%%%%%%%%%%%%%%

\section{In-Situ Resource Utilization}

Transporting every required resource from Earth is generally
uneconomical.

Suppose

\[
R_i
\]

denotes locally available resources.

Infrastructure planning seeks to maximize

\[
\boxed{
\sum_i
R_i,
}
\]

while minimizing imported mass.

Production operators therefore acquire increased importance within
planetary logistics.

%%%%%%%%%%%%%%%%%%%%%%%%%%%%%%%%%%%%%%%%%%%%%%%%%%%%%%%%%%%%%%%

\section{Energy Networks}

Extraterrestrial settlements require reliable energy production.

Possible generation sources include

solar arrays,

nuclear reactors,

geothermal systems,

and stored chemical fuels.

Energy distribution remains governed by the conservation and network
optimization principles developed in Chapter~\ref{chap:electrical-power-networks}.

%%%%%%%%%%%%%%%%%%%%%%%%%%%%%%%%%%%%%%%%%%%%%%%%%%%%%%%%%%%%%%%

\section{Closed Ecological Systems}

Habitats must recycle matter continuously.

Let

\[
M_{\rm air},
\qquad
M_{\rm water},
\qquad
M_{\rm food}
\]

represent essential resource inventories.

Their evolution satisfies

\[
\boxed{
\frac{dM}{dt}
=
P
-
C
+
R,
}
\]

where

\[
P
\]

is production,

\[
C
\]

consumption,

and

\[
R
\]

recycling.

Closed-loop operation becomes an optimization objective rather than an
engineering preference.

%%%%%%%%%%%%%%%%%%%%%%%%%%%%%%%%%%%%%%%%%%%%%%%%%%%%%%%%%%%%%%%

\section{Resilience}

Extraterrestrial infrastructure possesses limited repair capacity.

Failures may threaten the survival of an entire settlement.

Consequently,

redundancy,

fault isolation,

graceful degradation,

and recoverability become primary design objectives.

The resilience functionals introduced in Part~V therefore assume even
greater significance.

%%%%%%%%%%%%%%%%%%%%%%%%%%%%%%%%%%%%%%%%%%%%%%%%%%%%%%%%%%%%%%%

\section{Distributed Civilizations}

Future infrastructure may consist of numerous partially independent
settlements connected by transportation and communication networks.

Let

\[
G_C
\]

denote the civilization graph.

Vertices correspond to settlements,

while edges represent regular exchange of

people,

materials,

energy,

scientific knowledge,

and manufactured goods.

Long-term continuity depends upon maintaining sufficient connectivity
despite large physical separation.

%%%%%%%%%%%%%%%%%%%%%%%%%%%%%%%%%%%%%%%%%%%%%%%%%%%%%%%%%%%%%%%

\section{Unified Planetary Optimization}

The complete planning problem becomes

\[
\boxed{
\min_{\Gamma}
J_{\rm transport}
+
J_{\rm energy}
+
J_{\rm communication}
+
J_{\rm production}
+
J_{\rm resilience},
}
\]

subject to

orbital mechanics,

resource conservation,

launch constraints,

communication delays,

and environmental limitations.

This optimization differs in scale rather than mathematical structure
from the infrastructure problems considered throughout the remainder of
this volume.

%%%%%%%%%%%%%%%%%%%%%%%%%%%%%%%%%%%%%%%%%%%%%%%%%%%%%%%%%%%%%%%

\section{Discussion}

Planetary and interplanetary infrastructure demonstrate the scale
independence of the unified mathematical framework.

The same principles governing a municipal water network or an electrical
transmission grid also govern transportation between planets, provided
the governing physical constraints are modified appropriately.

Geometry,

topology,

conservation,

optimization,

control,

adaptation,

and resilience remain the fundamental organizing principles regardless
of spatial scale.

%%%%%%%%%%%%%%%%%%%%%%%%%%%%%%%%%%%%%%%%%%%%%%%%%%%%%%%%%%%%%%%

\section{Concluding Remarks}

The progression of this monograph has intentionally moved from local
geometric structures to infrastructure systems spanning entire
civilizations.

At every stage, increasingly complex applications have emerged by adding
new physical constraints while preserving a common mathematical core.

This observation motivates the central conclusion of the volume:

Infrastructure science is not fundamentally the study of roads, power
grids, pipelines, or communication systems individually.

It is the study of constrained transport, transformation, storage,
coordination, and continuation on embedded networks.

The engineering disciplines differ primarily in their constitutive laws;
their mathematical architecture is shared.

%%%%%%%%%%%%%%%%%%%%%%%%%%%%%%%%%%%%%%%%%%%%%%%%%%%%%%%%%%%%%%%%%%%%%%
\part{Foundations and Outlook}

%%%%%%%%%%%%%%%%%%%%%%%%%%%%%%%%%%%%%%%%%%%%%%%%%%%%%%%%%%%%%%%%%%%%%%
\chapter{Unifying Principles of Infrastructure Science}

\section{Introduction}

The preceding chapters have developed a mathematical theory of
infrastructure beginning with geometry, continuing through network
topology, conservation laws, optimization, control, uncertainty,
resilience, and finally large-scale integrated systems.

Although the subject matter has ranged from roads and pipelines to
planetary logistics, a remarkably small collection of mathematical
principles has appeared repeatedly.

The purpose of this chapter is to state those principles explicitly.

%%%%%%%%%%%%%%%%%%%%%%%%%%%%%%%%%%%%%%%%%%%%%%%%%%%%%%%%%%%%%%%

\section{Principle I: Geometry Precedes Infrastructure}

Every infrastructure system occupies space.

Distances,

curvature,

topography,

and physical embedding determine which connections are feasible before
optimization begins.

Infrastructure therefore cannot be separated from geometry.

Geometry is not an afterthought.

It is the domain upon which every subsequent optimization is defined.

%%%%%%%%%%%%%%%%%%%%%%%%%%%%%%%%%%%%%%%%%%%%%%%%%%%%%%%%%%%%%%%

\section{Principle II: Topology Determines Possibility}

Geometry determines where construction is possible.

Topology determines what can be connected.

Connectivity,

cycles,

cuts,

redundancy,

and accessibility arise from the topology of the network rather than
its physical realization.

Many operational questions therefore reduce first to topological
questions.

%%%%%%%%%%%%%%%%%%%%%%%%%%%%%%%%%%%%%%%%%%%%%%%%%%%%%%%%%%%%%%%

\section{Principle III: Conservation Governs Flow}

Every infrastructure transports something.

Electrical grids transport energy.

Hydraulic systems transport fluids.

Transportation networks transport passengers and freight.

Communication networks transport information.

Supply chains transport material.

Although the transported quantities differ,

each satisfies a conservation law modified only by sources,
sinks,
or transformation operators.

%%%%%%%%%%%%%%%%%%%%%%%%%%%%%%%%%%%%%%%%%%%%%%%%%%%%%%%%%%%%%%%

\section{Principle IV: Capacity Limits Throughput}

No infrastructure possesses infinite capacity.

Transmission lines overheat.

Roads become congested.

Pipelines experience pressure limits.

Communication channels saturate.

Warehouses become full.

Capacity therefore defines the feasible operating region of every
infrastructure system.

%%%%%%%%%%%%%%%%%%%%%%%%%%%%%%%%%%%%%%%%%%%%%%%%%%%%%%%%%%%%%%%

\section{Principle V: Optimization Allocates Scarce Resources}

Infrastructure planning is fundamentally an optimization problem.

Resources,

time,

space,

energy,

labor,

and capital are all finite.

Optimization therefore determines how limited resources should be
allocated while satisfying physical and operational constraints.

%%%%%%%%%%%%%%%%%%%%%%%%%%%%%%%%%%%%%%%%%%%%%%%%%%%%%%%%%%%%%%%

\section{Principle VI: Control Maintains Operation}

Infrastructure is never static.

Demand changes.

Equipment ages.

Weather varies.

Unexpected failures occur.

Control theory provides the mathematical machinery required to maintain
admissible operation despite continual disturbance.

%%%%%%%%%%%%%%%%%%%%%%%%%%%%%%%%%%%%%%%%%%%%%%%%%%%%%%%%%%%%%%%

\section{Principle VII: Uncertainty Is Fundamental}

No infrastructure planner possesses perfect information.

Future demand,

weather,

equipment failures,

economic conditions,

and human behavior remain uncertain.

Planning therefore requires probabilistic reasoning,

robust optimization,

or adaptive control rather than deterministic prediction alone.

%%%%%%%%%%%%%%%%%%%%%%%%%%%%%%%%%%%%%%%%%%%%%%%%%%%%%%%%%%%%%%%

\section{Principle VIII: Resilience Preserves Continuity}

Optimal operation is insufficient.

Infrastructure must continue functioning when optimization becomes
impossible.

Resilience therefore measures not efficiency under ideal conditions but
continuation under adverse conditions.

Recovery,

redundancy,

graceful degradation,

and repair are fundamental mathematical properties rather than merely
engineering considerations.

%%%%%%%%%%%%%%%%%%%%%%%%%%%%%%%%%%%%%%%%%%%%%%%%%%%%%%%%%%%%%%%

\section{Principle IX: Adaptation Outperforms Prediction}

No sufficiently complex infrastructure can be planned once and left
unchanged indefinitely.

Measurements continually update models.

Controllers continually revise decisions.

Infrastructure therefore evolves through repeated observation and
adaptation.

Planning is a process rather than a single computation.

%%%%%%%%%%%%%%%%%%%%%%%%%%%%%%%%%%%%%%%%%%%%%%%%%%%%%%%%%%%%%%%

\section{Principle X: Integration Creates Civilization}

Individual infrastructure systems rarely operate independently.

Energy,

water,

transportation,

communications,

and logistics continually influence one another.

Civilization emerges not from isolated networks but from their
coordination.

Infrastructure science therefore studies systems of systems rather than
individual technologies.

%%%%%%%%%%%%%%%%%%%%%%%%%%%%%%%%%%%%%%%%%%%%%%%%%%%%%%%%%%%%%%%

\section{A Unified Mathematical View}

The principles developed throughout the preceding chapters may be synthesized into a single mathematical optimization framework. Although each chapter has emphasized a different aspect of infrastructure science, they are ultimately components of a common variational problem whose objective is to identify infrastructure histories that remain feasible, efficient, and robust under the full collection of physical, informational, and operational constraints.

Let

\[
\mathcal H
\]

denote the set of all admissible infrastructure histories.

The central objective of infrastructure science is then to determine

\[
\boxed{
\Gamma^\star
=
\operatorname*{arg\,min}_{\Gamma\in\mathcal H}
J(\Gamma),
}
\]

where the admissible set is defined not by any single engineering criterion but by the simultaneous satisfaction of geometric, topological, conservation, capacity, control, uncertainty, and resilience constraints. The cost functional \(J\) therefore represents the aggregate performance of an infrastructure history over its lifetime, while the admissibility conditions determine which histories are physically, operationally, and mathematically realizable.

From this perspective, every chapter of the present volume has contributed one element of the broader optimization problem. Rather than constituting independent subjects, the individual mathematical formalisms collectively specify the constraints, objective functions, and governing principles required to characterize optimal infrastructure systems within a single unified framework.

%%%%%%%%%%%%%%%%%%%%%%%%%%%%%%%%%%%%%%%%%%%%%%%%%%%%%%%%%%%%%%%

\section{Closing Perspective}

Infrastructure is traditionally described in terms of its physical manifestations: roads, bridges, pipelines, electrical grids, ports, railways, satellites, and communication networks. While these systems differ substantially in their materials, operating environments, and engineering requirements, the mathematical perspective developed throughout this volume suggests that they should not be regarded as fundamentally distinct domains of study.

Instead, these diverse systems may be understood as alternative physical realizations of a common mathematical architecture. Each consists of interconnected components that transport matter, energy, information, or people across embedded networks while operating under constraints imposed by geometry, topology, conservation laws, finite capacity, control mechanisms, uncertainty, and continual adaptation. The apparent diversity of infrastructure therefore reflects differences in implementation rather than differences in underlying mathematical structure.

From this viewpoint, infrastructure science may be understood as the study of constrained flow on embedded networks evolving through time under physical law, operational control, uncertainty, and adaptive response. The specific engineering principles governing highways, electrical transmission systems, communication networks, water distribution systems, or transportation corridors necessarily differ, yet the mathematical language required to analyze their behavior is remarkably uniform.

The engineering details are domain-specific. The underlying mathematics is shared.

\newpage 
\section*{Appendices}

%%%%%%%%%%%%%%%%%%%%%%%%%%%%%%%%%%%%%%%%%%%%%%%%%%%%%%%%%%%%%%%%%%%%%%
\appendix
\part{Mathematical Foundations}

%%%%%%%%%%%%%%%%%%%%%%%%%%%%%%%%%%%%%%%%%%%%%%%%%%%%%%%%%%%%%%%%%%%%%%
\chapter{Foundations of Sets and Structures}

\section{Purpose}

Throughout this monograph, infrastructure systems are modeled using
graphs, manifolds, functions, dynamical systems, and optimization
problems. Each of these objects is ultimately constructed from the
language of sets and mappings.

This appendix establishes the mathematical conventions used throughout
the remainder of the volume. Rather than providing a comprehensive
introduction to set theory, the objective is to define the structures
that recur repeatedly in later chapters.

%%%%%%%%%%%%%%%%%%%%%%%%%%%%%%%%%%%%%%%%%%%%%%%%%%%%%%%%%%%%%%%

\section{Sets}

A set is a collection of distinct objects, called its elements.

If $x$ belongs to a set $A$, we write

\[
x\in A.
\]

If $x$ is not an element,

\[
x\notin A.
\]

Equality of sets is defined by extensionality.

\[
A=B
\iff
(\forall x)(x\in A \Leftrightarrow x\in B).
\]

%%%%%%%%%%%%%%%%%%%%%%%%%%%%%%%%%%%%%%%%%%%%%%%%%%%%%%%%%%%%%%%

\section{Subsets}

A set $A$ is a subset of $B$ if

\[
A\subseteq B
\]

whenever

\[
x\in A
\Longrightarrow
x\in B.
\]

Proper subsets satisfy

\[
A\subsetneq B.
\]

%%%%%%%%%%%%%%%%%%%%%%%%%%%%%%%%%%%%%%%%%%%%%%%%%%%%%%%%%%%%%%%

\section{Elementary Operations}

For sets $A$ and $B$

Union

\[
A\cup B
=
\{x:x\in A
\text{ or }
x\in B\}.
\]

Intersection

\[
A\cap B
=
\{x:x\in A
\text{ and }
x\in B\}.
\]

Difference

\[
A\setminus B
=
\{x:x\in A,x\notin B\}.
\]

Complement

\[
A^c
=
U\setminus A,
\]

where $U$ denotes the ambient universe.

%%%%%%%%%%%%%%%%%%%%%%%%%%%%%%%%%%%%%%%%%%%%%%%%%%%%%%%%%%%%%%%

\section{Cartesian Products}

Given sets

\[
A_1,\ldots,A_n,
\]

their Cartesian product is

\[
A_1\times\cdots\times A_n
=
\{
(a_1,\ldots,a_n)
:
a_i\in A_i
\}.
\]

Nearly every state space in this monograph is constructed as a Cartesian
product.

%%%%%%%%%%%%%%%%%%%%%%%%%%%%%%%%%%%%%%%%%%%%%%%%%%%%%%%%%%%%%%%

\section{Relations}

A binary relation on a set $X$ is a subset

\[
R
\subseteq
X\times X.
\]

We write

\[
xRy
\]

whenever

\[
(x,y)\in R.
\]

%%%%%%%%%%%%%%%%%%%%%%%%%%%%%%%%%%%%%%%%%%%%%%%%%%%%%%%%%%%%%%%

\section{Equivalence Relations}

A relation is an equivalence relation if it is

\begin{enumerate}
\item reflexive,
\item symmetric,
\item transitive.
\end{enumerate}

Equivalence relations partition sets into disjoint equivalence classes.

%%%%%%%%%%%%%%%%%%%%%%%%%%%%%%%%%%%%%%%%%%%%%%%%%%%%%%%%%%%%%%%

\begin{theorem}[Partition Theorem]

Every equivalence relation induces a unique partition of the underlying
set, and every partition determines an equivalence relation.

\end{theorem}

%%%%%%%%%%%%%%%%%%%%%%%%%%%%%%%%%%%%%%%%%%%%%%%%%%%%%%%%%%%%%%%

\section{Partial Orders}

A partial order satisfies

\begin{enumerate}
\item reflexivity,
\item antisymmetry,
\item transitivity.
\end{enumerate}

The notation

\[
x\preceq y
\]

denotes an ordered pair.

Many optimization problems naturally define partially ordered feasible
sets.

%%%%%%%%%%%%%%%%%%%%%%%%%%%%%%%%%%%%%%%%%%%%%%%%%%%%%%%%%%%%%%%

\section{Functions}

A function

\[
f:A\rightarrow B
\]

assigns every element of $A$ a unique element of $B$.

The set

\[
A
\]

is the domain,

while

\[
B
\]

is the codomain.

%%%%%%%%%%%%%%%%%%%%%%%%%%%%%%%%%%%%%%%%%%%%%%%%%%%%%%%%%%%%%%%

\section{Injective, Surjective, and Bijective Maps}

A function is injective if

\[
f(x)=f(y)
\Longrightarrow
x=y.
\]

It is surjective if

\[
f(A)=B.
\]

It is bijective if both properties hold.

%%%%%%%%%%%%%%%%%%%%%%%%%%%%%%%%%%%%%%%%%%%%%%%%%%%%%%%%%%%%%%%

\section{Composition}

If

\[
f:A\rightarrow B,
\qquad
g:B\rightarrow C,
\]

their composition is

\[
(g\circ f)(x)
=
g(f(x)).
\]

Composition is associative.

%%%%%%%%%%%%%%%%%%%%%%%%%%%%%%%%%%%%%%%%%%%%%%%%%%%%%%%%%%%%%%%

\begin{theorem}

For compatible functions,

\[
h\circ(g\circ f)
=
(h\circ g)\circ f.
\]

\end{theorem}

%%%%%%%%%%%%%%%%%%%%%%%%%%%%%%%%%%%%%%%%%%%%%%%%%%%%%%%%%%%%%%%

\section{Inverse Functions}

If

\[
f
\]

is bijective,

there exists a unique inverse

\[
f^{-1}
\]

satisfying

\[
f^{-1}(f(x))
=
x.
\]

%%%%%%%%%%%%%%%%%%%%%%%%%%%%%%%%%%%%%%%%%%%%%%%%%%%%%%%%%%%%%%%

\section{Families of Sets}

Infrastructure systems frequently involve indexed collections.

Let

\[
\{A_i\}_{i\in I}
\]

be a family of sets.

Their union is

\[
\bigcup_{i\in I}A_i,
\]

and their intersection is

\[
\bigcap_{i\in I}A_i.
\]

%%%%%%%%%%%%%%%%%%%%%%%%%%%%%%%%%%%%%%%%%%%%%%%%%%%%%%%%%%%%%%%

\section{Sequences}

A sequence is a function

\[
x:\mathbb N\rightarrow X.
\]

We denote

\[
x(n)
=
x_n.
\]

Sequences form the basis of iterative optimization algorithms.

%%%%%%%%%%%%%%%%%%%%%%%%%%%%%%%%%%%%%%%%%%%%%%%%%%%%%%%%%%%%%%%

\section{Metric Spaces}

A metric on a set $X$ is a function

\[
d:X\times X\rightarrow\mathbb R
\]

satisfying

\begin{enumerate}
\item positivity,
\item symmetry,
\item identity of indiscernibles,
\item triangle inequality.
\end{enumerate}

The pair

\[
(X,d)
\]

is called a metric space.

%%%%%%%%%%%%%%%%%%%%%%%%%%%%%%%%%%%%%%%%%%%%%%%%%%%%%%%%%%%%%%%

\section{Topological Spaces}

Every metric induces a topology.

The concepts of

continuity,

compactness,

connectedness,

and convergence

used throughout this volume are fundamentally topological.

%%%%%%%%%%%%%%%%%%%%%%%%%%%%%%%%%%%%%%%%%%%%%%%%%%%%%%%%%%%%%%%

\section{Conclusion}

The constructions introduced here—sets, relations, functions, products,
metrics, and topological structures—form the mathematical language from
which the remainder of this monograph is built. Subsequent appendices
develop increasingly specialized structures, including linear algebra,
differential geometry, graph theory, variational calculus, and
functional analysis, all of which rely on these foundational concepts.

%%%%%%%%%%%%%%%%%%%%%%%%%%%%%%%%%%%%%%%%%%%%%%%%%%%%%%%%%%%%%%%

\section{Functions Between Structured Sets}

Many mathematical objects encountered throughout this volume are more
than merely sets. They possess additional structure, such as topology,
algebraic operations, differentiable structure, or metric properties.

When mapping between such objects, one generally requires the mapping to
preserve the relevant structure.

Examples include

\begin{itemize}
\item continuous maps between topological spaces,
\item linear maps between vector spaces,
\item smooth maps between manifolds,
\item isometries between metric spaces,
\item graph homomorphisms between graphs.
\end{itemize}

Preservation of structure is one of the central themes of modern
mathematics and appears repeatedly throughout infrastructure science.

%%%%%%%%%%%%%%%%%%%%%%%%%%%%%%%%%%%%%%%%%%%%%%%%%%%%%%%%%%%%%%%

\section{Commutative Diagrams}

Relationships among functions are frequently summarized using
commutative diagrams.

Suppose

\[
f:X\rightarrow Y,
\qquad
g:X\rightarrow Z,
\qquad
h:Y\rightarrow Z.
\]

If

\[
g
=
h\circ f,
\]

then the diagram

\[
\begin{array}{ccc}
X & \xrightarrow{f} & Y\\
\downarrow g & & \downarrow h\\
Z & = & Z
\end{array}
\]

is said to commute.

Throughout this monograph, many optimization and projection operators
are naturally expressed by commutative diagrams.

%%%%%%%%%%%%%%%%%%%%%%%%%%%%%%%%%%%%%%%%%%%%%%%%%%%%%%%%%%%%%%%

\section{Products and Projections}

Given

\[
X
=
X_1\times\cdots\times X_n,
\]

the projection onto the

\[
i
\]

th coordinate is

\[
\pi_i:X\rightarrow X_i,
\]

defined by

\[
\pi_i(x_1,\ldots,x_n)
=
x_i.
\]

Conversely,

given functions

\[
f_i:Y\rightarrow X_i,
\]

there exists a unique product mapping

\[
f:Y\rightarrow
X_1\times\cdots\times X_n
\]

satisfying

\[
\pi_i\circ f=f_i.
\]

State vectors throughout this volume are constructed using precisely
this product structure.

%%%%%%%%%%%%%%%%%%%%%%%%%%%%%%%%%%%%%%%%%%%%%%%%%%%%%%%%%%%%%%%

\section{Quotient Sets}

Let

\[
\sim
\]

be an equivalence relation on

\[
X.
\]

The quotient set is

\[
X/\!\sim
=
\{
[x]
:
x\in X
\},
\]

where

\[
[x]
=
\{y:y\sim x\}
\]

denotes the equivalence class containing

\[
x.
\]

Quotient constructions arise naturally whenever one identifies objects
that are operationally equivalent.

%%%%%%%%%%%%%%%%%%%%%%%%%%%%%%%%%%%%%%%%%%%%%%%%%%%%%%%%%%%%%%%

\section{Cardinality}

The cardinality of a set

\[
A
\]

is denoted

\[
|A|.
\]

Finite sets satisfy

\[
|A|
=
n
\]

for some non-negative integer

\[
n.
\]

Infinite sets may be countably infinite,

such as

\[
\mathbb N,
\]

or uncountable,

such as

\[
\mathbb R.
\]

Many infrastructure graphs are finite,

whereas their associated state spaces are typically infinite-dimensional.

%%%%%%%%%%%%%%%%%%%%%%%%%%%%%%%%%%%%%%%%%%%%%%%%%%%%%%%%%%%%%%%

\section{Power Sets}

The power set of

\[
A
\]

is

\[
\mathcal P(A)
=
\{B:B\subseteq A\}.
\]

If

\[
|A|
=
n,
\]

then

\[
|\mathcal P(A)|
=
2^n.
\]

Power sets frequently appear when considering collections of feasible
subnetworks, admissible controls, or possible failure scenarios.

%%%%%%%%%%%%%%%%%%%%%%%%%%%%%%%%%%%%%%%%%%%%%%%%%%%%%%%%%%%%%%%

\section{Fixed Points}

A point

\[
x
\]

is called a fixed point of

\[
f:X\rightarrow X
\]

if

\[
f(x)=x.
\]

Many equilibrium states introduced throughout this monograph are fixed
points of dynamical systems.

%%%%%%%%%%%%%%%%%%%%%%%%%%%%%%%%%%%%%%%%%%%%%%%%%%%%%%%%%%%%%%%

\begin{theorem}[Banach Fixed-Point Theorem]

Let

\[
(X,d)
\]

be a complete metric space, and let

\[
f:X\rightarrow X
\]

be a contraction mapping.

Then

\begin{enumerate}
\item a unique fixed point exists,
\item repeated iteration
\[
x_{n+1}=f(x_n)
\]
converges to that fixed point.
\end{enumerate}

\end{theorem}

This theorem provides the mathematical justification for many iterative
numerical algorithms.

%%%%%%%%%%%%%%%%%%%%%%%%%%%%%%%%%%%%%%%%%%%%%%%%%%%%%%%%%%%%%%%

\section{Directed Systems}

Many infrastructure systems evolve through nested approximations.

A directed family

\[
X_1
\subseteq
X_2
\subseteq
X_3
\subseteq
\cdots
\]

represents increasing refinement.

Examples include

\begin{itemize}
\item progressively refined computational meshes,
\item expanding transportation networks,
\item growing communication systems,
\item successive optimization approximations.
\end{itemize}

%%%%%%%%%%%%%%%%%%%%%%%%%%%%%%%%%%%%%%%%%%%%%%%%%%%%%%%%%%%%%%%

\section{Universal Mathematical Language}

The concepts introduced in this appendix appear throughout modern
mathematics regardless of application domain.

Infrastructure science employs these structures because they provide a
precise language for describing

collections of objects,

relationships,

transformations,

constraints,

and admissible configurations.

Subsequent appendices build upon these foundations by introducing
algebraic, geometric, analytical, probabilistic, and variational
structures that together constitute the mathematical framework used
throughout this volume.

%%%%%%%%%%%%%%%%%%%%%%%%%%%%%%%%%%%%%%%%%%%%%%%%%%%%%%%%%%%%%%%

\section{Algebraic Structures}

Many mathematical objects possess internal operations in addition to
their underlying set structure.

These operations permit the construction of increasingly sophisticated
mathematical theories.

Throughout this appendix, binary operations are denoted by

\[
*:X\times X\rightarrow X.
\]

%%%%%%%%%%%%%%%%%%%%%%%%%%%%%%%%%%%%%%%%%%%%%%%%%%%%%%%%%%%%%%%

\subsection{Semigroups}

A semigroup is a set

\[
S
\]

equipped with an associative binary operation

\[
*:S\times S\rightarrow S
\]

satisfying

\[
(a*b)*c
=
a*(b*c)
\]

for every

\[
a,b,c\in S.
\]

Associativity allows repeated compositions to be written without
ambiguity.

%%%%%%%%%%%%%%%%%%%%%%%%%%%%%%%%%%%%%%%%%%%%%%%%%%%%%%%%%%%%%%%

\subsection{Monoids}

A monoid is a semigroup possessing an identity element

\[
e
\]

such that

\[
e*a
=
a*e
=
a
\]

for every

\[
a\in S.
\]

Many collections of transformations naturally form monoids under
composition.

%%%%%%%%%%%%%%%%%%%%%%%%%%%%%%%%%%%%%%%%%%%%%%%%%%%%%%%%%%%%%%%

\subsection{Groups}

A group is a monoid in which every element possesses an inverse.

For every

\[
a\in G
\]

there exists

\[
a^{-1}
\]

satisfying

\[
aa^{-1}
=
a^{-1}a
=
e.
\]

Groups arise naturally whenever reversible transformations are studied.

%%%%%%%%%%%%%%%%%%%%%%%%%%%%%%%%%%%%%%%%%%%%%%%%%%%%%%%%%%%%%%%

\begin{example}

Rigid rotations of three-dimensional space form a group under
composition.

Translations form another group.

Together they generate the Euclidean symmetry group.

\end{example}

%%%%%%%%%%%%%%%%%%%%%%%%%%%%%%%%%%%%%%%%%%%%%%%%%%%%%%%%%%%%%%%

\subsection{Abelian Groups}

A group is called Abelian whenever

\[
ab=ba
\]

for every pair of elements.

Vector addition is the standard example.

%%%%%%%%%%%%%%%%%%%%%%%%%%%%%%%%%%%%%%%%%%%%%%%%%%%%%%%%%%%%%%%

\section{Rings}

A ring is a set equipped with two binary operations,

usually written

\[
+
\]

and

\[
\cdot.
\]

Addition forms an Abelian group,

while multiplication is associative and distributes over addition.

The integers

\[
\mathbb Z
\]

form the prototypical ring.

%%%%%%%%%%%%%%%%%%%%%%%%%%%%%%%%%%%%%%%%%%%%%%%%%%%%%%%%%%%%%%%

\section{Fields}

A field is a commutative ring in which every nonzero element possesses a
multiplicative inverse.

The fields used throughout this monograph are

\[
\mathbb R
\]

and

\[
\mathbb C.
\]

These provide the scalar systems for vector spaces, matrices, and
differential equations.

%%%%%%%%%%%%%%%%%%%%%%%%%%%%%%%%%%%%%%%%%%%%%%%%%%%%%%%%%%%%%%%

\section{Vector Spaces}

Let

\[
F
\]

be a field.

A vector space over

\[
F
\]

consists of a set

\[
V
\]

equipped with

vector addition

and

scalar multiplication

satisfying the usual vector-space axioms.

Almost every state variable appearing in this volume belongs to some
finite- or infinite-dimensional vector space.

%%%%%%%%%%%%%%%%%%%%%%%%%%%%%%%%%%%%%%%%%%%%%%%%%%%%%%%%%%%%%%%

\section{Linear Independence}

Vectors

\[
v_1,\ldots,v_n
\]

are linearly independent if

\[
a_1v_1+\cdots+a_nv_n=0
\]

implies

\[
a_1=\cdots=a_n=0.
\]

Otherwise they are linearly dependent.

%%%%%%%%%%%%%%%%%%%%%%%%%%%%%%%%%%%%%%%%%%%%%%%%%%%%%%%%%%%%%%%

\section{Bases}

A basis is a linearly independent spanning set.

Every vector

\[
v
\]

may therefore be written uniquely as

\[
v
=
\sum_i
a_ie_i,
\]

where

\[
\{e_i\}
\]

is a basis.

%%%%%%%%%%%%%%%%%%%%%%%%%%%%%%%%%%%%%%%%%%%%%%%%%%%%%%%%%%%%%%%

\begin{theorem}

Every finite-dimensional vector space possesses a basis.

Furthermore,

every two bases have the same cardinality.

\end{theorem}

The common cardinality is called the dimension of the vector space.

%%%%%%%%%%%%%%%%%%%%%%%%%%%%%%%%%%%%%%%%%%%%%%%%%%%%%%%%%%%%%%%

\section{Linear Transformations}

A mapping

\[
T:V\rightarrow W
\]

is linear whenever

\[
T(av+bw)
=
aT(v)+bT(w)
\]

for every pair of vectors and every scalar.

Linear transformations preserve the algebraic structure of vector
spaces.

%%%%%%%%%%%%%%%%%%%%%%%%%%%%%%%%%%%%%%%%%%%%%%%%%%%%%%%%%%%%%%%

\section{Kernel and Image}

For a linear transformation

\[
T,
\]

its kernel is

\[
\ker(T)
=
\{v:T(v)=0\},
\]

while its image is

\[
\operatorname{im}(T)
=
\{T(v):v\in V\}.
\]

%%%%%%%%%%%%%%%%%%%%%%%%%%%%%%%%%%%%%%%%%%%%%%%%%%%%%%%%%%%%%%%

\begin{theorem}[Rank--Nullity]

If

\[
V
\]

is finite dimensional,

then

\[
\dim(V)
=
\dim(\ker T)
+
\dim(\operatorname{im}T).
\]

\end{theorem}

This theorem underlies much of network analysis, where kernels represent
conserved quantities and images represent attainable transformations.

%%%%%%%%%%%%%%%%%%%%%%%%%%%%%%%%%%%%%%%%%%%%%%%%%%%%%%%%%%%%%%%

\section{Mathematical Hierarchy}

The structures introduced throughout this appendix form a hierarchy.

\[
\text{Sets}
\subset
\text{Relations}
\subset
\text{Functions}
\subset
\text{Algebraic Structures}
\subset
\text{Vector Spaces}
\subset
\text{Topological Spaces}
\subset
\text{Manifolds}.
\]

Each successive level retains the properties of the previous levels
while introducing additional structure.

Throughout this monograph, increasing mathematical sophistication arises
primarily by enriching existing structures rather than replacing them.

%%%%%%%%%%%%%%%%%%%%%%%%%%%%%%%%%%%%%%%%%%%%%%%%%%%%%%%%%%%%%%%%%%%%%%
\chapter{Linear Algebra and Tensor Analysis}

\section{Purpose}

Linear algebra provides the algebraic language of modern mathematical
science.

State variables,

control inputs,

graph representations,

optimization problems,

differential equations,

and numerical algorithms are all naturally expressed using vectors,
matrices, and linear transformations.

This appendix develops the principal results required throughout this
volume.

%%%%%%%%%%%%%%%%%%%%%%%%%%%%%%%%%%%%%%%%%%%%%%%%%%%%%%%%%%%%%%%

\section{Vector Spaces}

Let

\[
F
\]

be a field.

A vector space over

\[
F
\]

is a set

\[
V
\]

equipped with

vector addition

and

scalar multiplication

satisfying the vector-space axioms.

Typical examples include

\[
\mathbb R^n,
\qquad
\mathbb C^n,
\]

spaces of continuous functions,

and solution spaces of differential equations.

%%%%%%%%%%%%%%%%%%%%%%%%%%%%%%%%%%%%%%%%%%%%%%%%%%%%%%%%%%%%%%%

\section{Subspaces}

A subset

\[
W\subseteq V
\]

is called a subspace if

\[
u,v\in W
\]

implies

\[
u+v\in W,
\]

and

\[
\alpha u\in W
\]

for every scalar

\[
\alpha.
\]

%%%%%%%%%%%%%%%%%%%%%%%%%%%%%%%%%%%%%%%%%%%%%%%%%%%%%%%%%%%%%%%

\section{Span}

The span of vectors

\[
v_1,\ldots,v_n
\]

is

\[
\operatorname{span}
\{v_1,\ldots,v_n\}
=
\left\{
\sum_i
a_i v_i
\right\}.
\]

It is the smallest subspace containing those vectors.

%%%%%%%%%%%%%%%%%%%%%%%%%%%%%%%%%%%%%%%%%%%%%%%%%%%%%%%%%%%%%%%

\section{Basis and Dimension}

A basis is a linearly independent spanning set.

\begin{theorem}

Every finite-dimensional vector space possesses a basis.

Every basis contains the same number of vectors.

\end{theorem}

The common cardinality defines

\[
\dim(V).
\]

%%%%%%%%%%%%%%%%%%%%%%%%%%%%%%%%%%%%%%%%%%%%%%%%%%%%%%%%%%%%%%%

\section{Coordinates}

Relative to a basis

\[
\mathcal B
=
\{e_1,\ldots,e_n\},
\]

every vector admits a unique representation

\[
v
=
\sum_i
x_i e_i.
\]

The coefficients

\[
(x_1,\ldots,x_n)
\]

are called the coordinates of

\[
v.
\]

%%%%%%%%%%%%%%%%%%%%%%%%%%%%%%%%%%%%%%%%%%%%%%%%%%%%%%%%%%%%%%%

\section{Matrices}

A matrix

\[
A
=
(a_{ij})
\]

represents a linear transformation relative to chosen bases.

Matrix multiplication corresponds to composition of linear maps.

\[
(AB)x
=
A(Bx).
\]

%%%%%%%%%%%%%%%%%%%%%%%%%%%%%%%%%%%%%%%%%%%%%%%%%%%%%%%%%%%%%%%

\section{Matrix Operations}

Define

addition,

scalar multiplication,

transpose,

trace,

and determinant.

Introduce

\[
A^T,
\qquad
\operatorname{tr}(A),
\qquad
\det(A).
\]

%%%%%%%%%%%%%%%%%%%%%%%%%%%%%%%%%%%%%%%%%%%%%%%%%%%%%%%%%%%%%%%

\section{Inverse Matrices}

A square matrix is invertible whenever

\[
\det(A)\neq0.
\]

Its inverse satisfies

\[
AA^{-1}
=
A^{-1}A
=
I.
\]

%%%%%%%%%%%%%%%%%%%%%%%%%%%%%%%%%%%%%%%%%%%%%%%%%%%%%%%%%%%%%%%

\section{Rank}

The rank of a matrix equals the dimension of its column space.

\[
\operatorname{rank}(A)
=
\dim(\operatorname{im}A).
\]

Rank measures the effective dimensionality of a linear mapping.

%%%%%%%%%%%%%%%%%%%%%%%%%%%%%%%%%%%%%%%%%%%%%%%%%%%%%%%%%%%%%%%

\section{Eigenvalues}

Suppose

\[
A:V\rightarrow V.
\]

An eigenvector satisfies

\[
Av
=
\lambda v.
\]

The scalar

\[
\lambda
\]

is called an eigenvalue.

%%%%%%%%%%%%%%%%%%%%%%%%%%%%%%%%%%%%%%%%%%%%%%%%%%%%%%%%%%%%%%%

\begin{theorem}

The eigenvalues of

\[
A
\]

are precisely the roots of

\[
\det(A-\lambda I)=0.
\]

\end{theorem}

%%%%%%%%%%%%%%%%%%%%%%%%%%%%%%%%%%%%%%%%%%%%%%%%%%%%%%%%%%%%%%%

\section{Diagonalization}

A matrix is diagonalizable whenever it possesses a complete basis of
eigenvectors.

Then

\[
A
=
PDP^{-1},
\]

where

\[
D
\]

is diagonal.

%%%%%%%%%%%%%%%%%%%%%%%%%%%%%%%%%%%%%%%%%%%%%%%%%%%%%%%%%%%%%%%

\section{Symmetric Matrices}

A matrix satisfying

\[
A=A^T
\]

is symmetric.

\begin{theorem}[Spectral Theorem]

Every real symmetric matrix possesses an orthonormal basis of
eigenvectors.

Its eigenvalues are real.

\end{theorem}

%%%%%%%%%%%%%%%%%%%%%%%%%%%%%%%%%%%%%%%%%%%%%%%%%%%%%%%%%%%%%%%

\section{Inner Products}

An inner product is a mapping

\[
\langle\cdot,\cdot\rangle
:
V\times V
\rightarrow
F
\]

satisfying

linearity,

symmetry,

and positive definiteness.

The Euclidean inner product is

\[
\langle x,y\rangle
=
x^Ty.
\]

%%%%%%%%%%%%%%%%%%%%%%%%%%%%%%%%%%%%%%%%%%%%%%%%%%%%%%%%%%%%%%%

\section{Norms}

The induced norm is

\[
\|x\|
=
\sqrt{\langle x,x\rangle}.
\]

Norms define distance,

angles,

and orthogonality.

%%%%%%%%%%%%%%%%%%%%%%%%%%%%%%%%%%%%%%%%%%%%%%%%%%%%%%%%%%%%%%%

\section{Orthogonal Projections}

Let

\[
W
\subseteq
V.
\]

Every vector decomposes uniquely as

\[
v
=
w+r,
\]

where

\[
w\in W,
\qquad
r\perp W.
\]

Orthogonal projection plays a central role in optimization and least
squares.

%%%%%%%%%%%%%%%%%%%%%%%%%%%%%%%%%%%%%%%%%%%%%%%%%%%%%%%%%%%%%%%

\section{Singular Value Decomposition}

Every matrix admits the factorization

\[
A
=
U\Sigma V^T.
\]

The diagonal entries of

\[
\Sigma
\]

are the singular values.

The SVD provides the optimal low-rank approximation of a matrix.

%%%%%%%%%%%%%%%%%%%%%%%%%%%%%%%%%%%%%%%%%%%%%%%%%%%%%%%%%%%%%%%

\section{Positive Definite Matrices}

A symmetric matrix is positive definite if

\[
x^TAx>0
\]

for every nonzero vector.

Positive definite matrices appear throughout optimization, graph
Laplacians, covariance matrices, and energy functionals.

%%%%%%%%%%%%%%%%%%%%%%%%%%%%%%%%%%%%%%%%%%%%%%%%%%%%%%%%%%%%%%%

\section{Tensor Products}

Given vector spaces

\[
V
\]

and

\[
W,
\]

their tensor product

\[
V\otimes W
\]

forms the space of bilinear combinations.

Tensors generalize vectors and matrices to multilinear mappings.

%%%%%%%%%%%%%%%%%%%%%%%%%%%%%%%%%%%%%%%%%%%%%%%%%%%%%%%%%%%%%%%

\section{Tensor Notation}

Introduce index notation,

Einstein summation,

covariant and contravariant indices,

tensor contraction,

and tensor products.

These conventions will be employed throughout the differential geometry
appendix.

%%%%%%%%%%%%%%%%%%%%%%%%%%%%%%%%%%%%%%%%%%%%%%%%%%%%%%%%%%%%%%%

\section{Concluding Remarks}

Linear algebra provides the common algebraic language for every major
construction developed in this monograph.

Graphs are represented by matrices,

state variables by vectors,

optimization by quadratic forms,

control systems by linear operators,

and geometric structures by tensors.

The following appendix builds upon these algebraic foundations to
develop the differential geometry required for infrastructure embedded
on curved manifolds.

%%%%%%%%%%%%%%%%%%%%%%%%%%%%%%%%%%%%%%%%%%%%%%%%%%%%%%%%%%%%%%%

\section{Matrix Norms}

Norms measure the size of matrices in a manner analogous to vector
norms.

A matrix norm

\[
\|\cdot\|
:
\mathbb R^{m\times n}
\rightarrow
\mathbb R
\]

satisfies

\begin{enumerate}
\item
\[
\|A\|\ge0,
\]

\item
\[
\|A\|=0
\Longleftrightarrow
A=0,
\]

\item
\[
\|\alpha A\|
=
|\alpha|
\|A\|,
\]

\item
\[
\|A+B\|
\le
\|A\|
+
\|B\|.
\]
\end{enumerate}

Common examples include

the Frobenius norm,

the spectral norm,

and induced operator norms.

%%%%%%%%%%%%%%%%%%%%%%%%%%%%%%%%%%%%%%%%%%%%%%%%%%%%%%%%%%%%%%%

\section{Operator Norms}

Suppose

\[
A:V\rightarrow W
\]

is linear.

Its induced norm is

\[
\boxed{
\|A\|
=
\sup_{x\neq0}
\frac{\|Ax\|}{\|x\|}.
}
\]

Operator norms quantify the largest possible amplification produced by
a linear transformation.

%%%%%%%%%%%%%%%%%%%%%%%%%%%%%%%%%%%%%%%%%%%%%%%%%%%%%%%%%%%%%%%

\section{Orthogonal Matrices}

A matrix

\[
Q
\]

is orthogonal whenever

\[
Q^TQ
=
QQ^T
=
I.
\]

Orthogonal matrices preserve

length,

angles,

and inner products.

%%%%%%%%%%%%%%%%%%%%%%%%%%%%%%%%%%%%%%%%%%%%%%%%%%%%%%%%%%%%%%%

\begin{theorem}

If

\[
Q
\]

is orthogonal,

then

\[
\|Qx\|
=
\|x\|
\]

for every vector

\[
x.
\]

\end{theorem}

Orthogonal transformations therefore represent rigid rotations and
reflections.

%%%%%%%%%%%%%%%%%%%%%%%%%%%%%%%%%%%%%%%%%%%%%%%%%%%%%%%%%%%%%%%

\section{Projection Matrices}

A matrix

\[
P
\]

is a projection whenever

\[
P^2=P.
\]

If

\[
P=P^T,
\]

then

\[
P
\]

is an orthogonal projection.

Projection operators appear throughout least-squares estimation,
optimization, and state estimation.

%%%%%%%%%%%%%%%%%%%%%%%%%%%%%%%%%%%%%%%%%%%%%%%%%%%%%%%%%%%%%%%

\section{Quadratic Forms}

Given a symmetric matrix

\[
A,
\]

the quadratic form associated with

\[
A
\]

is

\[
\boxed{
Q(x)
=
x^TAx.
}
\]

Quadratic forms describe

energy,

variance,

distance,

and optimization objectives.

%%%%%%%%%%%%%%%%%%%%%%%%%%%%%%%%%%%%%%%%%%%%%%%%%%%%%%%%%%%%%%%

\begin{theorem}

A symmetric matrix is positive definite if and only if

\[
Q(x)>0
\]

for every nonzero vector.

\end{theorem}

%%%%%%%%%%%%%%%%%%%%%%%%%%%%%%%%%%%%%%%%%%%%%%%%%%%%%%%%%%%%%%%

\section{Least-Squares Problems}

Many infrastructure models are overdetermined.

Given

\[
Ax\approx b,
\]

the least-squares solution minimizes

\[
\|Ax-b\|^2.
\]

The associated normal equations are

\[
\boxed{
A^TAx
=
A^Tb.
}
\]

%%%%%%%%%%%%%%%%%%%%%%%%%%%%%%%%%%%%%%%%%%%%%%%%%%%%%%%%%%%%%%%

\section{Moore--Penrose Pseudoinverse}

When

\[
A
\]

is not invertible,

one replaces the inverse by the Moore--Penrose pseudoinverse

\[
A^+.
\]

The minimum-norm least-squares solution becomes

\[
\boxed{
x
=
A^+b.
}
\]

The pseudoinverse is fundamental in inverse problems and state
reconstruction.

%%%%%%%%%%%%%%%%%%%%%%%%%%%%%%%%%%%%%%%%%%%%%%%%%%%%%%%%%%%%%%%

\section{Block Matrices}

Large infrastructure systems naturally decompose into interacting
subsystems.

Accordingly,

matrices frequently possess block structure,

for example

\[
A
=
\begin{bmatrix}
A_{11} & A_{12}\\
A_{21} & A_{22}
\end{bmatrix}.
\]

Block matrices simplify both theoretical analysis and numerical
computation.

%%%%%%%%%%%%%%%%%%%%%%%%%%%%%%%%%%%%%%%%%%%%%%%%%%%%%%%%%%%%%%%

\section{Kronecker Product}

Given matrices

\[
A
=
(a_{ij}),
\qquad
B,
\]

their Kronecker product is

\[
A\otimes B
=
\begin{bmatrix}
a_{11}B & \cdots & a_{1n}B\\
\vdots & \ddots & \vdots\\
a_{m1}B & \cdots & a_{mn}B
\end{bmatrix}.
\]

Kronecker products arise naturally in multidimensional discretizations
and coupled network models.

%%%%%%%%%%%%%%%%%%%%%%%%%%%%%%%%%%%%%%%%%%%%%%%%%%%%%%%%%%%%%%%

\section{Matrix Exponentials}

Linear dynamical systems satisfy

\[
\dot x
=
Ax.
\]

The unique solution is

\[
\boxed{
x(t)
=
e^{At}x(0),
}
\]

where

\[
e^{At}
=
\sum_{k=0}^{\infty}
\frac{(At)^k}{k!}.
\]

Matrix exponentials form the foundation of continuous-time linear system
theory.

%%%%%%%%%%%%%%%%%%%%%%%%%%%%%%%%%%%%%%%%%%%%%%%%%%%%%%%%%%%%%%%

\section{Jordan Canonical Form}

Not every matrix is diagonalizable.

Nevertheless,

every square matrix over

\[
\mathbb C
\]

is similar to a Jordan matrix,

whose blocks have the form

\[
\begin{bmatrix}
\lambda & 1 & 0 & \cdots\\
0 & \lambda & 1 & \cdots\\
\vdots & & \ddots & \ddots\\
0 & \cdots & 0 & \lambda
\end{bmatrix}.
\]

Jordan form characterizes the complete algebraic structure of linear
operators.

%%%%%%%%%%%%%%%%%%%%%%%%%%%%%%%%%%%%%%%%%%%%%%%%%%%%%%%%%%%%%%%

\section{Spectral Radius}

The spectral radius of a matrix is

\[
\boxed{
\rho(A)
=
\max_i
|\lambda_i|.
}
\]

Spectral radius governs the asymptotic behavior of repeated matrix
multiplication and appears throughout network dynamics and stability
analysis.

%%%%%%%%%%%%%%%%%%%%%%%%%%%%%%%%%%%%%%%%%%%%%%%%%%%%%%%%%%%%%%%

\section{Condition Numbers}

Numerical algorithms are influenced by matrix conditioning.

The condition number of an invertible matrix is

\[
\boxed{
\kappa(A)
=
\|A\|
\,
\|A^{-1}\|.
}
\]

Large condition numbers indicate sensitivity to perturbations.

Well-conditioned systems are numerically stable.

%%%%%%%%%%%%%%%%%%%%%%%%%%%%%%%%%%%%%%%%%%%%%%%%%%%%%%%%%%%%%%%

\section{Perron--Frobenius Theory}

Many infrastructure matrices contain only nonnegative entries.

%%%%%%%%%%%%%%%%%%%%%%%%%%%%%%%%%%%%%%%%%%%%%%%%%%%%%%%%%%%%%%%

\begin{theorem}[Perron--Frobenius]

If

\[
A
\]

is irreducible and nonnegative,

then

\begin{enumerate}
\item
its spectral radius is a positive eigenvalue,

\item
that eigenvalue possesses a strictly positive eigenvector,

\item
all other eigenvalues have smaller modulus.
\end{enumerate}

\end{theorem}

This theorem underlies ranking algorithms, population models,
transportation networks, and certain infrastructure equilibrium models.

%%%%%%%%%%%%%%%%%%%%%%%%%%%%%%%%%%%%%%%%%%%%%%%%%%%%%%%%%%%%%%%

\section{Summary}

Linear algebra provides the computational language of infrastructure
science.

Vectors represent states.

Matrices represent interactions.

Eigenvalues describe stability.

Quadratic forms quantify energy.

Orthogonal projections perform estimation.

Tensor products couple subsystems.

Every subsequent appendix builds upon these algebraic foundations.

The next appendix develops differential geometry, extending linear
concepts from flat vector spaces to curved manifolds embedded in
physical space.

%%%%%%%%%%%%%%%%%%%%%%%%%%%%%%%%%%%%%%%%%%%%%%%%%%%%%%%%%%%%%%%%%%%%%%
\chapter{Differential Geometry}

\section{Purpose}

Many infrastructure systems are embedded in physical space rather than
existing as abstract graphs alone.

Roads follow the Earth's surface.

Pipelines traverse terrain.

Transmission corridors cross mountains and rivers.

Communication satellites occupy orbital trajectories.

Consequently, optimization frequently occurs on curved geometric spaces
rather than Euclidean space.

Differential geometry provides the mathematical language required to
study such embedded infrastructure.

%%%%%%%%%%%%%%%%%%%%%%%%%%%%%%%%%%%%%%%%%%%%%%%%%%%%%%%%%%%%%%%

\section{Smooth Manifolds}

A manifold is a space that locally resembles Euclidean space.

%%%%%%%%%%%%%%%%%%%%%%%%%%%%%%%%%%%%%%%%%%%%%%%%%%%%%%%%%%%%%%%

\begin{definition}

An \(n\)-dimensional smooth manifold is a Hausdorff, second-countable
topological space in which every point possesses a neighborhood
diffeomorphic to an open subset of

\[
\mathbb R^n.
\]

\end{definition}

Although the global topology may be complicated, every sufficiently
small neighborhood admits ordinary Cartesian coordinates.

%%%%%%%%%%%%%%%%%%%%%%%%%%%%%%%%%%%%%%%%%%%%%%%%%%%%%%%%%%%%%%%

\section{Charts}

A chart consists of

\[
(U,\varphi),
\]

where

\[
U\subseteq M
\]

is open and

\[
\varphi:U\rightarrow\mathbb R^n
\]

is a homeomorphism onto its image.

The coordinates

\[
(x^1,\ldots,x^n)
\]

obtained from

\[
\varphi
\]

are called local coordinates.

%%%%%%%%%%%%%%%%%%%%%%%%%%%%%%%%%%%%%%%%%%%%%%%%%%%%%%%%%%%%%%%

\section{Atlases}

A collection of compatible charts covering

\[
M
\]

forms an atlas.

Compatibility requires that overlapping coordinate maps satisfy

\[
\boxed{
\varphi_\beta
\circ
\varphi_\alpha^{-1}
}
\]

being infinitely differentiable wherever defined.

%%%%%%%%%%%%%%%%%%%%%%%%%%%%%%%%%%%%%%%%%%%%%%%%%%%%%%%%%%%%%%%

\section{Smooth Maps}

Suppose

\[
M
\]

and

\[
N
\]

are smooth manifolds.

A mapping

\[
F:M\rightarrow N
\]

is smooth if every coordinate representation

\[
\psi
\circ
F
\circ
\varphi^{-1}
\]

is infinitely differentiable.

Smoothness therefore does not depend upon the particular coordinates
chosen.

%%%%%%%%%%%%%%%%%%%%%%%%%%%%%%%%%%%%%%%%%%%%%%%%%%%%%%%%%%%%%%%

\section{Curves}

A smooth curve is a mapping

\[
\gamma:I\rightarrow M,
\]

where

\[
I
\subseteq
\mathbb R.
\]

Infrastructure corridors,

vehicle trajectories,

communication paths,

and orbital transfers are all naturally represented by smooth curves.

%%%%%%%%%%%%%%%%%%%%%%%%%%%%%%%%%%%%%%%%%%%%%%%%%%%%%%%%%%%%%%%

\section{Tangent Vectors}

At each point

\[
p\in M,
\]

there exists an associated tangent space

\[
T_pM.
\]

Intuitively,

tangent vectors represent instantaneous directions of motion through
the manifold.

%%%%%%%%%%%%%%%%%%%%%%%%%%%%%%%%%%%%%%%%%%%%%%%%%%%%%%%%%%%%%%%

\section{The Tangent Bundle}

The collection of all tangent spaces forms the tangent bundle

\[
\boxed{
TM
=
\bigcup_{p\in M}
T_pM.
}
\]

A vector field assigns one tangent vector to every point of the
manifold.

%%%%%%%%%%%%%%%%%%%%%%%%%%%%%%%%%%%%%%%%%%%%%%%%%%%%%%%%%%%%%%%

\section{Vector Fields}

A smooth vector field is a mapping

\[
X:M\rightarrow TM
\]

satisfying

\[
X(p)\in T_pM
\]

for every

\[
p.
\]

Examples include

wind velocity,

fluid velocity,

traffic velocity,

and gradient directions.

%%%%%%%%%%%%%%%%%%%%%%%%%%%%%%%%%%%%%%%%%%%%%%%%%%%%%%%%%%%%%%%

\section{Differential of a Map}

Suppose

\[
F:M\rightarrow N.
\]

Its differential at

\[
p
\]

is the linear map

\[
\boxed{
dF_p
:
T_pM
\rightarrow
T_{F(p)}N.
}
\]

The differential generalizes the Jacobian matrix to arbitrary
manifolds.

%%%%%%%%%%%%%%%%%%%%%%%%%%%%%%%%%%%%%%%%%%%%%%%%%%%%%%%%%%%%%%%

\section{Cotangent Spaces}

Dual to each tangent space is the cotangent space

\[
T_p^*M,
\]

consisting of all linear functionals on

\[
T_pM.
\]

Elements of the cotangent bundle are called covectors.

Gradients naturally belong to cotangent spaces.

%%%%%%%%%%%%%%%%%%%%%%%%%%%%%%%%%%%%%%%%%%%%%%%%%%%%%%%%%%%%%%%

\section{Differential Forms}

A differential one-form assigns a covector to every point of the
manifold.

Higher-order differential forms arise through alternating tensor
products.

They provide the natural language for

line integrals,

surface integrals,

and volume integrals.

%%%%%%%%%%%%%%%%%%%%%%%%%%%%%%%%%%%%%%%%%%%%%%%%%%%%%%%%%%%%%%%

\section{Exterior Derivative}

The exterior derivative

\[
d
\]

maps

\[
k
\]

-forms into

\[
(k+1)
\]

-forms.

It satisfies

\[
\boxed{
d^2=0.
}
\]

This single identity unifies

gradient,

curl,

and divergence

within one mathematical operator.

%%%%%%%%%%%%%%%%%%%%%%%%%%%%%%%%%%%%%%%%%%%%%%%%%%%%%%%%%%%%%%%

\section{Pullbacks}

Suppose

\[
F:M\rightarrow N.
\]

The pullback operator

\[
F^*
\]

transfers differential forms from

\[
N
\]

back to

\[
M.
\]

Pullbacks permit measurements made in one coordinate system to be
interpreted in another.

%%%%%%%%%%%%%%%%%%%%%%%%%%%%%%%%%%%%%%%%%%%%%%%%%%%%%%%%%%%%%%%

\section{Pushforwards}

The differential

\[
dF
\]

is often called the pushforward.

It transports tangent vectors forward along smooth mappings,

complementing the pullback of differential forms.

%%%%%%%%%%%%%%%%%%%%%%%%%%%%%%%%%%%%%%%%%%%%%%%%%%%%%%%%%%%%%%%

\section{Summary}

The concepts introduced thus far establish the local differential
structure of smooth manifolds.

Charts provide coordinates.

Tangent spaces describe directions.

Cotangent spaces describe gradients.

Vector fields describe motion.

Differential forms describe integration.

The remaining sections develop metrics, curvature, geodesics, and
connections, providing the geometric machinery used throughout the main
text to model infrastructure embedded in curved physical space.

%%%%%%%%%%%%%%%%%%%%%%%%%%%%%%%%%%%%%%%%%%%%%%%%%%%%%%%%%%%%%%%

\section{Riemannian Metrics}

The tangent space at each point of a smooth manifold is a vector space.
To measure lengths, angles, and volumes, an additional geometric
structure is required.

\begin{definition}

A \emph{Riemannian metric} on a smooth manifold \(M\) is a smoothly
varying family of inner products

\[
g_p:T_pM\times T_pM\rightarrow\mathbb R,
\]

one for each point

\[
p\in M.
\]

\end{definition}

The pair

\[
(M,g)
\]

is called a \emph{Riemannian manifold}.

%%%%%%%%%%%%%%%%%%%%%%%%%%%%%%%%%%%%%%%%%%%%%%%%%%%%%%%%%%%%%%%

\section{Metric Components}

In local coordinates

\[
(x^1,\ldots,x^n),
\]

the metric is represented by the symmetric matrix

\[
g_{ij}
=
g\!\left(
\frac{\partial}{\partial x^i},
\frac{\partial}{\partial x^j}
\right).
\]

Accordingly,

\[
\boxed{
ds^2
=
g_{ij}\,
dx^i
dx^j.
}
\]

The metric tensor is positive definite,

so

\[
g_{ij}\xi^i\xi^j>0
\]

for every nonzero tangent vector.

%%%%%%%%%%%%%%%%%%%%%%%%%%%%%%%%%%%%%%%%%%%%%%%%%%%%%%%%%%%%%%%

\section{Length of Curves}

Let

\[
\gamma:[a,b]\rightarrow M
\]

be a smooth curve.

Its length is

\[
\boxed{
L(\gamma)
=
\int_a^b
\sqrt{
g_{ij}
\frac{dx^i}{dt}
\frac{dx^j}{dt}
}\,
dt.
}
\]

This generalizes the Euclidean arc-length formula.

%%%%%%%%%%%%%%%%%%%%%%%%%%%%%%%%%%%%%%%%%%%%%%%%%%%%%%%%%%%%%%%

\section{Distance}

The Riemannian distance between two points is defined as the minimum
length among all smooth curves joining them.

\[
\boxed{
d(p,q)
=
\inf_\gamma
L(\gamma).
}
\]

Distance therefore emerges from optimization over admissible paths.

%%%%%%%%%%%%%%%%%%%%%%%%%%%%%%%%%%%%%%%%%%%%%%%%%%%%%%%%%%%%%%%

\section{Examples}

In Euclidean space,

\[
g_{ij}
=
\delta_{ij},
\]

and

\[
ds^2
=
dx^2+dy^2+dz^2.
\]

On the unit sphere,

\[
ds^2
=
d\theta^2
+
\sin^2\theta\,d\phi^2.
\]

The latter geometry gives rise to great-circle (orthodromic) paths,
which play a central role in long-distance transportation and global
infrastructure.

%%%%%%%%%%%%%%%%%%%%%%%%%%%%%%%%%%%%%%%%%%%%%%%%%%%%%%%%%%%%%%%

\section{Covariant Derivatives}

Ordinary differentiation compares vectors located at different points.

Since tangent spaces at different points are distinct vector spaces,
ordinary subtraction is not intrinsically defined.

A covariant derivative provides the appropriate notion of
differentiation on a manifold.

%%%%%%%%%%%%%%%%%%%%%%%%%%%%%%%%%%%%%%%%%%%%%%%%%%%%%%%%%%%%%%%

\begin{definition}

A covariant derivative is an operator

\[
\nabla
\]

satisfying

\[
\nabla_XY
\]

for every pair of smooth vector fields

\[
X
\]

and

\[
Y,
\]

such that

\begin{enumerate}
\item it is linear in each argument,
\item it satisfies the Leibniz rule,
\item it reduces locally to ordinary differentiation plus correction
terms determined by the geometry.
\end{enumerate}

\end{definition}

%%%%%%%%%%%%%%%%%%%%%%%%%%%%%%%%%%%%%%%%%%%%%%%%%%%%%%%%%%%%%%%

\section{Christoffel Symbols}

In local coordinates,

the covariant derivative is represented by

\[
\boxed{
\nabla_iX^k
=
\frac{\partial X^k}{\partial x^i}
+
\Gamma^k_{ij}X^j.
}
\]

The quantities

\[
\Gamma^k_{ij}
\]

are called the Christoffel symbols.

Although they are not tensors,

they encode how the coordinate basis changes from point to point.

%%%%%%%%%%%%%%%%%%%%%%%%%%%%%%%%%%%%%%%%%%%%%%%%%%%%%%%%%%%%%%%

\begin{theorem}

For the unique metric-compatible torsion-free connection,

\[
\boxed{
\Gamma^k_{ij}
=
\frac12
g^{k\ell}
\left(
\frac{\partial g_{j\ell}}{\partial x^i}
+
\frac{\partial g_{i\ell}}{\partial x^j}
-
\frac{\partial g_{ij}}{\partial x^\ell}
\right).
}
\]

\end{theorem}

%%%%%%%%%%%%%%%%%%%%%%%%%%%%%%%%%%%%%%%%%%%%%%%%%%%%%%%%%%%%%%%

\section{Parallel Transport}

A vector field

\[
V(t)
\]

along a curve

\[
\gamma(t)
\]

is said to be parallel if

\[
\boxed{
\nabla_{\dot\gamma}V
=
0.
}
\]

Parallel transport preserves vector length and angle while moving the
vector along the curve.

Unlike Euclidean space,

the transported vector generally depends upon the path taken.

%%%%%%%%%%%%%%%%%%%%%%%%%%%%%%%%%%%%%%%%%%%%%%%%%%%%%%%%%%%%%%%

\section{Geodesics}

Geodesics generalize straight lines to curved manifolds.

%%%%%%%%%%%%%%%%%%%%%%%%%%%%%%%%%%%%%%%%%%%%%%%%%%%%%%%%%%%%%%%

\begin{definition}

A smooth curve

\[
\gamma(t)
\]

is a geodesic whenever its tangent vector transports parallel to itself,

\[
\boxed{
\nabla_{\dot\gamma}\dot\gamma
=
0.
}
\]

\end{definition}

In local coordinates,

this becomes

\[
\boxed{
\frac{d^2x^k}{dt^2}
+
\Gamma^k_{ij}
\frac{dx^i}{dt}
\frac{dx^j}{dt}
=
0.
}
\]

%%%%%%%%%%%%%%%%%%%%%%%%%%%%%%%%%%%%%%%%%%%%%%%%%%%%%%%%%%%%%%%

\section{Geodesic Interpretation}

Geodesics possess several equivalent interpretations.

They are

\begin{itemize}
\item locally shortest paths,
\item curves of stationary length,
\item curves with zero intrinsic acceleration,
\item trajectories determined entirely by the geometry of the manifold.
\end{itemize}

In Euclidean space,

geodesics are straight lines.

On the sphere,

they are great circles.

Throughout this volume, orthodromic infrastructure routes are modeled as
geodesics of the Earth's Riemannian geometry.

%%%%%%%%%%%%%%%%%%%%%%%%%%%%%%%%%%%%%%%%%%%%%%%%%%%%%%%%%%%%%%%

\section{Curvature}

The defining feature of a curved manifold is that vectors transported
around closed loops generally fail to return to their original
orientation.

This phenomenon is quantified by the curvature tensor.

%%%%%%%%%%%%%%%%%%%%%%%%%%%%%%%%%%%%%%%%%%%%%%%%%%%%%%%%%%%%%%%

\section{The Riemann Curvature Tensor}

Let

\[
X,\;Y,\;Z
\]

be smooth vector fields.

The Riemann curvature tensor is defined by

\[
\boxed{
R(X,Y)Z
=
\nabla_X\nabla_YZ
-
\nabla_Y\nabla_XZ
-
\nabla_{[X,Y]}Z,
}
\]

where

\[
[X,Y]
\]

denotes the Lie bracket.

Curvature therefore measures the noncommutativity of successive
covariant differentiation.

%%%%%%%%%%%%%%%%%%%%%%%%%%%%%%%%%%%%%%%%%%%%%%%%%%%%%%%%%%%%%%%

\section{Coordinate Representation}

In local coordinates,

the components of the curvature tensor are

\[
\boxed{
R^\ell{}_{ijk}
=
\frac{\partial\Gamma^\ell_{jk}}
{\partial x^i}
-
\frac{\partial\Gamma^\ell_{ik}}
{\partial x^j}
+
\Gamma^\ell_{im}\Gamma^m_{jk}
-
\Gamma^\ell_{jm}\Gamma^m_{ik}.
}
\]

These components depend only upon the metric and its first and second
derivatives.

%%%%%%%%%%%%%%%%%%%%%%%%%%%%%%%%%%%%%%%%%%%%%%%%%%%%%%%%%%%%%%%

\section{Symmetries of the Curvature Tensor}

The Riemann tensor possesses several important symmetries.

For every collection of indices,

\[
R_{ijkl}
=
-
R_{jikl},
\]

\[
R_{ijkl}
=
-
R_{ijlk},
\]

\[
R_{ijkl}
=
R_{klij},
\]

and

\[
R_{ijkl}
+
R_{iklj}
+
R_{iljk}
=
0.
\]

These identities substantially reduce the number of independent
components.

%%%%%%%%%%%%%%%%%%%%%%%%%%%%%%%%%%%%%%%%%%%%%%%%%%%%%%%%%%%%%%%

\section{Ricci Curvature}

Contracting the Riemann tensor produces the Ricci tensor,

\[
\boxed{
R_{ij}
=
R^k{}_{ikj}.
}
\]

Ricci curvature measures the average tendency of nearby geodesics to
converge or diverge.

Unlike the full Riemann tensor,

the Ricci tensor contains only second-order information about volume
distortion.

%%%%%%%%%%%%%%%%%%%%%%%%%%%%%%%%%%%%%%%%%%%%%%%%%%%%%%%%%%%%%%%

\section{Scalar Curvature}

A further contraction produces the scalar curvature,

\[
\boxed{
R
=
g^{ij}R_{ij}.
}
\]

Scalar curvature represents the average intrinsic curvature at a point.

Positive scalar curvature corresponds to locally sphere-like geometry,

negative curvature to saddle-like geometry,

and zero curvature to locally flat geometry.

%%%%%%%%%%%%%%%%%%%%%%%%%%%%%%%%%%%%%%%%%%%%%%%%%%%%%%%%%%%%%%%

\section{Sectional Curvature}

Given two linearly independent tangent vectors

\[
u,v\in T_pM,
\]

the sectional curvature of the plane they span is

\[
\boxed{
K(u,v)
=
\frac{
\langle
R(u,v)v,
u
\rangle
}
{
\|u\|^2\|v\|^2
-
\langle u,v\rangle^2
}.
}
\]

Sectional curvature generalizes Gaussian curvature to arbitrary
dimensions.

%%%%%%%%%%%%%%%%%%%%%%%%%%%%%%%%%%%%%%%%%%%%%%%%%%%%%%%%%%%%%%%

\section{Gaussian Curvature}

For a two-dimensional manifold,

the sectional curvature reduces to the Gaussian curvature,

\[
K.
\]

The Gaussian curvature depends solely upon intrinsic geometry and is
independent of the surrounding Euclidean space.

This remarkable fact forms the basis of Gauss's
\emph{Theorema Egregium}.

%%%%%%%%%%%%%%%%%%%%%%%%%%%%%%%%%%%%%%%%%%%%%%%%%%%%%%%%%%%%%%%

\begin{theorem}[Gauss's Theorema Egregium]

Gaussian curvature is an intrinsic invariant.

It may be computed entirely from the metric tensor without reference to
the manifold's embedding in any higher-dimensional space.

\end{theorem}

%%%%%%%%%%%%%%%%%%%%%%%%%%%%%%%%%%%%%%%%%%%%%%%%%%%%%%%%%%%%%%%

\section{Geodesic Deviation}

Curvature determines how nearby geodesics separate.

Let

\[
J
\]

denote the separation vector between neighboring geodesics.

Then

\[
\boxed{
\nabla_T\nabla_TJ
+
R(J,T)T
=
0,
}
\]

where

\[
T
\]

is the tangent vector of the reference geodesic.

This equation describes tidal effects, focusing phenomena, and the
stability of neighboring trajectories.

%%%%%%%%%%%%%%%%%%%%%%%%%%%%%%%%%%%%%%%%%%%%%%%%%%%%%%%%%%%%%%%

\section{Flat Manifolds}

A manifold is flat if

\[
\boxed{
R^\ell{}_{ijk}
=
0
}
\]

everywhere.

Euclidean space is flat.

In flat manifolds,

parallel transport is path independent,

geodesics are straight lines,

and Cartesian coordinates may be introduced globally.

%%%%%%%%%%%%%%%%%%%%%%%%%%%%%%%%%%%%%%%%%%%%%%%%%%%%%%%%%%%%%%%

\section{Spaces of Constant Curvature}

Three fundamental geometries occur repeatedly throughout mathematics.

\begin{center}
\begin{tabular}{ll}
\toprule
Curvature & Geometry\\
\midrule
$K>0$ & Spherical geometry\\
$K=0$ & Euclidean geometry\\
$K<0$ & Hyperbolic geometry\\
\bottomrule
\end{tabular}
\end{center}

Many optimization problems depend critically upon which of these
geometric settings is appropriate.

%%%%%%%%%%%%%%%%%%%%%%%%%%%%%%%%%%%%%%%%%%%%%%%%%%%%%%%%%%%%%%%

\section{Applications to Infrastructure}

Curvature influences infrastructure planning whenever the underlying
physical domain cannot be approximated as flat.

Examples include

\begin{itemize}
\item great-circle airline routes,
\item submarine communication cables,
\item intercontinental pipelines,
\item satellite constellations,
\item planetary transportation networks.
\end{itemize}

Ignoring curvature introduces systematic geometric errors that become
significant over continental or planetary scales.

%%%%%%%%%%%%%%%%%%%%%%%%%%%%%%%%%%%%%%%%%%%%%%%%%%%%%%%%%%%%%%%

\section{Summary}

Curvature measures the intrinsic deviation of a manifold from Euclidean
geometry.

The Riemann tensor encodes the full curvature structure.

The Ricci tensor and scalar curvature summarize its principal effects.

Sectional curvature characterizes the bending of individual tangent
planes, while geodesic deviation describes the evolution of neighboring
trajectories.

These concepts provide the geometric foundation for optimization and
transport on curved spaces developed throughout the main text.

%%%%%%%%%%%%%%%%%%%%%%%%%%%%%%%%%%%%%%%%%%%%%%%%%%%%%%%%%%%%%%%

\section{Lie Derivatives}

Many quantities evolve along the flow generated by a vector field.

The Lie derivative measures the intrinsic rate of change of geometric
objects under this motion.

Unlike the covariant derivative, which compares nearby tangent spaces,
the Lie derivative compares tensors transported by a flow.

%%%%%%%%%%%%%%%%%%%%%%%%%%%%%%%%%%%%%%%%%%%%%%%%%%%%%%%%%%%%%%%

\section{Flows of Vector Fields}

Let

\[
X
\]

be a smooth vector field on

\[
M.
\]

The associated flow is the one-parameter family of maps

\[
\Phi_t:M\rightarrow M,
\]

satisfying

\[
\frac{d}{dt}\Phi_t(p)
=
X(\Phi_t(p)),
\qquad
\Phi_0(p)=p.
\]

Every smooth vector field therefore generates a continuous family of
diffeomorphisms.

%%%%%%%%%%%%%%%%%%%%%%%%%%%%%%%%%%%%%%%%%%%%%%%%%%%%%%%%%%%%%%%

\section{Lie Derivative of Functions}

For a smooth scalar function

\[
f:M\rightarrow\mathbb R,
\]

the Lie derivative along

\[
X
\]

is simply the directional derivative,

\[
\boxed{
\mathcal L_Xf
=
X(f).
}
\]

In local coordinates,

\[
\mathcal L_Xf
=
X^i
\frac{\partial f}{\partial x^i}.
\]

%%%%%%%%%%%%%%%%%%%%%%%%%%%%%%%%%%%%%%%%%%%%%%%%%%%%%%%%%%%%%%%

\section{Lie Brackets}

Given vector fields

\[
X
\]

and

\[
Y,
\]

their Lie bracket is defined by

\[
\boxed{
[X,Y]
=
XY-YX.
}
\]

The Lie bracket measures the failure of successive infinitesimal flows
to commute.

%%%%%%%%%%%%%%%%%%%%%%%%%%%%%%%%%%%%%%%%%%%%%%%%%%%%%%%%%%%%%%%

\begin{theorem}

The collection of smooth vector fields on a manifold forms a Lie algebra
under the Lie bracket.

\end{theorem}

%%%%%%%%%%%%%%%%%%%%%%%%%%%%%%%%%%%%%%%%%%%%%%%%%%%%%%%%%%%%%%%

\section{Lie Derivative of Vector Fields}

The Lie derivative of one vector field along another is

\[
\boxed{
\mathcal L_XY
=
[X,Y].
}
\]

Thus,

vector fields remain unchanged along the flow precisely when their Lie
derivative vanishes.

%%%%%%%%%%%%%%%%%%%%%%%%%%%%%%%%%%%%%%%%%%%%%%%%%%%%%%%%%%%%%%%

\section{Lie Derivative of Tensor Fields}

If

\[
T
\]

is a tensor field,

its Lie derivative measures how

\[
T
\]

changes under infinitesimal deformations generated by

\[
X.
\]

For a tensor of type

\[
(r,s),
\]

the Lie derivative preserves tensor type and satisfies the Leibniz rule.

%%%%%%%%%%%%%%%%%%%%%%%%%%%%%%%%%%%%%%%%%%%%%%%%%%%%%%%%%%%%%%%

\section{Killing Vector Fields}

A vector field

\[
X
\]

is called a Killing field whenever

\[
\boxed{
\mathcal L_Xg
=
0,
}
\]

where

\[
g
\]

is the metric tensor.

Killing fields generate continuous symmetries of the geometry.

Examples include

translations,

rotations,

and other isometries.

%%%%%%%%%%%%%%%%%%%%%%%%%%%%%%%%%%%%%%%%%%%%%%%%%%%%%%%%%%%%%%%

\begin{theorem}

The collection of Killing vector fields forms a Lie algebra.

\end{theorem}

%%%%%%%%%%%%%%%%%%%%%%%%%%%%%%%%%%%%%%%%%%%%%%%%%%%%%%%%%%%%%%%

\section{Differential Forms Revisited}

Differential forms provide a coordinate-independent framework for
integration on manifolds.

A differential

\[
k
\]

-form assigns an alternating multilinear mapping to every tangent space.

The space of smooth

\[
k
\]

-forms is denoted

\[
\Omega^k(M).
\]

%%%%%%%%%%%%%%%%%%%%%%%%%%%%%%%%%%%%%%%%%%%%%%%%%%%%%%%%%%%%%%%

\section{The Wedge Product}

If

\[
\alpha
\]

is a

\[
p
\]

-form and

\[
\beta
\]

is a

\[
q
\]

-form,

their wedge product is

\[
\boxed{
\alpha\wedge\beta
\in
\Omega^{p+q}(M).
}
\]

The wedge product is bilinear and antisymmetric,

satisfying

\[
\alpha\wedge\beta
=
(-1)^{pq}
\beta\wedge\alpha.
\]

%%%%%%%%%%%%%%%%%%%%%%%%%%%%%%%%%%%%%%%%%%%%%%%%%%%%%%%%%%%%%%%

\section{Exterior Calculus}

The exterior derivative satisfies

\[
d:\Omega^k(M)\rightarrow\Omega^{k+1}(M),
\]

together with

\[
d^2=0.
\]

It obeys the graded Leibniz rule,

\[
\boxed{
d(\alpha\wedge\beta)
=
d\alpha\wedge\beta
+
(-1)^p
\alpha\wedge d\beta.
}
\]

Exterior calculus unifies gradient, curl, and divergence within a single
coordinate-free formalism.

%%%%%%%%%%%%%%%%%%%%%%%%%%%%%%%%%%%%%%%%%%%%%%%%%%%%%%%%%%%%%%%

\section{Interior Product}

For a vector field

\[
X,
\]

the interior product

\[
i_X
\]

contracts differential forms,

reducing degree by one.

It satisfies

\[
i_X:
\Omega^k(M)
\rightarrow
\Omega^{k-1}(M).
\]

%%%%%%%%%%%%%%%%%%%%%%%%%%%%%%%%%%%%%%%%%%%%%%%%%%%%%%%%%%%%%%%

\section{Cartan's Magic Formula}

One of the central identities of differential geometry is

\[
\boxed{
\mathcal L_X
=
d\,i_X
+
i_X\,d.
}
\]

This expresses the Lie derivative entirely in terms of the exterior
derivative and interior contraction.

Cartan's identity unifies symmetry, differentiation, and integration.

%%%%%%%%%%%%%%%%%%%%%%%%%%%%%%%%%%%%%%%%%%%%%%%%%%%%%%%%%%%%%%%

\section{Orientation}

Integration on manifolds requires a consistent choice of orientation.

An oriented manifold admits a globally consistent notion of
``positive'' volume.

Without an orientation,

integrals of differential forms cannot generally be defined globally.

%%%%%%%%%%%%%%%%%%%%%%%%%%%%%%%%%%%%%%%%%%%%%%%%%%%%%%%%%%%%%%%

\section{Volume Forms}

An oriented Riemannian manifold possesses a natural volume form,

\[
\boxed{
dV
=
\sqrt{\det(g)}
\,
dx^1
\wedge
\cdots
\wedge
dx^n.
}
\]

This form defines integration independently of coordinates.

Lengths,

areas,

volumes,

and total mass are all computed using the induced volume form.

%%%%%%%%%%%%%%%%%%%%%%%%%%%%%%%%%%%%%%%%%%%%%%%%%%%%%%%%%%%%%%%

\section{Integration of Differential Forms}

Let

\[
\omega
\]

be an

\[
n
\]

-form on an oriented

\[
n
\]

-dimensional manifold.

Its integral is written

\[
\boxed{
\int_M\omega.
}
\]

Unlike ordinary multiple integration,

this definition is invariant under smooth coordinate transformations.

%%%%%%%%%%%%%%%%%%%%%%%%%%%%%%%%%%%%%%%%%%%%%%%%%%%%%%%%%%%%%%%

\section{Change of Variables}

Suppose

\[
F:M\rightarrow N
\]

is an orientation-preserving diffeomorphism.

Then

\[
\boxed{
\int_N\omega
=
\int_MF^*\omega,
}
\]

where

\[
F^*
\]

denotes the pullback.

This theorem generalizes the classical Jacobian change-of-variables
formula.

%%%%%%%%%%%%%%%%%%%%%%%%%%%%%%%%%%%%%%%%%%%%%%%%%%%%%%%%%%%%%%%

\section{Summary}

Flows describe continuous motion on manifolds.

Lie derivatives measure intrinsic change.

Differential forms provide the natural language of integration.

The wedge product, exterior derivative, and Cartan's identity unify much
of modern differential geometry into a compact algebraic framework.

These ideas culminate in the generalized Stokes' theorem, developed in
the following section, which unifies the fundamental integral theorems
of vector calculus.

%%%%%%%%%%%%%%%%%%%%%%%%%%%%%%%%%%%%%%%%%%%%%%%%%%%%%%%%%%%%%%%

\section{Generalized Stokes' Theorem}

The generalized Stokes' theorem is one of the central results of
differential geometry.

It unifies the classical theorems of vector calculus into a single,
coordinate-independent statement.

%%%%%%%%%%%%%%%%%%%%%%%%%%%%%%%%%%%%%%%%%%%%%%%%%%%%%%%%%%%%%%%

\begin{theorem}[Generalized Stokes]

Let

\[
M
\]

be an oriented smooth manifold with boundary

\[
\partial M,
\]

and let

\[
\omega
\]

be a compactly supported differential

\[
(n-1)
\]

-form.

Then

\[
\boxed{
\int_M d\omega
=
\int_{\partial M}\omega.
}
\]

\end{theorem}

This theorem relates the behavior of a differential form within a
region to its behavior along the boundary of that region.

%%%%%%%%%%%%%%%%%%%%%%%%%%%%%%%%%%%%%%%%%%%%%%%%%%%%%%%%%%%%%%%

\section{Classical Special Cases}

Several familiar integral theorems arise as special cases of
generalized Stokes' theorem.

\subsection{Fundamental Theorem of Calculus}

For a smooth function

\[
f:[a,b]\rightarrow\mathbb R,
\]

\[
\boxed{
\int_a^b
f'(x)\,dx
=
f(b)-f(a).
}
\]

%%%%%%%%%%%%%%%%%%%%%%%%%%%%%%%%%%%%%%%%%%%%%%%%%%%%%%%%%%%%%%%

\subsection{Green's Theorem}

If

\[
D\subset\mathbb R^2,
\]

then

\[
\boxed{
\oint_{\partial D}
P\,dx+Q\,dy
=
\iint_D
\left(
\frac{\partial Q}{\partial x}
-
\frac{\partial P}{\partial y}
\right)
dx\,dy.
}
\]

%%%%%%%%%%%%%%%%%%%%%%%%%%%%%%%%%%%%%%%%%%%%%%%%%%%%%%%%%%%%%%%

\subsection{Stokes' Curl Theorem}

For an oriented surface

\[
S,
\]

\[
\boxed{
\oint_{\partial S}
\mathbf F\cdot d\mathbf r
=
\iint_S
(\nabla\times\mathbf F)\cdot\mathbf n\,dS.
}
\]

%%%%%%%%%%%%%%%%%%%%%%%%%%%%%%%%%%%%%%%%%%%%%%%%%%%%%%%%%%%%%%%

\subsection{Divergence Theorem}

For a compact region

\[
V,
\]

\[
\boxed{
\iiint_V
\nabla\cdot\mathbf F\,dV
=
\iint_{\partial V}
\mathbf F\cdot\mathbf n\,dS.
}
\]

Thus every major theorem of elementary vector calculus follows from the
single geometric identity established above.

%%%%%%%%%%%%%%%%%%%%%%%%%%%%%%%%%%%%%%%%%%%%%%%%%%%%%%%%%%%%%%%

\section{Conservation Laws}

Many physical laws express conservation of some quantity.

Let

\[
J
\]

denote a flux vector field.

The conservation law

\[
\boxed{
\nabla\cdot J
=
\sigma
}
\]

states that local accumulation equals local production.

Integrating over a region

\[
V
\]

and applying the divergence theorem gives

\[
\boxed{
\int_V\sigma\,dV
=
\int_{\partial V}
J\cdot n\,dS.
}
\]

The total production inside a region therefore equals the net outward
flux across its boundary.

%%%%%%%%%%%%%%%%%%%%%%%%%%%%%%%%%%%%%%%%%%%%%%%%%%%%%%%%%%%%%%%

\section{Applications to Infrastructure}

Generalized Stokes' theorem underlies many infrastructure models.

Examples include

\begin{itemize}
\item conservation of electrical current,
\item conservation of fluid mass,
\item conservation of traffic flow,
\item conservation of information,
\item conservation of economic commodities,
\item conservation of probability in stochastic systems.
\end{itemize}

In every case,

local differential equations are equivalent to global balance laws.

%%%%%%%%%%%%%%%%%%%%%%%%%%%%%%%%%%%%%%%%%%%%%%%%%%%%%%%%%%%%%%%

\section{The Hodge Star Operator}

A Riemannian metric induces an operator

\[
*
:
\Omega^k(M)
\rightarrow
\Omega^{n-k}(M),
\]

called the Hodge star.

The Hodge star converts differential forms into complementary forms,
allowing inner products to be defined on spaces of differential forms.

%%%%%%%%%%%%%%%%%%%%%%%%%%%%%%%%%%%%%%%%%%%%%%%%%%%%%%%%%%%%%%%

\section{The Codifferential}

Using the Hodge star,

one defines the codifferential

\[
\boxed{
\delta
=
(-1)^{nk+n+1}
*
d
*.
}
\]

The codifferential plays the role of divergence within exterior
calculus.

%%%%%%%%%%%%%%%%%%%%%%%%%%%%%%%%%%%%%%%%%%%%%%%%%%%%%%%%%%%%%%%

\section{The Laplace--Beltrami Operator}

The geometric Laplacian is defined by

\[
\boxed{
\Delta
=
d\delta
+
\delta d.
}
\]

On scalar functions,

this reduces locally to the familiar Laplacian,

while on differential forms it governs diffusion, harmonic analysis,
and many elliptic partial differential equations.

%%%%%%%%%%%%%%%%%%%%%%%%%%%%%%%%%%%%%%%%%%%%%%%%%%%%%%%%%%%%%%%

\section{Harmonic Forms}

A differential form

\[
\omega
\]

is called harmonic whenever

\[
\boxed{
\Delta\omega
=
0.
}
\]

Harmonic forms minimize natural energy functionals and play an important
role in topology, potential theory, and equilibrium transport problems.

%%%%%%%%%%%%%%%%%%%%%%%%%%%%%%%%%%%%%%%%%%%%%%%%%%%%%%%%%%%%%%%

\section{de Rham Cohomology}

Since

\[
d^2=0,
\]

every exact form is automatically closed.

The converse is generally false.

The resulting quotient space

\[
\boxed{
H^k(M)
=
\frac{\ker(d:\Omega^k\rightarrow\Omega^{k+1})}
{\operatorname{im}(d:\Omega^{k-1}\rightarrow\Omega^k)}
}
\]

is called the \(k\)-th de Rham cohomology group.

These groups encode the global topology of the manifold through the
behavior of differential forms.

%%%%%%%%%%%%%%%%%%%%%%%%%%%%%%%%%%%%%%%%%%%%%%%%%%%%%%%%%%%%%%%

\section{Hodge Decomposition}

One of the fundamental results of modern differential geometry states
that every differential form admits an orthogonal decomposition.

\begin{theorem}[Hodge Decomposition]

On a compact oriented Riemannian manifold,

every differential form can be written uniquely as

\[
\boxed{
\omega
=
d\alpha
+
\delta\beta
+
\gamma,
}
\]

where

\[
\gamma
\]

is harmonic.

\end{theorem}

This decomposition separates local variation from global topological
structure.

%%%%%%%%%%%%%%%%%%%%%%%%%%%%%%%%%%%%%%%%%%%%%%%%%%%%%%%%%%%%%%%

\section{Concluding Remarks}

Differential geometry extends the ideas of linear algebra from flat
vector spaces to curved manifolds.

Metrics measure distance.

Connections define differentiation.

Curvature quantifies intrinsic geometry.

Lie derivatives describe continuous symmetries.

Differential forms provide the natural language of integration, while
Stokes' theorem unifies the classical integral theorems of analysis.

Together, these concepts furnish the geometric framework required for
the study of infrastructure embedded in physical space and support the
optimization, control, and transport theories developed throughout this
volume.

The next appendix turns from geometry to variational calculus, where
optimization principles are formulated through functionals, Euler--
Lagrange equations, constrained extrema, and Hamiltonian methods.

%%%%%%%%%%%%%%%%%%%%%%%%%%%%%%%%%%%%%%%%%%%%%%%%%%%%%%%%%%%%%%%%%%%%%%
\chapter{Variational Calculus}

\section{Purpose}

Many of the optimization problems encountered throughout this volume are
most naturally formulated as variational problems.

Rather than optimizing finitely many variables,

one seeks an optimal function,

curve,

surface,

or field.

Examples include

\begin{itemize}
\item shortest transportation routes,
\item minimum-energy power distributions,
\item optimal pipeline layouts,
\item traffic equilibria,
\item geodesic construction,
\item distributed control,
\item resource allocation over time.
\end{itemize}

Variational calculus provides the mathematical framework for such
problems.

%%%%%%%%%%%%%%%%%%%%%%%%%%%%%%%%%%%%%%%%%%%%%%%%%%%%%%%%%%%%%%%

\section{Functionals}

Unlike an ordinary function,

which maps numbers to numbers,

a functional maps functions to numbers.

\begin{definition}

A functional is a mapping

\[
\mathcal J
:
\mathcal F
\rightarrow
\mathbb R,
\]

where

\[
\mathcal F
\]

is a space of admissible functions.

\end{definition}

Typical examples are

\[
\boxed{
\mathcal J[y]
=
\int_a^b
L(x,y,y')
\,dx,
}
\]

where

\[
L
\]

is called the Lagrangian.

%%%%%%%%%%%%%%%%%%%%%%%%%%%%%%%%%%%%%%%%%%%%%%%%%%%%%%%%%%%%%%%

\section{Admissible Functions}

Optimization takes place over an admissible class of functions.

Usually one prescribes boundary values,

such as

\[
y(a)=A,
\qquad
y(b)=B.
\]

Additional constraints,

including inequality or integral constraints,

may also be imposed.

%%%%%%%%%%%%%%%%%%%%%%%%%%%%%%%%%%%%%%%%%%%%%%%%%%%%%%%%%%%%%%%

\section{Variations}

A variation is a small perturbation of an admissible function.

Given

\[
y(x),
\]

define

\[
\boxed{
y_\varepsilon(x)
=
y(x)
+
\varepsilon
\eta(x),
}
\]

where

\[
\eta(a)
=
\eta(b)
=
0.
\]

The function

\[
\eta
\]

is called the variation.

%%%%%%%%%%%%%%%%%%%%%%%%%%%%%%%%%%%%%%%%%%%%%%%%%%%%%%%%%%%%%%%

\section{First Variation}

The first variation of

\[
\mathcal J
\]

is

\[
\boxed{
\delta\mathcal J
=
\left.
\frac{d}{d\varepsilon}
\mathcal J[y_\varepsilon]
\right|_{\varepsilon=0}.
}
\]

Necessary conditions for extrema require

\[
\delta\mathcal J=0.
\]

%%%%%%%%%%%%%%%%%%%%%%%%%%%%%%%%%%%%%%%%%%%%%%%%%%%%%%%%%%%%%%%

\section{Derivation of the Euler--Lagrange Equation}

Consider

\[
\mathcal J[y]
=
\int_a^b
L(x,y,y')
\,dx.
\]

Substituting

\[
y+\varepsilon\eta
\]

gives

\[
\delta\mathcal J
=
\int_a^b
\left(
\frac{\partial L}{\partial y}\eta
+
\frac{\partial L}{\partial y'}
\eta'
\right)
dx.
\]

Integrating the second term by parts,

and using

\[
\eta(a)
=
\eta(b)
=
0,
\]

yields

\[
\delta\mathcal J
=
\int_a^b
\left(
\frac{\partial L}{\partial y}
-
\frac{d}{dx}
\frac{\partial L}{\partial y'}
\right)
\eta
\,dx.
\]

Since

\[
\eta
\]

is arbitrary,

the Fundamental Lemma of the Calculus of Variations implies the
Euler--Lagrange equation.

%%%%%%%%%%%%%%%%%%%%%%%%%%%%%%%%%%%%%%%%%%%%%%%%%%%%%%%%%%%%%%%

\begin{theorem}[Euler--Lagrange Equation]

A necessary condition for an extremum is

\[
\boxed{
\frac{\partial L}{\partial y}
-
\frac{d}{dx}
\left(
\frac{\partial L}{\partial y'}
\right)
=
0.
}
\]

\end{theorem}

This equation forms the foundation of modern optimization,
classical mechanics,
field theory,
and optimal infrastructure design.

%%%%%%%%%%%%%%%%%%%%%%%%%%%%%%%%%%%%%%%%%%%%%%%%%%%%%%%%%%%%%%%

\section{Multiple Variables}

Suppose

\[
y=(y^1,\ldots,y^m).
\]

Then each component satisfies

\[
\boxed{
\frac{\partial L}{\partial y^i}
-
\frac{d}{dx}
\left(
\frac{\partial L}{\partial y^{i\,\prime}}
\right)
=
0,
\qquad
i=1,\ldots,m.
}
\]

The resulting equations are generally coupled.

%%%%%%%%%%%%%%%%%%%%%%%%%%%%%%%%%%%%%%%%%%%%%%%%%%%%%%%%%%%%%%%

\section{Several Independent Variables}

If

\[
u=u(x^1,\ldots,x^n),
\]

the functional becomes

\[
\mathcal J[u]
=
\int_\Omega
L
(x,u,\nabla u)
\,dV.
\]

The Euler--Lagrange equation generalizes to

\[
\boxed{
\frac{\partial L}{\partial u}
-
\sum_i
\frac{\partial}{\partial x^i}
\left(
\frac{\partial L}
{\partial u_{x^i}}
\right)
=
0.
}
\]

Many partial differential equations arise from precisely this
construction.

%%%%%%%%%%%%%%%%%%%%%%%%%%%%%%%%%%%%%%%%%%%%%%%%%%%%%%%%%%%%%%%

\section{Natural Boundary Conditions}

When endpoint values are not prescribed,

additional boundary conditions arise naturally from the integration by
parts.

These are called natural or transversality conditions and form an
essential part of many free-boundary optimization problems.

%%%%%%%%%%%%%%%%%%%%%%%%%%%%%%%%%%%%%%%%%%%%%%%%%%%%%%%%%%%%%%%

\section{Summary}

Variational calculus transforms optimization over functions into
differential equations.

The central concepts are

functionals,

variations,

and stationarity.

The Euler--Lagrange equation provides the necessary condition for an
extremum and serves as the starting point for much of modern mechanics,
optimal control, continuum physics, and infrastructure optimization.

The following sections develop constrained variational principles,
Hamiltonian formulations, and optimal control theory.

%%%%%%%%%%%%%%%%%%%%%%%%%%%%%%%%%%%%%%%%%%%%%%%%%%%%%%%%%%%%%%%

\section{Constrained Variational Problems}

Many optimization problems require the admissible functions to satisfy
additional constraints.

These constraints may be algebraic,

differential,

integral,

or geometric.

Typical examples include

\begin{itemize}
\item conservation of mass,
\item fixed total resource budgets,
\item prescribed network capacities,
\item boundary constraints,
\item geometric embedding conditions.
\end{itemize}

%%%%%%%%%%%%%%%%%%%%%%%%%%%%%%%%%%%%%%%%%%%%%%%%%%%%%%%%%%%%%%%

\section{Lagrange Multipliers}

Suppose one wishes to extremize

\[
\mathcal J[y]
\]

subject to the constraint

\[
\mathcal G[y]=0.
\]

Introduce a multiplier

\[
\lambda,
\]

and define the augmented functional

\[
\boxed{
\widetilde{\mathcal J}
=
\mathcal J
+
\lambda
\mathcal G.
}
\]

Stationarity is then imposed upon

\[
\widetilde{\mathcal J}
\]

rather than upon

\[
\mathcal J
\]

alone.

%%%%%%%%%%%%%%%%%%%%%%%%%%%%%%%%%%%%%%%%%%%%%%%%%%%%%%%%%%%%%%%

\begin{theorem}

Necessary conditions for constrained extrema are obtained by requiring

\[
\delta\widetilde{\mathcal J}
=
0.
\]

The resulting Euler--Lagrange equations are supplemented by the original
constraint equations.

\end{theorem}

%%%%%%%%%%%%%%%%%%%%%%%%%%%%%%%%%%%%%%%%%%%%%%%%%%%%%%%%%%%%%%%

\section{Isoperimetric Problems}

A classical constrained variational problem seeks to extremize

\[
\mathcal J[y]
=
\int_a^b
L(x,y,y')
\,dx
\]

subject to

\[
\boxed{
\int_a^b
G(x,y,y')
\,dx
=
C,
}
\]

where

\[
C
\]

is prescribed.

Introducing a constant multiplier

\[
\lambda,
\]

one instead extremizes

\[
\boxed{
\int_a^b
\left(
L
+
\lambda G
\right)
dx.
}
\]

%%%%%%%%%%%%%%%%%%%%%%%%%%%%%%%%%%%%%%%%%%%%%%%%%%%%%%%%%%%%%%%

\section{Integral Constraints}

Infrastructure optimization frequently imposes limits upon total
resources.

Examples include

\[
\int_\Omega
\rho(x)
\,dV
=
M,
\]

representing

fixed material,

fixed investment,

or prescribed energy budgets.

Such constraints naturally enter the variational formulation through
Lagrange multipliers.

%%%%%%%%%%%%%%%%%%%%%%%%%%%%%%%%%%%%%%%%%%%%%%%%%%%%%%%%%%%%%%%

\section{Variational Principles in Mechanics}

Classical mechanics is governed by Hamilton's Principle.

Let

\[
L(q,\dot q,t)
=
T-V
\]

denote the Lagrangian,

where

\[
T
\]

is kinetic energy and

\[
V
\]

is potential energy.

The action is

\[
\boxed{
S[q]
=
\int_{t_1}^{t_2}
L(q,\dot q,t)
\,dt.
}
\]

The physical trajectory renders the action stationary.

%%%%%%%%%%%%%%%%%%%%%%%%%%%%%%%%%%%%%%%%%%%%%%%%%%%%%%%%%%%%%%%

\begin{theorem}[Hamilton's Principle]

The equations of motion satisfy

\[
\boxed{
\delta S
=
0,
}
\]

which is equivalent to the Euler--Lagrange equations.

\end{theorem}

%%%%%%%%%%%%%%%%%%%%%%%%%%%%%%%%%%%%%%%%%%%%%%%%%%%%%%%%%%%%%%%

\section{Conservation Laws}

Suppose the Lagrangian contains no explicit dependence upon one of its
variables.

Then a corresponding conserved quantity exists.

For example,

if

\[
L
\]

contains no explicit dependence upon

\[
x,
\]

then

\[
\boxed{
L
-
y'
\frac{\partial L}{\partial y'}
=
\text{constant}.
}
\]

This identity is known as the Beltrami identity.

%%%%%%%%%%%%%%%%%%%%%%%%%%%%%%%%%%%%%%%%%%%%%%%%%%%%%%%%%%%%%%%

\section{Noether's Theorem}

One of the deepest results in variational calculus relates symmetry to
conservation.

\begin{theorem}[Noether]

Every continuous symmetry of the action functional corresponds to a
conserved quantity.

\end{theorem}

Examples include

\begin{center}
\begin{tabular}{ll}
\toprule
Symmetry & Conserved Quantity\\
\midrule
Time translation & Energy\\
Spatial translation & Linear momentum\\
Rotation & Angular momentum\\
Gauge symmetry & Charge\\
\bottomrule
\end{tabular}
\end{center}

This theorem unifies geometry, optimization, and conservation.

%%%%%%%%%%%%%%%%%%%%%%%%%%%%%%%%%%%%%%%%%%%%%%%%%%%%%%%%%%%%%%%

\section{Hamiltonian Formulation}

Define the generalized momentum

\[
\boxed{
p
=
\frac{\partial L}{\partial\dot q}.
}
\]

The Hamiltonian is obtained through the Legendre transformation,

\[
\boxed{
H(q,p)
=
p\dot q
-
L.
}
\]

The equations of motion become

\[
\boxed{
\dot q
=
\frac{\partial H}{\partial p},
\qquad
\dot p
=
-
\frac{\partial H}{\partial q}.
}
\]

Hamiltonian systems often possess advantageous geometric and numerical
properties.

%%%%%%%%%%%%%%%%%%%%%%%%%%%%%%%%%%%%%%%%%%%%%%%%%%%%%%%%%%%%%%%

\section{Canonical Transformations}

A transformation

\[
(q,p)
\mapsto
(Q,P)
\]

is canonical if it preserves Hamilton's equations.

Canonical transformations simplify many optimization and dynamical
problems while preserving their essential structure.

%%%%%%%%%%%%%%%%%%%%%%%%%%%%%%%%%%%%%%%%%%%%%%%%%%%%%%%%%%%%%%%

\section{Variational Inequalities}

Not every optimization problem possesses equality constraints.

Frequently one seeks

\[
u
\]

satisfying

\[
\boxed{
\langle
F(u),
v-u
\rangle
\ge
0
\qquad
\forall v\in K,
}
\]

where

\[
K
\]

is a convex admissible set.

Variational inequalities arise naturally in

traffic equilibrium,

contact mechanics,

network congestion,

and capacity-limited infrastructure.

%%%%%%%%%%%%%%%%%%%%%%%%%%%%%%%%%%%%%%%%%%%%%%%%%%%%%%%%%%%%%%%

\section{Existence of Minimizers}

The existence of a minimizing function cannot be assumed.

A typical existence theorem requires

\begin{itemize}
\item coercivity,
\item lower semicontinuity,
\item weak compactness of the admissible set.
\end{itemize}

Under these hypotheses,

the direct method of the calculus of variations guarantees the existence
of a minimizer.

%%%%%%%%%%%%%%%%%%%%%%%%%%%%%%%%%%%%%%%%%%%%%%%%%%%%%%%%%%%%%%%

\section{Summary}

Variational methods extend optimization from finite-dimensional vectors
to infinite-dimensional function spaces.

Constraints are incorporated through Lagrange multipliers.

Hamilton's principle connects stationary action with physical dynamics.

Noether's theorem reveals that conservation laws arise from continuous
symmetries.

Hamiltonian formulations provide an equivalent first-order description,
while variational inequalities extend the theory to constrained and
capacity-limited systems.

These ideas form the mathematical foundation for the optimal control
theory developed in the following appendix.

%%%%%%%%%%%%%%%%%%%%%%%%%%%%%%%%%%%%%%%%%%%%%%%%%%%%%%%%%%%%%%%%%%%%%%
\chapter{Optimal Control Theory}

\section{Purpose}

Variational calculus determines optimal trajectories when the entire
trajectory is chosen simultaneously.

Optimal control extends this framework to dynamical systems whose
evolution may be influenced continuously through time.

Applications throughout this volume include

\begin{itemize}
\item traffic signal coordination,
\item electrical grid management,
\item water distribution,
\item transportation routing,
\item inventory regulation,
\item autonomous infrastructure,
\item resilient network recovery.
\end{itemize}

%%%%%%%%%%%%%%%%%%%%%%%%%%%%%%%%%%%%%%%%%%%%%%%%%%%%%%%%%%%%%%%

\section{Control Systems}

A control system is described by a differential equation

\[
\boxed{
\dot{x}(t)
=
f(x(t),u(t),t),
}
\]

where

\[
x(t)\in\mathbb{R}^n
\]

is the state vector and

\[
u(t)\in\mathbb{R}^m
\]

is the control input.

The function

\[
f
\]

determines the system dynamics.

%%%%%%%%%%%%%%%%%%%%%%%%%%%%%%%%%%%%%%%%%%%%%%%%%%%%%%%%%%%%%%%

\section{Admissible Controls}

Not every control function is physically realizable.

The admissible controls satisfy constraints such as

\[
u(t)\in U,
\]

where

\[
U
\]

is a prescribed feasible set.

Typical constraints include

\begin{itemize}
\item actuator limits,
\item capacity bounds,
\item safety constraints,
\item resource limitations,
\item regulatory requirements.
\end{itemize}

%%%%%%%%%%%%%%%%%%%%%%%%%%%%%%%%%%%%%%%%%%%%%%%%%%%%%%%%%%%%%%%

\section{Performance Functionals}

The objective is to minimize a cost functional

\[
\boxed{
J[u]
=
\Phi(x(T))
+
\int_{0}^{T}
L(x,u,t)\,dt,
}
\]

where

\begin{itemize}
\item
\[
\Phi
\]
is the terminal cost,

\item
\[
L
\]
is the running cost.
\end{itemize}

Examples include fuel consumption,

travel time,

energy expenditure,

or operational cost.

%%%%%%%%%%%%%%%%%%%%%%%%%%%%%%%%%%%%%%%%%%%%%%%%%%%%%%%%%%%%%%%

\section{The Hamiltonian}

Introduce the costate variable

\[
\lambda(t).
\]

The optimal-control Hamiltonian is

\[
\boxed{
H(x,u,\lambda,t)
=
L(x,u,t)
+
\lambda^{T}f(x,u,t).
}
\]

Unlike the Hamiltonian of classical mechanics,

this Hamiltonian incorporates both dynamics and optimization.

%%%%%%%%%%%%%%%%%%%%%%%%%%%%%%%%%%%%%%%%%%%%%%%%%%%%%%%%%%%%%%%

\section{Pontryagin's Maximum Principle}

The fundamental theorem of optimal control provides necessary conditions
for optimality.

\begin{theorem}[Pontryagin Maximum Principle]

If

\[
u^*(t)
\]

is optimal,

then there exists a costate trajectory

\[
\lambda(t)
\]

such that

\[
\boxed{
\dot{x}
=
\frac{\partial H}{\partial\lambda},
}
\]

\[
\boxed{
\dot{\lambda}
=
-
\frac{\partial H}{\partial x},
}
\]

and

\[
\boxed{
H(x,u^*,\lambda,t)
=
\min_{u\in U}
H(x,u,\lambda,t).
}
\]

\end{theorem}

%%%%%%%%%%%%%%%%%%%%%%%%%%%%%%%%%%%%%%%%%%%%%%%%%%%%%%%%%%%%%%%

\section{Boundary Conditions}

Typically,

the initial state is prescribed,

\[
x(0)=x_0,
\]

while the terminal cost determines the final costate,

\[
\boxed{
\lambda(T)
=
\nabla\Phi(x(T)).
}
\]

These conditions define a two-point boundary-value problem.

%%%%%%%%%%%%%%%%%%%%%%%%%%%%%%%%%%%%%%%%%%%%%%%%%%%%%%%%%%%%%%%

\section{Linear Quadratic Regulator}

One of the most important optimal-control problems is the Linear
Quadratic Regulator (LQR).

The dynamics are

\[
\dot{x}
=
Ax
+
Bu,
\]

and the objective is

\[
\boxed{
J
=
\int_0^\infty
\left(
x^TQx
+
u^TRu
\right)
dt.
}
\]

Here,

\[
Q
\]

and

\[
R
\]

are symmetric positive-definite weighting matrices.

%%%%%%%%%%%%%%%%%%%%%%%%%%%%%%%%%%%%%%%%%%%%%%%%%%%%%%%%%%%%%%%

\section{Optimal Feedback}

The optimal control law is linear,

\[
\boxed{
u
=
-
Kx,
}
\]

where

\[
K
\]

is obtained from the algebraic Riccati equation.

Feedback control continuously adjusts the system in response to its
current state, providing robustness against disturbances.

%%%%%%%%%%%%%%%%%%%%%%%%%%%%%%%%%%%%%%%%%%%%%%%%%%%%%%%%%%%%%%%

\section{The Riccati Equation}

The feedback matrix is determined by the symmetric matrix

\[
P,
\]

which satisfies

\[
\boxed{
A^{T}P
+
PA
-
PBR^{-1}B^{T}P
+
Q
=
0.
}
\]

The resulting controller minimizes the quadratic performance index while
stabilizing the system.

%%%%%%%%%%%%%%%%%%%%%%%%%%%%%%%%%%%%%%%%%%%%%%%%%%%%%%%%%%%%%%%

\section{Dynamic Programming}

An alternative approach to optimal control is Bellman's Principle of
Optimality.

\begin{quote}
An optimal policy has the property that every remaining portion of the
policy must itself be optimal.
\end{quote}

This recursive principle transforms global optimization into a sequence
of local optimization problems.

%%%%%%%%%%%%%%%%%%%%%%%%%%%%%%%%%%%%%%%%%%%%%%%%%%%%%%%%%%%%%%%

\section{The Value Function}

Define

\[
V(x,t)
\]

as the minimum achievable future cost beginning from state

\[
x
\]

at time

\[
t.
\]

The value function completely characterizes the optimal control problem.

%%%%%%%%%%%%%%%%%%%%%%%%%%%%%%%%%%%%%%%%%%%%%%%%%%%%%%%%%%%%%%%

\section{Hamilton--Jacobi--Bellman Equation}

The value function satisfies the nonlinear partial differential equation

\[
\boxed{
-\frac{\partial V}{\partial t}
=
\min_u
\left[
L(x,u,t)
+
\nabla V\cdot f(x,u,t)
\right].
}
\]

The Hamilton--Jacobi--Bellman equation provides both necessary and
sufficient conditions for optimality under appropriate regularity
assumptions.

%%%%%%%%%%%%%%%%%%%%%%%%%%%%%%%%%%%%%%%%%%%%%%%%%%%%%%%%%%%%%%%

\section{Comparison of Approaches}

Pontryagin's Maximum Principle and Dynamic Programming solve the same
optimization problem from complementary perspectives.

\begin{center}
\begin{tabular}{ll}
\toprule
Pontryagin & Dynamic Programming\\
\midrule
Boundary-value formulation & Value-function formulation\\
Necessary conditions & Necessary and sufficient conditions\\
Costate variables & Value function\\
Ordinary differential equations & Partial differential equation\\
\bottomrule
\end{tabular}
\end{center}

Each framework has computational advantages depending upon the structure
of the problem.

%%%%%%%%%%%%%%%%%%%%%%%%%%%%%%%%%%%%%%%%%%%%%%%%%%%%%%%%%%%%%%%%%%%%%%
\section{Discrete-Time Optimal Control}

Many engineering systems evolve in discrete time rather than
continuously.

Let

\[
x_{k+1}
=
f(x_k,u_k),
\]

where

\[
k=0,1,\ldots,N.
\]

The performance functional becomes

\[
\boxed{
J
=
\Phi(x_N)
+
\sum_{k=0}^{N-1}
L(x_k,u_k).
}
\]

Discrete formulations are fundamental for digital controllers,
computer simulations, and optimization algorithms.

%%%%%%%%%%%%%%%%%%%%%%%%%%%%%%%%%%%%%%%%%%%%%%%%%%%%%%%%%%%%%%%

\section{Model Predictive Control}

Model Predictive Control (MPC) repeatedly solves a finite-horizon
optimization problem while applying only the first control action.

At each time step,

\begin{enumerate}
\item measure the current state,
\item predict future evolution,
\item solve a constrained optimization problem,
\item apply the first control input,
\item repeat using updated measurements.
\end{enumerate}

The optimization problem is

\[
\boxed{
\min_{u_0,\ldots,u_{N-1}}
\left[
\Phi(x_N)
+
\sum_{k=0}^{N-1}
L(x_k,u_k)
\right]
}
\]

subject to

\[
x_{k+1}=f(x_k,u_k),
\]

together with all operational constraints.

MPC naturally incorporates changing demand,
capacity limits,
and uncertainty.

%%%%%%%%%%%%%%%%%%%%%%%%%%%%%%%%%%%%%%%%%%%%%%%%%%%%%%%%%%%%%%%

\section{State Constraints}

Infrastructure systems frequently require state variables to remain
inside admissible regions.

Examples include

\[
x(t)\in X,
\]

where

\[
X
\]

may represent

\begin{itemize}
\item safe operating temperatures,
\item reservoir capacities,
\item voltage limits,
\item traffic density limits,
\item environmental regulations.
\end{itemize}

Such constraints convert unconstrained optimization into constrained
optimal control.

%%%%%%%%%%%%%%%%%%%%%%%%%%%%%%%%%%%%%%%%%%%%%%%%%%%%%%%%%%%%%%%

\section{Control Constraints}

Similarly,

control actions satisfy

\[
u_{\min}
\le
u
\le
u_{\max}.
\]

Typical examples include

maximum pump rates,

generator output limits,

speed limits,

fuel restrictions,

and actuator saturation.

%%%%%%%%%%%%%%%%%%%%%%%%%%%%%%%%%%%%%%%%%%%%%%%%%%%%%%%%%%%%%%%

\section{Stochastic Optimal Control}

Many systems evolve under uncertainty.

The dynamics become

\[
dx
=
f(x,u)\,dt
+
G(x,u)\,dW_t,
\]

where

\[
W_t
\]

denotes a Wiener process.

The objective is generally to minimize expected cost,

\[
\boxed{
\min
\;
\mathbb E[J].
}
\]

Stochastic control combines optimization with probability theory.

%%%%%%%%%%%%%%%%%%%%%%%%%%%%%%%%%%%%%%%%%%%%%%%%%%%%%%%%%%%%%%%

\section{Robust Control}

Instead of assuming probabilistic uncertainty,

robust control seeks satisfactory performance under all admissible
disturbances.

One commonly studies

\[
\boxed{
\min_u
\;
\max_{w\in\mathcal W}
J(u,w),
}
\]

where

\[
w
\]

represents unknown disturbances.

Robust methods emphasize guaranteed performance rather than average
performance.

%%%%%%%%%%%%%%%%%%%%%%%%%%%%%%%%%%%%%%%%%%%%%%%%%%%%%%%%%%%%%%%

\section{Distributed Control}

Large infrastructure systems consist of many interacting subsystems.

Rather than a single centralized controller,

one considers

\[
N
\]

controllers,

each operating with partial information.

The overall objective becomes

\[
\boxed{
J
=
\sum_{i=1}^{N}
J_i
+
J_{\mathrm{coupling}},
}
\]

where the coupling term represents interactions among subsystems.

Distributed control is particularly important for

\begin{itemize}
\item electrical grids,
\item transportation systems,
\item communication networks,
\item autonomous vehicle fleets,
\item water distribution systems.
\end{itemize}

%%%%%%%%%%%%%%%%%%%%%%%%%%%%%%%%%%%%%%%%%%%%%%%%%%%%%%%%%%%%%%%

\section{Numerical Solution Methods}

Most optimal-control problems admit no closed-form solution.

Common numerical approaches include

\begin{itemize}
\item shooting methods,
\item multiple shooting,
\item direct collocation,
\item direct transcription,
\item sequential quadratic programming,
\item interior-point optimization,
\item dynamic programming,
\item policy iteration,
\item value iteration.
\end{itemize}

The choice of method depends upon the dimensionality,
constraints,
and smoothness of the problem.

%%%%%%%%%%%%%%%%%%%%%%%%%%%%%%%%%%%%%%%%%%%%%%%%%%%%%%%%%%%%%%%

\section{Applications to Infrastructure}

Optimal control appears throughout infrastructure science.

Representative examples include

\begin{center}
\begin{tabular}{ll}
\toprule
System & Control Variable\\
\midrule
Electrical grid & Generator dispatch\\
Transportation & Signal timing\\
Water network & Pump scheduling\\
Communication network & Routing policies\\
Supply chain & Inventory replenishment\\
Satellite constellation & Orbital maneuvers\\
Emergency response & Resource allocation\\
\bottomrule
\end{tabular}
\end{center}

Although the physical systems differ,

their mathematical formulations share the same underlying optimal-control
structure.

%%%%%%%%%%%%%%%%%%%%%%%%%%%%%%%%%%%%%%%%%%%%%%%%%%%%%%%%%%%%%%%

\section{Concluding Remarks}

Optimal control extends variational calculus from static optimization to
dynamic decision making.

State equations describe system evolution.

Control variables influence future trajectories.

Performance functionals quantify competing objectives.

Pontryagin's Maximum Principle and Dynamic Programming provide
complementary theoretical foundations,

while modern numerical algorithms enable practical computation for
large-scale systems.

Throughout this volume,

optimal control serves as the mathematical bridge connecting geometric
models,
network dynamics,
optimization,
and adaptive infrastructure management.

The following appendix develops the functional-analytic foundations upon
which infinite-dimensional optimization, partial differential equations,
and operator theory are built.

%%%%%%%%%%%%%%%%%%%%%%%%%%%%%%%%%%%%%%%%%%%%%%%%%%%%%%%%%%%%%%%%%%%%%%
\chapter{Functional Analysis}

\section{Purpose}

Many problems in optimization, partial differential equations, control
theory, and continuum mechanics are naturally formulated on
infinite-dimensional spaces.

Finite-dimensional linear algebra provides the necessary intuition, but
many of its fundamental results require careful extension.

Functional analysis studies vector spaces of functions together with
their topological and geometric structure.

Throughout this appendix we develop the principal concepts required for
the analysis of infinite-dimensional optimization and differential
equations.

%%%%%%%%%%%%%%%%%%%%%%%%%%%%%%%%%%%%%%%%%%%%%%%%%%%%%%%%%%%%%%%

\section{Normed Vector Spaces}

A normed vector space is a vector space

\[
V
\]

equipped with a norm

\[
\|\cdot\|:V\rightarrow\mathbb R,
\]

satisfying

\[
\|x\|\ge0,
\]

\[
\|x\|=0
\Longleftrightarrow
x=0,
\]

\[
\|\alpha x\|
=
|\alpha|
\|x\|,
\]

and

\[
\|x+y\|
\le
\|x\|
+
\|y\|.
\]

The norm induces a metric,

\[
d(x,y)
=
\|x-y\|,
\]

thereby providing a notion of convergence.

%%%%%%%%%%%%%%%%%%%%%%%%%%%%%%%%%%%%%%%%%%%%%%%%%%%%%%%%%%%%%%%

\section{Examples}

Common normed spaces include

\[
\mathbb R^n,
\]

with the Euclidean norm,

the space

\[
C([a,b])
\]

of continuous functions,

and sequence spaces such as

\[
\ell^p.
\]

Many state spaces appearing throughout this volume belong naturally to
one of these examples.

%%%%%%%%%%%%%%%%%%%%%%%%%%%%%%%%%%%%%%%%%%%%%%%%%%%%%%%%%%%%%%%

\section{Banach Spaces}

Completeness is one of the central concepts of functional analysis.

\begin{definition}

A Banach space is a normed vector space in which every Cauchy sequence
converges to a point within the space.

\end{definition}

Completeness guarantees that limits of convergent approximation schemes
remain admissible.

%%%%%%%%%%%%%%%%%%%%%%%%%%%%%%%%%%%%%%%%%%%%%%%%%%%%%%%%%%%%%%%

\section{Hilbert Spaces}

A Hilbert space is a complete inner-product space.

Its norm is induced by the inner product,

\[
\boxed{
\|x\|
=
\sqrt{\langle x,x\rangle}.
}
\]

Hilbert spaces generalize Euclidean geometry to infinitely many
dimensions.

%%%%%%%%%%%%%%%%%%%%%%%%%%%%%%%%%%%%%%%%%%%%%%%%%%%%%%%%%%%%%%%

\section{Examples of Hilbert Spaces}

Important examples include

\[
L^2(\Omega),
\]

the space of square-integrable functions,

and

\[
\ell^2,
\]

the space of square-summable sequences.

These spaces arise naturally in signal processing,
quantum mechanics,
diffusion,
and variational formulations of partial differential equations.

%%%%%%%%%%%%%%%%%%%%%%%%%%%%%%%%%%%%%%%%%%%%%%%%%%%%%%%%%%%%%%%

\section{Convergence}

Several notions of convergence are distinguished in functional
analysis.

A sequence

\[
x_n
\]

converges strongly to

\[
x
\]

if

\[
\boxed{
\|x_n-x\|
\rightarrow
0.
}
\]

Strong convergence corresponds to convergence in norm.

%%%%%%%%%%%%%%%%%%%%%%%%%%%%%%%%%%%%%%%%%%%%%%%%%%%%%%%%%%%%%%%

\section{Weak Convergence}

A sequence

\[
x_n
\]

converges weakly to

\[
x
\]

whenever

\[
\boxed{
f(x_n)
\rightarrow
f(x)
}
\]

for every continuous linear functional

\[
f.
\]

Weak convergence is generally easier to establish than strong
convergence and plays a fundamental role in existence theorems.

%%%%%%%%%%%%%%%%%%%%%%%%%%%%%%%%%%%%%%%%%%%%%%%%%%%%%%%%%%%%%%%

\section{Linear Operators}

A linear operator between Banach spaces is a mapping

\[
T:X\rightarrow Y
\]

satisfying

\[
T(ax+by)
=
aT(x)
+
bT(y).
\]

Operators generalize matrices to infinite-dimensional spaces.

%%%%%%%%%%%%%%%%%%%%%%%%%%%%%%%%%%%%%%%%%%%%%%%%%%%%%%%%%%%%%%%

\section{Bounded Operators}

A linear operator is bounded whenever there exists

\[
C>0
\]

such that

\[
\boxed{
\|Tx\|
\le
C
\|x\|
}
\]

for every

\[
x.
\]

Boundedness is equivalent to continuity for linear operators.

%%%%%%%%%%%%%%%%%%%%%%%%%%%%%%%%%%%%%%%%%%%%%%%%%%%%%%%%%%%%%%%

\section{Operator Norm}

The norm of a bounded operator is

\[
\boxed{
\|T\|
=
\sup_{x\neq0}
\frac{\|Tx\|}{\|x\|}.
}
\]

Operator norms measure the maximal amplification produced by the
operator.

%%%%%%%%%%%%%%%%%%%%%%%%%%%%%%%%%%%%%%%%%%%%%%%%%%%%%%%%%%%%%%%

\section{Dual Spaces}

The dual space of

\[
X
\]

is

\[
\boxed{
X^*
=
\{
f:X\rightarrow\mathbb R
\mid
f
\text{ is linear and continuous}
\}.
}
\]

Elements of the dual space are called continuous linear functionals.

Duality is fundamental throughout optimization and variational analysis.

%%%%%%%%%%%%%%%%%%%%%%%%%%%%%%%%%%%%%%%%%%%%%%%%%%%%%%%%%%%%%%%

\section{The Riesz Representation Theorem}

Hilbert spaces possess a particularly elegant duality theory.

\begin{theorem}[Riesz Representation]

Let

\[
H
\]

be a Hilbert space.

Every continuous linear functional

\[
f\in H^*
\]

may be written uniquely as

\[
\boxed{
f(x)
=
\langle x,y\rangle
}
\]

for a unique element

\[
y\in H.
\]

\end{theorem}

Thus the dual of a Hilbert space is naturally identified with the space
itself.

%%%%%%%%%%%%%%%%%%%%%%%%%%%%%%%%%%%%%%%%%%%%%%%%%%%%%%%%%%%%%%%

\section{Orthogonal Complements}

For a subspace

\[
M
\subseteq
H,
\]

its orthogonal complement is

\[
\boxed{
M^\perp
=
\{
x\in H
:
\langle x,m\rangle=0
\quad
\forall m\in M
\}.
}
\]

Orthogonal decomposition underlies least-squares approximation,
projection methods, and Galerkin techniques.

%%%%%%%%%%%%%%%%%%%%%%%%%%%%%%%%%%%%%%%%%%%%%%%%%%%%%%%%%%%%%%%

\section{Projection Theorem}

\begin{theorem}

Every closed convex subset of a Hilbert space possesses a unique nearest
point to every element of the space.

Equivalently,

orthogonal projections onto closed subspaces always exist and are
unique.

\end{theorem}

Projection operators form the mathematical basis of approximation
methods used throughout numerical optimization and finite-element
analysis.

%%%%%%%%%%%%%%%%%%%%%%%%%%%%%%%%%%%%%%%%%%%%%%%%%%%%%%%%%%%%%%%

\section{Summary}

Functional analysis extends linear algebra from finite-dimensional
vector spaces to spaces of functions.

Norms define convergence,

Banach spaces ensure completeness,

Hilbert spaces generalize Euclidean geometry,

bounded operators replace matrices,

and dual spaces provide the natural setting for optimization and
variational principles.

These concepts provide the analytical foundation for the study of
partial differential equations developed in the following appendix.

%%%%%%%%%%%%%%%%%%%%%%%%%%%%%%%%%%%%%%%%%%%%%%%%%%%%%%%%%%%%%%%

\section{Compact Operators}

Compact operators play a role in infinite-dimensional analysis similar
to that of matrices in finite-dimensional linear algebra.

\begin{definition}

A bounded linear operator

\[
T:X\rightarrow Y
\]

is called compact if every bounded sequence

\[
\{x_n\}
\]

contains a subsequence for which

\[
\{Tx_n\}
\]

converges.

\end{definition}

Compact operators frequently arise from integral equations,
elliptic partial differential equations,
and inverse problems.

%%%%%%%%%%%%%%%%%%%%%%%%%%%%%%%%%%%%%%%%%%%%%%%%%%%%%%%%%%%%%%%

\section{The Spectrum of an Operator}

Let

\[
T:X\rightarrow X
\]

be a bounded linear operator.

The spectrum of

\[
T
\]

is

\[
\boxed{
\sigma(T)
=
\{
\lambda\in\mathbb C
:
T-\lambda I
\text{ is not invertible}
\}.
}
\]

The spectrum generalizes the notion of eigenvalues to
infinite-dimensional spaces.

%%%%%%%%%%%%%%%%%%%%%%%%%%%%%%%%%%%%%%%%%%%%%%%%%%%%%%%%%%%%%%%

\section{Self-Adjoint Operators}

Let

\[
H
\]

be a Hilbert space.

An operator

\[
T:H\rightarrow H
\]

is self-adjoint whenever

\[
\boxed{
\langle Tx,y\rangle
=
\langle x,Ty\rangle
}
\]

for every

\[
x,y\in H.
\]

Self-adjoint operators generalize symmetric matrices.

%%%%%%%%%%%%%%%%%%%%%%%%%%%%%%%%%%%%%%%%%%%%%%%%%%%%%%%%%%%%%%%

\begin{theorem}[Spectral Theorem]

Every bounded self-adjoint operator admits a spectral decomposition.

Its spectrum is entirely real.

\end{theorem}

This theorem forms the mathematical foundation of quantum mechanics,
harmonic analysis, and many elliptic differential equations.

%%%%%%%%%%%%%%%%%%%%%%%%%%%%%%%%%%%%%%%%%%%%%%%%%%%%%%%%%%%%%%%

\section{Compactness}

Compactness generalizes the notion of finite bounded sets.

A subset

\[
K
\]

of a metric space is compact whenever every open cover possesses a
finite subcover.

In Euclidean spaces,

compactness is equivalent to

closedness together with boundedness.

%%%%%%%%%%%%%%%%%%%%%%%%%%%%%%%%%%%%%%%%%%%%%%%%%%%%%%%%%%%%%%%

\section{Weak Compactness}

Infinite-dimensional spaces exhibit behavior quite different from
finite-dimensional spaces.

Closed and bounded sets need not be compact.

However,

many important spaces possess weak compactness properties that allow
minimizing sequences to converge weakly.

Weak compactness is fundamental to the direct method of the calculus of
variations.

%%%%%%%%%%%%%%%%%%%%%%%%%%%%%%%%%%%%%%%%%%%%%%%%%%%%%%%%%%%%%%%

\section{Reflexive Spaces}

A Banach space

\[
X
\]

is called reflexive if the canonical embedding into its double dual,

\[
X^{**},
\]

is an isomorphism.

Hilbert spaces,

together with

\[
L^p
\]

spaces for

\[
1<p<\infty,
\]

are reflexive.

Reflexivity provides many of the compactness properties required in
existence proofs.

%%%%%%%%%%%%%%%%%%%%%%%%%%%%%%%%%%%%%%%%%%%%%%%%%%%%%%%%%%%%%%%

\section{Sobolev Spaces}

Many differential equations require functions whose derivatives exist
only in a weak sense.

For an integer

\[
k\ge0,
\]

the Sobolev space

\[
\boxed{
H^k(\Omega)
=
\{
u
:
D^\alpha u
\in
L^2(\Omega),
\ |\alpha|\le k
\}
}
\]

consists of square-integrable functions whose weak derivatives up to
order

\[
k
\]

are also square integrable.

Sobolev spaces provide the natural setting for weak formulations of
partial differential equations.

%%%%%%%%%%%%%%%%%%%%%%%%%%%%%%%%%%%%%%%%%%%%%%%%%%%%%%%%%%%%%%%

\section{Weak Derivatives}

A function need not possess classical derivatives to satisfy a
differential equation.

Instead,

one defines derivatives through integration by parts.

A function

\[
v
\]

is the weak derivative of

\[
u
\]

if

\[
\boxed{
\int_\Omega
u
\frac{\partial\phi}{\partial x_i}
\,dx
=
-
\int_\Omega
v
\phi
\,dx
}
\]

for every smooth compactly supported test function

\[
\phi.
\]

Weak derivatives substantially enlarge the class of admissible
solutions.

%%%%%%%%%%%%%%%%%%%%%%%%%%%%%%%%%%%%%%%%%%%%%%%%%%%%%%%%%%%%%%%

\section{Weak Solutions}

Rather than satisfying a differential equation pointwise,

a weak solution satisfies an integrated variational identity.

Weak formulations possess several important advantages.

\begin{itemize}
\item Existence is easier to establish.
\item Lower regularity is required.
\item Numerical approximation becomes natural.
\item Variational methods apply directly.
\end{itemize}

Most modern finite-element methods compute weak rather than classical
solutions.

%%%%%%%%%%%%%%%%%%%%%%%%%%%%%%%%%%%%%%%%%%%%%%%%%%%%%%%%%%%%%%%

\section{The Lax--Milgram Theorem}

One of the central existence theorems for elliptic problems is the
Lax--Milgram theorem.

\begin{theorem}[Lax--Milgram]

Let

\[
H
\]

be a Hilbert space,

and let

\[
a(\cdot,\cdot)
\]

be a continuous coercive bilinear form.

Then for every continuous linear functional

\[
f,
\]

there exists a unique element

\[
u\in H
\]

such that

\[
\boxed{
a(u,v)
=
f(v)
\qquad
\forall v\in H.
}
\]

\end{theorem}

This theorem establishes existence and uniqueness for a large class of
boundary-value problems.

%%%%%%%%%%%%%%%%%%%%%%%%%%%%%%%%%%%%%%%%%%%%%%%%%%%%%%%%%%%%%%%

\section{The Hahn--Banach Theorem}

The Hahn--Banach theorem is one of the fundamental structural results of
functional analysis.

\begin{theorem}[Hahn--Banach]

Every bounded linear functional defined on a subspace of a normed vector
space admits an extension to the entire space without increasing its
norm.

\end{theorem}

This theorem underlies separation theorems, duality theory, and convex
optimization.

%%%%%%%%%%%%%%%%%%%%%%%%%%%%%%%%%%%%%%%%%%%%%%%%%%%%%%%%%%%%%%%

\section{Applications}

Functional analysis provides the mathematical framework for

\begin{itemize}
\item variational calculus,
\item partial differential equations,
\item optimal control,
\item inverse problems,
\item finite-element methods,
\item continuum mechanics,
\item signal processing,
\item infinite-dimensional optimization.
\end{itemize}

Nearly every modern analytical method employed throughout this volume is
naturally formulated within Banach or Hilbert spaces.

%%%%%%%%%%%%%%%%%%%%%%%%%%%%%%%%%%%%%%%%%%%%%%%%%%%%%%%%%%%%%%%

\section{Concluding Remarks}

Functional analysis unifies geometry, topology, and linear algebra in
the infinite-dimensional setting.

Compactness enables existence proofs.

Duality characterizes optimization.

Weak convergence permits stable approximation.

Sobolev spaces accommodate generalized derivatives.

Operator theory extends matrix methods to spaces of functions.

Together these concepts provide the analytical foundation for the study
of partial differential equations, which form the subject of the next
appendix.

%%%%%%%%%%%%%%%%%%%%%%%%%%%%%%%%%%%%%%%%%%%%%%%%%%%%%%%%%%%%%%%%%%%%%%
\chapter{Partial Differential Equations}

\section{Purpose}

Many infrastructure systems evolve continuously in both space and time.

Examples include

\begin{itemize}
\item electrical potentials,
\item fluid flow,
\item heat transport,
\item traffic density,
\item pollutant dispersion,
\item groundwater movement,
\item communication signal propagation.
\end{itemize}

The governing equations for these systems are partial differential
equations (PDEs).

This appendix develops the principal classes of PDEs together with their
analytical structure and applications.

%%%%%%%%%%%%%%%%%%%%%%%%%%%%%%%%%%%%%%%%%%%%%%%%%%%%%%%%%%%%%%%

\section{Partial Derivatives}

Let

\[
u=u(x^1,\ldots,x^n)
\]

be a scalar field.

The partial derivative with respect to

\[
x^i
\]

is

\[
\boxed{
\frac{\partial u}{\partial x^i}.
}
\]

Higher-order derivatives are written

\[
\frac{\partial^2u}
{\partial x^i\partial x^j},
\]

and collectively form the Hessian matrix.

%%%%%%%%%%%%%%%%%%%%%%%%%%%%%%%%%%%%%%%%%%%%%%%%%%%%%%%%%%%%%%%

\section{General Form of a PDE}

A partial differential equation relates an unknown function to its
partial derivatives.

Its general form is

\[
\boxed{
F
\left(
x,
u,
\nabla u,
\nabla^2u,
\ldots
\right)
=
0.
}
\]

Unlike ordinary differential equations,

PDEs involve several independent variables.

%%%%%%%%%%%%%%%%%%%%%%%%%%%%%%%%%%%%%%%%%%%%%%%%%%%%%%%%%%%%%%%

\section{Linear and Nonlinear Equations}

A PDE is linear if the unknown function and all of its derivatives
appear linearly.

Otherwise,

the equation is nonlinear.

Linear equations often admit superposition,

whereas nonlinear equations may exhibit shocks,

instabilities,

pattern formation,

and multiple equilibria.

%%%%%%%%%%%%%%%%%%%%%%%%%%%%%%%%%%%%%%%%%%%%%%%%%%%%%%%%%%%%%%%

\section{Order}

The order of a PDE is determined by its highest derivative.

Examples include

\[
\begin{array}{ll}
\text{First order} &
u_t+c\,u_x=0,
\\[1ex]
\text{Second order} &
u_{tt}-c^2u_{xx}=0.
\end{array}
\]

Second-order equations dominate continuum physics and infrastructure
applications.

%%%%%%%%%%%%%%%%%%%%%%%%%%%%%%%%%%%%%%%%%%%%%%%%%%%%%%%%%%%%%%%

\section{Classification of Second-Order PDEs}

Consider the equation

\[
Au_{xx}
+
2Bu_{xy}
+
Cu_{yy}
+\cdots
=0.
\]

The discriminant

\[
\boxed{
\Delta
=
B^2-AC
}
\]

determines the equation's type.

\begin{center}
\begin{tabular}{ll}
\toprule
Condition & Type\\
\midrule
$\Delta<0$ & Elliptic\\
$\Delta=0$ & Parabolic\\
$\Delta>0$ & Hyperbolic\\
\bottomrule
\end{tabular}
\end{center}

This classification governs both qualitative behavior and analytical
techniques.

%%%%%%%%%%%%%%%%%%%%%%%%%%%%%%%%%%%%%%%%%%%%%%%%%%%%%%%%%%%%%%%

\section{Boundary Conditions}

Solutions are determined not only by the differential equation but also
by boundary conditions.

Common examples include

Dirichlet conditions,

\[
\boxed{
u=g
\quad
\text{on }\partial\Omega,
}
\]

Neumann conditions,

\[
\boxed{
\frac{\partial u}{\partial n}
=h,
}
\]

and Robin conditions,

\[
\boxed{
a\,u
+
b\,
\frac{\partial u}{\partial n}
=
c.
}
\]

%%%%%%%%%%%%%%%%%%%%%%%%%%%%%%%%%%%%%%%%%%%%%%%%%%%%%%%%%%%%%%%

\section{Initial Conditions}

Time-dependent equations additionally require initial data,

\[
\boxed{
u(x,0)
=
u_0(x).
}
\]

For second-order evolution equations,

one also specifies

\[
u_t(x,0).
\]

Together,

initial and boundary conditions determine a well-posed problem.

%%%%%%%%%%%%%%%%%%%%%%%%%%%%%%%%%%%%%%%%%%%%%%%%%%%%%%%%%%%%%%%

\section{The Laplace Equation}

One of the most fundamental elliptic equations is

\[
\boxed{
\Delta u
=
0.
}
\]

Solutions are called harmonic functions.

They describe steady-state phenomena including

electrostatic potential,

steady groundwater flow,

stationary temperature fields,

and equilibrium pressure distributions.

%%%%%%%%%%%%%%%%%%%%%%%%%%%%%%%%%%%%%%%%%%%%%%%%%%%%%%%%%%%%%%%

\section{The Poisson Equation}

Introducing a source term yields

\[
\boxed{
\Delta u
=
f.
}
\]

The function

\[
f
\]

represents distributed production or consumption.

Poisson equations arise whenever equilibrium is maintained in the
presence of internal sources.

%%%%%%%%%%%%%%%%%%%%%%%%%%%%%%%%%%%%%%%%%%%%%%%%%%%%%%%%%%%%%%%

\section{The Heat Equation}

The diffusion equation is

\[
\boxed{
u_t
=
D
\Delta u,
}
\]

where

\[
D
\]

is the diffusion coefficient.

Heat,

pollutants,

information,

traffic density,

and many other quantities evolve according to diffusion-like dynamics.

%%%%%%%%%%%%%%%%%%%%%%%%%%%%%%%%%%%%%%%%%%%%%%%%%%%%%%%%%%%%%%%

\section{The Wave Equation}

Wave propagation is governed by

\[
\boxed{
u_{tt}
=
c^2
\Delta u.
}
\]

Solutions propagate disturbances with finite speed.

Applications include

acoustics,

elasticity,

communication,

and seismic waves.

%%%%%%%%%%%%%%%%%%%%%%%%%%%%%%%%%%%%%%%%%%%%%%%%%%%%%%%%%%%%%%%

\section{Summary}

Partial differential equations describe the evolution of fields across
space and time.

Elliptic equations characterize equilibrium,

parabolic equations govern diffusion,

and hyperbolic equations describe wave propagation.

Together they provide the mathematical language for continuum models
throughout infrastructure science.

The following sections develop conservation laws, weak formulations,
existence theory, and numerical approximation methods for these
equations.

%%%%%%%%%%%%%%%%%%%%%%%%%%%%%%%%%%%%%%%%%%%%%%%%%%%%%%%%%%%%%%%

\section{Conservation Laws}

Many partial differential equations arise from the principle of local
conservation.

Let

\[
u(x,t)
\]

denote the density of a conserved quantity,

and let

\[
J(x,t)
\]

denote its flux.

Local conservation is expressed by

\[
\boxed{
\frac{\partial u}{\partial t}
+
\nabla\cdot J
=
S,
}
\]

where

\[
S
\]

represents distributed sources or sinks.

When

\[
S=0,
\]

the total quantity changes only through boundary flux.

%%%%%%%%%%%%%%%%%%%%%%%%%%%%%%%%%%%%%%%%%%%%%%%%%%%%%%%%%%%%%%%

\section{The Continuity Equation}

If the flux is transported by a velocity field

\[
v,
\]

then

\[
J
=
uv,
\]

and the conservation law becomes

\[
\boxed{
\frac{\partial u}{\partial t}
+
\nabla\cdot(uv)
=
S.
}
\]

This equation governs

mass conservation,

traffic flow,

charge transport,

population dynamics,

and commodity movement.

%%%%%%%%%%%%%%%%%%%%%%%%%%%%%%%%%%%%%%%%%%%%%%%%%%%%%%%%%%%%%%%

\section{Advection--Diffusion Equation}

Many physical systems exhibit both transport and diffusion.

The governing equation is

\[
\boxed{
\frac{\partial u}{\partial t}
+
v\cdot\nabla u
=
D\Delta u
+
S.
}
\]

The advection term transports information,

while the diffusion term smooths spatial gradients.

%%%%%%%%%%%%%%%%%%%%%%%%%%%%%%%%%%%%%%%%%%%%%%%%%%%%%%%%%%%%%%%

\section{Reaction--Diffusion Equations}

If local reactions occur simultaneously with diffusion,

one obtains

\[
\boxed{
u_t
=
D\Delta u
+
R(u),
}
\]

where

\[
R(u)
\]

describes local production or consumption.

Reaction--diffusion systems model

chemical kinetics,

biological growth,

ecological spread,

and numerous pattern-forming phenomena.

%%%%%%%%%%%%%%%%%%%%%%%%%%%%%%%%%%%%%%%%%%%%%%%%%%%%%%%%%%%%%%%

\section{Systems of Partial Differential Equations}

Many applications involve several interacting fields.

A general coupled system may be written

\[
\boxed{
\frac{\partial U}{\partial t}
=
F
\left(
U,
\nabla U,
\nabla^2U
\right),
}
\]

where

\[
U
=
(u_1,\ldots,u_m).
\]

Electrical,

hydraulic,

transportation,

and communication networks all produce coupled systems rather than single
equations.

%%%%%%%%%%%%%%%%%%%%%%%%%%%%%%%%%%%%%%%%%%%%%%%%%%%%%%%%%%%%%%%

\section{Weak Formulation}

Classical differentiability is often unnecessarily restrictive.

Instead,

multiply the governing equation by a test function

\[
\phi,
\]

integrate over the domain,

and transfer derivatives using integration by parts.

For Poisson's equation,

\[
-\Delta u=f,
\]

one obtains

\[
\boxed{
\int_\Omega
\nabla u\cdot\nabla\phi
\,dV
=
\int_\Omega
f\phi
\,dV.
}
\]

Weak formulations enlarge the admissible solution space while preserving
the underlying physical principles.

%%%%%%%%%%%%%%%%%%%%%%%%%%%%%%%%%%%%%%%%%%%%%%%%%%%%%%%%%%%%%%%

\section{Existence and Uniqueness}

A mathematical model is useful only if its solutions exist and are
uniquely determined.

A boundary-value problem is said to be well posed when

\begin{enumerate}
\item a solution exists,
\item the solution is unique,
\item the solution depends continuously upon the data.
\end{enumerate}

These criteria were formalized by Hadamard and serve as the standard
definition of a physically meaningful PDE problem.

%%%%%%%%%%%%%%%%%%%%%%%%%%%%%%%%%%%%%%%%%%%%%%%%%%%%%%%%%%%%%%%

\section{Maximum Principles}

Elliptic and parabolic equations satisfy remarkable comparison
properties.

\begin{theorem}[Maximum Principle]

Let

\[
u
\]

satisfy

\[
\Delta u
\ge0
\]

within a bounded domain.

Then

\[
u
\]

cannot attain a strict interior maximum unless it is constant.

\end{theorem}

Maximum principles provide powerful uniqueness results and useful
qualitative information about solutions.

%%%%%%%%%%%%%%%%%%%%%%%%%%%%%%%%%%%%%%%%%%%%%%%%%%%%%%%%%%%%%%%

\section{Energy Methods}

Many PDEs admit natural energy functionals.

For the heat equation,

one defines

\[
\boxed{
E(t)
=
\frac12
\int_\Omega
u^2\,dV.
}
\]

Differentiating with respect to time yields

\[
\boxed{
\frac{dE}{dt}
=
-
D
\int_\Omega
|\nabla u|^2
\,dV
\le
0,
}
\]

assuming homogeneous boundary conditions.

Thus diffusion monotonically decreases the total energy.

%%%%%%%%%%%%%%%%%%%%%%%%%%%%%%%%%%%%%%%%%%%%%%%%%%%%%%%%%%%%%%%

\section{Finite Difference Methods}

Finite differences approximate derivatives using nearby grid values.

For example,

\[
\boxed{
\frac{\partial u}{\partial x}
\approx
\frac{
u_{i+1}-u_i
}
{\Delta x},
}
\]

and

\[
\boxed{
\frac{\partial^2u}{\partial x^2}
\approx
\frac{
u_{i+1}
-
2u_i
+
u_{i-1}
}
{\Delta x^2}.
}
\]

These approximations convert PDEs into algebraic systems suitable for
computer solution.

%%%%%%%%%%%%%%%%%%%%%%%%%%%%%%%%%%%%%%%%%%%%%%%%%%%%%%%%%%%%%%%

\section{Finite Element Methods}

Rather than approximating derivatives directly,

the finite-element method approximates the solution itself by piecewise
basis functions.

One writes

\[
\boxed{
u_h
=
\sum_i
c_i
\phi_i,
}
\]

where

\[
\phi_i
\]

are locally supported basis functions.

Substitution into the weak formulation produces a sparse linear system
for the coefficients

\[
c_i.
\]

Finite-element methods are particularly effective for irregular
geometries and complex infrastructure domains.

%%%%%%%%%%%%%%%%%%%%%%%%%%%%%%%%%%%%%%%%%%%%%%%%%%%%%%%%%%%%%%%

\section{Finite Volume Methods}

Finite-volume methods integrate the conservation law over small control
volumes.

The numerical update is obtained directly from flux balance across cell
boundaries.

Consequently,

global conservation is preserved exactly at the discrete level,

making finite-volume methods especially valuable for transport and flow
problems.

%%%%%%%%%%%%%%%%%%%%%%%%%%%%%%%%%%%%%%%%%%%%%%%%%%%%%%%%%%%%%%%

\section{Applications to Infrastructure}

Partial differential equations provide mathematical models for

\begin{itemize}
\item groundwater flow,
\item electrical potential,
\item traffic density,
\item heat transport,
\item pollutant dispersion,
\item elastic deformation,
\item fluid mechanics,
\item communication signal propagation,
\item environmental diffusion.
\end{itemize}

Although the governing equations differ,

they share a common analytical framework based upon conservation,
boundary conditions,
weak formulations,
and numerical approximation.

%%%%%%%%%%%%%%%%%%%%%%%%%%%%%%%%%%%%%%%%%%%%%%%%%%%%%%%%%%%%%%%

\section{Concluding Remarks}

Partial differential equations extend ordinary differential equations to
fields distributed over space and time.

Conservation laws provide their physical foundation.

Weak formulations permit generalized solutions.

Energy methods establish stability,

while finite-difference,
finite-element,
and finite-volume methods enable practical numerical computation.

Together these ideas complete the analytical framework needed for the
continuum models developed throughout this volume.

The following appendix develops probability and stochastic processes,
allowing uncertainty, noise, and random disturbances to be incorporated
systematically into infrastructure analysis.

%%%%%%%%%%%%%%%%%%%%%%%%%%%%%%%%%%%%%%%%%%%%%%%%%%%%%%%%%%%%%%%%%%%%%%
\chapter{Probability and Stochastic Processes}

\section{Purpose}

Infrastructure systems operate under uncertainty.

Demand fluctuates,

weather varies,

equipment fails,

traffic changes,

communication channels experience noise,

and human behavior remains only partially predictable.

Probability theory provides the mathematical framework for describing
uncertainty, while stochastic processes describe its evolution through
time.

%%%%%%%%%%%%%%%%%%%%%%%%%%%%%%%%%%%%%%%%%%%%%%%%%%%%%%%%%%%%%%%

\section{Probability Spaces}

The mathematical foundation of probability consists of three objects,

\[
(\Omega,\mathcal F,P),
\]

where

\begin{itemize}
\item
\[
\Omega
\]
is the sample space,

\item
\[
\mathcal F
\]
is a sigma-algebra of measurable events,

\item
\[
P
\]
is a probability measure satisfying

\[
P(\Omega)=1.
\]
\end{itemize}

Every probabilistic model begins with this structure.

%%%%%%%%%%%%%%%%%%%%%%%%%%%%%%%%%%%%%%%%%%%%%%%%%%%%%%%%%%%%%%%

\section{Events}

An event is a measurable subset

\[
A\subseteq\Omega.
\]

Its probability satisfies

\[
0
\le
P(A)
\le
1.
\]

If

\[
A
\]

and

\[
B
\]

are disjoint,

then

\[
\boxed{
P(A\cup B)
=
P(A)+P(B).
}
\]

%%%%%%%%%%%%%%%%%%%%%%%%%%%%%%%%%%%%%%%%%%%%%%%%%%%%%%%%%%%%%%%

\section{Conditional Probability}

When additional information becomes available,

probabilities are updated through conditioning.

For

\[
P(B)>0,
\]

the conditional probability is

\[
\boxed{
P(A|B)
=
\frac{P(A\cap B)}
{P(B)}.
}
\]

Conditional probability forms the basis of Bayesian inference and
stochastic control.

%%%%%%%%%%%%%%%%%%%%%%%%%%%%%%%%%%%%%%%%%%%%%%%%%%%%%%%%%%%%%%%

\section{Bayes' Theorem}

Reversing conditional probabilities gives

\[
\boxed{
P(A|B)
=
\frac{
P(B|A)P(A)
}
{P(B)}.
}
\]

Bayes' theorem provides the mathematical rule for updating beliefs using
new observations.

%%%%%%%%%%%%%%%%%%%%%%%%%%%%%%%%%%%%%%%%%%%%%%%%%%%%%%%%%%%%%%%

\section{Random Variables}

A random variable is a measurable mapping

\[
X:\Omega\rightarrow\mathbb R.
\]

Rather than studying elementary outcomes directly,

probability theory usually analyzes random variables derived from those
outcomes.

%%%%%%%%%%%%%%%%%%%%%%%%%%%%%%%%%%%%%%%%%%%%%%%%%%%%%%%%%%%%%%%

\section{Probability Distributions}

A random variable is characterized by its probability distribution.

For continuous variables,

the probability density function

\[
f(x)
\]

satisfies

\[
\boxed{
P(a\le X\le b)
=
\int_a^b
f(x)\,dx.
}
\]

The cumulative distribution function is

\[
F(x)
=
P(X\le x).
\]

%%%%%%%%%%%%%%%%%%%%%%%%%%%%%%%%%%%%%%%%%%%%%%%%%%%%%%%%%%%%%%%

\section{Expectation}

The expected value of a random variable is

\[
\boxed{
\mathbb E[X]
=
\int_\Omega
X\,dP.
}
\]

For continuous variables,

\[
\mathbb E[X]
=
\int_{-\infty}^{\infty}
x
f(x)
\,dx.
\]

Expectation represents the probabilistic average of repeated
observations.

%%%%%%%%%%%%%%%%%%%%%%%%%%%%%%%%%%%%%%%%%%%%%%%%%%%%%%%%%%%%%%%

\section{Variance}

The variance measures dispersion about the mean.

It is defined by

\[
\boxed{
\operatorname{Var}(X)
=
\mathbb E
[(X-\mathbb E[X])^2].
}
\]

The standard deviation is

\[
\sigma
=
\sqrt{\operatorname{Var}(X)}.
\]

%%%%%%%%%%%%%%%%%%%%%%%%%%%%%%%%%%%%%%%%%%%%%%%%%%%%%%%%%%%%%%%

\section{Covariance}

Given two random variables,

their covariance is

\[
\boxed{
\operatorname{Cov}(X,Y)
=
\mathbb E
[(X-\mu_X)(Y-\mu_Y)].
}
\]

Covariance quantifies statistical dependence.

%%%%%%%%%%%%%%%%%%%%%%%%%%%%%%%%%%%%%%%%%%%%%%%%%%%%%%%%%%%%%%%

\section{Correlation}

The correlation coefficient is

\[
\boxed{
\rho
=
\frac{
\operatorname{Cov}(X,Y)
}
{
\sigma_X\sigma_Y
}.
}
\]

Its values lie between

\[
-1
\]

and

\[
1.
\]

Correlation measures linear association independently of scale.

%%%%%%%%%%%%%%%%%%%%%%%%%%%%%%%%%%%%%%%%%%%%%%%%%%%%%%%%%%%%%%%

\section{Important Probability Distributions}

Several probability distributions occur repeatedly in engineering and
scientific applications.

\begin{center}
\begin{tabular}{ll}
\toprule
Distribution & Typical Application\\
\midrule
Bernoulli & Binary events\\
Binomial & Repeated trials\\
Poisson & Rare events\\
Exponential & Waiting times\\
Gaussian & Measurement noise\\
Gamma & Service times\\
Weibull & Reliability analysis\\
Lognormal & Multiplicative growth\\
\bottomrule
\end{tabular}
\end{center}

%%%%%%%%%%%%%%%%%%%%%%%%%%%%%%%%%%%%%%%%%%%%%%%%%%%%%%%%%%%%%%%

\section{The Gaussian Distribution}

The normal distribution has density

\[
\boxed{
f(x)
=
\frac{1}
{\sqrt{2\pi\sigma^2}}
\exp
\left(
-
\frac{(x-\mu)^2}
{2\sigma^2}
\right).
}
\]

Its importance follows from the Central Limit Theorem,

which explains its widespread appearance in natural and engineered
systems.

%%%%%%%%%%%%%%%%%%%%%%%%%%%%%%%%%%%%%%%%%%%%%%%%%%%%%%%%%%%%%%%

\section{Independence}

Random variables

\[
X
\]

and

\[
Y
\]

are independent whenever

\[
\boxed{
P(X,Y)
=
P(X)P(Y).
}
\]

Independence simplifies both theoretical analysis and computational
modeling,

although it is rarely exact in large infrastructure systems.

\newpage

\begin{thebibliography}{99}

\bibitem{alexandrov2005}
Alexandrov, P. S.
\emph{Combinatorial Topology}.
Dover Publications, 2005.

\bibitem{berge1985}
Berge, C.
\emph{Graphs and Hypergraphs}.
North-Holland, 1985.

\bibitem{boyd2004}
Boyd, S., and Vandenberghe, L.
\emph{Convex Optimization}.
Cambridge University Press, 2004.

\bibitem{deberg2008}
de Berg, M., Cheong, O., van Kreveld, M., and Overmars, M.
\emph{Computational Geometry: Algorithms and Applications}.
3rd ed., Springer, 2008.

\bibitem{dijkstra1959}
Dijkstra, E. W.
``A Note on Two Problems in Connexion with Graphs.''
\emph{Numerische Mathematik},
1 (1959), 269--271.

\bibitem{edelsbrunner2001}
Edelsbrunner, H.
\emph{Geometry and Topology for Mesh Generation}.
Cambridge University Press, 2001.

\bibitem{edelsbrunner2010}
Edelsbrunner, H., and Harer, J.
\emph{Computational Topology: An Introduction}.
American Mathematical Society, 2010.

\bibitem{ford1956}
Ford, L. R., and Fulkerson, D. R.
``Maximal Flow Through a Network.''
\emph{Canadian Journal of Mathematics},
8 (1956), 399--404.

\bibitem{gallager2013}
Gallager, R. G.
\emph{Stochastic Processes: Theory for Applications}.
Cambridge University Press, 2013.

\bibitem{grunbaum2003}
Grünbaum, B.
\emph{Convex Polytopes}.
2nd ed., Springer, 2003.

\bibitem{gross2006}
Gross, J. L., and Yellen, J.
\emph{Graph Theory and Its Applications}.
2nd ed., Chapman \& Hall/CRC, 2006.

\bibitem{lee2013}
Lee, J. M.
\emph{Introduction to Smooth Manifolds}.
2nd ed., Springer, 2013.

\bibitem{okabe2000}
Okabe, A., Boots, B., Sugihara, K., and Chiu, S. N.
\emph{Spatial Tessellations: Concepts and Applications of Voronoi Diagrams}.
2nd ed., Wiley, 2000.

\bibitem{papadimitriou1998}
Papadimitriou, C. H., and Steiglitz, K.
\emph{Combinatorial Optimization: Algorithms and Complexity}.
Dover Publications, 1998.

\bibitem{press2007}
Press, W. H., Teukolsky, S. A., Vetterling, W. T., and Flannery, B. P.
\emph{Numerical Recipes: The Art of Scientific Computing}.
3rd ed., Cambridge University Press, 2007.

\bibitem{spivak1979}
Spivak, M.
\emph{A Comprehensive Introduction to Differential Geometry}.
Vols. I--V.
Publish or Perish, 1979.

\bibitem{strang1988}
Strang, G.
\emph{Linear Algebra and Its Applications}.
3rd ed., Brooks/Cole, 1988.

\bibitem{thurston1997}
Thurston, W. P.
\emph{Three-Dimensional Geometry and Topology}.
Vol. 1.
Princeton University Press, 1997.

\bibitem{west2001}
West, D. B.
\emph{Introduction to Graph Theory}.
2nd ed., Prentice Hall, 2001.

\bibitem{zwillinger2018}
Zwillinger, D. (ed.)
\emph{CRC Standard Mathematical Tables and Formulae}.
33rd ed., CRC Press, 2018.

\end{thebibliography}

\end{document}
